\chapter*{Possible misprints, not verified}
\addcontentsline{toc}{chapter}{Possible misprints, not verified}
These proposals remain open. No correction is made on the authority of these notices; the agreed text is retained. Readings and reasons are quoted in the language of the source records.
\begingroup\small\raggedright\sloppy\setlength{\parskip}{4pt}
\par\noindent\textbf{OM-OPUS-0004 -- t1 I p.29 -- OPUS}\par
first line, 'si et seulement si \ensuremath{\Psi}(K)'\par
\emph{Reported reading}: \ensuremath{\Psi}(\ensuremath{\underline{K}}) with an underlined K\par
\emph{Proposal}: \$\textbackslash{}Psi(K)\$\par
The change from script K to plain K requires a verified identification of the presheaf and the underlined printed glyph.\par
\par\noindent\textbf{OM-OPUS-0006 -- t1 I p.56 -- OPUS}\par
same formula, base of the fibre product\par
\emph{Reported reading}: \ensuremath{\phi}\_\{i(Y)\}\par
\emph{Proposal}: \textbackslash{}varphi\_i(Y)\par
The intended scope of the subscript in phi\_i(Y) requires checking the definition of the family phi\_i.\par
\par\noindent\textbf{OM-OPUS-0018 -- t1 I p.190 -- OPUS}\par
first paragraph\par
\emph{Reported reading}: X et X' sont des éléments de U munis\par
\emph{Proposal}: Si \$X\$ et \$X'\$ sont des éléments\par
The proposed Si repairs a possible missing antecedent, but the full grammatical clause must establish that exact insertion.\par
\par\noindent\textbf{OM-OPUS-0019 -- t1 I p.196 -- OPUS}\par
(FI.2), « et si card(I) < d »\par
\emph{Reported reading}: card(I) < d (plain d, not the double-underlined cardinal)\par
\emph{Proposal}: \$\textbackslash{}card(I)<\textbackslash{}sheaf\{\textbackslash{}sheaf\{d\}\}\$\par
The record reports a typographical distinction between plain and underlined d; the symbol identity requires the surrounding cardinal convention.\par
\par\noindent\textbf{OM-OPUS-0022 -- t1 I p.199 -- OPUS}\par
(U.IV), displayed formula\par
\emph{Reported reading}: I =\allowbreak  \ensuremath{\cup}\_\{\ensuremath{\lambda} \ensuremath{\alpha} K\} [\ensuremath{\varepsilon}, \ensuremath{\alpha}(\ensuremath{\lambda})]\par
\emph{Proposal}: \textbackslash{}bigcup\_\{\textbackslash{}lambda\textbackslash{}in K\}\par
The glyph alpha versus membership is a scan-dependent reading; preserve the agreed source pending that comparison.\par
\par\noindent\textbf{OM-OPUS-0023 -- t1 I p.210 -- OPUS}\par
"(début du \$\{\textbackslash{}rm n\textasciicircum{}\{\textbackslash{}circ\}\}\$)"\par
\emph{Reported reading}: (début du n°) — no number\par
\emph{Proposal}: (début du \$\{\textbackslash{}rm n\textasciicircum{}\{\textbackslash{}circ\}\}\$\textasciitilde{}6)\par
The omitted paragraph number is not supplied by the sentence; 6 remains a conjectured target.\par
\par\noindent\textbf{OM-OPUS-0024 -- t1 I p.211 -- OPUS}\par
Lemme 3, second sequence\par
\emph{Reported reading}: (E\_n)\_\{n\ensuremath{\geq}1\} (\ensuremath{\phi}\_n)\_\{n\ensuremath{\geq}1\} d'ensembles finis E\_n et d'applications f\_n\par
\emph{Proposal}: \$(f\_n)\_\{n \textbackslash{}geq 1\}\$\par
The reported Greek glyph and the Latin f require the scan to identify the printed symbol precisely.\par
\par\noindent\textbf{OM-OPUS-0027 -- t1 II p.255 -- OPUS}\par
end of proof of 5.4\par
\emph{Reported reading}: pr\_2 : H\ensuremath{\times}\_K X \ensuremath{\to} X (twice)\par
\emph{Proposal}: \$\textbackslash{}pr\_2:H\textbackslash{}times\_K X\_\{\textbackslash{}lambda\}\textbackslash{}rightarrow X\_\{\textbackslash{}lambda\}\$\par
Restoring lambda changes both the pullback and its codomain; the projection's exact source must be checked.\par
\par\noindent\textbf{OM-OPUS-0029 -- t1 II p.261 -- OPUS}\par
formulas (*), (**), (***)\par
\emph{Reported reading}: i\_a(Ker(u)), a\_A coker(i\_a(u)), a\_A coim(i\_a(u))\par
\emph{Proposal}: i\_A(\textbackslash{}Ker(u)),\textbackslash{} a\_A\textbackslash{}coker(i\_A(u)),\textbackslash{} a\_A\textbackslash{}coim(i\_A(u))\par
The subscript A/a requires identifying the inclusion functor in all three formulas, not just regularizing its letter case.\par
\par\noindent\textbf{OM-OPUS-0033 -- t1 III p.294 -- OPUS}\par
Preuve of 5.2, item 4)\par
\emph{Reported reading}: Soit (G)\_i, i\ensuremath{\in}I\par
\emph{Proposal}: \$(G\_i)\$, \$i\textbackslash{}in I\$\par
The proposal moves i across a parenthesis in the family notation (G)\_i; the supplied notation alone does not establish an intended change in the indexed family.\par
\par\noindent\textbf{OM-OPUS-0045 -- t1 IV p.340 -- OPUS}\par
«le choix d'un morphisme (4.3.1)»\par
\emph{Reported reading}: (4.3.1)\par
\emph{Proposal}: (4.3.1), with the display f:E\textbackslash{}rightarrow P on p. 339 numbered (4.3.1)\par
The reference is present but the proposed new displayed number needs identification of the intended equation.\par
\par\noindent\textbf{OM-OPUS-0061 -- t1 IV p.418 -- OPUS}\par
7.8 b), 'Si \ensuremath{𝒪}' est un ensemble ordonné'\par
\emph{Reported reading}: Si \ensuremath{𝒪} est un ensemble ordonné satisfaisant (no prime)\par
\emph{Proposal}: \$\textbackslash{}Oo'\$\par
The prime distinguishes two ordered sets; the antecedent of this quantified set requires verification.\par
\par\noindent\textbf{OM-OPUS-0064 -- t1 IV p.436 -- OPUS}\par
b), end\par
\emph{Reported reading}: g* transforme objets de S en isomorphismes\par
\emph{Proposal}: transforme flèches de \$S\$ en isomorphismes\par
Replacing objects by arrows changes the described action of g*; the definition of S must settle it.\par
\par\noindent\textbf{OM-OPUS-0066 -- t1 IV p.452 -- OPUS}\par
Prop. 9.2.4 b), reference\par
\emph{Reported reading}: (1.2.7)\par
\emph{Proposal}: (\textbackslash{}hyperrefext[I][I.def2.7]\{2.7\})\par
The proposed reference is plausible but must be compared with the stated definition and its label.\par
\par\noindent\textbf{OM-OPUS-0067 -- t1 IV p.458 -- OPUS}\par
diagram after 'considérons le diagramme de topos'\par
\emph{Reported reading}: bottom arrow E\_\{/U\} \ensuremath{\to} E labelled i\par
\emph{Proposal}: \textbackslash{}ar[r]\textasciicircum{}\{j\}\par
The arrow label needs the definitions and endpoints of i and j in the gluing diagram.\par
\par\noindent\textbf{OM-OPUS-0068 -- t1 IV p.466 -- OPUS}\par
top diagrams (X —w— Y over j\_*j\textasciicircum{}*X ... j\_*j\textasciicircum{}*Y and i\_*i\textasciicircum{}*X ... i\_*i\textasciicircum{}*Y)\par
\emph{Reported reading}: only the left vertical arrows from X are drawn; there is no vertical arrow from Y and no arrow along the bottom rows (only the labels j\_*m', i\_*m)\par
\emph{Proposal}: complete squares as in the TeX\par
The record explicitly allows loss of faint strokes in the scan; no diagram completion is certified.\par
\par\noindent\textbf{OM-OPUS-0074 -- t1 IV p.487 -- OPUS}\par
e), homomorphisms of presheaves and of sheaves\par
\emph{Reported reading}: f\_*(Quot\_\{X'\}) \ensuremath{\longrightarrow} Quot\_X … f\_*(Q\_\{X'\}) \ensuremath{\longrightarrow} Q\_X d'où f\textasciicircum{}*(Q\_X) \ensuremath{\longrightarrow} Q\_\{X'\}\par
\emph{Proposal}: \$f\_\{\textbackslash{}ast\}(\textbackslash{}Quot\_\{X'\})\textbackslash{}leftarrow \textbackslash{}Quot\_\{X\}\$, \$f\_\{\textbackslash{}ast\}(Q\_\{X'\})\textbackslash{}leftarrow Q\_X\$\par
Reversing both arrows requires the constructed map of relations and the adjunction; the snippet alone does not verify that construction.\par
\par\noindent\textbf{OM-OPUS-0076 -- t1 IV p.495 -- OPUS}\par
first line, middle term\par
\emph{Reported reading}: \ensuremath{\Pi}\_i \ensuremath{𝓗}om(v*(e'),v*Y)\par
\emph{Proposal}: \textbackslash{}Pi\_I\textbackslash{}scrHom(v\textasciicircum{}\{\textbackslash{}ast\}(e'), v\textasciicircum{}\{\textbackslash{}ast\}Y)\par
The product index may be a bound index or the indexing object; the formula defining Pi must settle the capital.\par
\par\noindent\textbf{OM-OPUS-0078 -- t1 IV p.502 -- OPUS}\par
chain of isomorphisms proving b) (rows printed on p. 503), labels '(12.2.1)'\par
\emph{Reported reading}: (12.2.1)\par
\emph{Proposal}: (12.1.2)\par
The claimed replacement cross-reference requires comparison with both numbered adjunction formulas.\par
\par\noindent\textbf{OM-OPUS-0079 -- t1 IV p.502 -- OPUS}\par
chain proving b), third row (printed on p. 503)\par
\emph{Reported reading}: Hom\_Λ(M,j\_\{X*\}…)\par
\emph{Proposal}: \textbackslash{}Hom\_A(M, j\_\{X\textbackslash{}ast\}\par
Changing the ring subscript requires the module structures on both arguments of Hom.\par
\par\noindent\textbf{OM-OPUS-0080 -- t1 IV p.502 -- OPUS}\par
chain proving b), third row, right-hand side (printed on p. 503)\par
\emph{Reported reading}: Hom\_\{A/X\}(j\textasciicircum{}*\_X M,\ensuremath{𝓗}om\_E(j\textasciicircum{}*\_X N))\par
\emph{Proposal}: \textbackslash{}scrHom\_E (Y,j\textasciicircum{}\{\textbackslash{}ast\}\_XN)\par
The missing first argument Y requires checking the internal-Hom adjunction used in the chain.\par
\par\noindent\textbf{OM-OPUS-0081 -- t1 IV p.502 -- OPUS}\par
chain proving b), fourth row, end (printed on p. 503)\par
\emph{Reported reading}: Hom(Y,\ensuremath{𝓗}om\_\{A|X\}(j\textasciicircum{}*\_X M,j\textasciicircum{}*\_X M))\par
\emph{Proposal}: \textbackslash{}scrHom\_\{A| X\}(j\textasciicircum{}\{\textbackslash{}ast\}\_X M, j\textasciicircum{}\{\textbackslash{}ast\}\_XN))\par
The second internal-Hom argument must be traced through the chain before replacing M by N.\par
\par\noindent\textbf{OM-OPUS-0084 -- t1 IV p.515 -- OPUS}\par
diagram (14.2.1), left object\par
\emph{Reported reading}: (\_\{A|U\} \ensuremath{𝒰}) — a handwritten script U\par
\emph{Proposal}: (\{\}\_\{A|U\}U)\par
The script/plain U distinction needs the local notational identification and the printed glyph.\par
\par\noindent\textbf{OM-OPUS-0085 -- t2 V p.12 -- OPUS}\par
1.6.2, x : P \ensuremath{\to} E\par
\emph{Reported reading}: x : P \ensuremath{\longmapsto} E (mapsto arrow)\par
\emph{Proposal}: x : P \textbackslash{}to E\par
The arrow-style change requires establishing a misprint rather than a printed notation convention.\par
\par\noindent\textbf{OM-OPUS-0087 -- t2 V p.21 -- OPUS}\par
item g), 'acycliques pour lim\_C' and 'le complexe lim\_C(Z\_\{ES(C)\}\ensuremath{\otimes}M)'\par
\emph{Reported reading}: \ensuremath{\ell}im with a left arrow under it, twice\par
\emph{Proposal}: \textbackslash{}varinjlim\_\{C\}\par
Replacing two projective limits by inductive limits requires the homological construction of this complex.\par
\par\noindent\textbf{OM-OPUS-0088 -- t2 V p.27 -- OPUS}\par
(3.4.3), (3.4.4)\par
\emph{Reported reading}: Φ\textasciicircum{}q(M) : \ensuremath{\mathcal{H}}\textasciicircum{}q(M) \ensuremath{\to} \ensuremath{\mathcal{H}}\textasciicircum{}q(M) ; Φ\textasciicircum{}q\_X(M) : H\textasciicircum{}q(X,M) \ensuremath{\to} H\textasciicircum{}q(X,M)\par
\emph{Proposal}: \textbackslash{}check\{\textbackslash{}Hh\}\textasciicircum{}q(M) \textbackslash{}to \textbackslash{}Hh\textasciicircum{}q(M)\textbackslash{} ;\textbackslash{} \textbackslash{}check\{H\}\textasciicircum{}q(X, M) \textbackslash{}to H\textasciicircum{}q(X, M)\par
The source of both edge morphisms requires its Čech-cohomology definition and the degree conventions.\par
\par\noindent\textbf{OM-OPUS-0094 -- t2 V p.71 -- OPUS}\par
8.1.10 (ii)\par
\emph{Reported reading}: f : V,U,\ensuremath{\phi},\ensuremath{\lambda}) \ensuremath{\longrightarrow} (V',\ensuremath{\phi}',U',\ensuremath{\lambda}') est un triple (f\_1,f\_2,F\_3)\par
\emph{Proposal}: \$f: (V, U, \textbackslash{}phi, \textbackslash{}lambda) \textbackslash{}to (V', U', \textbackslash{}phi', \textbackslash{}lambda')\$ ... \$(f\_1, f\_2, f\_3)\$\par
This record bundles a tuple permutation and glyph repairs; each component needs a separate exact placement and definition check.\par
\par\noindent\textbf{OM-OPUS-0095 -- t2 V p.72 -- OPUS}\par
same formula, and after the diagram\par
\emph{Reported reading}: q\textasciicircum{}x U\_2 ; (\ensuremath{\lambda}\_1χ\_j , h\_2χ\_j) ; pour que h, χ\_j =\allowbreak  h\_2 χ\_j\par
\emph{Proposal}: \$q\textasciicircum{}\{\textbackslash{}ast\} U\textasciicircum{}2\$ ; \$(h\_1 \textbackslash{}chi\_j, h\_2 \textbackslash{}chi\_j)\$ ; \$h\_1 \textbackslash{}chi\_j =\allowbreak  h\_2 \textbackslash{}chi\_j\$\par
Three independent proposed repairs are bundled; their index and morphism identifications require the full diagram.\par
\par\noindent\textbf{OM-OPUS-0096 -- t2 V p.72 -- OPUS}\par
(i)\par
\emph{Reported reading}: \ensuremath{𝒮}\textasciicircum{}1 et \ensuremath{𝒮}\textasciicircum{}2 dans \ensuremath{𝒮}\textasciicircum{}xU ... \ensuremath{\mathcal{F}}(\ensuremath{𝒮}\textasciicircum{}1) et \ensuremath{\mathcal{F}}(\ensuremath{𝒮}\textasciicircum{}2)\par
\emph{Proposal}: \$\textbackslash{}lambda\textasciicircum{}1\$ et \$\textbackslash{}lambda\textasciicircum{}2\$ dans \$\textbackslash{}Ss\textasciicircum{}\{\textbackslash{}ast\}\_U\$ ... \$\textbackslash{}Ff(\textbackslash{}lambda\textasciicircum{}1)\$ et \$\textbackslash{}Ff(\textbackslash{}lambda\textasciicircum{}2)\$\par
The identification of the handwritten script glyph with lambda is expressly a scan reading.\par
\par\noindent\textbf{OM-OPUS-0097 -- t2 V p.73 -- OPUS}\par
(i) and (ii), several places\par
\emph{Reported reading}: \ensuremath{\mathcal{F}}(\ensuremath{𝒮}\textasciicircum{}1); \ensuremath{\mu}\_k \ensuremath{\to} 2; avec \ensuremath{𝒮}\textasciicircum{}2 remplacé par \ensuremath{𝒮}\textasciicircum{}1; diagram \ensuremath{\lambda}\_j \ensuremath{\to} \ensuremath{𝒮}, lower \ensuremath{\mu}; \ensuremath{\mathcal{F}}(\ensuremath{𝒮}\_j); \ensuremath{\lambda}\_j par \ensuremath{𝒮}\_k; \ensuremath{\mu}' \ensuremath{\in} Ob \textasciicircum{}x\_U\par
\emph{Proposal}: \$\textbackslash{}Ff(\textbackslash{}lambda\textasciicircum{}1)\$; \$\textbackslash{}lambda\textasciicircum{}2\$; \$\textbackslash{}lambda\textasciicircum{}2\$ remplacé par \$\textbackslash{}lambda\textasciicircum{}1\$; \$\textbackslash{}lambda\$ (and \$\textbackslash{}mu'\$ for the lower \$\textbackslash{}mu\$); \$\textbackslash{}Ff(\textbackslash{}lambda\_j)\$; \$\textbackslash{}lambda\_k\$; \$\textbackslash{}mu' \textbackslash{}in \textbackslash{}Ob \textbackslash{}Ss\textasciicircum{}\{\textbackslash{}ast\}\_U\$\par
The record bundles several handwritten-glyph readings and diagram changes; those readings cannot be certified without the scan.\par
\par\noindent\textbf{OM-OPUS-0099 -- t2 V p.80 -- OPUS}\par
diagram, indices\par
\emph{Reported reading}: Hom(\ensuremath{𝔉}\_1,\ensuremath{𝒢}), Hom(\ensuremath{𝔉}\_2,\ensuremath{𝒢})\par
\emph{Proposal}: \textbackslash{}Hhom(\textbackslash{}Ff\_0, \textbackslash{}Gg) ... \textbackslash{}Hhom(\textbackslash{}Ff\_1,\textbackslash{}Gg)\par
The presentation suggests indices 0 and 1, but the two diagram terms must be matched with its arrows before replacement.\par
\par\noindent\textbf{OM-OPUS-0102 -- t2 Vbis p.97 -- OPUS}\par
1.3.8, first line\par
\emph{Reported reading}: nous fournirent\par
\emph{Proposal}: nous fournissent\par
Both past and present verb forms are grammatical; the reported tense change needs the whole sentence.\par
\par\noindent\textbf{OM-OPUS-0106 -- t2 Vbis p.124 -- OPUS}\par
Proposition 3.1.1, \ensuremath{\eta}\textasciicircum{}f\_G, target\par
\emph{Reported reading}: L\textasciicircum{}+ ū\textasciicircum{}* (no index G)\par
\emph{Proposal}: \textbackslash{}LL\textasciicircum{}\{+\}\_G \textbackslash{}bar\{u\}\textasciicircum{}\textbackslash{}ast\par
The proposed corrected field and quoted TeX put G on different operators; that ambiguity prevents an exact repair.\par
\par\noindent\textbf{OM-OPUS-0111 -- t2 Vbis p.154 -- OPUS}\par
5.2.1, first complex\par
\emph{Reported reading}: \ensuremath{\underline{R}}⁺ ū\_+(\ensuremath{\Omega}\textasciicircum{}\{Zar.\}\_\{X/k\})\par
\emph{Proposal}: \textbackslash{}bar\{u\}\_\{\textbackslash{}ast\}\par
The proposed star is described as presumptive; the notation for the three direct-image complexes must settle it.\par
\par\noindent\textbf{OM-OPUS-0112 -- t2 VI p.169 -- OPUS}\par
proof of 1.4, formula X\_i =\allowbreak  X\_i \ensuremath{\times}\_X \ensuremath{\coprod}\_\{j=\allowbreak J\} X\_j =\allowbreak  ...\par
\emph{Reported reading}: \ensuremath{\coprod} with index j=\allowbreak J\par
\emph{Proposal}: \textbackslash{}amalg\_\{j\textbackslash{}in J\}\par
The equality/membership mark in the coproduct index requires checking the printed glyph and the indexing convention.\par
\par\noindent\textbf{OM-OPUS-0114 -- t2 VI p.190 -- OPUS}\par
same sentence\par
\emph{Reported reading}: donc X\_o est limite inductive filtrante des Coker(R\_i \ensuremath{\rightrightarrows} X\_o) =\allowbreak  X\_i\par
\emph{Proposal}: donc \$X\$ est limite inductive filtrante des \$\textbackslash{}Coker(R\_i \textbackslash{}rightrightarrows X\_0) =\allowbreak  X'\_i\$\par
The record combines a missing clause and two variable changes; the absent clause cannot be reconstructed from this snippet.\par
\par\noindent\textbf{OM-OPUS-0126 -- t2 VI p.243 -- OPUS}\par
Remarque 4.6.1\par
\emph{Reported reading}: sous-objets quasi-compacts U\_\ensuremath{\alpha} ... des sous-topos fermés E''\_i complémentaire des E\_\{/U\_i\}\par
\emph{Proposal}: \$U\_i\$ (or \$E''\_\textbackslash{}alpha\$, \$E\_\{/U\_\textbackslash{}alpha\}\$)\par
The record offers two alternative index repairs and hence does not determine one exact corrected reading.\par
\par\noindent\textbf{OM-OPUS-0128 -- t2 VI p.252 -- OPUS}\par
c), first sentence 'pour tout objet X'' de F''\par
\emph{Reported reading}: F"\par
\emph{Proposal}: \$E''\$\par
The target topos must be matched with the preceding morphism before changing F double-prime to E double-prime.\par
\par\noindent\textbf{OM-OPUS-0129 -- t2 VI p.252 -- OPUS}\par
'un morphisme d'un (X'\_i,u\_i)'\par
\emph{Reported reading}: (X'\_i,u\_i)\par
\emph{Proposal}: \$(X'\_\{\textbackslash{}alpha\}, u\_\{\textbackslash{}alpha\})\$\par
The change of family index requires tracking the corresponding transition maps.\par
\par\noindent\textbf{OM-OPUS-0136 -- t2 VI p.297 -- OPUS}\par
7.4.13.6, 'dans le site total D\_\{/(e\_i)\}'\par
\emph{Reported reading}: D\_\{/ (e\_i)\}\par
\emph{Proposal}: D\_\{/e\_i\}\par
Parentheses around the slice index may group an indexed family; the supplied snippet does not uniquely establish a misprint.\par
\par\noindent\textbf{OM-OPUS-0138 -- t2 VI p.299 -- OPUS}\par
diagram (8.1.2.1)\par
\emph{Reported reading}: vertical m \ensuremath{\times} id\_I; diagonal m\_1\par
\emph{Proposal}: vertical n\textbackslash{}times \textbackslash{}id\_I, diagonal m\_0\par
The labels require matching all endpoints with the definition of gamma; the record alone does not prove the diagram.\par
\par\noindent\textbf{OM-OPUS-0144 -- t2 VII p.346 -- OPUS}\par
1.6, i\_\{X/Y\} =\allowbreak  X\_ét \ensuremath{\to} Y\_ét\par
\emph{Reported reading}: i\_\{X/Y\} =\allowbreak  X\_ét \ensuremath{\to} Y\_ét\par
\emph{Proposal}: \$i\_\{X/Y\} : X\_\{\textbackslash{}text\{ét\}\} \textbackslash{}to Y\_\{\textbackslash{}text\{ét\}\}\$\par
Replacing an equals sign by a colon requires the local convention for defining a named arrow.\par
\par\noindent\textbf{OM-OPUS-0147 -- t2 VIII p.392 -- OPUS}\par
proof of 6.3, second half\par
\emph{Reported reading}: G \ensuremath{\to} i\_*i*G ... l'hypothèse faite sur G ... i*(G) \ensuremath{\to} i*(i\_* i*(G))\par
\emph{Proposal}: F \textbackslash{}to i\_\{\textbackslash{}ast\} i\textasciicircum{}\{\textbackslash{}ast\} F\par
The proposed renaming runs through an entire proof; every bound occurrence and antecedent must be tracked.\par
\par\noindent\textbf{OM-OPUS-0148 -- t2 VIII p.402 -- OPUS}\par
Prop. 8.1, spectral sequence\par
\emph{Reported reading}: H*(X,F) ⇐ E\_2\textasciicircum{}\{pq\} =\allowbreak  H\textasciicircum{}p(X/Y, \ensuremath{𝓗}\textasciicircum{}q(F))\par
\emph{Proposal}: H\textasciicircum{}\{\textbackslash{}ast\} (Y, F) \textbackslash{}Longleftarrow\par
The abutment requires checking the site and the sheaf in the spectral-sequence construction.\par
\par\noindent\textbf{OM-OPUS-0149 -- t3 IX p.1 -- OPUS}\par
sur le corps \$K\$ alg. clos ... variété abélienne sur \$k\$\par
\emph{Reported reading}: corps K ... sur k\par
\emph{Proposal}: le corps \$k\$ alg. clos\par
The two field letters may encode an implicit identification; the record does not determine which occurrence to replace.\par
\par\noindent\textbf{OM-OPUS-0150 -- t3 IX p.3 -- OPUS}\par
Prop. 1.2 (v)\par
\emph{Reported reading}: si F est un faisceau de p-torsion sur T, alors les R\textasciicircum{}qu\_*F' sont des faisceaux de p-torsion sur T\par
\emph{Proposal}: les \$R\textasciicircum{}qu\_\{\textbackslash{}ast\}F\$ sont des faisceaux de \$p\$-torsion sur \$T'\$\par
Both the sheaf and its base change; the direct-image morphism and its source/target must be checked together.\par
\par\noindent\textbf{OM-OPUS-0151 -- t3 IX p.4 -- OPUS}\par
(iii) « La deuxième assertion est une conséquence de »\par
\emph{Reported reading}: sentence breaks off after « de »\par
\emph{Proposal}: La deuxième assertion est une conséquence de \textbackslash{}ldots\par
The sentence ends at de and supplies no unique missing reference or completion.\par
\par\noindent\textbf{OM-OPUS-0155 -- t3 IX p.16 -- OPUS}\par
la réunion des \$T\_i\$ est \$X\$ puisque \$f\$ est surjectif, d'où s'ensuit que \$X\$ est réunion d'un nombre fini des \$T\_i\$\par
\emph{Reported reading}: la réunion des T\_i est X puisque f est surjectif, d'où s'ensuit que X est réunion d'un nombre fini des T\_i\par
\emph{Proposal}: la réunion des \$T\_i\$ est \$Y\$ puisque \$f\$ est surjectif, d'où s'ensuit que \$Y\$ est réunion d'un nombre fini des \$T\_i\$\par
The two replacements depend on the ambient sets of T\_i and the surjectivity argument, which need explicit checking.\par
\par\noindent\textbf{OM-OPUS-0159 -- t3 X p.48 -- OPUS}\par
morphisme cobord, target\par
\emph{Reported reading}: H\textasciicircum{}\{n+1\}(Ḡ, \ensuremath{\mu}\_\ensuremath{\ell})\par
\emph{Proposal}: H\textasciicircum{}\{n+1\}(G, \ensuremath{𝛍}\_\{\textbackslash{}ell\})\par
This finding intersects Orgogozo's \textbackslash{}sisi material at morphisme cobord, target; preserve that human entry unchanged. The proposed reading H\textasciicircum{}\{n+1\}(G, \ensuremath{𝛍}\_\{\textbackslash{}ell\}) at morphisme cobord, target changes a substantive mathematical assertion or bibliographic identification; no independent derivation of this exact replacement is established in this review.\par
\par\noindent\textbf{OM-OPUS-0160 -- t3 X p.49 -- OPUS}\par
diagramme, right entry\par
\emph{Reported reading}: H\textasciicircum{}\{n+1\}(Ḡ, \ensuremath{\mu}\_\ensuremath{\ell})\par
\emph{Proposal}: H\textasciicircum{}\{n+1\}(G,\ensuremath{𝛍}\_\{\textbackslash{}ell\})\par
This finding intersects Orgogozo's \textbackslash{}sisi material at diagramme, right entry; preserve that human entry unchanged. The proposed reading H\textasciicircum{}\{n+1\}(G,\ensuremath{𝛍}\_\{\textbackslash{}ell\}) at diagramme, right entry changes a substantive mathematical assertion or bibliographic identification; no independent derivation of this exact replacement is established in this review.\par
\par\noindent\textbf{OM-OPUS-0161 -- t3 X p.57 -- OPUS}\par
Corollaire 4.4, first line\par
\emph{Reported reading}: Soit X noéthérien\par
\emph{Proposal}: noethérien\par
The supplied record establishes a spelling preference but gives no textual test of the author's intended accent.\par
\par\noindent\textbf{OM-OPUS-0162 -- t3 X p.62 -- OPUS}\par
'au-dessus de 0[i]'\par
\emph{Reported reading}: 0[i]\par
\emph{Proposal}: \$\textbackslash{}QQ[i]\$\par
The record offers 0 or O as the printed glyph; that printed reading cannot be fixed without the scan.\par
\par\noindent\textbf{OM-OPUS-0163 -- t3 X p.63 -- OPUS}\par
bibliography [5]\par
\emph{Reported reading}: cohomologique de groupes profinis\par
\emph{Proposal}: des\par
This finding intersects Orgogozo's \textbackslash{}sisi material at bibliography [5]; preserve that human entry unchanged. The proposed reading des at bibliography [5] changes a substantive mathematical assertion or bibliographic identification; no independent derivation of this exact replacement is established in this review.\par
\par\noindent\textbf{OM-OPUS-0164 -- t3 X p.63 -- OPUS}\par
bibliography [6]\par
\emph{Reported reading}: Théorie des Faisceaux, Paris (1958).\par
\emph{Proposal}: Topologie algébrique et théorie des faisceaux\par
This finding intersects Orgogozo's \textbackslash{}sisi material at bibliography [6]; preserve that human entry unchanged. The proposed reading Topologie algébrique et théorie des faisceaux at bibliography [6] changes a substantive mathematical assertion or bibliographic identification; no independent derivation of this exact replacement is established in this review.\par
\par\noindent\textbf{OM-OPUS-0170 -- t3 XII p.92 -- OPUS}\par
(puisque H\textasciicircum{}1(\ensuremath{\bar{X}},Ḡ) =\allowbreak  lim H\textasciicircum{}1(X',C) =\allowbreak  0 ...)\par
\emph{Reported reading}: H\textasciicircum{}1(\textbackslash{}bar X,\textbackslash{}bar G)\par
\emph{Proposal}: H\textasciicircum{}1(\textbackslash{}bar\{X\},\textbackslash{}bar\{C\})\par
The cohomology coefficient must be traced through the chosen Godement resolution.\par
\par\noindent\textbf{OM-OPUS-0174 -- t3 XII p.100 -- OPUS}\par
diagramme de la démonstration de 6.2, ligne du haut\par
\emph{Reported reading}: T\textasciicircum{}n(C) \ensuremath{\to} R\textasciicircum{}\{n+1\}(A) \ensuremath{\to} T\textasciicircum{}\{n+1\}(B)\par
\emph{Proposal}: T\textasciicircum{}\{n+1\}(A)\par
The connecting morphism must be checked against the named delta-functor before changing R to T.\par
\par\noindent\textbf{OM-OPUS-0176 -- t3 XII p.116 -- OPUS}\par
lines 4 and display (2): H*(R\textasciicircum{}1(p\_*(\ensuremath{\bar{F}})) \ensuremath{\to} ... and H*(Rp\_*(\ensuremath{\bar{F}})) \ensuremath{\to} Rp\_\{Y*\}(\ensuremath{\bar{h}}*(\ensuremath{\bar{F}}))\par
\emph{Reported reading}: H*(R\textasciicircum{}1(p\_*(\ensuremath{\bar{F}})) and H*(Rp\_*(\ensuremath{\bar{F}})), with a capital H\par
\emph{Proposal}: h\textasciicircum{}\{\textbackslash{}ast\}(R\textasciicircum{}1p\_\{\textbackslash{}ast\}(\textbackslash{}bar\{F\})) and h\textasciicircum{}\{\textbackslash{}ast\}(Rp\_\{\textbackslash{}ast\}(\textbackslash{}bar\{F\}))\par
The proposed lower-case h requires the exact base-change square and functor types in both occurrences.\par
\par\noindent\textbf{OM-OPUS-0177 -- t3 XII p.123 -- OPUS}\par
Remarques 6.13\par
\emph{Reported reading}: où X est le spectre d'un anneau noethérien, séparé et complet ... où Y =\allowbreak  Spec (A/J)\par
\emph{Proposal}: le spectre d'un anneau noethérien \$A\$, séparé\par
Naming the ring A is plausible, but the record combines that insertion with an unclosed parenthesis without a unique complete repair.\par
\par\noindent\textbf{OM-OPUS-0180 -- t3 XII p.131 -- OPUS}\par
diagram, middle column, both rows: H\textasciicircum{}2(X, \_n) and H\textasciicircum{}2(X\_0, \_n)\par
\emph{Reported reading}: H\textasciicircum{}2(X, \_n) and H\textasciicircum{}2(X\_o, \_n): there is a blank where the handwritten \ensuremath{\mu} should stand. It is not filled in, unlike the handwritten \ensuremath{\mu} in the text above.\par
\emph{Proposal}: H\textasciicircum{}\{2\}(X, \ensuremath{𝛍}\_n) ... H\textasciicircum{}2(X\_0, \ensuremath{𝛍}\_n)\par
The record explicitly cannot distinguish an omitted mu from a faint lost stroke; retain the printed blank.\par
\par\noindent\textbf{OM-OPUS-0181 -- t3 XIII p.137 -- OPUS}\par
first line, «exprimer la condition (**) ainsi»\par
\emph{Reported reading}: exprimer la condition (**) ainsi\par
\emph{Proposal}: \$(\textbackslash{}ast)\$\par
The star count requires comparison with the two numbered conditions and their stated equivalence.\par
\par\noindent\textbf{OM-OPUS-0183 -- t3 XIV p.154 -- OPUS}\par
« qui induise Y \ensuremath{\to} X\_s \ensuremath{\otimes}\_K K' sur les fibres »\par
\emph{Reported reading}: Y \ensuremath{\longrightarrow} X\_S \ensuremath{\otimes}\_K K' (subscript looks like a capital S)\par
\emph{Proposal}: Y\textbackslash{}rightarrow X\_s \textbackslash{}otimes\_K K'\par
The record explicitly describes the printed case as unclear; the scan must resolve it.\par
\par\noindent\textbf{OM-OPUS-0184 -- t3 XIV p.156 -- OPUS}\par
fibre of R\textasciicircum{}2f\_*(Z/\ensuremath{\ell})\par
\emph{Reported reading}: au-dessus de s S est isomorphe\par
\emph{Proposal}: \$s \textbackslash{}in S\$\par
The record explicitly allows a lost handwritten membership sign; the printed reading remains unverified.\par
\par\noindent\textbf{OM-OPUS-0185 -- t3 XV p.172 -- OPUS}\par
proof of 1.5, pour tout y Y\par
\emph{Reported reading}: pour tout y Y (blank where \ensuremath{\in} belongs)\par
\emph{Proposal}: \$y \textbackslash{}in Y\$\par
The supposed blank may be a lost handwritten membership sign; the printed reading remains unverified.\par
\par\noindent\textbf{OM-OPUS-0186 -- t3 XV p.172 -- OPUS}\par
Proposition 1.6, n N\par
\emph{Reported reading}: et n N un entier naturel (blank where \ensuremath{\in} belongs)\par
\emph{Proposal}: \$n \textbackslash{}in \textbackslash{}NN\$\par
The supposed missing membership sign needs a scan check before a diplomatic variant can be specified.\par
\par\noindent\textbf{OM-OPUS-0188 -- t3 XVII p.292 -- OPUS}\par
Cas triangulé, second triangle\par
\emph{Reported reading}: et soit (f*Y', f*Y).\par
\emph{Proposal}: \$(f\textasciicircum{}\{\textbackslash{}ast\}Y',f\textasciicircum{}\{\textbackslash{}ast\}Y, X'')\$ un triangle distingué dans \$\textbackslash{}Aa\_x\$\par
This finding intersects Orgogozo's \textbackslash{}sisi material at Cas triangulé, second triangle; preserve that human entry unchanged. The record documents a lacuna at Cas triangulé, second triangle; it does not uniquely establish the proposed completion \$(f\textasciicircum{}\{\textbackslash{}ast\}Y',f\textasciicircum{}\{\textbackslash{}ast\}Y, X'')\$ un triangle distingué dans \$\textbackslash{}Aa\_x\$.\par
\par\noindent\textbf{OM-OPUS-0189 -- t3 XVII p.299 -- OPUS}\par
Proposition 2.4.1 (*), x \ensuremath{\to} y \ensuremath{\to} Z\par
\emph{Reported reading}: y \textbackslash{}xrightarrow\{g\} Z\par
\emph{Proposal}: z\par
This finding intersects Orgogozo's \textbackslash{}sisi material at Proposition 2.4.1 (*), x \ensuremath{\to} y \ensuremath{\to} Z; preserve that human entry unchanged. The replacement z at Proposition 2.4.1 (*), x \ensuremath{\to} y \ensuremath{\to} Z identifies a particular indexed object or symbol; that identification is not independently established by the quoted local expression.\par
\par\noindent\textbf{OM-OPUS-0190 -- t3 XVII p.300 -- OPUS}\par
« trouver B\textasciicircum{}1, v et s \ensuremath{\in} S\textasciicircum{}1 » and diagram (2.4.1.4)\par
\emph{Reported reading}: B\textasciicircum{}1, v et s \ensuremath{\in} S\textasciicircum{}1; B\_1, A\_1, F\_*A\_1 in the diagram\par
\emph{Proposal}: B', s\textbackslash{}in S\textasciicircum{}1 et u; A\textasciicircum{}1, F\_\textbackslash{}ast A\textasciicircum{}1, B\textasciicircum{}1\par
This finding intersects Orgogozo's \textbackslash{}sisi material at « trouver B\textasciicircum{}1, v et s \ensuremath{\in} S\textasciicircum{}1 » and diagram (2.4.1.4); preserve that human entry unchanged. The replacement B', s\textbackslash{}in S\textasciicircum{}1 et u; A\textasciicircum{}1, F\_\textbackslash{}ast A\textasciicircum{}1, B\textasciicircum{}1 at « trouver B\textasciicircum{}1, v et s \ensuremath{\in} S\textasciicircum{}1 » and diagram (2.4.1.4) identifies a particular indexed object or symbol; that identification is not independently established by the quoted local expression.\par
\par\noindent\textbf{OM-OPUS-0191 -- t3 XVII p.303 -- OPUS}\par
'(cf. 5.3 dans le cas des schémas'\par
\emph{Reported reading}: 5.3\par
\emph{Proposal}: 6.1\par
This finding intersects Orgogozo's \textbackslash{}sisi material at '(cf. 5.3 dans le cas des schémas'; preserve that human entry unchanged. The proposed reference 6.1 at '(cf. 5.3 dans le cas des schémas' needs comparison with the statement actually invoked; that target has not been independently established here.\par
\par\noindent\textbf{OM-OPUS-0196 -- t3 XVII p.371 -- OPUS}\par
introduction to 5.5, symmetric power (X\_*)\textasciicircum{}n/…\par
\emph{Reported reading}: (X\_*)\textasciicircum{}n/G\_n\par
\emph{Proposal}: (X\_\{\textbackslash{}ast\})\textasciicircum{}\{n\}/\textbackslash{}Sf\_\{n\}\par
The symmetric-group letter may be a font/handwriting issue; the scan must establish the exact printed glyph.\par
\par\noindent\textbf{OM-OPUS-0201 -- t3 XVII p.383 -- OPUS}\par
5.5.7.1, first two lines: «Un point x de Sym\_S(X) … point géométrique de Sym\_S(X)»\par
\emph{Reported reading}: Un point x de Sym\_S(X) … point géométrique de Sym\_S(X)\par
\emph{Proposal}: \$\textbackslash{}Sym\textasciicircum{}\{k\}\_\{S\}(X)\$\par
The degree k must be checked against all symmetric powers in the paragraph before inserting two exponents.\par
\par\noindent\textbf{OM-OPUS-0203 -- t3 XVII p.401 -- OPUS}\par
5.5.22.5, «telle que x soit dans l'image»\par
\emph{Reported reading}: telle que x soit dans l'image de Sym\textasciicircum{}n\_k(Y')\par
\emph{Proposal}: \$y\$\par
This finding intersects Orgogozo's \textbackslash{}sisi material at 5.5.22.5, «telle que x soit dans l'image»; preserve that human entry unchanged. The proposed reading \$y\$ at 5.5.22.5, «telle que x soit dans l'image» changes a substantive mathematical assertion or bibliographic identification; no independent derivation of this exact replacement is established in this review.\par
\par\noindent\textbf{OM-OPUS-0206 -- t3 XVII p.404 -- OPUS}\par
5.5.25.4, «Soit donc F un faisceau»\par
\emph{Reported reading}: Soit donc F un faisceau de A-vectoriels\par
\emph{Proposal}: \$K\$\par
This finding intersects Orgogozo's \textbackslash{}sisi material at 5.5.25.4, «Soit donc F un faisceau»; preserve that human entry unchanged. The proposed reading \$K\$ at 5.5.25.4, «Soit donc F un faisceau» changes a substantive mathematical assertion or bibliographic identification; no independent derivation of this exact replacement is established in this review.\par
\par\noindent\textbf{OM-OPUS-0207 -- t3 XVII p.404 -- OPUS}\par
5.5.25.5, «le cas où F est de la forme j\_!A\_\{X'\}»\par
\emph{Reported reading}: le cas où F est de la forme j\_! A\_\{X'\}\par
\emph{Proposal}: \$K\$\par
The proposed replacement requires tracking the object being reduced throughout the dévissage.\par
\par\noindent\textbf{OM-OPUS-0208 -- t3 XVII p.405 -- OPUS}\par
«Un ultime dévissage … en les faisceaux Z/\ensuremath{\ell}\_X»\par
\emph{Reported reading}: en les faisceaux Z/\ensuremath{\ell}\_X , pour X' composante connexe\par
\emph{Proposal}: \$(\textbackslash{}ZZ/\textbackslash{}ell)\_\{X'\}\$\par
This finding intersects Orgogozo's \textbackslash{}sisi material at «Un ultime dévissage … en les faisceaux Z/\ensuremath{\ell}\_X»; preserve that human entry unchanged. The replacement \$(\textbackslash{}ZZ/\textbackslash{}ell)\_\{X'\}\$ at «Un ultime dévissage … en les faisceaux Z/\ensuremath{\ell}\_X» identifies a particular indexed object or symbol; that identification is not independently established by the quoted local expression.\par
\par\noindent\textbf{OM-OPUS-0210 -- t3 XVII p.421 -- OPUS}\par
Prop. 6.2.2 (ii)\par
\emph{Reported reading}: définissent un isomorphismes ... \ensuremath{\to} f\_!f\textasciicircum{}*\ensuremath{𝔉} .\par
\emph{Proposal}: un isomorphisme ... f\_\{!\}f\textasciicircum{}\{\textbackslash{}ast\}\textbackslash{}Ff|U\par
The singular correction is grammatical, but this record also inserts the restriction |U in a morphism target; that second assertion is not independently derived here. Both changes already occur in Orgogozo's human note, which is preserved unchanged.\par
\par\noindent\textbf{OM-OPUS-0213 -- t3 XVII p.457 -- OPUS}\par
6.3.25.4, displayed isomorphism\par
\emph{Reported reading}: d*K\_n \ensuremath{\simeq} pr\_1\textasciicircum{}* K\_\{n\_1\} + pr\_2\textasciicircum{}* K\_\{n\_2\}\par
\emph{Proposal}: \textbackslash{}alpha\textasciicircum{}\{\textbackslash{}ast\}K\_\{n\}\par
The d/alpha reading concerns a handwritten symbol whose exact printed identity needs verification.\par
\par\noindent\textbf{OM-OPUS-0214 -- t3 XVII p.459 -- OPUS}\par
Proposition 6.3.29\par
\emph{Reported reading}: et c : A \ensuremath{\to} B* le morphisme structural ... \ensuremath{\varepsilon}* soit fidèle\par
\emph{Proposal}: \$\textbackslash{}epsilon:A\textbackslash{}rightarrow B\textasciicircum{}\{\textbackslash{}ast\}\$\par
The record explicitly calls the replacement probable; the structural morphism and pullback need an identification.\par
\par\noindent\textbf{OM-OPUS-0215 -- t3 XVII p.473 -- OPUS}\par
proof of 7.1.6, first sentence\par
\emph{Reported reading}: R\textasciicircum{}+fg\_! --c\_fg--> Rf\_! o Rg\_!\par
\emph{Proposal}: R\textasciicircum{}\{+\}f\_\{!\}\textbackslash{}circ R\textasciicircum{}\{+\}g\_\{!\}\par
The plus signs specify derived-functor domains; the source and target categories must be checked.\par
\par\noindent\textbf{OM-OPUS-0216 -- t3 XVII p.476 -- OPUS}\par
Remarque 7.2.3\par
\emph{Reported reading}: (cf.(3.7))\par
\emph{Proposal}: \textbackslash{}cf 7.3.7\par
This finding intersects Orgogozo's \textbackslash{}sisi material at Remarque 7.2.3; preserve that human entry unchanged. The proposed reading \textbackslash{}cf 7.3.7 at Remarque 7.2.3 changes a substantive mathematical assertion or bibliographic identification; no independent derivation of this exact replacement is established in this review.\par
\par\noindent\textbf{OM-OPUS-0217 -- t3 XVII p.478 -- OPUS}\par
(7.3.1.2)\par
\emph{Reported reading}: X\textasciicircum{}1(F) \ensuremath{\to} ... X\textasciicircum{}\{2d-1\}(F)\par
\emph{Proposal}: \textbackslash{}ldots\textbackslash{}to\textbackslash{}Xx\textasciicircum{}\{2d-1\}(F)\par
This finding intersects Orgogozo's \textbackslash{}sisi material at (7.3.1.2); preserve that human entry unchanged. The replacement \textbackslash{}ldots\textbackslash{}to\textbackslash{}Xx\textasciicircum{}\{2d-1\}(F) at (7.3.1.2) requires a complete check of the maps, their endpoints and the composition convention; this identification is not independently verified.\par
\par\noindent\textbf{OM-OPUS-0218 -- t3 XVII p.479 -- OPUS}\par
7.3.5\par
\emph{Reported reading}: grâce à 7.2.3\par
\emph{Proposal}: 7.3.3\par
This finding intersects Orgogozo's \textbackslash{}sisi material at 7.3.5; preserve that human entry unchanged. The proposed reference 7.3.3 at 7.3.5 needs comparison with the statement actually invoked; that target has not been independently established here.\par
\par\noindent\textbf{OM-OPUS-0219 -- t3 XVII p.479 -- OPUS}\par
7.3.5\par
\emph{Reported reading}: En vertu de 7.2.4\par
\emph{Proposal}: 7.3.4\par
This finding intersects Orgogozo's \textbackslash{}sisi material at 7.3.5; preserve that human entry unchanged. The proposed reference 7.3.4 at 7.3.5 needs comparison with the statement actually invoked; that target has not been independently established here.\par
\par\noindent\textbf{OM-OPUS-0220 -- t3 XVII p.480 -- OPUS}\par
Dold-Puppe entry\par
\emph{Reported reading}: nicht additiver\par
\emph{Proposal}: nicht-additiver\par
This finding intersects Orgogozo's \textbackslash{}sisi material at Dold-Puppe entry; preserve that human entry unchanged. The proposed reading nicht-additiver at Dold-Puppe entry changes a substantive mathematical assertion or bibliographic identification; no independent derivation of this exact replacement is established in this review.\par
\par\noindent\textbf{OM-OPUS-0222 -- t3 XVIII p.489 -- OPUS}\par
opening diagram and the following sentence\par
\emph{Reported reading}: Pic(X)/n, H\textasciicircum{}2\_c(X,...), Tr\_u, Pic(Y)/n, H\textasciicircum{}2\_c(Y,...); « sur X »\par
\emph{Proposal}: \textbackslash{}bar\{X\}, \textbackslash{}bar\{Y\}, \textbackslash{}Tr\_\{\textbackslash{}bar\{u\}\}\par
The record combines several bar insertions; the compactified schemes and their maps must be tracked separately.\par
\par\noindent\textbf{OM-OPUS-0223 -- t3 XVIII p.490 -- OPUS}\par
first paragraph\par
\emph{Reported reading}: dans P\textasciicircum{}1\_s\par
\emph{Proposal}: \textbackslash{}PP\textasciicircum{}\{1\}\_S\par
The printed index only looks lowercase according to the record; the scan must resolve it.\par
\par\noindent\textbf{OM-OPUS-0225 -- t3 XVIII p.495 -- OPUS}\par
Hom\_s(W\_n,G) and Hom\_s(W\_n,G\_n)\par
\emph{Reported reading}: Hom\_s (lower-case s)\par
\emph{Proposal}: \textbackslash{}Hhom\_\{S\}\par
The case of the base subscript needs verification against the printed symbols and the base-scheme convention.\par
\par\noindent\textbf{OM-OPUS-0226 -- t3 XVIII p.500 -- OPUS}\par
after Variantes 1.2.15, the two proof paragraphs\par
\emph{Reported reading}: Dans le cas (i), les flèches (1.2.2.1) sont des isomorphismes d'après le théorème de changement de base ... Dans le cas (ii), que les flèches (1.2.2.1) soient des isomorphismes est standard ... Zariski.\par
\emph{Proposal}: Dans le cas (i), que les flèches (1.2.2.1) soient des isomorphismes est standard, compte tenu de ce que la cohomologie fppf de \$G\$ coïncide avec sa cohomologie pour la topologie de Zariski. Dans le cas (ii), les flèches (1.2.2.1) sont des isomorphismes d'après le théorème de changement de base pour un morphisme propre et la simple connexité de \$\textbackslash{}PP\textasciicircum{}\{r\}\_\{k\}\$ pour \$k\$ algébriquement clos.\par
Swapping whole proof paragraphs requires proving which case each argument establishes; no paraphrased replacement is accepted.\par
\par\noindent\textbf{OM-OPUS-0227 -- t3 XVIII p.503 -- OPUS}\par
display after « se réécrit », second line\par
\emph{Reported reading}: p\_12\textasciicircum{}*\ensuremath{\langle}\ensuremath{\mathcal{L}}1,K] + … (no relation sign at the start of the line)\par
\emph{Proposal}: \textbackslash{}simeq p\textasciicircum{}\{\textbackslash{}ast\}\_\{12\}\textbackslash{}langle \textbackslash{}Ll\_\{1\},K]\par
The record proposes an isomorphism sign while reporting that TeX inserts equality; the exact relation is unresolved.\par
\par\noindent\textbf{OM-OPUS-0231 -- t3 XVIII p.517 -- OPUS}\par
«Pour D et E disjoints, on a …»\par
\emph{Reported reading}: O\_D \ensuremath{\otimes}\textasciicircum{}L O\_E =\allowbreak  O (script O)\par
\emph{Proposal}: \textbackslash{}Oo\_\{D\}\textbackslash{}mathop\{\textbackslash{}otimes\}\textasciicircum{}\{\textbackslash{}LL\}\textbackslash{}Oo\_\{E\}=\allowbreak 0\par
Disjoint supports motivate a zero object, but the reported script O/zero reading still requires the exact printed glyph.\par
\par\noindent\textbf{OM-OPUS-0232 -- t3 XVIII p.519 -- OPUS}\par
a\_\{g,h\} and the middle term of the top row of 1.4.1.1\par
\emph{Reported reading}: \ensuremath{\underline{F}}(F(g),\ensuremath{\underline{F}}(h)); F(F(f),\ensuremath{\underline{F}}(gh))\par
\emph{Proposal}: F(\ensuremath{\mathcal{F}}(g),\ensuremath{\mathcal{F}}(h)), F(\ensuremath{\mathcal{F}}(f),\ensuremath{\mathcal{F}}(gh))\par
The functor F and the object assignment script F must be identified throughout the coherence formula.\par
\par\noindent\textbf{OM-OPUS-0233 -- t3 XVIII p.523 -- OPUS}\par
1.4.8 a), last words of the page\par
\emph{Reported reading}: vérifie les compatibilités de 1.4.3.\par
\emph{Proposal}: \textbackslash{}hyperref[XVIII.subsubsec1.4.6]\{1.4.6\}\par
The target cross-reference requires comparison with the additive-functor compatibility definitions.\par
\par\noindent\textbf{OM-OPUS-0236 -- t3 XVIII p.535 -- OPUS}\par
diagram (1.5.3.6), bottom left\par
\emph{Reported reading}: \ensuremath{\tau}\_\ensuremath{\leq}1 RHom(\ensuremath{\tau}\_<0 Rf\_2* G\_m[1], G)\par
\emph{Proposal}: \textbackslash{}tau\_\{\textbackslash{}leq 0\}\par
The truncation bound changes the complex; the degrees and the displayed comparison require independent verification.\par
\par\noindent\textbf{OM-OPUS-0237 -- t3 XVIII p.539 -- OPUS}\par
Si E est un diviseur relatif …\par
\emph{Reported reading}: m'\_1 =\allowbreak  m\_i \ensuremath{\otimes} 1\par
\emph{Proposal}: m'\_\{i\}=\allowbreak m\_\{i\}\textbackslash{}otimes 1\par
The index i/1 reading requires the chosen family of local equations and the printed symbol.\par
\par\noindent\textbf{OM-OPUS-0239 -- t3 XVIII p.539 -- OPUS}\par
(II) L|f\textasciicircum{}\{-1\}(U)\par
\emph{Reported reading}: L|f\textasciicircum{}\{-1\}(U)\par
\emph{Proposal}: \textbackslash{}bar\{f\}\par
The preimage map depends on the compactification; the two morphisms must be identified.\par
\par\noindent\textbf{OM-OPUS-0240 -- t3 XVIII p.540 -- OPUS}\par
1.5.13 and 1.5.13.1\par
\emph{Reported reading}: f\_* G\_m(m)\par
\emph{Proposal}: \textbackslash{}bar\{f\}\_\{\textbackslash{}ast\}\par
The direct-image functor depends on the compactification and must be checked against the defined domain.\par
\par\noindent\textbf{OM-OPUS-0241 -- t3 XVIII p.541 -- OPUS}\par
throughout the page (f\_* G\_m(m), R\textasciicircum{}1 f\_* G\_m(m), e : f\_* ch(G[1]) …)\par
\emph{Reported reading}: f\_* G\_m (m), R\textasciicircum{}1 f\_* G\_m (m)\par
\emph{Proposal}: \textbackslash{}bar\{f\}\_\{\textbackslash{}ast\}\par
The record proposes multiple bar insertions across the page, requiring a tracked occurrence list.\par
\par\noindent\textbf{OM-OPUS-0242 -- t3 XVIII p.542 -- OPUS}\par
diagrams (1.5.15.1) and (1.5.15.2)\par
\emph{Reported reading}: Ext\textasciicircum{}1(R\textasciicircum{}1f\_\{1*\} G\_m (m), G); Ext\textasciicircum{}1(u\textasciicircum{}*, Id); Ext\textasciicircum{}1(N\_\{X\_1/X\_2\}, Id)\par
\emph{Proposal}: \textbackslash{}Ext\textasciicircum{}\{1\}(R\textasciicircum{}1\textbackslash{}bar\{f\}\_\{1\textbackslash{}ast\}\textbackslash{}GG\_m(\textbackslash{}mathfrak\{m\}\_1),G),\textbackslash{} \textbackslash{}Ext\textasciicircum{}\{1\}(\textbackslash{}bar\{u\}\textasciicircum{}\{\textbackslash{}ast\},\textbackslash{}Id),\textbackslash{} \textbackslash{}Ext\textasciicircum{}\{1\}(N\_\{\textbackslash{}bar\{X\}\_1/\textbackslash{}bar\{X\}\_2\},\textbackslash{}Id)\par
The bundled bar and modulus-index insertions require a separate type check for every diagram entry.\par
\par\noindent\textbf{OM-OPUS-0243 -- t3 XVIII p.546 -- OPUS}\par
Les faisceaux étales R\textasciicircum{}1f\_1(Z/n), R\textasciicircum{}1f'\_1(Z/n)\par
\emph{Reported reading}: R\textasciicircum{}1f\_1(\textbackslash{}ZZ/n), R\textasciicircum{}1f'\_1(\textbackslash{}ZZ/n) (no star)\par
\emph{Proposal}: R\textasciicircum{}\{1\}f\_\{1\textbackslash{}ast\}(\textbackslash{}ZZ/n), R\textasciicircum{}\{1\}f'\_\{1\textbackslash{}ast\}(\textbackslash{}ZZ/n)\par
The omitted stars require checking which derived direct-image functors the sentence defines.\par
\par\noindent\textbf{OM-OPUS-0245 -- t3 XVIII p.558 -- OPUS}\par
after (2.9.2)\par
\emph{Reported reading}: le diagramme de morphismes trace (2.8.1) et XVII 6.2. )\par
\emph{Proposal}: le diagramme de morphismes trace (cf. (2.8.1) et XVII 6.2)\par
The printed closing parenthesis lacks a matching opener, but the proposed correction also adds cf.; the exact completion is an editorial reconstruction rather than a recovered reading.\par
\par\noindent\textbf{OM-OPUS-0249 -- t3 XIX p.600 -- OPUS}\par
Pour démontrer que c' est surjectif\par
\emph{Reported reading}: c' est surjectif\par
\emph{Proposal}: \$c\$\par
The prime distinguishes maps in the argument; the exact surjectivity claim must identify the map.\par
\par\noindent\textbf{OM-OPUS-0250 -- t3 XIX p.602 -- OPUS}\par
première ligne\par
\emph{Reported reading}: la fibre fermée de A[1/x\_r] est non-vide\par
\emph{Proposal}: \$A\_\{\textbackslash{}alpha\}[1/x\_\{r\}]\$\par
The record describes A\_alpha as presumptive, so the omitted index is not determined.\par
\par\noindent\textbf{OM-OPUS-0251 -- t3 XIX p.609 -- OPUS}\par
(3.3.2)\par
\emph{Reported reading}: (R\textasciicircum{}* i\textasciicircum{}!)\ensuremath{\mu}\_n\par
\emph{Proposal}: (R\textasciicircum{}\{2\}i\textasciicircum{}\{!\})\ensuremath{𝛍}\_\{n\}\par
Replacing the cohomological star by 2 requires verifying the purity degree and the meaning of the displayed graded notation.\par
\par\noindent\textbf{OM-OPUS-0258 -- t3 XIX p.638 -- OPUS}\par
bibliography [1], year of the Annals paper\par
\emph{Reported reading}: vol. 63 (1963)\par
\emph{Proposal}: vol. 63 (1956)\par
The year requires an external bibliographic authority; the source record alone does not establish it.\par
\par\noindent\textbf{OM-SOL-0001 -- t1 I p.2 -- SOL}\par
Troisième propriété, définition du couple de Kuratowski\par
\emph{Reported reading}: \$(x,y)=\allowbreak \textbackslash{}\{\textbackslash{}\{x,y\textbackslash{}\},x\textbackslash{}\}\$\par
\emph{Proposal}: \$(x,y)=\allowbreak \textbackslash{}\{\textbackslash{}\{x,y\textbackslash{}\},\textbackslash{}\{x\textbackslash{}\}\textbackslash{}\}\$\par
The proposed reading \$(x,y)=\allowbreak \textbackslash{}\{\textbackslash{}\{x,y\textbackslash{}\},\textbackslash{}\{x\textbackslash{}\}\textbackslash{}\}\$ at Troisième propriété, définition du couple de Kuratowski changes a substantive mathematical assertion or bibliographic identification; no independent derivation of this exact replacement is established in this review.\par
\par\noindent\textbf{OM-SOL-0004 -- t1 I p.4 -- SOL}\par
Deuxième formule isolée\par
\emph{Reported reading}: \$U\_c\textbackslash{}in U'\_\{c'\}\$\par
\emph{Proposal}: \$\ensuremath{𝒰}\_c\textbackslash{}in \ensuremath{𝒰}\_\{c'\}\$\par
The replacement \$\ensuremath{𝒰}\_c\textbackslash{}in \ensuremath{𝒰}\_\{c'\}\$ at Deuxième formule isolée identifies a particular indexed object or symbol; that identification is not independently established by the quoted local expression.\par
\par\noindent\textbf{OM-SOL-0005 -- t1 I p.6 -- SOL}\par
After Definition 1.2, sentence beginning «Lorsque C est U-petite»\par
\emph{Reported reading}: La catégorie C\_\{\textbackslash{}underline\{U\}\}, without a hat on C.\par
\emph{Proposal}: \textbackslash{}widehat\{C\}\_\{\ensuremath{𝒰}\}\par
The replacement \textbackslash{}widehat\{C\}\_\{\ensuremath{𝒰}\} at After Definition 1.2, sentence beginning «Lorsque C est U-petite» identifies a particular indexed object or symbol; that identification is not independently established by the quoted local expression.\par
\par\noindent\textbf{OM-SOL-0006 -- t1 I p.7 -- SOL}\par
Construction-definition 1.3, display defining h\_U(x)(u), final inverse factor\par
\emph{Reported reading}: \textbackslash{}varphi(y,x)\textasciicircum{}\{-1\}\par
\emph{Proposal}: \textbackslash{}varphi(y',x)\textasciicircum{}\{\{\}-1\}\par
The replacement \textbackslash{}varphi(y',x)\textasciicircum{}\{\{\}-1\} at Construction-definition 1.3, display defining h\_U(x)(u), final inverse factor identifies a particular indexed object or symbol; that identification is not independently established by the quoted local expression.\par
\par\noindent\textbf{OM-SOL-0009 -- t1 I p.12 -- SOL}\par
2.4.2, caractérisation de l’exactitude\par
\emph{Reported reading}: pour qu’un foncteur soit exact\par
\emph{Proposal}: pour qu'un foncteur soit exact à gauche\par
This finding intersects Orgogozo's \textbackslash{}sisi material at 2.4.2, caractérisation de l’exactitude; preserve that human entry unchanged. The proposed reading pour qu'un foncteur soit exact à gauche at 2.4.2, caractérisation de l’exactitude changes a substantive mathematical assertion or bibliographic identification; no independent derivation of this exact replacement is established in this review.\par
\par\noindent\textbf{OM-SOL-0012 -- t1 I p.15 -- SOL}\par
Définition 2.7, flèches vers l’objet c\par
\emph{Reported reading}: des flèches a c et b c\par
\emph{Proposal}: \$a\textbackslash{}rightarrow c\$ et \$b\textbackslash{}rightarrow c\$\par
This finding intersects Orgogozo's \textbackslash{}sisi material at Définition 2.7, flèches vers l’objet c; preserve that human entry unchanged. The replacement \$a\textbackslash{}rightarrow c\$ et \$b\textbackslash{}rightarrow c\$ at Définition 2.7, flèches vers l’objet c requires a complete check of the maps, their endpoints and the composition convention; this identification is not independently verified.\par
\par\noindent\textbf{OM-SOL-0015 -- t1 I p.23 -- SOL}\par
Diagram (*), horizontal arrows\par
\emph{Reported reading}: \$I\textasciicircum{}Y\_u\textbackslash{}longrightarrow I\textasciicircum{}f\_u\textbackslash{}longrightarrow I\textasciicircum{}\{Y'\}\_u\$\par
\emph{Proposal}: \$I\textasciicircum{}Y\_u\textbackslash{}xleftarrow\{I\textasciicircum{}f\_u\}I\textasciicircum{}\{Y'\}\_u\$\par
The replacement \$I\textasciicircum{}Y\_u\textbackslash{}xleftarrow\{I\textasciicircum{}f\_u\}I\textasciicircum{}\{Y'\}\_u\$ at Diagram (*), horizontal arrows requires a complete check of the maps, their endpoints and the composition convention; this identification is not independently verified.\par
\par\noindent\textbf{OM-SOL-0016 -- t1 I p.27 -- SOL}\par
Corollary 5.5.1, final parenthesis\par
\emph{Reported reading}: (1.4, 3))\par
\emph{Proposal}: (\textbackslash{}hyperref[I.prop5.4]\{5.4\}.3))\par
The proposed reference (\textbackslash{}hyperref[I.prop5.4]\{5.4\}.3)) at Corollary 5.5.1, final parenthesis needs comparison with the statement actually invoked; that target has not been independently established here.\par
\par\noindent\textbf{OM-SOL-0017 -- t1 I p.28 -- SOL}\par
Proof of Proposition 5.6, middle-left term of the diagram\par
\emph{Reported reading}: \textbackslash{}Hom\_\{\textbackslash{}widehat\{C\}\}(u\_\{!\}H,u\_\{\textbackslash{}ast\}K)\par
\emph{Proposal}: \textbackslash{}Hom\_\{\textbackslash{}widehat\{C'\}\}(u\_\{!\}H,u\_\{\textbackslash{}ast\}K)\par
The replacement \textbackslash{}Hom\_\{\textbackslash{}widehat\{C'\}\}(u\_\{!\}H,u\_\{\textbackslash{}ast\}K) at Proof of Proposition 5.6, middle-left term of the diagram identifies a particular indexed object or symbol; that identification is not independently established by the quoted local expression.\par
\par\noindent\textbf{OM-SOL-0018 -- t1 I p.29 -- SOL}\par
Remark 5.7(b), reference following « démonstration de »\par
\emph{Reported reading}: démonstration de 1.4\par
\emph{Proposal}: démonstration de \textbackslash{}hyperref[I.prop5.4]\{5.4\}\par
The proposed reference démonstration de \textbackslash{}hyperref[I.prop5.4]\{5.4\} at Remark 5.7(b), reference following « démonstration de » needs comparison with the statement actually invoked; that target has not been independently established here.\par
\par\noindent\textbf{OM-SOL-0019 -- t1 I p.29 -- SOL}\par
Remark 5.7(b), reference following « hypothèses générales de »\par
\emph{Reported reading}: hypothèses générales de 1.1\par
\emph{Proposal}: hypothèses générales de \textbackslash{}hyperref[I.prop5.1]\{5.1\}\par
The proposed reference hypothèses générales de \textbackslash{}hyperref[I.prop5.1]\{5.1\} at Remark 5.7(b), reference following « hypothèses générales de » needs comparison with the statement actually invoked; that target has not been independently established here.\par
\par\noindent\textbf{OM-SOL-0022 -- t1 I p.34 -- SOL}\par
Diagramme (5.13.2), sommet supérieur droit\par
\emph{Reported reading}: \$(C'/u\textasciicircum{}\{\textbackslash{}cdot\}H)\textasciicircum{}\{\textbackslash{}widehat\{\textbackslash{}phantom\{,\}\}\}\$\par
\emph{Proposal}: (C'/u\_\{!\}H)\textasciicircum{}\{\textbackslash{}widehat\{\textbackslash{}phantom\{,\}\}\}\par
The replacement (C'/u\_\{!\}H)\textasciicircum{}\{\textbackslash{}widehat\{\textbackslash{}phantom\{,\}\}\} at Diagramme (5.13.2), sommet supérieur droit identifies a particular indexed object or symbol; that identification is not independently established by the quoted local expression.\par
\par\noindent\textbf{OM-SOL-0023 -- t1 I p.36 -- SOL}\par
Preuve 1), item b), initial morphism\par
\emph{Reported reading}: m\_0=\allowbreak \textbackslash{}mathrm\{id\}\_X\par
\emph{Proposal}: m\_\textbackslash{}circ=\allowbreak \textbackslash{}id\_\{uX\}\par
The replacement m\_\textbackslash{}circ=\allowbreak \textbackslash{}id\_\{uX\} at Preuve 1), item b), initial morphism identifies a particular indexed object or symbol; that identification is not independently established by the quoted local expression.\par
\par\noindent\textbf{OM-SOL-0024 -- t1 I p.40 -- SOL}\par
Proposition 6.2(v), final category\par
\emph{Reported reading}: dans E\_i\par
\emph{Proposal}: dans \$F\_i\$\par
The proposed reading dans \$F\_i\$ at Proposition 6.2(v), final category changes a substantive mathematical assertion or bibliographic identification; no independent derivation of this exact replacement is established in this review.\par
\par\noindent\textbf{OM-SOL-0025 -- t1 I p.40 -- SOL}\par
Démonstration (iv), conclusion of the argument\par
\emph{Reported reading}: ce sont des monomorphismes d'après ce qui précède, donc des isomorphismes\par
\emph{Proposal}: \$u\$ est un monomorphisme d'après ce qui précède, donc un isomorphisme\par
The replacement \$u\$ est un monomorphisme d'après ce qui précède, donc un isomorphisme at Démonstration (iv), conclusion of the argument requires a complete check of the maps, their endpoints and the composition convention; this identification is not independently verified.\par
\par\noindent\textbf{OM-SOL-0028 -- t1 I p.41 -- SOL}\par
Proposition 6.4, critère de fidélité de l'adjoint à droite\par
\emph{Reported reading}: le morphisme d'adjonction \$X\textbackslash{}rightarrow\textbackslash{}psi\textbackslash{}varphi(X)\$ soit un monomorphisme\par
\emph{Proposal}: le morphisme d'adjonction \$X\textbackslash{}rightarrow \textbackslash{}psi\textbackslash{}varphi(X)\$ (\textbackslash{}resp \$\textbackslash{}varphi\textbackslash{}psi(Y)\textbackslash{}rightarrow Y\$) soit un monomorphisme (\textbackslash{}resp un épimorphisme).\par
The proposed reading le morphisme d'adjonction \$X\textbackslash{}rightarrow \textbackslash{}psi\textbackslash{}varphi(X)\$ (\textbackslash{}resp \$\textbackslash{}varphi\textbackslash{}psi(Y)\textbackslash{}rightarrow Y\$) soit un monomorphisme (\textbackslash{}resp un épimorphisme). at Proposition 6.4, critère de fidélité de l'adjoint à droite changes a substantive mathematical assertion or bibliographic identification; no independent derivation of this exact replacement is established in this review.\par
\par\noindent\textbf{OM-SOL-0029 -- t1 I p.43 -- SOL}\par
Démonstration de 6.5.2 (ii), dernière conclusion de la page\par
\emph{Reported reading}: donc \$E\_X\$ est ordonnée\par
\emph{Proposal}: donc \$E'\_X\$ est ordonnée.\par
The replacement donc \$E'\_X\$ est ordonnée. at Démonstration de 6.5.2 (ii), dernière conclusion de la page identifies a particular indexed object or symbol; that identification is not independently established by the quoted local expression.\par
\par\noindent\textbf{OM-SOL-0031 -- t1 I p.46 -- SOL}\par
Formule (7.1.3.1)\par
\emph{Reported reading}: \textbackslash{}widehat\{C\}=\allowbreak \textbackslash{}underline\{\textbackslash{}mathrm\{Hom\}\}(C,\textbackslash{}underline\{U\}\textbackslash{}text\{-\}\textbackslash{}mathrm\{Ens\})\par
\emph{Proposal}: \textbackslash{}widehat\{C\}=\allowbreak \textbackslash{}Hhom(C\textasciicircum{}\textbackslash{}circ, \ensuremath{𝒰}\textbackslash{}text\{-\}\textbackslash{}Ens)\par
The replacement \textbackslash{}widehat\{C\}=\allowbreak \textbackslash{}Hhom(C\textasciicircum{}\textbackslash{}circ, \ensuremath{𝒰}\textbackslash{}text\{-\}\textbackslash{}Ens) at Formule (7.1.3.1) requires a complete check of the maps, their endpoints and the composition convention; this identification is not independently verified.\par
\par\noindent\textbf{OM-SOL-0034 -- t1 I p.47 -- SOL}\par
7.2 (ii), diagramme d'implications, branche vers 2)\par
\emph{Reported reading}: Le terme 2) est isolé : aucune flèche ne relie 1) à 2).\par
\emph{Proposal}: pleinement fidèle)\}\} \textbackslash{}ar@\{=\allowbreak >\}[ld] \textbackslash{}ar@\{=\allowbreak >\}[rd]\&\textbackslash{}\textbackslash{}\par
The proposed reading pleinement fidèle)\}\} \textbackslash{}ar@\{=\allowbreak >\}[ld] \textbackslash{}ar@\{=\allowbreak >\}[rd]\&\textbackslash{}\textbackslash{} at 7.2 (ii), diagramme d'implications, branche vers 2) changes a substantive mathematical assertion or bibliographic identification; no independent derivation of this exact replacement is established in this review.\par
\par\noindent\textbf{OM-SOL-0035 -- t1 I p.48 -- SOL}\par
Corollaire 7.3, première phrase\par
\emph{Reported reading}: « Toutes les notions envisagées dans 6.1 »\par
\emph{Proposal}: Toutes les notions envisagées dans \textbackslash{}hyperref[I.def7.1]\{7.1\}\par
The proposed reference Toutes les notions envisagées dans \textbackslash{}hyperref[I.def7.1]\{7.1\} at Corollaire 7.3, première phrase needs comparison with the statement actually invoked; that target has not been independently established here.\par
\par\noindent\textbf{OM-SOL-0036 -- t1 I p.49 -- SOL}\par
Démonstration du corollaire 7.3, dernier renvoi\par
\emph{Reported reading}: « au diagramme d'implications 6.2 (ii) »\par
\emph{Proposal}: diagramme d'implications \textbackslash{}hyperref[I.prop7.2]\{7.2\} (ii)\par
The proposed reference diagramme d'implications \textbackslash{}hyperref[I.prop7.2]\{7.2\} (ii) at Démonstration du corollaire 7.3, dernier renvoi needs comparison with the statement actually invoked; that target has not been independently established here.\par
\par\noindent\textbf{OM-SOL-0038 -- t1 I p.49 -- SOL}\par
Dernier paragraphe, les deux catégories d'indices précédant la formule finale\par
\emph{Reported reading}: \textbackslash{}widehat\{C\}\_\{/\textbackslash{}varphi(X)\}, aux deux occurrences\par
\emph{Proposal}: C\_\{/\textbackslash{}varphi(X)\}\par
The proposed reading C\_\{/\textbackslash{}varphi(X)\} at Dernier paragraphe, les deux catégories d'indices précédant la formule finale changes a substantive mathematical assertion or bibliographic identification; no independent derivation of this exact replacement is established in this review.\par
\par\noindent\textbf{OM-SOL-0039 -- t1 I p.54 -- SOL}\par
Proof of Lemma 7.6.2, opening sentence\par
\emph{Reported reading}: la famille des X\_i \ensuremath{\times}\_X T \ensuremath{\to} est épimorphique\par
\emph{Proposal}: \textbackslash{}fprod\{X\_i\}\{T\}\{X\}\textbackslash{}rightarrow T\par
The proposed reading \textbackslash{}fprod\{X\_i\}\{T\}\{X\}\textbackslash{}rightarrow T at Proof of Lemma 7.6.2, opening sentence changes a substantive mathematical assertion or bibliographic identification; no independent derivation of this exact replacement is established in this review.\par
\par\noindent\textbf{OM-SOL-0040 -- t1 I p.54 -- SOL}\par
Proof of Lemma 7.6.2, generating assumption\par
\emph{Reported reading}: comme C est génératrice par épimorphismes stricts\par
\emph{Proposal}: génératrice par épimorphismes\par
Deleting stricts changes a generating hypothesis; the supplied snippet does not establish the exact hypothesis used in the proof.\par
\par\noindent\textbf{OM-SOL-0042 -- t1 I p.55 -- SOL}\par
Proposition 7.7, existence of a detecting functor\par
\emph{Reported reading}: un i \ensuremath{\in} J tel que\par
\emph{Proposal}: i\textbackslash{}in I\par
The replacement i\textbackslash{}in I at Proposition 7.7, existence of a detecting functor identifies a particular indexed object or symbol; that identification is not independently established by the quoted local expression.\par
\par\noindent\textbf{OM-SOL-0043 -- t1 I p.57 -- SOL}\par
Proof of 7.8, « Soit inversement v »\par
\emph{Reported reading}: \$v:F|D\textbackslash{}rightarrow G|D\$\par
\emph{Proposal}: \$v:F|C\textbackslash{}rightarrow G|C\$\par
This finding intersects Orgogozo's \textbackslash{}sisi material at Proof of 7.8, « Soit inversement v »; preserve that human entry unchanged. The proposed reading \$v:F|C\textbackslash{}rightarrow G|C\$ at Proof of 7.8, « Soit inversement v » changes a substantive mathematical assertion or bibliographic identification; no independent derivation of this exact replacement is established in this review.\par
\par\noindent\textbf{OM-SOL-0045 -- t1 I p.57 -- SOL}\par
Last sentence of the proof of 7.8, restrictions induced by u\par
\emph{Reported reading}: \$F|C'\$ dans \$G|C'\$\par
\emph{Proposal}: \$F|C\$ dans \$G|C\$\par
This finding intersects Orgogozo's \textbackslash{}sisi material at Last sentence of the proof of 7.8, restrictions induced by u; preserve that human entry unchanged. The replacement \$F|C\$ dans \$G|C\$ at Last sentence of the proof of 7.8, restrictions induced by u identifies a particular indexed object or symbol; that identification is not independently established by the quoted local expression.\par
\par\noindent\textbf{OM-SOL-0047 -- t1 I p.61 -- SOL}\par
Définition 8.1.1, catégorie contenant les limites inductives\par
\emph{Reported reading}: \textbackslash{}Hhom(G(\ensuremath{𝒱}\textbackslash{}text\{-\}\textbackslash{}Ens))\textasciicircum{}\textbackslash{}circ\par
\emph{Proposal}: \textbackslash{}Hhom(C,(\ensuremath{𝒱}\textbackslash{}text\{-\}\textbackslash{}Ens))\textasciicircum{}\textbackslash{}circ\par
The proposed reading \textbackslash{}Hhom(C,(\ensuremath{𝒱}\textbackslash{}text\{-\}\textbackslash{}Ens))\textasciicircum{}\textbackslash{}circ at Définition 8.1.1, catégorie contenant les limites inductives changes a substantive mathematical assertion or bibliographic identification; no independent derivation of this exact replacement is established in this review.\par
\par\noindent\textbf{OM-SOL-0048 -- t1 I p.63 -- SOL}\par
Démonstration de 8.1.3, application de F 1) à la double flèche\par
\emph{Reported reading}: on peut supposer \$k'\$ de la forme \$\textbackslash{}varphi(i)\$\par
\emph{Proposal}: on peut supposer \$j'\$ de la forme \$\textbackslash{}varphi(i)\$\par
The proposed reading on peut supposer \$j'\$ de la forme \$\textbackslash{}varphi(i)\$ at Démonstration de 8.1.3, application de F 1) à la double flèche changes a substantive mathematical assertion or bibliographic identification; no independent derivation of this exact replacement is established in this review.\par
\par\noindent\textbf{OM-SOL-0049 -- t1 I p.66 -- SOL}\par
3°, opening sentence\par
\emph{Reported reading}: \$D,D'\textbackslash{}subset E\$\par
\emph{Proposal}: \$D,D'\textbackslash{}in E\$\par
The proposed reading \$D,D'\textbackslash{}in E\$ at 3°, opening sentence changes a substantive mathematical assertion or bibliographic identification; no independent derivation of this exact replacement is established in this review.\par
\par\noindent\textbf{OM-SOL-0050 -- t1 I p.66 -- SOL}\par
3°, final equation\par
\emph{Reported reading}: \$ff\_x=\allowbreak f'f'\_\{x'\}\$\par
\emph{Proposal}: \$ff\_x=\allowbreak f'f'\_x\$\par
The proposed reading \$ff\_x=\allowbreak f'f'\_x\$ at 3°, final equation changes a substantive mathematical assertion or bibliographic identification; no independent derivation of this exact replacement is established in this review.\par
\par\noindent\textbf{OM-SOL-0052 -- t1 I p.70 -- SOL}\par
After (8.2.4.3), definition of composition\par
\emph{Reported reading}: On définit la composition des ind-objets de \$\textbackslash{}underline\{C\}\$\par
\emph{Proposal}: On définit la composition des morphismes des ind-objets de \$\ensuremath{𝒞}\$\par
The replacement On définit la composition des morphismes des ind-objets de \$\ensuremath{𝒞}\$ at After (8.2.4.3), definition of composition requires a complete check of the maps, their endpoints and the composition convention; this identification is not independently verified.\par
\par\noindent\textbf{OM-SOL-0055 -- t1 I p.74 -- SOL}\par
Property 1) of the functor G\par
\emph{Reported reading}: F(s)\par
\emph{Proposal}: \$G(s)\$\par
The proposed reading \$G(s)\$ at Property 1) of the functor G changes a substantive mathematical assertion or bibliographic identification; no independent derivation of this exact replacement is established in this review.\par
\par\noindent\textbf{OM-SOL-0056 -- t1 I p.76 -- SOL}\par
Théorème 8.3.3, (iii bis), dernière flèche\par
\emph{Reported reading}: \$X\_i\textbackslash{}rightarrow X\$\par
\emph{Proposal}: \$X\textbackslash{}rightarrow X\_i\$\par
The replacement \$X\textbackslash{}rightarrow X\_i\$ at Théorème 8.3.3, (iii bis), dernière flèche requires a complete check of the maps, their endpoints and the composition convention; this identification is not independently verified.\par
\par\noindent\textbf{OM-SOL-0057 -- t1 I p.76 -- SOL}\par
Théorème 8.3.3, hypothèse de (v)\par
\emph{Reported reading}: les limites inductives filtrantes\par
\emph{Proposal}: les limites inductives finies\par
The finiteness qualifier changes the theorem; equivalence of the two listed conditions needs proof from the surrounding statement.\par
\par\noindent\textbf{OM-SOL-0058 -- t1 I p.77 -- SOL}\par
Démonstration de (ii) \ensuremath{\Rightarrow} (i), foncteur canonique affiché\par
\emph{Reported reading}: \$\textbackslash{}varphi:I=\allowbreak C\_\{/F\}\textbackslash{}longrightarrow F\$\par
\emph{Proposal}: \$\textbackslash{}varphi:I=\allowbreak C\_\{/F\}\textbackslash{}longrightarrow C\$\par
The replacement \$\textbackslash{}varphi:I=\allowbreak C\_\{/F\}\textbackslash{}longrightarrow C\$ at Démonstration de (ii) \ensuremath{\Rightarrow} (i), foncteur canonique affiché requires a complete check of the maps, their endpoints and the composition convention; this identification is not independently verified.\par
\par\noindent\textbf{OM-SOL-0059 -- t1 I p.78 -- SOL}\par
Remarque 8.3.4, parenthèse « considéré comme préfaisceau »\par
\emph{Reported reading}: préfaisceau sur \$C\textasciicircum{}\textbackslash{}circ\$\par
\emph{Proposal}: préfaisceau sur \$C\$\par
The replacement préfaisceau sur \$C\$ at Remarque 8.3.4, parenthèse « considéré comme préfaisceau » identifies a particular indexed object or symbol; that identification is not independently established by the quoted local expression.\par
\par\noindent\textbf{OM-SOL-0064 -- t1 I p.83 -- SOL}\par
Exercice 8.5.5 a), condition (i)\par
\emph{Reported reading}: Dans C les petites limites inductives sont représentables\par
\emph{Proposal}: Dans \$C\$ les petites limites inductives filtrantes sont représentables,\par
The proposed reading Dans \$C\$ les petites limites inductives filtrantes sont représentables, at Exercice 8.5.5 a), condition (i) changes a substantive mathematical assertion or bibliographic identification; no independent derivation of this exact replacement is established in this review.\par
\par\noindent\textbf{OM-SOL-0065 -- t1 I p.83 -- SOL}\par
Exercice 8.5.5 a), condition (ii)\par
\emph{Reported reading}: Dans C les petites limites inductives sont représentables\par
\emph{Proposal}: Dans \$C\$ les petites limites inductives filtrantes sont représentables,\par
The proposed reading Dans \$C\$ les petites limites inductives filtrantes sont représentables, at Exercice 8.5.5 a), condition (ii) changes a substantive mathematical assertion or bibliographic identification; no independent derivation of this exact replacement is established in this review.\par
\par\noindent\textbf{OM-SOL-0068 -- t1 I p.89 -- SOL}\par
8.7.1.7, dernière ligne, domaine de la variable X\par
\emph{Reported reading}: \$X\textbackslash{}in \textbackslash{}ob \textbackslash{}Ind(C)'\$\par
\emph{Proposal}: \$X\textbackslash{}in \textbackslash{}ob C\$\par
The proposed reading \$X\textbackslash{}in \textbackslash{}ob C\$ at 8.7.1.7, dernière ligne, domaine de la variable X changes a substantive mathematical assertion or bibliographic identification; no independent derivation of this exact replacement is established in this review.\par
\par\noindent\textbf{OM-SOL-0069 -- t1 I p.92 -- SOL}\par
Formule (8.7.4.2), après la flèche d'application\par
\emph{Reported reading}: \textbackslash{}lim\_i f(X\_i), sans flèche de limite inductive\par
\emph{Proposal}: \textbackslash{}varinjlim\_i f(X\_i)\par
The replacement \textbackslash{}varinjlim\_i f(X\_i) at Formule (8.7.4.2), après la flèche d'application requires a complete check of the maps, their endpoints and the composition convention; this identification is not independently verified.\par
\par\noindent\textbf{OM-SOL-0070 -- t1 I p.93 -- SOL}\par
Démonstration de 8.7.5 a), dernière phrase\par
\emph{Reported reading}: deux ind-objets quelconques de E\par
\emph{Proposal}: deux ind-objets quelconques de \$C\$\par
The replacement deux ind-objets quelconques de \$C\$ at Démonstration de 8.7.5 a), dernière phrase requires a complete check of the maps, their endpoints and the composition convention; this identification is not independently verified.\par
\par\noindent\textbf{OM-SOL-0071 -- t1 I p.94 -- SOL}\par
Première formule, cible de u\_j\par
\emph{Reported reading}: \textbackslash{}Hom(\textbackslash{}bar\{f\}(\textbackslash{}underline\{X\}\_j),\textbackslash{}underline\{Y\})\par
\emph{Proposal}: \textbackslash{}Hom(\textbackslash{}bar\{f\}(\ensuremath{𝒳}\_j),\textbackslash{}bar\{f\}(\ensuremath{𝒴}))\par
The replacement \textbackslash{}Hom(\textbackslash{}bar\{f\}(\ensuremath{𝒳}\_j),\textbackslash{}bar\{f\}(\ensuremath{𝒴})) at Première formule, cible de u\_j identifies a particular indexed object or symbol; that identification is not independently established by the quoted local expression.\par
\par\noindent\textbf{OM-SOL-0072 -- t1 I p.97 -- SOL}\par
Premier paragraphe, factorisation de l'identité de X\par
\emph{Reported reading}: l'homomorphisme canonique X\_i \ensuremath{\to} X admet un inverse à gauche\par
\emph{Proposal}: \$X\_i\textbackslash{}rightarrow X\$ admet un inverse à droite\par
The replacement \$X\_i\textbackslash{}rightarrow X\$ admet un inverse à droite at Premier paragraphe, factorisation de l'identité de X requires a complete check of the maps, their endpoints and the composition convention; this identification is not independently verified.\par
\par\noindent\textbf{OM-SOL-0073 -- t1 I p.107 -- SOL}\par
Exercice 8.8.7 b), catégorie de restriction\par
\emph{Reported reading}: I\_\{/i\_\textbackslash{}circ\}\par
\emph{Proposal}: I\_\{i\_\textbackslash{}circ/\}\par
The replacement I\_\{i\_\textbackslash{}circ/\} at Exercice 8.8.7 b), catégorie de restriction requires a complete check of the maps, their endpoints and the composition convention; this identification is not independently verified.\par
\par\noindent\textbf{OM-SOL-0074 -- t1 I p.107 -- SOL}\par
Exercice 8.8.7 c), identification de C tilde\par
\emph{Reported reading}: C\textasciicircum{}\{\textbackslash{}sim\} est isomorphe à \textbackslash{}NN\textasciicircum{}\textbackslash{}circ\textbackslash{}times\textbackslash{}NN\textasciicircum{}\{\textbackslash{}sim\}\par
\emph{Proposal}: C\textasciicircum{}\{\textbackslash{}widetilde\{\textbackslash{}phantom\{,\}\}\}\textbackslash{}simeq\textbackslash{}\{(j,i)\textbackslash{}in\textbackslash{}NN\textasciicircum{}\textbackslash{}circ\textbackslash{}times\textbackslash{}NN\textasciicircum{}\{\textbackslash{}widetilde\{\textbackslash{}phantom\{,\}\}\}\textbackslash{}mid i\textbackslash{}geq j\textbackslash{}\}\par
The replacement C\textasciicircum{}\{\textbackslash{}widetilde\{\textbackslash{}phantom\{,\}\}\}\textbackslash{}simeq\textbackslash{}\{(j,i)\textbackslash{}in\textbackslash{}NN\textasciicircum{}\textbackslash{}circ\textbackslash{}times\textbackslash{}NN\textasciicircum{}\{\textbackslash{}widetilde\{\textbackslash{}phantom\{,\}\}\}\textbackslash{}mid i\textbackslash{}geq j\textbackslash{}\} at Exercice 8.8.7 c), identification de C tilde identifies a particular indexed object or symbol; that identification is not independently established by the quoted local expression.\par
\par\noindent\textbf{OM-SOL-0077 -- t1 I p.109 -- SOL}\par
Démonstration a), après « fait de F une »\par
\emph{Reported reading}: limite projective des F\_j\par
\emph{Proposal}: limite projective des \$F\_i\$\par
The replacement limite projective des \$F\_i\$ at Démonstration a), après « fait de F une » identifies a particular indexed object or symbol; that identification is not independently established by the quoted local expression.\par
\par\noindent\textbf{OM-SOL-0079 -- t1 I p.110 -- SOL}\par
Démonstration, lettres des deux premiers alinéas\par
\emph{Reported reading}: c) Résulte formellement ... ; d) Il est bien connu ...\par
\emph{Proposal}: d) Résulte formellement ... ; c) Il est bien connu ...\par
The proposed reading d) Résulte formellement ... ; c) Il est bien connu ... at Démonstration, lettres des deux premiers alinéas changes a substantive mathematical assertion or bibliographic identification; no independent derivation of this exact replacement is established in this review.\par
\par\noindent\textbf{OM-SOL-0080 -- t1 I p.110 -- SOL}\par
Début du paragraphe « Supposons d’abord »\par
\emph{Reported reading}: Supposons d'abord I fini.\par
\emph{Proposal}: Supposons d'abord \$J\$ fini.\par
The replacement Supposons d'abord \$J\$ fini. at Début du paragraphe « Supposons d’abord » identifies a particular indexed object or symbol; that identification is not independently established by the quoted local expression.\par
\par\noindent\textbf{OM-SOL-0082 -- t1 I p.110 -- SOL}\par
Numéro du corollaire\par
\emph{Reported reading}: Corollaire 8.8.2\par
\emph{Proposal}: Corollaire 8.9.2\par
This finding intersects Orgogozo's \textbackslash{}sisi material at Numéro du corollaire; preserve that human entry unchanged. The proposed reading Corollaire 8.9.2 at Numéro du corollaire changes a substantive mathematical assertion or bibliographic identification; no independent derivation of this exact replacement is established in this review.\par
\par\noindent\textbf{OM-SOL-0083 -- t1 I p.111 -- SOL}\par
Fin de la preuve de 8.9.2, renvoi après « prouvé dans »\par
\emph{Reported reading}: 8.9.4 a)\par
\emph{Proposal}: \textbackslash{}hyperref[I.prop8.9.5]\{8.9.5\} a)\par
The proposed reference \textbackslash{}hyperref[I.prop8.9.5]\{8.9.5\} a) at Fin de la preuve de 8.9.2, renvoi après « prouvé dans » needs comparison with the statement actually invoked; that target has not been independently established here.\par
\par\noindent\textbf{OM-SOL-0084 -- t1 I p.112 -- SOL}\par
Remarques 8.9.4, « alors il en est de même de… »\par
\emph{Reported reading}: \textbackslash{}varprojlim\_i\textbackslash{}varphi\par
\emph{Proposal}: \textbackslash{}varprojlim\_j\textbackslash{}varphi\par
The replacement \textbackslash{}varprojlim\_j\textbackslash{}varphi at Remarques 8.9.4, « alors il en est de même de… » identifies a particular indexed object or symbol; that identification is not independently established by the quoted local expression.\par
\par\noindent\textbf{OM-SOL-0088 -- t1 I p.117 -- SOL}\par
Exercice 8.9.10 b), expression après « chacune des familles »\par
\emph{Reported reading}: familles \textbackslash{}varprojlim\_J X\_\{i\_\textbackslash{}alpha\}\par
\emph{Proposal}: familles \$(X\_\{i\_\textbackslash{}alpha\})\_\{\textbackslash{}alpha\textbackslash{}in J\}\$\par
The replacement familles \$(X\_\{i\_\textbackslash{}alpha\})\_\{\textbackslash{}alpha\textbackslash{}in J\}\$ at Exercice 8.9.10 b), expression après « chacune des familles » identifies a particular indexed object or symbol; that identification is not independently established by the quoted local expression.\par
\par\noindent\textbf{OM-SOL-0089 -- t1 I p.117 -- SOL}\par
Exercice 8.9.11 a), dernière ligne\par
\emph{Reported reading}: morphismes de but X\_i\par
\emph{Proposal}: morphismes de but \$X\$\par
The proposed reading morphismes de but \$X\$ at Exercice 8.9.11 a), dernière ligne changes a substantive mathematical assertion or bibliographic identification; no independent derivation of this exact replacement is established in this review.\par
\par\noindent\textbf{OM-SOL-0091 -- t1 I p.123 -- SOL}\par
8.11.1, object assignment before (8.11.1.2)\par
\emph{Reported reading}: F\textbackslash{}longrightarrow F\textbackslash{}circ f\par
\emph{Proposal}: F\textbackslash{}longmapsto F\textbackslash{}circ f\par
This finding intersects Orgogozo's \textbackslash{}sisi material at 8.11.1, object assignment before (8.11.1.2); preserve that human entry unchanged. The replacement F\textbackslash{}longmapsto F\textbackslash{}circ f at 8.11.1, object assignment before (8.11.1.2) requires a complete check of the maps, their endpoints and the composition convention; this identification is not independently verified.\par
\par\noindent\textbf{OM-SOL-0092 -- t1 I p.123 -- SOL}\par
8.11.1, « la sous-catégorie pleine de C′ formée des foncteurs ind-représentables »\par
\emph{Reported reading}: sous-catégorie pleine de C′\par
\emph{Proposal}: sous-catégorie pleine de \$\textbackslash{}widehat\{C\}\$ formée des foncteurs\par
The replacement sous-catégorie pleine de \$\textbackslash{}widehat\{C\}\$ formée des foncteurs at 8.11.1, « la sous-catégorie pleine de C′ formée des foncteurs ind-représentables » identifies a particular indexed object or symbol; that identification is not independently established by the quoted local expression.\par
\par\noindent\textbf{OM-SOL-0093 -- t1 I p.123 -- SOL}\par
8.11.1, « f* applique objets de C »\par
\emph{Reported reading}: applique objets de C\par
\emph{Proposal}: applique objets de \$C'\$\par
The replacement applique objets de \$C'\$ at 8.11.1, « f* applique objets de C » identifies a particular indexed object or symbol; that identification is not independently established by the quoted local expression.\par
\par\noindent\textbf{OM-SOL-0094 -- t1 I p.123 -- SOL}\par
8.11.1, parenthesis following « applique objets de C »\par
\emph{Reported reading}: identifié à une sous-catégorie pleine de \$\textbackslash{}widehat\{C\}\$\par
\emph{Proposal}: sous-catégorie pleine de \$\textbackslash{}widehat\{C\}'\$)\par
The replacement sous-catégorie pleine de \$\textbackslash{}widehat\{C\}'\$) at 8.11.1, parenthesis following « applique objets de C » identifies a particular indexed object or symbol; that identification is not independently established by the quoted local expression.\par
\par\noindent\textbf{OM-SOL-0095 -- t1 I p.124 -- SOL}\par
Immediately after (8.11.1.5), quantification of Y′\par
\emph{Reported reading}: Ordinary, non-underlined Y′ in « Y′ \ensuremath{\in} Ob Ind(C′) »\par
\emph{Proposal}: \$\ensuremath{𝒴}'\textbackslash{}in \textbackslash{}Ob \textbackslash{}Ind(C')\$\par
The replacement \$\ensuremath{𝒴}'\textbackslash{}in \textbackslash{}Ob \textbackslash{}Ind(C')\$ at Immediately after (8.11.1.5), quantification of Y′ identifies a particular indexed object or symbol; that identification is not independently established by the quoted local expression.\par
\par\noindent\textbf{OM-SOL-0096 -- t1 I p.124 -- SOL}\par
Paragraph before (8.11.1.8), « prolongement canonique de f »\par
\emph{Reported reading}: prolongement canonique de f aux ind-objets\par
\emph{Proposal}: prolongement canonique de \$f\textasciicircum{}\{\textbackslash{}ad\}\$ aux ind-objets\par
The replacement prolongement canonique de \$f\textasciicircum{}\{\textbackslash{}ad\}\$ aux ind-objets at Paragraph before (8.11.1.8), « prolongement canonique de f » identifies a particular indexed object or symbol; that identification is not independently established by the quoted local expression.\par
\par\noindent\textbf{OM-SOL-0097 -- t1 I p.125 -- SOL}\par
Proof of Proposition 8.11.2, immediately before (8.11.2.1)\par
\emph{Reported reading}: en les pro-objets\par
\emph{Proposal}: en les ind-objets\par
Changing pro-objets to ind-objets requires checking the variance and the displayed adjunction.\par
\par\noindent\textbf{OM-SOL-0098 -- t1 I p.125 -- SOL}\par
Proposition 8.11.4, necessary exactness condition\par
\emph{Reported reading}: exact à gauche\par
\emph{Proposal}: exact à droite\par
This finding intersects Orgogozo's \textbackslash{}sisi material at Proposition 8.11.4, necessary exactness condition; preserve that human entry unchanged. The proposed reading exact à droite at Proposition 8.11.4, necessary exactness condition changes a substantive mathematical assertion or bibliographic identification; no independent derivation of this exact replacement is established in this review.\par
\par\noindent\textbf{OM-SOL-0099 -- t1 I p.126 -- SOL}\par
Immediately after (8.11.5.1)\par
\emph{Reported reading}: induit par F\par
\emph{Proposal}: induit par \$f\$\par
This finding intersects Orgogozo's \textbackslash{}sisi material at Immediately after (8.11.5.1); preserve that human entry unchanged. The proposed reading induit par \$f\$ at Immediately after (8.11.5.1) changes a substantive mathematical assertion or bibliographic identification; no independent derivation of this exact replacement is established in this review.\par
\par\noindent\textbf{OM-SOL-0100 -- t1 I p.126 -- SOL}\par
(8.11.5.2), domain and codomain of g\par
\emph{Reported reading}: g:\textbackslash{}Pro(C)\textbackslash{}longrightarrow\textbackslash{}Pro(C')\par
\emph{Proposal}: g:\textbackslash{}Pro(C')\textbackslash{}longrightarrow\textbackslash{}Pro(C)\par
This finding intersects Orgogozo's \textbackslash{}sisi material at (8.11.5.2), domain and codomain of g; preserve that human entry unchanged. The proposed reading g:\textbackslash{}Pro(C')\textbackslash{}longrightarrow\textbackslash{}Pro(C) at (8.11.5.2), domain and codomain of g changes a substantive mathematical assertion or bibliographic identification; no independent derivation of this exact replacement is established in this review.\par
\par\noindent\textbf{OM-SOL-0102 -- t1 I p.126 -- SOL}\par
Paragraph after (8.11.5.3), “Pro(f) … admet un adjoint”\par
\emph{Reported reading}: admet un adjoint à droite\par
\emph{Proposal}: admet un adjoint à gauche\par
This finding intersects Orgogozo's \textbackslash{}sisi material at Paragraph after (8.11.5.3), “Pro(f) … admet un adjoint”; preserve that human entry unchanged. The replacement admet un adjoint à gauche at Paragraph after (8.11.5.3), “Pro(f) … admet un adjoint” requires a complete check of the maps, their endpoints and the composition convention; this identification is not independently verified.\par
\par\noindent\textbf{OM-SOL-0103 -- t1 I p.126 -- SOL}\par
Same paragraph, “le foncteur g précédent est un tel adjoint”\par
\emph{Reported reading}: est un tel adjoint à droite\par
\emph{Proposal}: est un tel adjoint à gauche\par
This finding intersects Orgogozo's \textbackslash{}sisi material at Same paragraph, “le foncteur g précédent est un tel adjoint”; preserve that human entry unchanged. The proposed reading est un tel adjoint à gauche at Same paragraph, “le foncteur g précédent est un tel adjoint” changes a substantive mathematical assertion or bibliographic identification; no independent derivation of this exact replacement is established in this review.\par
\par\noindent\textbf{OM-SOL-0104 -- t1 I p.126 -- SOL}\par
Same paragraph, immediately after the second “adjoint à droite”\par
\emph{Reported reading}: de Pro(C)\par
\emph{Proposal}: de \$\textbackslash{}Pro(f)\$\par
This finding intersects Orgogozo's \textbackslash{}sisi material at Same paragraph, immediately after the second “adjoint à droite”; preserve that human entry unchanged. The proposed reading de \$\textbackslash{}Pro(f)\$ at Same paragraph, immediately after the second “adjoint à droite” changes a substantive mathematical assertion or bibliographic identification; no independent derivation of this exact replacement is established in this review.\par
\par\noindent\textbf{OM-SOL-0105 -- t1 I p.131 -- SOL}\par
Démonstration, b), nécessité\par
\emph{Reported reading}: foncteurs exacts à droite\par
\emph{Proposal}: foncteurs exacts à gauche\par
Reversing exactitude à droite to à gauche requires checking the hypothesis and the argument, beyond this quoted phrase.\par
\par\noindent\textbf{OM-SOL-0106 -- t1 I p.134 -- SOL}\par
Démonstration du lemme 8.12.8.1, catégorie dite filtrante\par
\emph{Reported reading}: \$(\{\}\_F\textbackslash{}backslash C\textasciicircum{}\textbackslash{}circ)\$\par
\emph{Proposal}: \$(\{\}\_F\textbackslash{}backslash C)\textasciicircum{}\textbackslash{}circ\$\par
The replacement \$(\{\}\_F\textbackslash{}backslash C)\textasciicircum{}\textbackslash{}circ\$ at Démonstration du lemme 8.12.8.1, catégorie dite filtrante identifies a particular indexed object or symbol; that identification is not independently established by the quoted local expression.\par
\par\noindent\textbf{OM-SOL-0107 -- t1 I p.136 -- SOL}\par
8.13.2, démonstration b), catégorie où chercher une sous-catégorie cofinale\par
\emph{Reported reading}: \textbackslash{}operatorname\{Ob\}(\{\}\_F\textbackslash{}backslash C)\textasciicircum{}\textbackslash{}circ\par
\emph{Proposal}: (\{\}\_F\textbackslash{}backslash C)\textasciicircum{}\textbackslash{}circ\par
The proposed reading (\{\}\_F\textbackslash{}backslash C)\textasciicircum{}\textbackslash{}circ at 8.13.2, démonstration b), catégorie où chercher une sous-catégorie cofinale changes a substantive mathematical assertion or bibliographic identification; no independent derivation of this exact replacement is established in this review.\par
\par\noindent\textbf{OM-SOL-0108 -- t1 I p.136 -- SOL}\par
8.13.2, démonstration b), indice de la catégorie comma\par
\emph{Reported reading}: \{\}\_F\textbackslash{}backslash C\par
\emph{Proposal}: \{\}\_f\textbackslash{}backslash C\par
The proposed reading \{\}\_f\textbackslash{}backslash C at 8.13.2, démonstration b), indice de la catégorie comma changes a substantive mathematical assertion or bibliographic identification; no independent derivation of this exact replacement is established in this review.\par
\par\noindent\textbf{OM-SOL-0109 -- t1 I p.137 -- SOL}\par
Premier paragraphe, fin de la démonstration de 8.13.2 b)\par
\emph{Reported reading}: F(X), F(X\_i), \{\}\_F\textbackslash{}backslash C, \{\}\_F\textbackslash{}backslash C'\par
\emph{Proposal}: f(X),\textbackslash{}quad f(X\_i),\textbackslash{}quad \{\}\_f\textbackslash{}backslash C,\textbackslash{}quad \{\}\_f\textbackslash{}backslash C'\par
The proposed reading f(X),\textbackslash{}quad f(X\_i),\textbackslash{}quad \{\}\_f\textbackslash{}backslash C,\textbackslash{}quad \{\}\_f\textbackslash{}backslash C' at Premier paragraphe, fin de la démonstration de 8.13.2 b) changes a substantive mathematical assertion or bibliographic identification; no independent derivation of this exact replacement is established in this review.\par
\par\noindent\textbf{OM-SOL-0110 -- t1 I p.138 -- SOL}\par
Preuve de 8.13.4, les deux occurrences de « Y' \ensuremath{\in} Ob C' »\par
\emph{Reported reading}: Y'\textbackslash{}in \textbackslash{}operatorname\{Ob\} C'\par
\emph{Proposal}: Y'\textbackslash{}in \textbackslash{}Ob C'\_1\par
The replacement Y'\textbackslash{}in \textbackslash{}Ob C'\_1 at Preuve de 8.13.4, les deux occurrences de « Y' \ensuremath{\in} Ob C' » identifies a particular indexed object or symbol; that identification is not independently established by the quoted local expression.\par
\par\noindent\textbf{OM-SOL-0116 -- t1 I p.146 -- SOL}\par
Proof, paragraph beginning « Le fait que »\par
\emph{Reported reading}: tout objet de E soit accessible\par
\emph{Proposal}: tout objet de \$E\$ soit préaccessible\par
The proposed reading tout objet de \$E\$ soit préaccessible at Proof, paragraph beginning « Le fait que » changes a substantive mathematical assertion or bibliographic identification; no independent derivation of this exact replacement is established in this review.\par
\par\noindent\textbf{OM-SOL-0117 -- t1 I p.147 -- SOL}\par
Lemma 9.11.2, final hypothesis\par
\emph{Reported reading}: \ensuremath{\pi}-préadmissible\par
\emph{Proposal}: \$\textbackslash{}Pi\$-préaccessible\par
The proposed reading \$\textbackslash{}Pi\$-préaccessible at Lemma 9.11.2, final hypothesis changes a substantive mathematical assertion or bibliographic identification; no independent derivation of this exact replacement is established in this review.\par
\par\noindent\textbf{OM-SOL-0120 -- t1 I p.149 -- SOL}\par
Definition 9.12 b), reference after L\_\{\ensuremath{\Pi}₀\}\par
\emph{Reported reading}: (9.3)\par
\emph{Proposal}: (\textbackslash{}hyperref[I.def9.1]\{9.1\})\par
The proposed reference (\textbackslash{}hyperref[I.def9.1]\{9.1\}) at Definition 9.12 b), reference after L\_\{\ensuremath{\Pi}₀\} needs comparison with the statement actually invoked; that target has not been independently established here.\par
\par\noindent\textbf{OM-SOL-0121 -- t1 I p.150 -- SOL}\par
End of Definition 9.12 c), additional cardinal bound\par
\emph{Reported reading}: si X\textbackslash{}in\textbackslash{}ob\textbackslash{}Filt\textasciicircum{}\{\textbackslash{}Pi\}(E), \textbackslash{}Pi'\textbackslash{}geq\textbackslash{}Pi\par
\emph{Proposal}: si \$X\textbackslash{}in \textbackslash{}ob \textbackslash{}Filt\textasciicircum{}\{\textbackslash{}Pi'\}(E)\$, \$\textbackslash{}Pi'\textbackslash{}geq \textbackslash{}Pi\$\par
The replacement si \$X\textbackslash{}in \textbackslash{}ob \textbackslash{}Filt\textasciicircum{}\{\textbackslash{}Pi'\}(E)\$, \$\textbackslash{}Pi'\textbackslash{}geq \textbackslash{}Pi\$ at End of Definition 9.12 c), additional cardinal bound identifies a particular indexed object or symbol; that identification is not independently established by the quoted local expression.\par
\par\noindent\textbf{OM-SOL-0122 -- t1 I p.152 -- SOL}\par
Parenthetical application immediately before Lemma 9.13.1\par
\emph{Reported reading}: I=\allowbreak C\_\{/X\}\par
\emph{Proposal}: J=\allowbreak C\_\{/X\}\par
The proposed reading J=\allowbreak C\_\{/X\} at Parenthetical application immediately before Lemma 9.13.1 changes a substantive mathematical assertion or bibliographic identification; no independent derivation of this exact replacement is established in this review.\par
\par\noindent\textbf{OM-SOL-0126 -- t1 I p.157 -- SOL}\par
Proof of Corollaire 9.19, domain of the functor defined by the inductive limit\par
\emph{Reported reading}: \$\textbackslash{}Ind(E\_\{\textbackslash{}Pi\})\$\par
\emph{Proposal}: \$\textbackslash{}Ind(E\_\{\textbackslash{}Pi\})\_\{\textbackslash{}Pi\}\$\par
The replacement \$\textbackslash{}Ind(E\_\{\textbackslash{}Pi\})\_\{\textbackslash{}Pi\}\$ at Proof of Corollaire 9.19, domain of the functor defined by the inductive limit identifies a particular indexed object or symbol; that identification is not independently established by the quoted local expression.\par
\par\noindent\textbf{OM-SOL-0127 -- t1 I p.158 -- SOL}\par
Exercice 9.20.1, condition on the initial cardinal\par
\emph{Reported reading}: \$\textbackslash{}Pi\_\textbackslash{}circ\textbackslash{}in \textbackslash{}card F\textbackslash{}ell\textbackslash{},C\$\par
\emph{Proposal}: \$\textbackslash{}Pi\_\textbackslash{}circ\textbackslash{}geq \textbackslash{}card F\textbackslash{}ell\textbackslash{},C\$\par
The proposed reading \$\textbackslash{}Pi\_\textbackslash{}circ\textbackslash{}geq \textbackslash{}card F\textbackslash{}ell\textbackslash{},C\$ at Exercice 9.20.1, condition on the initial cardinal changes a substantive mathematical assertion or bibliographic identification; no independent derivation of this exact replacement is established in this review.\par
\par\noindent\textbf{OM-SOL-0128 -- t1 I p.164 -- SOL}\par
Première ligne, « dans les E »\par
\emph{Reported reading}: des systèmes inductifs dans les \$E\$\par
\emph{Proposal}: dans les \$E\_\{\textbackslash{}alpha\}\$\par
The replacement dans les \$E\_\{\textbackslash{}alpha\}\$ at Première ligne, « dans les E » identifies a particular indexed object or symbol; that identification is not independently established by the quoted local expression.\par
\par\noindent\textbf{OM-SOL-0130 -- t1 I p.166 -- SOL}\par
Après « En vertu de (9.21.4 a)) »\par
\emph{Reported reading}: X(\textbackslash{}alpha)\textbackslash{}in \textbackslash{}ob C\_\{\textbackslash{}alpha\}\par
\emph{Proposal}: X(\textbackslash{}alpha)\_i\textbackslash{}in \textbackslash{}ob C\_\{\textbackslash{}alpha\}\par
The replacement X(\textbackslash{}alpha)\_i\textbackslash{}in \textbackslash{}ob C\_\{\textbackslash{}alpha\} at Après « En vertu de (9.21.4 a)) » requires a complete check of the maps, their endpoints and the composition convention; this identification is not independently verified.\par
\par\noindent\textbf{OM-SOL-0131 -- t1 I p.166 -- SOL}\par
« On trouve ainsi, pour i\ensuremath{\in}I et f\ensuremath{\in}… »\par
\emph{Reported reading}: \textbackslash{}mathrm\{F\}\{\}l\textbackslash{}, F\par
\emph{Proposal}: \textbackslash{}mathrm\{F\}\{\}l\textbackslash{}, B\par
The replacement \textbackslash{}mathrm\{F\}\{\}l\textbackslash{}, B at « On trouve ainsi, pour i\ensuremath{\in}I et f\ensuremath{\in}… » requires a complete check of the maps, their endpoints and the composition convention; this identification is not independently verified.\par
\par\noindent\textbf{OM-SOL-0132 -- t1 I p.167 -- SOL}\par
Premier paragraphe, « Donc, comme … est c-accessible »\par
\emph{Reported reading}: X(\textbackslash{}alpha)\_\{\textbackslash{}varphi\textasciicircum{}2(i)\}\par
\emph{Proposal}: X(\textbackslash{}alpha)\_i\par
The replacement X(\textbackslash{}alpha)\_i at Premier paragraphe, « Donc, comme … est c-accessible » identifies a particular indexed object or symbol; that identification is not independently established by the quoted local expression.\par
\par\noindent\textbf{OM-SOL-0133 -- t1 I p.167 -- SOL}\par
Même phrase, présentation de la limite inductive\par
\emph{Reported reading}: X(\textbackslash{}alpha)\textbackslash{}simeq \textbackslash{}varinjlim\_\{j\textbackslash{}in I\}X(\textbackslash{}alpha)\_j\par
\emph{Proposal}: X(\textbackslash{}beta)\textbackslash{}simeq \textbackslash{}varinjlim\_\{j\textbackslash{}in I\}X(\textbackslash{}beta)\_j\par
The replacement X(\textbackslash{}beta)\textbackslash{}simeq \textbackslash{}varinjlim\_\{j\textbackslash{}in I\}X(\textbackslash{}beta)\_j at Même phrase, présentation de la limite inductive requires a complete check of the maps, their endpoints and the composition convention; this identification is not independently verified.\par
\par\noindent\textbf{OM-SOL-0137 -- t1 I p.171 -- SOL}\par
Lemma 9.21.17 (i), « X(\ensuremath{\alpha}) soit un objet d-accessible de E »\par
\emph{Reported reading}: de E\par
\emph{Proposal}: de \$E\_\{\textbackslash{}alpha\}\$\par
The replacement de \$E\_\{\textbackslash{}alpha\}\$ at Lemma 9.21.17 (i), « X(\ensuremath{\alpha}) soit un objet d-accessible de E » identifies a particular indexed object or symbol; that identification is not independently established by the quoted local expression.\par
\par\noindent\textbf{OM-SOL-0139 -- t1 I p.173 -- SOL}\par
First sentence of the proof of Lemma 9.21.19\par
\emph{Reported reading}: on procède comme dans 9.21.3\par
\emph{Proposal}: on procède comme dans 9.21.13\par
The proposed reference on procède comme dans 9.21.13 at First sentence of the proof of Lemma 9.21.19 needs comparison with the statement actually invoked; that target has not been independently established here.\par
\par\noindent\textbf{OM-SOL-0140 -- t1 I p.174 -- SOL}\par
Conclusion immediately before 9.21.21\par
\emph{Reported reading}: \$X\textbackslash{}in\textbackslash{}ob F'\$\par
\emph{Proposal}: X\_J\textbackslash{}in \textbackslash{}ob F'\par
The proposed reading X\_J\textbackslash{}in \textbackslash{}ob F' at Conclusion immediately before 9.21.21 changes a substantive mathematical assertion or bibliographic identification; no independent derivation of this exact replacement is established in this review.\par
\par\noindent\textbf{OM-SOL-0141 -- t1 I p.176 -- SOL}\par
Premier paragraphe, égalité pour C\_alpha\par
\emph{Reported reading}: \textbackslash{}Filt\textasciicircum{}C(E\_\{\textbackslash{}alpha\})\par
\emph{Proposal}: \textbackslash{}Filt\textasciicircum{}c(E\_\{\textbackslash{}alpha\})\par
The replacement \textbackslash{}Filt\textasciicircum{}c(E\_\{\textbackslash{}alpha\}) at Premier paragraphe, égalité pour C\_alpha identifies a particular indexed object or symbol; that identification is not independently established by the quoted local expression.\par
\par\noindent\textbf{OM-SOL-0142 -- t1 I p.179 -- SOL}\par
10.1, dernière phrase\par
\emph{Reported reading}: catégorie opposée à \textbackslash{}underline\{A\}\par
\emph{Proposal}: \$C\$\par
This finding intersects Orgogozo's \textbackslash{}sisi material at 10.1, dernière phrase; preserve that human entry unchanged. The proposed reading \$C\$ at 10.1, dernière phrase changes a substantive mathematical assertion or bibliographic identification; no independent derivation of this exact replacement is established in this review.\par
\par\noindent\textbf{OM-SOL-0143 -- t1 I p.180 -- SOL}\par
Formule (10.3.1), les deux indices de Hom\par
\emph{Reported reading}: \textbackslash{}Hom\_\{\textbackslash{}underline\{A\}\}(B,C)\textbackslash{}longrightarrow\textbackslash{}prod\_i\textbackslash{}Hom\_\{\textbackslash{}underline\{A\}\}(A\_i,C)\par
\emph{Proposal}: \textbackslash{}Hom\_C(B,C)\textbackslash{}longrightarrow \textbackslash{}prod\_i \textbackslash{}Hom\_C(A\_i,C)\par
The replacement \textbackslash{}Hom\_C(B,C)\textbackslash{}longrightarrow \textbackslash{}prod\_i \textbackslash{}Hom\_C(A\_i,C) at Formule (10.3.1), les deux indices de Hom identifies a particular indexed object or symbol; that identification is not independently established by the quoted local expression.\par
\par\noindent\textbf{OM-SOL-0145 -- t1 I p.181 -- SOL}\par
10.4, famille de flèches définie comme monomorphique\par
\emph{Reported reading}: \$f\_i:A\_i\textbackslash{}rightarrow B\$\par
\emph{Proposal}: \$f\_i:B\textbackslash{}rightarrow A\_i\$\par
The proposed reading \$f\_i:B\textbackslash{}rightarrow A\_i\$ at 10.4, famille de flèches définie comme monomorphique changes a substantive mathematical assertion or bibliographic identification; no independent derivation of this exact replacement is established in this review.\par
\par\noindent\textbf{OM-SOL-0148 -- t1 I p.196 -- SOL}\par
Formula (2), explanatory description of A\_\{n+1\}\par
\emph{Reported reading}: ensemble des éléments de A\_n\par
\emph{Proposal}: ensemble des éléments des éléments de \$A\_n\$.\par
The proposed reading ensemble des éléments des éléments de \$A\_n\$. at Formula (2), explanatory description of A\_\{n+1\} changes a substantive mathematical assertion or bibliographic identification; no independent derivation of this exact replacement is established in this review.\par
\par\noindent\textbf{OM-SOL-0149 -- t1 I p.211 -- SOL}\par
Lemme 3, notation de la seconde suite\par
\emph{Reported reading}: La seconde suite porte un glyphe grec arrondi, différent du f employé ensuite pour les applications.\par
\emph{Proposal}: \$(f\_n)\_\{n\textbackslash{}geq 1\}\$\par
The proposed reading \$(f\_n)\_\{n\textbackslash{}geq 1\}\$ involves a printed glyph or a missing stroke at Lemme 3, notation de la seconde suite; the textual record does not independently establish that glyph.\par
\par\noindent\textbf{OM-SOL-0151 -- t1 II p.234 -- SOL}\par
Remarque 3.3, final sentence\par
\emph{Reported reading}: les formules on retrouve la construction de [\par
\emph{Proposal}: [\textbackslash{}text\{référence à compléter\}]\par
The proposed reference [\textbackslash{}text\{référence à compléter\}] at Remarque 3.3, final sentence needs comparison with the statement actually invoked; that target has not been independently established here.\par
\par\noindent\textbf{OM-SOL-0152 -- t1 II p.242 -- SOL}\par
Après le diagramme (4.6.2), renvoi au lemme 3.1\par
\emph{Reported reading}: 3.1 1)\par
\emph{Proposal}: \textbackslash{}hyperref[II.lem3.1]\{3.1\} 2)\par
The proposed reference \textbackslash{}hyperref[II.lem3.1]\{3.1\} 2) at Après le diagramme (4.6.2), renvoi au lemme 3.1 needs comparison with the statement actually invoked; that target has not been independently established here.\par
\par\noindent\textbf{OM-SOL-0153 -- t1 II p.243 -- SOL}\par
Première phrase, renvoi au lemme 3.1\par
\emph{Reported reading}: 3.1 2)\par
\emph{Proposal}: \textbackslash{}hyperref[II.lem3.1]\{3.1\} 3)\par
The proposed reading \textbackslash{}hyperref[II.lem3.1]\{3.1\} 3) at Première phrase, renvoi au lemme 3.1 changes a substantive mathematical assertion or bibliographic identification; no independent derivation of this exact replacement is established in this review.\par
\par\noindent\textbf{OM-SOL-0154 -- t1 II p.243 -- SOL}\par
Corollaire 4.4.4, condition définissant T''\par
\emph{Reported reading}: (resp. des préfaisceaux)\par
\emph{Proposal}: (\textbackslash{}resp des préfaisceaux séparés)\par
The proposed reading (\textbackslash{}resp des préfaisceaux séparés) at Corollaire 4.4.4, condition définissant T'' changes a substantive mathematical assertion or bibliographic identification; no independent derivation of this exact replacement is established in this review.\par
\par\noindent\textbf{OM-SOL-0155 -- t1 II p.246 -- SOL}\par
Corollaire 4.6.1, preuve de 2)\par
\emph{Reported reading}: Y est un objet initial de X\par
\emph{Proposal}: \$Y\$ est un objet initial de \$C\$\par
The proposed reading \$Y\$ est un objet initial de \$C\$ at Corollaire 4.6.1, preuve de 2) changes a substantive mathematical assertion or bibliographic identification; no independent derivation of this exact replacement is established in this review.\par
\par\noindent\textbf{OM-SOL-0156 -- t1 II p.249 -- SOL}\par
Proposition 4.7, continuation of ii) \ensuremath{\Rightarrow} i), displayed application\par
\emph{Reported reading}: \textbackslash{}Hom(X,F)\textbackslash{}longrightarrow \textbackslash{}Hom(S,F)\par
\emph{Proposal}: \textbackslash{}Hom(Y,F)\textbackslash{}longrightarrow \textbackslash{}Hom(S,F)\par
The proposed reading \textbackslash{}Hom(Y,F)\textbackslash{}longrightarrow \textbackslash{}Hom(S,F) at Proposition 4.7, continuation of ii) \ensuremath{\Rightarrow} i), displayed application changes a substantive mathematical assertion or bibliographic identification; no independent derivation of this exact replacement is established in this review.\par
\par\noindent\textbf{OM-SOL-0158 -- t1 II p.250 -- SOL}\par
4.9, «Alors la catégorie C des U-faisceaux»\par
\emph{Reported reading}: Alors la catégorie \$C\$ des \$\ensuremath{𝒰}\$-faisceaux sur \$C\$\par
\emph{Proposal}: Alors la catégorie \$\textbackslash{}widetilde\{C\}\$ des \$\ensuremath{𝒰}\$-faisceaux sur \$C\$\par
The replacement Alors la catégorie \$\textbackslash{}widetilde\{C\}\$ des \$\ensuremath{𝒰}\$-faisceaux sur \$C\$ at 4.9, «Alors la catégorie C des U-faisceaux» identifies a particular indexed object or symbol; that identification is not independently established by the quoted local expression.\par
\par\noindent\textbf{OM-SOL-0159 -- t1 II p.250 -- SOL}\par
4.9, «famille génératrice dans C coïncident»\par
\emph{Reported reading}: famille génératrice dans \$C\$ coïncident\par
\emph{Proposal}: \textbackslash{}emph\{famille génératrice\} dans \$\textbackslash{}widetilde\{C\}\$ coïncident\par
The replacement \textbackslash{}emph\{famille génératrice\} dans \$\textbackslash{}widetilde\{C\}\$ coïncident at 4.9, «famille génératrice dans C coïncident» identifies a particular indexed object or symbol; that identification is not independently established by the quoted local expression.\par
\par\noindent\textbf{OM-SOL-0160 -- t1 II p.250 -- SOL}\par
Proposition 4.10, final sentence\par
\emph{Reported reading}: d'objets de \$C\$ est génératrice\par
\emph{Proposal}: \$(\textbackslash{}epsilon(X))\_\{X\textbackslash{}in \textbackslash{}ob C\}\$ d'objets de \$\textbackslash{}widetilde\{C\}\$ est génératrice.\par
The replacement \$(\textbackslash{}epsilon(X))\_\{X\textbackslash{}in \textbackslash{}ob C\}\$ d'objets de \$\textbackslash{}widetilde\{C\}\$ est génératrice. at Proposition 4.10, final sentence identifies a particular indexed object or symbol; that identification is not independently established by the quoted local expression.\par
\par\noindent\textbf{OM-SOL-0161 -- t1 II p.251 -- SOL}\par
Preuve du corollaire 4.11, « D'autre part »\par
\emph{Reported reading}: \$F,F'\textbackslash{}in\textbackslash{}ob C\$\par
\emph{Proposal}: \$F\$, \$F'\textbackslash{}in \textbackslash{}ob \textbackslash{}widetilde\{C\}\$\par
The replacement \$F\$, \$F'\textbackslash{}in \textbackslash{}ob \textbackslash{}widetilde\{C\}\$ at Preuve du corollaire 4.11, « D'autre part » identifies a particular indexed object or symbol; that identification is not independently established by the quoted local expression.\par
\par\noindent\textbf{OM-SOL-0163 -- t1 II p.252 -- SOL}\par
Preuve de la proposition 5.1, formule suivant le renvoi à 4.4\par
\emph{Reported reading}: \$u:(H\textbackslash{}times\_K X)\textbackslash{}rightarrow\textbackslash{}underline\{a\}(X)\$\par
\emph{Proposal}: \$u:\textbackslash{}underline\{a\}(H\textbackslash{}times\_K X)\textbackslash{}rightarrow \textbackslash{}underline\{a\}(X)\$\par
The replacement \$u:\textbackslash{}underline\{a\}(H\textbackslash{}times\_K X)\textbackslash{}rightarrow \textbackslash{}underline\{a\}(X)\$ at Preuve de la proposition 5.1, formule suivant le renvoi à 4.4 requires a complete check of the maps, their endpoints and the composition convention; this identification is not independently verified.\par
\par\noindent\textbf{OM-SOL-0164 -- t1 II p.258 -- SOL}\par
Remarques 6.2 (1), isomorphisme affiché\par
\emph{Reported reading}: \textbackslash{}Hom\_D(S,F(X))\textbackslash{}xrightarrow\{\textbackslash{}sim\}\textbackslash{}varprojlim\_\{C/R\}F(.)\par
\emph{Proposal}: \textbackslash{}Hom\_D(S,F(X))\textbackslash{}xrightarrow\{\textbackslash{}sim\}\textbackslash{}varprojlim\_\{C/R\}\textbackslash{}Hom\_D(S,F(.))\par
The proposed reading \textbackslash{}Hom\_D(S,F(X))\textbackslash{}xrightarrow\{\textbackslash{}sim\}\textbackslash{}varprojlim\_\{C/R\}\textbackslash{}Hom\_D(S,F(.)) at Remarques 6.2 (1), isomorphisme affiché changes a substantive mathematical assertion or bibliographic identification; no independent derivation of this exact replacement is established in this review.\par
\par\noindent\textbf{OM-SOL-0165 -- t1 II p.260 -- SOL}\par
Proposition 6.5, source de esj et indice du second Hom\par
\emph{Reported reading}: \textbackslash{}gamma-\ensuremath{𝒰}\par
\emph{Proposal}: \ensuremath{𝒰}-\textbackslash{}gamma\par
The proposed reading \ensuremath{𝒰}-\textbackslash{}gamma at Proposition 6.5, source de esj et indice du second Hom changes a substantive mathematical assertion or bibliographic identification; no independent derivation of this exact replacement is established in this review.\par
\par\noindent\textbf{OM-SOL-0166 -- t1 II p.261 -- SOL}\par
Formule (*)\par
\emph{Reported reading}: i\_a(\textbackslash{}Ker(u))\par
\emph{Proposal}: i\_A(\textbackslash{}Ker(u))\par
The replacement i\_A(\textbackslash{}Ker(u)) at Formule (*) identifies a particular indexed object or symbol; that identification is not independently established by the quoted local expression.\par
\par\noindent\textbf{OM-SOL-0167 -- t1 II p.261 -- SOL}\par
Formule (**), argument de coker à droite\par
\emph{Reported reading}: \textbackslash{}coker(i\_a(u))\par
\emph{Proposal}: \textbackslash{}coker(i\_A(u))\par
The replacement \textbackslash{}coker(i\_A(u)) at Formule (**), argument de coker à droite identifies a particular indexed object or symbol; that identification is not independently established by the quoted local expression.\par
\par\noindent\textbf{OM-SOL-0168 -- t1 II p.261 -- SOL}\par
Formule (***), argument de coim à droite\par
\emph{Reported reading}: a\_A\textbackslash{}coim(i\_a(u))\par
\emph{Proposal}: a\_A\textbackslash{}coim(i\_A(u))\par
The replacement a\_A\textbackslash{}coim(i\_A(u)) at Formule (***), argument de coim à droite identifies a particular indexed object or symbol; that identification is not independently established by the quoted local expression.\par
\par\noindent\textbf{OM-SOL-0170 -- t1 III p.266 -- SOL}\par
Proposition 1.2 ii), arrow after u\_!R\par
\emph{Reported reading}: u\_\{!\}R\textbackslash{}hookrightarrow u(X)\par
\emph{Proposal}: u\_\{!\}R\textbackslash{}rightarrow u(X)\par
The replacement u\_\{!\}R\textbackslash{}rightarrow u(X) at Proposition 1.2 ii), arrow after u\_!R requires a complete check of the maps, their endpoints and the composition convention; this identification is not independently verified.\par
\par\noindent\textbf{OM-SOL-0171 -- t1 III p.269 -- SOL}\par
Opening sentence, reference for commutativity of the upper square\par
\emph{Reported reading}: 1.4 3)\par
\emph{Proposal}: \textbackslash{}hyperrefext[I][I.prop5.4]\{5.4\} 3)\par
The proposed reference \textbackslash{}hyperrefext[I][I.prop5.4]\{5.4\} 3) at Opening sentence, reference for commutativity of the upper square needs comparison with the statement actually invoked; that target has not been independently established here.\par
\par\noindent\textbf{OM-SOL-0175 -- t1 III p.274 -- SOL}\par
Proposition 1.7, source and target of u\_s\textasciicircum{}\ensuremath{\gamma}\par
\emph{Reported reading}: u\textasciicircum{}\{\textbackslash{}gamma\}\_s:C\textasciicircum{}\{\textbackslash{}sim\}\_\{\textbackslash{}gamma\}\textbackslash{}rightarrow \{C'\}\textasciicircum{}\{\textbackslash{}sim\}\_\{\textbackslash{}gamma\}\par
\emph{Proposal}: u\textasciicircum{}\{\textbackslash{}gamma\}\_s:\{C'\}\textasciicircum{}\{\textbackslash{}sim\}\_\{\textbackslash{}gamma\}\textbackslash{}rightarrow C\textasciicircum{}\{\textbackslash{}sim\}\_\{\textbackslash{}gamma\}\par
The proposed reading u\textasciicircum{}\{\textbackslash{}gamma\}\_s:\{C'\}\textasciicircum{}\{\textbackslash{}sim\}\_\{\textbackslash{}gamma\}\textbackslash{}rightarrow C\textasciicircum{}\{\textbackslash{}sim\}\_\{\textbackslash{}gamma\} at Proposition 1.7, source and target of u\_s\textasciicircum{}\ensuremath{\gamma} changes a substantive mathematical assertion or bibliographic identification; no independent derivation of this exact replacement is established in this review.\par
\par\noindent\textbf{OM-SOL-0176 -- t1 III p.275 -- SOL}\par
Proposition 1.7, item 2), right-hand side of the isomorphism\par
\emph{Reported reading}: \textbackslash{}underline\{a\}'\_\{\textbackslash{}gamma\}u'\_\{\textbackslash{}gamma!\}\par
\emph{Proposal}: \textbackslash{}underline\{a\}'\_\{\textbackslash{}gamma\} u\_\{\textbackslash{}gamma !\}\par
The replacement \textbackslash{}underline\{a\}'\_\{\textbackslash{}gamma\} u\_\{\textbackslash{}gamma !\} at Proposition 1.7, item 2), right-hand side of the isomorphism identifies a particular indexed object or symbol; that identification is not independently established by the quoted local expression.\par
\par\noindent\textbf{OM-SOL-0177 -- t1 III p.275 -- SOL}\par
Final paragraph of the proposition, « l'adjoint à gauche de u\_s »\par
\emph{Reported reading}: u\_s\par
\emph{Proposal}: u\textasciicircum{}\{\textbackslash{}gamma\}\_s\par
The replacement u\textasciicircum{}\{\textbackslash{}gamma\}\_s at Final paragraph of the proposition, « l'adjoint à gauche de u\_s » identifies a particular indexed object or symbol; that identification is not independently established by the quoted local expression.\par
\par\noindent\textbf{OM-SOL-0178 -- t1 III p.279 -- SOL}\par
Last pullback formula in the proof of 2.2\par
\emph{Reported reading}: u\textasciicircum{}\{\textbackslash{}ast\}R\textbackslash{}times\_\{\textbackslash{}widehat\{u\}\textasciicircum{}\{\textbackslash{}ast\}S\}Y\textbackslash{}hookrightarrow Y\par
\emph{Proposal}: \textbackslash{}widehat\{u\}\textasciicircum{}\{\textbackslash{}ast\}R\textbackslash{}times\_\{\textbackslash{}widehat\{u\}\textasciicircum{}\{\textbackslash{}ast\}S\}Y\textbackslash{}hookrightarrow Y\par
The replacement \textbackslash{}widehat\{u\}\textasciicircum{}\{\textbackslash{}ast\}R\textbackslash{}times\_\{\textbackslash{}widehat\{u\}\textasciicircum{}\{\textbackslash{}ast\}S\}Y\textbackslash{}hookrightarrow Y at Last pullback formula in the proof of 2.2 identifies a particular indexed object or symbol; that identification is not independently established by the quoted local expression.\par
\par\noindent\textbf{OM-SOL-0179 -- t1 III p.283 -- SOL}\par
Suite de l'exemple 2.7, avant les trois foncteurs adjoints\par
\emph{Reported reading}: Le foncteur i est adjoint à droite au foncteur Φ.\par
\emph{Proposal}: Le foncteur \$i\textasciicircum{}\{-1\}\$ est adjoint à droite au foncteur \$\textbackslash{}Phi\$.\par
The replacement Le foncteur \$i\textasciicircum{}\{-1\}\$ est adjoint à droite au foncteur \$\textbackslash{}Phi\$. at Suite de l'exemple 2.7, avant les trois foncteurs adjoints identifies a particular indexed object or symbol; that identification is not independently established by the quoted local expression.\par
\par\noindent\textbf{OM-SOL-0181 -- t1 III p.287 -- SOL}\par
Proof of Proposition 3.7, final description of the defining family of presheaves\par
\emph{Reported reading}: \$F\textbackslash{}in \textbackslash{}ob C\_i\$\par
\emph{Proposal}: \$F\textbackslash{}in \textbackslash{}ob \textbackslash{}widetilde\{C\_i\}\$\par
The replacement \$F\textbackslash{}in \textbackslash{}ob \textbackslash{}widetilde\{C\_i\}\$ at Proof of Proposition 3.7, final description of the defining family of presheaves identifies a particular indexed object or symbol; that identification is not independently established by the quoted local expression.\par
\par\noindent\textbf{OM-SOL-0183 -- t1 III p.289 -- SOL}\par
First step, sentence introducing the two morphisms p and q\par
\emph{Reported reading}: d'un objet \$Z\$ de \$C'\$ dans \$H'\$\par
\emph{Proposal}: d'un objet \$Z\$ de \$C'\$ dans \$u\_\{!\}u\textasciicircum{}\{\textbackslash{}ast\}H\$\par
The proposed reading d'un objet \$Z\$ de \$C'\$ dans \$u\_\{!\}u\textasciicircum{}\{\textbackslash{}ast\}H\$ at First step, sentence introducing the two morphisms p and q changes a substantive mathematical assertion or bibliographic identification; no independent derivation of this exact replacement is established in this review.\par
\par\noindent\textbf{OM-SOL-0184 -- t1 III p.290 -- SOL}\par
Opening paragraph, covering sieve obtained by base change\par
\emph{Reported reading}: \$R\textbackslash{}times\_\{u(Y)\}u(X)\$ est un crible couvrant \$X\$\par
\emph{Proposal}: \$R\textbackslash{}times\_\{u(Y)\}u(X)\$ est un crible couvrant \$u(X)\$\par
The proposed reading \$R\textbackslash{}times\_\{u(Y)\}u(X)\$ est un crible couvrant \$u(X)\$ at Opening paragraph, covering sieve obtained by base change changes a substantive mathematical assertion or bibliographic identification; no independent derivation of this exact replacement is established in this review.\par
\par\noindent\textbf{OM-SOL-0186 -- t1 III p.292 -- SOL}\par
4.4, phrase « De plus »\par
\emph{Reported reading}: \$i\textasciicircum{}\{\textbackslash{}ast\}:\textbackslash{}widetilde\{C\}'\textbackslash{}rightarrow \textbackslash{}widetilde\{G\}\$\par
\emph{Proposal}: \$i\textasciicircum{}\{\textbackslash{}ast\}:\textbackslash{}widetilde\{C\}\textbackslash{}rightarrow \textbackslash{}widetilde\{G\}\$\par
The replacement \$i\textasciicircum{}\{\textbackslash{}ast\}:\textbackslash{}widetilde\{C\}\textbackslash{}rightarrow \textbackslash{}widetilde\{G\}\$ at 4.4, phrase « De plus » identifies a particular indexed object or symbol; that identification is not independently established by the quoted local expression.\par
\par\noindent\textbf{OM-SOL-0187 -- t1 III p.294 -- SOL}\par
5.3, référence après « trois foncteurs adjoints »\par
\emph{Reported reading}: (4.3.2.)\par
\emph{Proposal}: (I\textasciitilde{}4.3.2)\par
The proposed reference (I\textasciitilde{}4.3.2) at 5.3, référence après « trois foncteurs adjoints » needs comparison with the statement actually invoked; that target has not been independently established here.\par
\par\noindent\textbf{OM-SOL-0188 -- t1 IV p.303 -- SOL}\par
Théorème 1.2, condition i′)\par
\emph{Reported reading}: les limites projectives soient représentables dans C\par
\emph{Proposal}: limites projectives finies soient représentables dans \$C\$\par
The proposed reading limites projectives finies soient représentables dans \$C\$ at Théorème 1.2, condition i′) changes a substantive mathematical assertion or bibliographic identification; no independent derivation of this exact replacement is established in this review.\par
\par\noindent\textbf{OM-SOL-0190 -- t1 IV p.307 -- SOL}\par
Remarque 1.3, renvoi final\par
\emph{Reported reading}: (VIII 1)\par
\emph{Proposal}: (VII\textasciitilde{}1)\par
This finding intersects Orgogozo's \textbackslash{}sisi material at Remarque 1.3, renvoi final; preserve that human entry unchanged. The proposed reference (VII\textasciitilde{}1) at Remarque 1.3, renvoi final needs comparison with the statement actually invoked; that target has not been independently established here.\par
\par\noindent\textbf{OM-SOL-0191 -- t1 IV p.308 -- SOL}\par
Topologie induite sur C\par
\emph{Reported reading}: (III 4.5)\par
\emph{Proposal}: (III\textasciitilde{}3.5)\par
This finding intersects Orgogozo's \textbackslash{}sisi material at Topologie induite sur C; preserve that human entry unchanged. The proposed reference (III\textasciitilde{}3.5) at Topologie induite sur C needs comparison with the statement actually invoked; that target has not been independently established here.\par
\par\noindent\textbf{OM-SOL-0192 -- t1 IV p.308 -- SOL}\par
Définition du foncteur canonique J\_E, avant le diagramme (*)\par
\emph{Reported reading}: \$J\_E:E\textbackslash{}rightarrow E\_\{\textbackslash{}underline\{V\}\}\$\par
\emph{Proposal}: \$J\_E:E\textbackslash{}rightarrow E\textasciicircum{}\{\textbackslash{}widetilde\{\textbackslash{}phantom\{,\}\}\}\_\{\ensuremath{𝒱}\}\$\par
The replacement \$J\_E:E\textbackslash{}rightarrow E\textasciicircum{}\{\textbackslash{}widetilde\{\textbackslash{}phantom\{,\}\}\}\_\{\ensuremath{𝒱}\}\$ at Définition du foncteur canonique J\_E, avant le diagramme (*) identifies a particular indexed object or symbol; that identification is not independently established by the quoted local expression.\par
\par\noindent\textbf{OM-SOL-0193 -- t1 IV p.310 -- SOL}\par
Démonstration du corollaire 1.7, suffisance\par
\emph{Reported reading}: III 2.6\par
\emph{Proposal}: \textbackslash{}hyperrefext[III][III.prop1.6]\{1.6\}\par
This finding intersects Orgogozo's \textbackslash{}sisi material at Démonstration du corollaire 1.7, suffisance; preserve that human entry unchanged. The proposed reference \textbackslash{}hyperrefext[III][III.prop1.6]\{1.6\} at Démonstration du corollaire 1.7, suffisance needs comparison with the statement actually invoked; that target has not been independently established here.\par
\par\noindent\textbf{OM-SOL-0194 -- t1 IV p.314 -- SOL}\par
2.3, assertion about subobjects of the final object of B\_G\par
\emph{Reported reading}: le seul sous-objet de l'objet final e du topos B\_G est e ou le faisceau vide \ensuremath{\varnothing}\par
\emph{Proposal}: les seuls sous-objets de l'objet final \$e\$ du topos \$B\_G\$ sont \$e\$ et le faisceau vide \$\ensuremath{\varnothing}\$\par
The proposed reading les seuls sous-objets de l'objet final \$e\$ du topos \$B\_G\$ sont \$e\$ et le faisceau vide \$\ensuremath{\varnothing}\$ at 2.3, assertion about subobjects of the final object of B\_G changes a substantive mathematical assertion or bibliographic identification; no independent derivation of this exact replacement is established in this review.\par
\par\noindent\textbf{OM-SOL-0195 -- t1 IV p.316 -- SOL}\par
First paragraph, reference after “la prétopologie”\par
\emph{Reported reading}: (I 1.3)\par
\emph{Proposal}: II\textasciitilde{}1.3\par
This finding intersects Orgogozo's \textbackslash{}sisi material at First paragraph, reference after “la prétopologie”; preserve that human entry unchanged. The proposed reading II\textasciitilde{}1.3 at First paragraph, reference after “la prétopologie” changes a substantive mathematical assertion or bibliographic identification; no independent derivation of this exact replacement is established in this review.\par
\par\noindent\textbf{OM-SOL-0198 -- t1 IV p.320 -- SOL}\par
§2.7.1, same sentence, “une opération de G”\par
\emph{Reported reading}: une opération de \$G\$ sur \$E\$\par
\emph{Proposal}: une opération de \$\ensuremath{𝒢}\$ sur \$E\$\par
The replacement une opération de \$\ensuremath{𝒢}\$ sur \$E\$ at §2.7.1, same sentence, “une opération de G” identifies a particular indexed object or symbol; that identification is not independently established by the quoted local expression.\par
\par\noindent\textbf{OM-SOL-0199 -- t1 IV p.321 -- SOL}\par
Exercise 2.7.5, footnote\par
\emph{Reported reading}: cf. 8.7 1).\par
\emph{Proposal}: \textbackslash{}cf \textbackslash{}hyperref[IV.exer8.7]\{8.7\} l).\par
This finding intersects Orgogozo's \textbackslash{}sisi material at Exercise 2.7.5, footnote; preserve that human entry unchanged. The proposed reference \textbackslash{}cf \textbackslash{}hyperref[IV.exer8.7]\{8.7\} l). at Exercise 2.7.5, footnote needs comparison with the statement actually invoked; that target has not been independently established here.\par
\par\noindent\textbf{OM-SOL-0203 -- t1 IV p.329 -- SOL}\par
Last paragraph, immediately after the reference 3.1\par
\emph{Reported reading}: garder la notion "fine" 3.1 de la notion de morphisme entre morphismes de topos\par
\emph{Proposal}: garder la notion « fine » \textbackslash{}hyperref[IV.def3.1]\{3.1\} et la notion de morphisme entre morphismes de topos\par
The replacement garder la notion « fine » \textbackslash{}hyperref[IV.def3.1]\{3.1\} et la notion de morphisme entre morphismes de topos at Last paragraph, immediately after the reference 3.1 requires a complete check of the maps, their endpoints and the composition convention; this identification is not independently verified.\par
\par\noindent\textbf{OM-SOL-0204 -- t1 IV p.331 -- SOL}\par
3.4, identities corresponding to the composites uv and vu\par
\emph{Reported reading}: uv et vu ... respectivement au morphisme identique de E et de E'\par
\emph{Proposal}: \$vu\$ et \$uv\$ soient isomorphes respectivement au morphisme identique de \$E\$ et de \$E'\$\par
The replacement \$vu\$ et \$uv\$ soient isomorphes respectivement au morphisme identique de \$E\$ et de \$E'\$ at 3.4, identities corresponding to the composites uv and vu requires a complete check of the maps, their endpoints and the composition convention; this identification is not independently verified.\par
\par\noindent\textbf{OM-SOL-0205 -- t1 IV p.333 -- SOL}\par
4.1.1, first description of the adjoint of f\_*\par
\emph{Reported reading}: admettant un adjoint à droite f\textasciicircum{}*\par
\emph{Proposal}: admettant un adjoint à gauche \$f\textasciicircum{}\{\textbackslash{}ast\}\$\par
The proposed reading admettant un adjoint à gauche \$f\textasciicircum{}\{\textbackslash{}ast\}\$ at 4.1.1, first description of the adjoint of f\_* changes a substantive mathematical assertion or bibliographic identification; no independent derivation of this exact replacement is established in this review.\par
\par\noindent\textbf{OM-SOL-0206 -- t1 IV p.333 -- SOL}\par
4.1.1, sentence about choosing the adjoint\par
\emph{Reported reading}: du choix de l'adjoint à droite f\textasciicircum{}* de f\_*\par
\emph{Proposal}: choix de l'adjoint à gauche \$f\textasciicircum{}\{\textbackslash{}ast\}\$ de \$f\_\{\textbackslash{}ast\}\$\par
The proposed reading choix de l'adjoint à gauche \$f\textasciicircum{}\{\textbackslash{}ast\}\$ de \$f\_\{\textbackslash{}ast\}\$ at 4.1.1, sentence about choosing the adjoint changes a substantive mathematical assertion or bibliographic identification; no independent derivation of this exact replacement is established in this review.\par
\par\noindent\textbf{OM-SOL-0207 -- t1 IV p.339 -- SOL}\par
4.3, formule donnant l'unique morphisme vers le topos final\par
\emph{Reported reading}: \$f:E\textbackslash{}longrightarrow P\$, sans numéro\par
\emph{Proposal}: \textbackslash{}begin\{equation*\} f:E\textbackslash{}rightarrow P\textbackslash{}tag\{4.3.1\} \textbackslash{}end\{equation*\}\par
The proposed reference \textbackslash{}begin\{equation*\} f:E\textbackslash{}rightarrow P\textbackslash{}tag\{4.3.1\} \textbackslash{}end\{equation*\} at 4.3, formule donnant l'unique morphisme vers le topos final needs comparison with the statement actually invoked; that target has not been independently established here.\par
\par\noindent\textbf{OM-SOL-0208 -- t1 IV p.340 -- SOL}\par
Justification de la commutation du foncteur \$I\textbackslash{}mapsto T\textbackslash{}times I\$ aux limites inductives\par
\emph{Reported reading}: adjoint à droite\par
\emph{Proposal}: adjoint à gauche\par
This finding intersects Orgogozo's \textbackslash{}sisi material at Justification de la commutation du foncteur \$I\textbackslash{}mapsto T\textbackslash{}times I\$ aux limites inductives; preserve that human entry unchanged. The proposed reading adjoint à gauche at Justification de la commutation du foncteur \$I\textbackslash{}mapsto T\textbackslash{}times I\$ aux limites inductives changes a substantive mathematical assertion or bibliographic identification; no independent derivation of this exact replacement is established in this review.\par
\par\noindent\textbf{OM-SOL-0209 -- t1 IV p.344 -- SOL}\par
4.5.1, explanation of the direct-image formula\par
\emph{Reported reading}: où \$X\$ est un objet de \$X\$\par
\emph{Proposal}: où \$X\$ est un objet de \$E\$\par
This finding intersects Orgogozo's \textbackslash{}sisi material at 4.5.1, explanation of the direct-image formula; preserve that human entry unchanged. The proposed reading où \$X\$ est un objet de \$E\$ at 4.5.1, explanation of the direct-image formula changes a substantive mathematical assertion or bibliographic identification; no independent derivation of this exact replacement is established in this review.\par
\par\noindent\textbf{OM-SOL-0210 -- t1 IV p.344 -- SOL}\par
4.5.1, definition of the G-action on H\_s\par
\emph{Reported reading}: déduites de \$F\$\par
\emph{Proposal}: déduites de \$f\$\par
This finding intersects Orgogozo's \textbackslash{}sisi material at 4.5.1, definition of the G-action on H\_s; preserve that human entry unchanged. The replacement déduites de \$f\$ at 4.5.1, definition of the G-action on H\_s requires a complete check of the maps, their endpoints and the composition convention; this identification is not independently verified.\par
\par\noindent\textbf{OM-SOL-0211 -- t1 IV p.344 -- SOL}\par
4.5.1, arbitrary limits in the justification for f\_!\par
\emph{Reported reading}: aux \$\textbackslash{}varinjlim\$ quelconques\par
\emph{Proposal}: aux \$\textbackslash{}varprojlim\$ quelconques\par
This finding intersects Orgogozo's \textbackslash{}sisi material at 4.5.1, arbitrary limits in the justification for f\_!; preserve that human entry unchanged. The replacement aux \$\textbackslash{}varprojlim\$ quelconques at 4.5.1, arbitrary limits in the justification for f\_! requires a complete check of the maps, their endpoints and the composition convention; this identification is not independently verified.\par
\par\noindent\textbf{OM-SOL-0212 -- t1 IV p.344 -- SOL}\par
4.5.1, finite limits in the same parenthesis\par
\emph{Reported reading}: aux \$\textbackslash{}varinjlim\$ finies\par
\emph{Proposal}: aux \$\textbackslash{}varprojlim\$ finies\par
This finding intersects Orgogozo's \textbackslash{}sisi material at 4.5.1, finite limits in the same parenthesis; preserve that human entry unchanged. The replacement aux \$\textbackslash{}varprojlim\$ finies at 4.5.1, finite limits in the same parenthesis requires a complete check of the maps, their endpoints and the composition convention; this identification is not independently verified.\par
\par\noindent\textbf{OM-SOL-0213 -- t1 IV p.350 -- SOL}\par
4.7, action formula following « ou f\_* »\par
\emph{Reported reading}: \textbackslash{}widehat\{f\}\_\{\textbackslash{}ast\}\textbackslash{}text\{ ou \}f\_\{\textbackslash{}ast\}:F\textbackslash{}mapsto F\textbackslash{}circ f\textbackslash{}text\{ de \}\textbackslash{}widehat\{C\}\textbackslash{}text\{ dans \}\textbackslash{}widehat\{C\}'\par
\emph{Proposal}: le foncteur \$\textbackslash{}widehat\{f\}\_\{\textbackslash{}ast\}\$ ou \$f\_\{\textbackslash{}ast\}\$ de \$\textbackslash{}widehat\{C\}\$ dans \$\textbackslash{}widehat\{C\}'\$\par
The proposed reading le foncteur \$\textbackslash{}widehat\{f\}\_\{\textbackslash{}ast\}\$ ou \$f\_\{\textbackslash{}ast\}\$ de \$\textbackslash{}widehat\{C\}\$ dans \$\textbackslash{}widehat\{C\}'\$ at 4.7, action formula following « ou f\_* » changes a substantive mathematical assertion or bibliographic identification; no independent derivation of this exact replacement is established in this review.\par
\par\noindent\textbf{OM-SOL-0215 -- t1 IV p.354 -- SOL}\par
4.9.1, phrase « On fera attention qu'en général… »\par
\emph{Reported reading}: \$f\textasciicircum{}s\$ n'est pas exact à droite\par
\emph{Proposal}: \$f\textasciicircum{}s\$ n'est pas exact à gauche\par
This finding intersects Orgogozo's \textbackslash{}sisi material at 4.9.1, phrase « On fera attention qu'en général… »; preserve that human entry unchanged. The proposed reading \$f\textasciicircum{}s\$ n'est pas exact à gauche at 4.9.1, phrase « On fera attention qu'en général… » changes a substantive mathematical assertion or bibliographic identification; no independent derivation of this exact replacement is established in this review.\par
\par\noindent\textbf{OM-SOL-0216 -- t1 IV p.356 -- SOL}\par
4.9.3, direction of the morphisms of sites\par
\emph{Reported reading}: Morsite(C',C) de C dans C'\par
\emph{Proposal}: \$\textbackslash{}Mmorsite(C',C)\$ de \$C'\$ dans \$C\$\par
The replacement \$\textbackslash{}Mmorsite(C',C)\$ de \$C'\$ dans \$C\$ at 4.9.3, direction of the morphisms of sites requires a complete check of the maps, their endpoints and the composition convention; this identification is not independently verified.\par
\par\noindent\textbf{OM-SOL-0217 -- t1 IV p.360 -- SOL}\par
Opening paragraph, morphism before (4.10.3.1)\par
\emph{Reported reading}: \$u\_\{\textbackslash{}ast\}(\textbackslash{}underline\{F\}\_\{X''\})\textbackslash{}longrightarrow \textbackslash{}underline\{F\}\_\{X'\}\$\par
\emph{Proposal}: \$u\_\{\textbackslash{}ast\}(\textbackslash{}underline\{F\}\_\{X''\})\textbackslash{}longleftarrow \textbackslash{}underline\{F\}\_\{X'\}\$\par
The replacement \$u\_\{\textbackslash{}ast\}(\textbackslash{}underline\{F\}\_\{X''\})\textbackslash{}longleftarrow \textbackslash{}underline\{F\}\_\{X'\}\$ at Opening paragraph, morphism before (4.10.3.1) requires a complete check of the maps, their endpoints and the composition convention; this identification is not independently verified.\par
\par\noindent\textbf{OM-SOL-0218 -- t1 IV p.361 -- SOL}\par
4.10.4, reference after « morphismes de transitivité »\par
\emph{Reported reading}: 4.11.3\par
\emph{Proposal}: \textbackslash{}eqhyperref[IV.eq4.10.3.1]\{4.10.3.1\}\par
This finding intersects Orgogozo's \textbackslash{}sisi material at 4.10.4, reference after « morphismes de transitivité »; preserve that human entry unchanged. The proposed reference \textbackslash{}eqhyperref[IV.eq4.10.3.1]\{4.10.3.1\} at 4.10.4, reference after « morphismes de transitivité » needs comparison with the statement actually invoked; that target has not been independently established here.\par
\par\noindent\textbf{OM-SOL-0219 -- t1 IV p.362 -- SOL}\par
4.10.5 c), reference immediately after \$R\textasciicircum{}i g\_\{\textbackslash{}ast\}\$\par
\emph{Reported reading}: (V )\par
\emph{Proposal}: \textbackslash{}hyperrefext[V][V.sec5]\{5\}\par
This finding intersects Orgogozo's \textbackslash{}sisi material at 4.10.5 c), reference immediately after \$R\textasciicircum{}i g\_\{\textbackslash{}ast\}\$; preserve that human entry unchanged. The proposed reference \textbackslash{}hyperrefext[V][V.sec5]\{5\} at 4.10.5 c), reference immediately after \$R\textasciicircum{}i g\_\{\textbackslash{}ast\}\$ needs comparison with the statement actually invoked; that target has not been independently established here.\par
\par\noindent\textbf{OM-SOL-0220 -- t1 IV p.362 -- SOL}\par
4.10.5 c), reference immediately after « donc »\par
\emph{Reported reading}: (V )\par
\emph{Proposal}: \textbackslash{}hyperrefext[V][V.prop5.4]\{5.4\}\par
This finding intersects Orgogozo's \textbackslash{}sisi material at 4.10.5 c), reference immediately after « donc »; preserve that human entry unchanged. The proposed reference \textbackslash{}hyperrefext[V][V.prop5.4]\{5.4\} at 4.10.5 c), reference immediately after « donc » needs comparison with the statement actually invoked; that target has not been independently established here.\par
\par\noindent\textbf{OM-SOL-0221 -- t1 IV p.363 -- SOL}\par
Exercice 4.10.6 a), reference after « quarrables »\par
\emph{Reported reading}: (I 10.3)\par
\emph{Proposal}: \textbackslash{}hyperrefext[I][I.sec10.7]\{10.7\}\par
This finding intersects Orgogozo's \textbackslash{}sisi material at Exercice 4.10.6 a), reference after « quarrables »; preserve that human entry unchanged. The proposed reference \textbackslash{}hyperrefext[I][I.sec10.7]\{10.7\} at Exercice 4.10.6 a), reference after « quarrables » needs comparison with the statement actually invoked; that target has not been independently established here.\par
\par\noindent\textbf{OM-SOL-0222 -- t1 IV p.368 -- SOL}\par
5.5, après (5.5.1)\par
\emph{Reported reading}: Appliquant la construction de 5.1\par
\emph{Proposal}: construction de \textbackslash{}hyperref[IV.subsec5.2]\{5.2\}\par
The proposed reference construction de \textbackslash{}hyperref[IV.subsec5.2]\{5.2\} at 5.5, après (5.5.1) needs comparison with the statement actually invoked; that target has not been independently established here.\par
\par\noindent\textbf{OM-SOL-0223 -- t1 IV p.370 -- SOL}\par
5.7, dernière phrase\par
\emph{Reported reading}: réserver dans 5.1 le terme de « morphisme d'inclusion »\par
\emph{Proposal}: dans \textbackslash{}hyperref[IV.subsec5.2]\{5.2\}\par
The proposed reference dans \textbackslash{}hyperref[IV.subsec5.2]\{5.2\} at 5.7, dernière phrase needs comparison with the statement actually invoked; that target has not been independently established here.\par
\par\noindent\textbf{OM-SOL-0224 -- t1 IV p.375 -- SOL}\par
Exercice 5.9 b), cible de l'équivalence\par
\emph{Reported reading}: \textbackslash{}Ttors(E,q\textasciicircum{}\{\textbackslash{}ast\}(G))\par
\emph{Proposal}: \textbackslash{}Ttors(E',q\textasciicircum{}\{\textbackslash{}ast\}(G))\par
The proposed reading \textbackslash{}Ttors(E',q\textasciicircum{}\{\textbackslash{}ast\}(G)) at Exercice 5.9 b), cible de l'équivalence changes a substantive mathematical assertion or bibliographic identification; no independent derivation of this exact replacement is established in this review.\par
\par\noindent\textbf{OM-SOL-0225 -- t1 IV p.375 -- SOL}\par
Exercice 5.9 b), dernière phrase\par
\emph{Reported reading}: des \$q\textasciicircum{}\{\textbackslash{}ast\}(G)\$-torseurs (à droite) sur \$E\$\par
\emph{Proposal}: \$q\textasciicircum{}\{\textbackslash{}ast\}(G)\$-torseurs (à droite) sur \$E'\$.\par
The proposed reading \$q\textasciicircum{}\{\textbackslash{}ast\}(G)\$-torseurs (à droite) sur \$E'\$. at Exercice 5.9 b), dernière phrase changes a substantive mathematical assertion or bibliographic identification; no independent derivation of this exact replacement is established in this review.\par
\par\noindent\textbf{OM-SOL-0227 -- t1 IV p.378 -- SOL}\par
Proposition 5.11, direction of \ensuremath{\beta}\par
\emph{Reported reading}: \textbackslash{}beta:j\_X\textbackslash{}circ g\_1\textbackslash{}rightarrow f\textbackslash{}circ g\_2\par
\emph{Proposal}: \textbackslash{}beta:f\textbackslash{}circ g\_2\textbackslash{}rightarrow j\_X\textbackslash{}circ g\_1\par
The replacement \textbackslash{}beta:f\textbackslash{}circ g\_2\textbackslash{}rightarrow j\_X\textbackslash{}circ g\_1 at Proposition 5.11, direction of \ensuremath{\beta} requires a complete check of the maps, their endpoints and the composition convention; this identification is not independently verified.\par
\par\noindent\textbf{OM-SOL-0228 -- t1 IV p.389 -- SOL}\par
6.4.2, dernière phrase\par
\emph{Reported reading}: Enfin, pour un intéressant théorème d'existence (dû à P. Deligne) de suffisamment de points, couvrant tous les cas rencontrés en géométrie algébrique.\par
\emph{Proposal}: Enfin, pour un intéressant théorème d'existence (dû à P. Deligne) de suffisamment de points, couvrant tous les cas rencontrés en géométrie algébrique, voir \textbackslash{}hyperrefext[VI][VI.sec9]\{9\}.\par
This finding intersects Orgogozo's \textbackslash{}sisi material at 6.4.2, dernière phrase; preserve that human entry unchanged. The proposed reading Enfin, pour un intéressant théorème d'existence (dû à P. Deligne) de suffisamment de points, couvrant tous les cas rencontrés en géométrie algébrique, voir \textbackslash{}hyperrefext[VI][VI.sec9]\{9\}. involves a printed glyph or a missing stroke at 6.4.2, dernière phrase; the textual record does not independently establish that glyph.\par
\par\noindent\textbf{OM-SOL-0230 -- t1 IV p.395 -- SOL}\par
Remarque 6.7.4, ensemble indexant les points au-dessus des points donnés\par
\emph{Reported reading}: l'ensemble somme des \$E\_\{p\_\{\textbackslash{}alpha\}\}\$\par
\emph{Proposal}: l'ensemble somme des \$X\_\{p\_\{\textbackslash{}alpha\}\}\$\par
The replacement l'ensemble somme des \$X\_\{p\_\{\textbackslash{}alpha\}\}\$ at Remarque 6.7.4, ensemble indexant les points au-dessus des points donnés identifies a particular indexed object or symbol; that identification is not independently established by the quoted local expression.\par
\par\noindent\textbf{OM-SOL-0231 -- t1 IV p.395 -- SOL}\par
6.8, catégorie ambiante dans la parenthèse\par
\emph{Reported reading}: \$(E\textasciicircum{}\textbackslash{}circ)\textasciicircum{}\{\textbackslash{}widehat\{\textbackslash{}phantom\{,\}\}\}\_\{/\textbackslash{}varphi\_x\}\$\par
\emph{Proposal}: \$(E\textasciicircum{}\textbackslash{}circ)\textasciicircum{}\{\textbackslash{}widehat\{\textbackslash{}phantom\{,\}\}\}\_\{/\textbackslash{}varphi\_p\}\$\par
The replacement \$(E\textasciicircum{}\textbackslash{}circ)\textasciicircum{}\{\textbackslash{}widehat\{\textbackslash{}phantom\{,\}\}\}\_\{/\textbackslash{}varphi\_p\}\$ at 6.8, catégorie ambiante dans la parenthèse identifies a particular indexed object or symbol; that identification is not independently established by the quoted local expression.\par
\par\noindent\textbf{OM-SOL-0232 -- t1 IV p.396 -- SOL}\par
First paragraph, reference after « y commute »\par
\emph{Reported reading}: (I )\par
\emph{Proposal}: I\textasciitilde{}8.10.10\par
This finding intersects Orgogozo's \textbackslash{}sisi material at First paragraph, reference after « y commute »; preserve that human entry unchanged. The proposed reference I\textasciitilde{}8.10.10 at First paragraph, reference after « y commute » needs comparison with the statement actually invoked; that target has not been independently established here.\par
\par\noindent\textbf{OM-SOL-0233 -- t1 IV p.398 -- SOL}\par
End of the proof of (6.8.4), enlargement of universes\par
\emph{Reported reading}: formation du préfaisceau associé\par
\emph{Proposal}: formation du faisceau associé\par
The distinction between the associated presheaf and sheaf functor requires the definition of the functor a in this argument.\par
\par\noindent\textbf{OM-SOL-0234 -- t1 IV p.398 -- SOL}\par
6.8.6, after « topologie chaotique »\par
\emph{Reported reading}: \$C=\allowbreak \textbackslash{}widehat\{C\}\$\par
\emph{Proposal}: \$C\textasciicircum{}\{\textbackslash{}sim\}=\allowbreak \textbackslash{}widehat\{C\}\$\par
The replacement \$C\textasciicircum{}\{\textbackslash{}sim\}=\allowbreak \textbackslash{}widehat\{C\}\$ at 6.8.6, after « topologie chaotique » identifies a particular indexed object or symbol; that identification is not independently established by the quoted local expression.\par
\par\noindent\textbf{OM-SOL-0235 -- t1 IV p.401 -- SOL}\par
c), indice du pro-objet définissant le foncteur fibre\par
\emph{Reported reading}: (X\_i)\_\{i\textbackslash{}in I'\}\par
\emph{Proposal}: (X\_i)\_\{i\textbackslash{}in P(I')\}\par
The replacement (X\_i)\_\{i\textbackslash{}in P(I')\} at c), indice du pro-objet définissant le foncteur fibre identifies a particular indexed object or symbol; that identification is not independently established by the quoted local expression.\par
\par\noindent\textbf{OM-SOL-0239 -- t1 IV p.411 -- SOL}\par
7.3, premier paragraphe, référence après les foncteurs de limite inductive\par
\emph{Reported reading}: (I et )\par
\emph{Proposal}: \textbackslash{}hyperrefext[I][I.prop2.8]\{2.8\}\par
This finding intersects Orgogozo's \textbackslash{}sisi material at 7.3, premier paragraphe, référence après les foncteurs de limite inductive; preserve that human entry unchanged. The proposed reference \textbackslash{}hyperrefext[I][I.prop2.8]\{2.8\} at 7.3, premier paragraphe, référence après les foncteurs de limite inductive needs comparison with the statement actually invoked; that target has not been independently established here.\par
\par\noindent\textbf{OM-SOL-0241 -- t1 IV p.411 -- SOL}\par
7.3, premier paragraphe, numéro de l'exercice cité\par
\emph{Reported reading}: 7.1.10\par
\emph{Proposal}: \textbackslash{}hyperref[IV.exer7.1.11]\{7.1.11\}\par
This finding intersects Orgogozo's \textbackslash{}sisi material at 7.3, premier paragraphe, numéro de l'exercice cité; preserve that human entry unchanged. The proposed reference \textbackslash{}hyperref[IV.exer7.1.11]\{7.1.11\} at 7.3, premier paragraphe, numéro de l'exercice cité needs comparison with the statement actually invoked; that target has not been independently established here.\par
\par\noindent\textbf{OM-SOL-0242 -- t1 IV p.412 -- SOL}\par
7.4, définition des familles couvrantes\par
\emph{Reported reading}: d'éléments \$E\_i\$ de \$E\$ majorés par \$E\$\par
\emph{Proposal}: d'éléments \$E\_i\$ de \$U\$ majorés par \$E\$\par
The proposed reading d'éléments \$E\_i\$ de \$U\$ majorés par \$E\$ at 7.4, définition des familles couvrantes changes a substantive mathematical assertion or bibliographic identification; no independent derivation of this exact replacement is established in this review.\par
\par\noindent\textbf{OM-SOL-0244 -- t1 IV p.418 -- SOL}\par
7.8 b), dernier morphisme d'ensembles ordonnés\par
\emph{Reported reading}: \$\textbackslash{}Oo\textbackslash{}longrightarrow\textbackslash{}Oo'\$\par
\emph{Proposal}: \$\textbackslash{}Oo'\textbackslash{}rightarrow\textbackslash{}Oo\$\par
The replacement \$\textbackslash{}Oo'\textbackslash{}rightarrow\textbackslash{}Oo\$ at 7.8 b), dernier morphisme d'ensembles ordonnés requires a complete check of the maps, their endpoints and the composition convention; this identification is not independently verified.\par
\par\noindent\textbf{OM-SOL-0246 -- t1 IV p.419 -- SOL}\par
7.8 e), premier morphisme de topos affiché\par
\emph{Reported reading}: \$\textbackslash{}Ouv(E)\textasciicircum{}\{\textbackslash{}sim\}\textbackslash{}longrightarrow\textbackslash{}Top(X')\$\par
\emph{Proposal}: \$\textbackslash{}Ouv(E)\textasciicircum{}\{\textbackslash{}sim\}\textbackslash{}longleftarrow\textbackslash{}Top(X')\$\par
The replacement \$\textbackslash{}Ouv(E)\textasciicircum{}\{\textbackslash{}sim\}\textbackslash{}longleftarrow\textbackslash{}Top(X')\$ at 7.8 e), premier morphisme de topos affiché requires a complete check of the maps, their endpoints and the composition convention; this identification is not independently verified.\par
\par\noindent\textbf{OM-SOL-0247 -- t1 IV p.422 -- SOL}\par
8.4.3, renvoi après \$B\_G\$\par
\emph{Reported reading}: (2.4)\par
\emph{Proposal}: \textbackslash{}hyperref[IV.subsec2.6]\{2.6\}\par
This finding intersects Orgogozo's \textbackslash{}sisi material at 8.4.3, renvoi après \$B\_G\$; preserve that human entry unchanged. The proposed reference \textbackslash{}hyperref[IV.subsec2.6]\{2.6\} at 8.4.3, renvoi après \$B\_G\$ needs comparison with the statement actually invoked; that target has not been independently established here.\par
\par\noindent\textbf{OM-SOL-0248 -- t1 IV p.429 -- SOL}\par
Exercice 8.7 j), référence avant « est bijectif »\par
\emph{Reported reading}: (8.7.5)\par
\emph{Proposal}: \textbackslash{}eqhyperref[IV.eq8.7.6]\{8.7.6\}\par
The proposed reference \textbackslash{}eqhyperref[IV.eq8.7.6]\{8.7.6\} at Exercice 8.7 j), référence avant « est bijectif » needs comparison with the statement actually invoked; that target has not been independently established here.\par
\par\noindent\textbf{OM-SOL-0249 -- t1 IV p.432 -- SOL}\par
Remarque 9.1.2 d), « Avec les notations de… »\par
\emph{Reported reading}: « Avec les notations de b) »\par
\emph{Proposal}: Avec les notations de c)\par
This finding intersects Orgogozo's \textbackslash{}sisi material at Remarque 9.1.2 d), « Avec les notations de… »; preserve that human entry unchanged. The proposed reference Avec les notations de c) at Remarque 9.1.2 d), « Avec les notations de… » needs comparison with the statement actually invoked; that target has not been independently established here.\par
\par\noindent\textbf{OM-SOL-0250 -- t1 IV p.435 -- SOL}\par
Exercice 9.1.7 a), dernière ligne, après Ob\par
\emph{Reported reading}: \$X\textbackslash{}in\textbackslash{}Ob\textbackslash{}not\{E\}\$\par
\emph{Proposal}: \$X\textbackslash{}in\textbackslash{}Ob(E)\$\par
This finding intersects Orgogozo's \textbackslash{}sisi material at Exercice 9.1.7 a), dernière ligne, après Ob; preserve that human entry unchanged. The replacement \$X\textbackslash{}in\textbackslash{}Ob(E)\$ at Exercice 9.1.7 a), dernière ligne, après Ob identifies a particular indexed object or symbol; that identification is not independently established by the quoted local expression.\par
\par\noindent\textbf{OM-SOL-0253 -- t1 IV p.437 -- SOL}\par
Exercice 9.1.7.2, début du premier alinéa\par
\emph{Reported reading}: Soit \$f:F\textbackslash{}longrightarrow E\$\par
\emph{Proposal}: a) Soit \$f:F\textbackslash{}longrightarrow E\$\par
Adding an item label is structurally plausible, but the proposed absent printed label cannot be verified from this source record alone.\par
\par\noindent\textbf{OM-SOL-0257 -- t1 IV p.444 -- SOL}\par
9.1.11 b), ensemble ordonné des sous-topos\par
\emph{Reported reading}: avec l'ensemble ordonné des sous-topos de X\par
\emph{Proposal}: avec l'ensemble ordonné des sous-topos de \$E\$.\par
The proposed reading avec l'ensemble ordonné des sous-topos de \$E\$. at 9.1.11 b), ensemble ordonné des sous-topos changes a substantive mathematical assertion or bibliographic identification; no independent derivation of this exact replacement is established in this review.\par
\par\noindent\textbf{OM-SOL-0258 -- t1 IV p.446 -- SOL}\par
Fin de 9.1.11 g), compacité locale\par
\emph{Reported reading}: \$X\$ est localement compact\par
\emph{Proposal}: \$X\_\{\textbackslash{}cons\}\$ est localement compact\par
This finding intersects Orgogozo's \textbackslash{}sisi material at Fin de 9.1.11 g), compacité locale; preserve that human entry unchanged. The proposed reading \$X\_\{\textbackslash{}cons\}\$ est localement compact at Fin de 9.1.11 g), compacité locale changes a substantive mathematical assertion or bibliographic identification; no independent derivation of this exact replacement is established in this review.\par
\par\noindent\textbf{OM-SOL-0259 -- t1 IV p.446 -- SOL}\par
Fin de 9.1.11 g), topos noethérien\par
\emph{Reported reading}: le topos \$\textbackslash{}Top(E)\$ est noethérien\par
\emph{Proposal}: le topos \$\textbackslash{}Top(X)\$ est noethérien\par
This finding intersects Orgogozo's \textbackslash{}sisi material at Fin de 9.1.11 g), topos noethérien; preserve that human entry unchanged. The proposed reading le topos \$\textbackslash{}Top(X)\$ est noethérien at Fin de 9.1.11 g), topos noethérien changes a substantive mathematical assertion or bibliographic identification; no independent derivation of this exact replacement is established in this review.\par
\par\noindent\textbf{OM-SOL-0261 -- t1 IV p.451 -- SOL}\par
9.2.1, renvoi après « l'est également »\par
\emph{Reported reading}: (1.5.7 a))\par
\emph{Proposal}: \textbackslash{}hyperrefext[I][I.rem5.7]\{5.7\}\par
This finding intersects Orgogozo's \textbackslash{}sisi material at 9.2.1, renvoi après « l'est également »; preserve that human entry unchanged. The proposed reference \textbackslash{}hyperrefext[I][I.rem5.7]\{5.7\} at 9.2.1, renvoi après « l'est également » needs comparison with the statement actually invoked; that target has not been independently established here.\par
\par\noindent\textbf{OM-SOL-0262 -- t1 IV p.454 -- SOL}\par
9.3.1, description des objets de Top(Y)\par
\emph{Reported reading}: dont la restriction à U est le topos final de U\par
\emph{Proposal}: l'objet final de \$U\$\par
This finding intersects Orgogozo's \textbackslash{}sisi material at 9.3.1, description des objets de Top(Y); preserve that human entry unchanged. The proposed reading l'objet final de \$U\$ at 9.3.1, description des objets de Top(Y) changes a substantive mathematical assertion or bibliographic identification; no independent derivation of this exact replacement is established in this review.\par
\par\noindent\textbf{OM-SOL-0263 -- t1 IV p.456 -- SOL}\par
Proposition 9.3.4, définition de la sous-catégorie\par
\emph{Reported reading}: définie par les sous-objets \$X\$ de \$E\$\par
\emph{Proposal}: définie par les objets \$X\$ de \$E\$\par
The term sous-objets may refer to a specified object or subcategory; its replacement needs the full construction.\par
\par\noindent\textbf{OM-SOL-0265 -- t1 IV p.458 -- SOL}\par
Démonstration de 9.4.1, exactitude du foncteur de recollement\par
\emph{Reported reading}: Les foncteurs \$i\_\{\textbackslash{}ast\}\$ et \$j\_\{\textbackslash{}ast\}\$ sont exacts à gauche\par
\emph{Proposal}: \$i\textasciicircum{}\{\textbackslash{}ast\}\$ et \$j\_\{\textbackslash{}ast\}\$ sont exacts à gauche\par
The replacement \$i\textasciicircum{}\{\textbackslash{}ast\}\$ et \$j\_\{\textbackslash{}ast\}\$ sont exacts à gauche at Démonstration de 9.4.1, exactitude du foncteur de recollement identifies a particular indexed object or symbol; that identification is not independently established by the quoted local expression.\par
\par\noindent\textbf{OM-SOL-0266 -- t1 IV p.461 -- SOL}\par
9.4.7 b), première formule\par
\emph{Reported reading}: \$(C\_\{/S\})\textasciicircum{}\{\textbackslash{}sim\}\textbackslash{}approxeq E\textasciicircum{}\{\textbackslash{}sim\}\_\{/\textbackslash{}varepsilon(S)\}\$\par
\emph{Proposal}: \$(C\_\{/S\})\textasciicircum{}\{\textbackslash{}sim\}\textbackslash{}approxeq E\_\{/\textbackslash{}varepsilon(S)\}\$\par
The replacement \$(C\_\{/S\})\textasciicircum{}\{\textbackslash{}sim\}\textbackslash{}approxeq E\_\{/\textbackslash{}varepsilon(S)\}\$ at 9.4.7 b), première formule identifies a particular indexed object or symbol; that identification is not independently established by the quoted local expression.\par
\par\noindent\textbf{OM-SOL-0267 -- t1 IV p.462 -- SOL}\par
9.4.8 c), explicitation de « sommande direct »\par
\emph{Reported reading}: un ouvert V […] tel qu'il existe un sous-objet V de e avec e\ensuremath{\simeq}U\ensuremath{\amalg}V\par
\emph{Proposal}: tel qu'il existe un sous-objet \$U\$ de \$e\$ avec \$e\textbackslash{}simeq U\textbackslash{}amalg V\$\par
The proposed reading tel qu'il existe un sous-objet \$U\$ de \$e\$ avec \$e\textbackslash{}simeq U\textbackslash{}amalg V\$ at 9.4.8 c), explicitation de « sommande direct » changes a substantive mathematical assertion or bibliographic identification; no independent derivation of this exact replacement is established in this review.\par
\par\noindent\textbf{OM-SOL-0268 -- t1 IV p.462 -- SOL}\par
9.4.8 d), identification du sous-topos induit\par
\emph{Reported reading}: \$F'\textbackslash{}simeq F\_\{/S'\}\$\par
\emph{Proposal}: \$F'\textbackslash{}simeq E'\_\{/S'\}\$\par
The proposed reading \$F'\textbackslash{}simeq E'\_\{/S'\}\$ at 9.4.8 d), identification du sous-topos induit changes a substantive mathematical assertion or bibliographic identification; no independent derivation of this exact replacement is established in this review.\par
\par\noindent\textbf{OM-SOL-0270 -- t1 IV p.464 -- SOL}\par
9.5.3, renvoi après \ensuremath{\rho}=\allowbreak i*j\_*\par
\emph{Reported reading}: (9.10)\par
\emph{Proposal}: (9.3.6.3)\par
This finding intersects Orgogozo's \textbackslash{}sisi material at 9.5.3, renvoi après \ensuremath{\rho}=\allowbreak i*j\_*; preserve that human entry unchanged. The proposed reference (9.3.6.3) at 9.5.3, renvoi après \ensuremath{\rho}=\allowbreak i*j\_* needs comparison with the statement actually invoked; that target has not been independently established here.\par
\par\noindent\textbf{OM-SOL-0272 -- t1 IV p.465 -- SOL}\par
Note (*) de 9.5.6, renvoi « plus général que 9.5.5 »\par
\emph{Reported reading}: Pour un énoncé de recollement plus général que 9.5.5, cf. exercice.\par
\emph{Proposal}: énoncé de recollement plus général que \textbackslash{}hyperref[IV.theo9.5.4]\{9.5.4\}\par
The proposed reference énoncé de recollement plus général que \textbackslash{}hyperref[IV.theo9.5.4]\{9.5.4\} at Note (*) de 9.5.6, renvoi « plus général que 9.5.5 » needs comparison with the statement actually invoked; that target has not been independently established here.\par
\par\noindent\textbf{OM-SOL-0273 -- t1 IV p.466 -- SOL}\par
9.5.7, reference in the proof heading\par
\emph{Reported reading}: Démonstration de 9.5.4.1)\par
\emph{Proposal}: 9.5.4 b)\par
This finding intersects Orgogozo's \textbackslash{}sisi material at 9.5.7, reference in the proof heading; preserve that human entry unchanged. The proposed reference 9.5.4 b) at 9.5.7, reference in the proof heading needs comparison with the statement actually invoked; that target has not been independently established here.\par
\par\noindent\textbf{OM-SOL-0274 -- t1 IV p.467 -- SOL}\par
c), lower pair of arrows in the displayed diagram\par
\emph{Reported reading}: f(u), f(v)\par
\emph{Proposal}: \textbackslash{}rho(u),\textbackslash{}quad\textbackslash{}rho(v)\par
The proposed reading \textbackslash{}rho(u),\textbackslash{}quad\textbackslash{}rho(v) at c), lower pair of arrows in the displayed diagram changes a substantive mathematical assertion or bibliographic identification; no independent derivation of this exact replacement is established in this review.\par
\par\noindent\textbf{OM-SOL-0276 -- t1 IV p.468 -- SOL}\par
Opening paragraph, immediately after Ob C''\par
\emph{Reported reading}: A small raised x follows C'', before i\ensuremath{\in}I(X'').\par
\emph{Proposal}: X''\textbackslash{}in\textbackslash{}Ob C'',\textbackslash{}quad i\textbackslash{}in I(X'')\par
The proposed reading X''\textbackslash{}in\textbackslash{}Ob C'',\textbackslash{}quad i\textbackslash{}in I(X'') at Opening paragraph, immediately after Ob C'' changes a substantive mathematical assertion or bibliographic identification; no independent derivation of this exact replacement is established in this review.\par
\par\noindent\textbf{OM-SOL-0278 -- t1 IV p.469 -- SOL}\par
9.5.8.5, parenthetical explanation of the initial object\par
\emph{Reported reading}: (\ensuremath{\phi} objet initial de )\par
\emph{Proposal}: (\$\textbackslash{}phi\$ objet initial de \$\textbackslash{}mathcal\{F\}\$)\par
The record documents a lacuna at 9.5.8.5, parenthetical explanation of the initial object; it does not uniquely establish the proposed completion (\$\textbackslash{}phi\$ objet initial de \$\textbackslash{}mathcal\{F\}\$).\par
\par\noindent\textbf{OM-SOL-0280 -- t1 IV p.469 -- SOL}\par
9.5.8.8, lower-left object of the diagram\par
\emph{Reported reading}: \textbackslash{}rho(y)\par
\emph{Proposal}: \textbackslash{}rho(Y)\par
The replacement \textbackslash{}rho(Y) at 9.5.8.8, lower-left object of the diagram identifies a particular indexed object or symbol; that identification is not independently established by the quoted local expression.\par
\par\noindent\textbf{OM-SOL-0281 -- t1 IV p.469 -- SOL}\par
Exercise 9.5.9 a), second reference in « 1.8 ou 1.8 »\par
\emph{Reported reading}: 1.8 ou 1.8\par
\emph{Proposal}: I\textasciitilde{}8\par
This finding intersects Orgogozo's \textbackslash{}sisi material at Exercise 9.5.9 a), second reference in « 1.8 ou 1.8 »; preserve that human entry unchanged. The proposed reference I\textasciitilde{}8 at Exercise 9.5.9 a), second reference in « 1.8 ou 1.8 » needs comparison with the statement actually invoked; that target has not been independently established here.\par
\par\noindent\textbf{OM-SOL-0282 -- t1 IV p.470 -- SOL}\par
Exercise 9.5.10 a), references following « cribles C' de C »\par
\emph{Reported reading}: (1.4.1, 4.2)\par
\emph{Proposal}: (I\textasciitilde{}4.1, 8.4.4)\par
The proposed reference (I\textasciitilde{}4.1, 8.4.4) at Exercise 9.5.10 a), references following « cribles C' de C » needs comparison with the statement actually invoked; that target has not been independently established here.\par
\par\noindent\textbf{OM-SOL-0283 -- t1 IV p.471 -- SOL}\par
Exercice 9.5.11 a), signature du foncteur\par
\emph{Reported reading}: \textbackslash{}rho\textbackslash{}longrightarrow\textbackslash{}mathcal\{F\}\par
\emph{Proposal}: \textbackslash{}rho:\textbackslash{}mathcal\{U\}\textbackslash{}rightarrow\textbackslash{}mathcal\{F\}\par
This finding intersects Orgogozo's \textbackslash{}sisi material at Exercice 9.5.11 a), signature du foncteur; preserve that human entry unchanged. The proposed reading \textbackslash{}rho:\textbackslash{}mathcal\{U\}\textbackslash{}rightarrow\textbackslash{}mathcal\{F\} at Exercice 9.5.11 a), signature du foncteur changes a substantive mathematical assertion or bibliographic identification; no independent derivation of this exact replacement is established in this review.\par
\par\noindent\textbf{OM-SOL-0284 -- t1 IV p.472 -- SOL}\par
Exercice 9.5.11 a), équivalence de catégories\par
\emph{Reported reading}: (\textbackslash{}bar\{\textbackslash{}mathcal\{U\}\},\textbackslash{}bar\{\textbackslash{}mathcal\{F\}\},\textbackslash{}bar\{\textbackslash{}rho\})\par
\emph{Proposal}: (\textbackslash{}bar\{\textbackslash{}mathcal\{F\}\},\textbackslash{}bar\{\textbackslash{}mathcal\{U\}\},\textbackslash{}bar\{\textbackslash{}rho\})\par
The proposed reading (\textbackslash{}bar\{\textbackslash{}mathcal\{F\}\},\textbackslash{}bar\{\textbackslash{}mathcal\{U\}\},\textbackslash{}bar\{\textbackslash{}rho\}) at Exercice 9.5.11 a), équivalence de catégories changes a substantive mathematical assertion or bibliographic identification; no independent derivation of this exact replacement is established in this review.\par
\par\noindent\textbf{OM-SOL-0286 -- t1 IV p.473 -- SOL}\par
Exercice 9.5.12 b), signature de \textbackslash{}rho\par
\emph{Reported reading}: \textbackslash{}rho:\textbackslash{}Top(X)\textbackslash{}longrightarrow\textbackslash{}Top(Y)\par
\emph{Proposal}: \textbackslash{}rho:\textbackslash{}Top(Y)\textbackslash{}rightarrow\textbackslash{}Top(X)\par
The proposed reading \textbackslash{}rho:\textbackslash{}Top(Y)\textbackslash{}rightarrow\textbackslash{}Top(X) at Exercice 9.5.12 b), signature de \textbackslash{}rho changes a substantive mathematical assertion or bibliographic identification; no independent derivation of this exact replacement is established in this review.\par
\par\noindent\textbf{OM-SOL-0287 -- t1 IV p.474 -- SOL}\par
Premier paragraphe, référence après « topologie canonique »\par
\emph{Reported reading}: (1.1.2 (iii))\par
\emph{Proposal}: \textbackslash{}hyperref[IV.theo1.2]\{1.2\} (iii)\par
This finding intersects Orgogozo's \textbackslash{}sisi material at Premier paragraphe, référence après « topologie canonique »; preserve that human entry unchanged. The proposed reference \textbackslash{}hyperref[IV.theo1.2]\{1.2\} (iii) at Premier paragraphe, référence après « topologie canonique » needs comparison with the statement actually invoked; that target has not been independently established here.\par
\par\noindent\textbf{OM-SOL-0289 -- t1 IV p.475 -- SOL}\par
Formule (9.7.2.3), terme intermédiaire\par
\emph{Reported reading}: \textbackslash{}Hom(\textbackslash{}Pro(\textbackslash{}rho)(p),q)\par
\emph{Proposal}: \textbackslash{}Hom(u(p),v(q))\textbackslash{}simeq \textbackslash{}operatorname\{Nat\}(q\textasciicircum{}\{\textbackslash{}ast\}\textbackslash{}circ\textbackslash{}rho,p\textasciicircum{}\{\textbackslash{}ast\})\textbackslash{}simeq \textbackslash{}Hom(p,\textbackslash{}sigma(q))\par
The proposed reading \textbackslash{}Hom(u(p),v(q))\textbackslash{}simeq \textbackslash{}operatorname\{Nat\}(q\textasciicircum{}\{\textbackslash{}ast\}\textbackslash{}circ\textbackslash{}rho,p\textasciicircum{}\{\textbackslash{}ast\})\textbackslash{}simeq \textbackslash{}Hom(p,\textbackslash{}sigma(q)) at Formule (9.7.2.3), terme intermédiaire changes a substantive mathematical assertion or bibliographic identification; no independent derivation of this exact replacement is established in this review.\par
\par\noindent\textbf{OM-SOL-0290 -- t1 IV p.476 -- SOL}\par
Démonstration, b), les quatre arguments des foncteurs image inverse\par
\emph{Reported reading}: \textbackslash{}mathcal\{U\} dans les quatre expressions p\textasciicircum{}\{\textbackslash{}ast\}(\textbackslash{}mathcal\{U\}) ou \{p'\}\textasciicircum{}\{\textbackslash{}ast\}(\textbackslash{}mathcal\{U\})\par
\emph{Proposal}: e\_\{\textbackslash{}mathcal\{U\}\}\par
This finding intersects Orgogozo's \textbackslash{}sisi material at Démonstration, b), les quatre arguments des foncteurs image inverse; preserve that human entry unchanged. The proposed reading e\_\{\textbackslash{}mathcal\{U\}\} at Démonstration, b), les quatre arguments des foncteurs image inverse changes a substantive mathematical assertion or bibliographic identification; no independent derivation of this exact replacement is established in this review.\par
\par\noindent\textbf{OM-SOL-0291 -- t1 IV p.477 -- SOL}\par
Exercice 9.7.5, définition de P(F)\par
\emph{Reported reading}: points de E qui se factorisent par E\par
\emph{Proposal}: points de \$E\$ qui se factorisent par \$F\$\par
This finding intersects Orgogozo's \textbackslash{}sisi material at Exercice 9.7.5, définition de P(F); preserve that human entry unchanged. The proposed reading points de \$E\$ qui se factorisent par \$F\$ involves a printed glyph or a missing stroke at Exercice 9.7.5, définition de P(F); the textual record does not independently establish that glyph.\par
\par\noindent\textbf{OM-SOL-0292 -- t1 IV p.478 -- SOL}\par
Exercice 9.7.5, première phrase\par
\emph{Reported reading}: ouverte (resp. fermée, resp. localement fermée) de E\par
\emph{Proposal}: \textbackslash{}Point(E)\par
This finding intersects Orgogozo's \textbackslash{}sisi material at Exercice 9.7.5, première phrase; preserve that human entry unchanged. The proposed reading \textbackslash{}Point(E) at Exercice 9.7.5, première phrase changes a substantive mathematical assertion or bibliographic identification; no independent derivation of this exact replacement is established in this review.\par
\par\noindent\textbf{OM-SOL-0293 -- t1 IV p.479 -- SOL}\par
Exercice 9.8.2 a), définition de la donnée de descente\par
\emph{Reported reading}: relativement au diagramme (9.8.6.1)\par
\emph{Proposal}: (9.8.2.1)\par
This finding intersects Orgogozo's \textbackslash{}sisi material at Exercice 9.8.2 a), définition de la donnée de descente; preserve that human entry unchanged. The proposed reading (9.8.2.1) at Exercice 9.8.2 a), définition de la donnée de descente changes a substantive mathematical assertion or bibliographic identification; no independent derivation of this exact replacement is established in this review.\par
\par\noindent\textbf{OM-SOL-0294 -- t1 IV p.481 -- SOL}\par
e), condition (i'), appartenance de S\par
\emph{Reported reading}: \$S\textbackslash{}in\textbackslash{}Ob E\$\par
\emph{Proposal}: S\textbackslash{}in\textbackslash{}Ob T\par
The proposed reading S\textbackslash{}in\textbackslash{}Ob T at e), condition (i'), appartenance de S changes a substantive mathematical assertion or bibliographic identification; no independent derivation of this exact replacement is established in this review.\par
\par\noindent\textbf{OM-SOL-0295 -- t1 IV p.482 -- SOL}\par
g), « Déterminer Point(J) »\par
\emph{Reported reading}: \$\textbackslash{}Ppoint(J)\$\par
\emph{Proposal}: \textbackslash{}Ppoint(T)\par
The proposed reading \textbackslash{}Ppoint(T) at g), « Déterminer Point(J) » changes a substantive mathematical assertion or bibliographic identification; no independent derivation of this exact replacement is established in this review.\par
\par\noindent\textbf{OM-SOL-0296 -- t1 IV p.484 -- SOL}\par
Hint, fin de la première phrase\par
\emph{Reported reading}: \$T\_\{/p\}\$\par
\emph{Proposal}: T\_\{/P\}\par
This finding intersects Orgogozo's \textbackslash{}sisi material at Hint, fin de la première phrase; preserve that human entry unchanged. The proposed reading T\_\{/P\} at Hint, fin de la première phrase changes a substantive mathematical assertion or bibliographic identification; no independent derivation of this exact replacement is established in this review.\par
\par\noindent\textbf{OM-SOL-0297 -- t1 IV p.484 -- SOL}\par
j), énumération suivant « où on prendra pour X respectivement »\par
\emph{Reported reading}: une variété différentiable, un espace analytique complexe ou un espace analytique réel\par
\emph{Proposal}: une variété différentiable, un espace analytique réel ou un espace analytique complexe\par
Reordering the analytic cases requires matching all three cases with the preceding definition.\par
\par\noindent\textbf{OM-SOL-0299 -- t1 IV p.486 -- SOL}\par
c), phrase finale sur la correspondance biunivoque\par
\emph{Reported reading}: relations d'équivalence globalement cohérentes \$R\$ sur \$X\$\par
\emph{Proposal}: cohérentes \$r\$ sur \$X\$\par
This finding intersects Orgogozo's \textbackslash{}sisi material at c), phrase finale sur la correspondance biunivoque; preserve that human entry unchanged. The proposed reading cohérentes \$r\$ sur \$X\$ at c), phrase finale sur la correspondance biunivoque changes a substantive mathematical assertion or bibliographic identification; no independent derivation of this exact replacement is established in this review.\par
\par\noindent\textbf{OM-SOL-0301 -- t1 IV p.487 -- SOL}\par
e), homomorphisme de préfaisceaux\par
\emph{Reported reading}: \$f\_\{\textbackslash{}ast\}(\textbackslash{}Quot\_\{X'\})\textbackslash{}longrightarrow \textbackslash{}Quot\_X\$\par
\emph{Proposal}: \$f\_\{\textbackslash{}ast\}(\textbackslash{}Quot\_\{X'\})\textbackslash{}leftarrow \textbackslash{}Quot\_X\$\par
The replacement \$f\_\{\textbackslash{}ast\}(\textbackslash{}Quot\_\{X'\})\textbackslash{}leftarrow \textbackslash{}Quot\_X\$ at e), homomorphisme de préfaisceaux requires a complete check of the maps, their endpoints and the composition convention; this identification is not independently verified.\par
\par\noindent\textbf{OM-SOL-0302 -- t1 IV p.487 -- SOL}\par
e), homomorphisme de faisceaux\par
\emph{Reported reading}: \$f\_\{\textbackslash{}ast\}(Q\_\{X'\})\textbackslash{}longrightarrow Q\_X\$\par
\emph{Proposal}: \$f\_\{\textbackslash{}ast\}(Q\_\{X'\})\textbackslash{}leftarrow Q\_X\$\par
The replacement \$f\_\{\textbackslash{}ast\}(Q\_\{X'\})\textbackslash{}leftarrow Q\_X\$ at e), homomorphisme de faisceaux requires a complete check of the maps, their endpoints and the composition convention; this identification is not independently verified.\par
\par\noindent\textbf{OM-SOL-0307 -- t1 IV p.497 -- SOL}\par
11.2.1, renvoi après « morphisme de localisation »\par
\emph{Reported reading}: (IV 8)\par
\emph{Proposal}: (IV\textasciitilde{}5.2)\par
This finding intersects Orgogozo's \textbackslash{}sisi material at 11.2.1, renvoi après « morphisme de localisation »; preserve that human entry unchanged. The proposed reference (IV\textasciitilde{}5.2) at 11.2.1, renvoi après « morphisme de localisation » needs comparison with the statement actually invoked; that target has not been independently established here.\par
\par\noindent\textbf{OM-SOL-0313 -- t1 IV p.503 -- SOL}\par
Première chaîne, preuve de b), première ligne, renvoi à droite\par
\emph{Reported reading}: (12.2.1)\par
\emph{Proposal}: \textbackslash{}eqhyperref[IV.eq12.1.2]\{12.1.2\}\par
The proposed reference \textbackslash{}eqhyperref[IV.eq12.1.2]\{12.1.2\} at Première chaîne, preuve de b), première ligne, renvoi à droite needs comparison with the statement actually invoked; that target has not been independently established here.\par
\par\noindent\textbf{OM-SOL-0314 -- t1 IV p.503 -- SOL}\par
Première chaîne, preuve de b), deuxième ligne, membre gauche\par
\emph{Reported reading}: \textbackslash{}Hom(M,\textbackslash{}scrHom\_E(j\_\{X!\}Y,N))\par
\emph{Proposal}: \textbackslash{}Hom\_A(M,\textbackslash{}scrHom\_E(j\_\{X!\}Y,N))\par
The replacement \textbackslash{}Hom\_A(M,\textbackslash{}scrHom\_E(j\_\{X!\}Y,N)) at Première chaîne, preuve de b), deuxième ligne, membre gauche requires a complete check of the maps, their endpoints and the composition convention; this identification is not independently verified.\par
\par\noindent\textbf{OM-SOL-0315 -- t1 IV p.503 -- SOL}\par
Première chaîne, preuve de b), troisième ligne, indice du Hom de gauche\par
\emph{Reported reading}: \textbackslash{}Hom\_\{\textbackslash{}Lambda\}(M,j\_\{X\textbackslash{}ast\}\textbackslash{}scrHom\_\{E\_\{/X\}\}(Y,j\textasciicircum{}\{\textbackslash{}ast\}\_XN))\par
\emph{Proposal}: \textbackslash{}Hom\_A(M,j\_\{X\textbackslash{}ast\}\textbackslash{}scrHom\_\{E\_\{/X\}\}(Y,j\textasciicircum{}\{\textbackslash{}ast\}\_XN))\par
The replacement \textbackslash{}Hom\_A(M,j\_\{X\textbackslash{}ast\}\textbackslash{}scrHom\_\{E\_\{/X\}\}(Y,j\textasciicircum{}\{\textbackslash{}ast\}\_XN)) at Première chaîne, preuve de b), troisième ligne, indice du Hom de gauche requires a complete check of the maps, their endpoints and the composition convention; this identification is not independently verified.\par
\par\noindent\textbf{OM-SOL-0316 -- t1 IV p.503 -- SOL}\par
Première chaîne, preuve de b), troisième ligne, membre droit\par
\emph{Reported reading}: \textbackslash{}Hom\_\{A/X\}(j\textasciicircum{}\{\textbackslash{}ast\}\_XM,\textbackslash{}scrHom\_E(j\textasciicircum{}\{\textbackslash{}ast\}\_XN))\par
\emph{Proposal}: \textbackslash{}Hom\_\{A/X\}(j\textasciicircum{}\{\textbackslash{}ast\}\_XM,\textbackslash{}scrHom\_\{E\_\{/X\}\}(Y,j\textasciicircum{}\{\textbackslash{}ast\}\_XN))\par
The replacement \textbackslash{}Hom\_\{A/X\}(j\textasciicircum{}\{\textbackslash{}ast\}\_XM,\textbackslash{}scrHom\_\{E\_\{/X\}\}(Y,j\textasciicircum{}\{\textbackslash{}ast\}\_XN)) at Première chaîne, preuve de b), troisième ligne, membre droit identifies a particular indexed object or symbol; that identification is not independently established by the quoted local expression.\par
\par\noindent\textbf{OM-SOL-0317 -- t1 IV p.503 -- SOL}\par
Première chaîne, preuve de b), quatrième ligne, Hom interne du membre gauche\par
\emph{Reported reading}: \textbackslash{}scrHom\_E(Y,j\textasciicircum{}\{\textbackslash{}ast\}\_XN)\par
\emph{Proposal}: \textbackslash{}scrHom\_\{E\_\{/X\}\}(Y,j\textasciicircum{}\{\textbackslash{}ast\}\_XN)\par
The replacement \textbackslash{}scrHom\_\{E\_\{/X\}\}(Y,j\textasciicircum{}\{\textbackslash{}ast\}\_XN) at Première chaîne, preuve de b), quatrième ligne, Hom interne du membre gauche identifies a particular indexed object or symbol; that identification is not independently established by the quoted local expression.\par
\par\noindent\textbf{OM-SOL-0318 -- t1 IV p.503 -- SOL}\par
Première chaîne, preuve de b), quatrième ligne, renvoi à droite\par
\emph{Reported reading}: (12.2.1)\par
\emph{Proposal}: \textbackslash{}eqhyperref[IV.eq12.1.2]\{12.1.2\}\par
The proposed reference \textbackslash{}eqhyperref[IV.eq12.1.2]\{12.1.2\} at Première chaîne, preuve de b), quatrième ligne, renvoi à droite needs comparison with the statement actually invoked; that target has not been independently established here.\par
\par\noindent\textbf{OM-SOL-0319 -- t1 IV p.503 -- SOL}\par
Démonstration de c), première ligne, premier j du membre droit\par
\emph{Reported reading}: j\_X\textasciicircum{}\{X\}Y\par
\emph{Proposal}: j\textasciicircum{}\{\textbackslash{}ast\}\_XY\par
The replacement j\textasciicircum{}\{\textbackslash{}ast\}\_XY at Démonstration de c), première ligne, premier j du membre droit identifies a particular indexed object or symbol; that identification is not independently established by the quoted local expression.\par
\par\noindent\textbf{OM-SOL-0320 -- t1 IV p.503 -- SOL}\par
Démonstration de c), deuxième ligne, renvoi à droite\par
\emph{Reported reading}: (12.2.1)\par
\emph{Proposal}: \textbackslash{}eqhyperref[IV.eq12.1.2]\{12.1.2\}\par
The proposed reference \textbackslash{}eqhyperref[IV.eq12.1.2]\{12.1.2\} at Démonstration de c), deuxième ligne, renvoi à droite needs comparison with the statement actually invoked; that target has not been independently established here.\par
\par\noindent\textbf{OM-SOL-0321 -- t1 IV p.503 -- SOL}\par
Démonstration de c), cinquième ligne, renvoi à droite\par
\emph{Reported reading}: (12.2.1)\par
\emph{Proposal}: \textbackslash{}eqhyperref[IV.eq12.1.2]\{12.1.2\}\par
The proposed reference \textbackslash{}eqhyperref[IV.eq12.1.2]\{12.1.2\} at Démonstration de c), cinquième ligne, renvoi à droite needs comparison with the statement actually invoked; that target has not been independently established here.\par
\par\noindent\textbf{OM-SOL-0322 -- t1 IV p.503 -- SOL}\par
Démonstration de d), première ligne, renvoi à droite\par
\emph{Reported reading}: (12.2.1)\par
\emph{Proposal}: \textbackslash{}eqhyperref[IV.eq12.1.2]\{12.1.2\}\par
The proposed reference \textbackslash{}eqhyperref[IV.eq12.1.2]\{12.1.2\} at Démonstration de d), première ligne, renvoi à droite needs comparison with the statement actually invoked; that target has not been independently established here.\par
\par\noindent\textbf{OM-SOL-0323 -- t1 IV p.503 -- SOL}\par
Démonstration de d), quatrième ligne, renvoi à droite\par
\emph{Reported reading}: (12.2.1)\par
\emph{Proposal}: \textbackslash{}eqhyperref[IV.eq12.1.2]\{12.1.2\}\par
The proposed reference \textbackslash{}eqhyperref[IV.eq12.1.2]\{12.1.2\} at Démonstration de d), quatrième ligne, renvoi à droite needs comparison with the statement actually invoked; that target has not been independently established here.\par
\par\noindent\textbf{OM-SOL-0327 -- t1 IV p.506 -- SOL}\par
Proposition 12.10, préfaisceau de l'énoncé\par
\emph{Reported reading}: M'(X)\textbackslash{}otimes\_\{A(X)\}N'(X)\par
\emph{Proposal}: M'(X)\textbackslash{}otimes\_\{A'(X)\}N'(X)\par
The replacement M'(X)\textbackslash{}otimes\_\{A'(X)\}N'(X) at Proposition 12.10, préfaisceau de l'énoncé identifies a particular indexed object or symbol; that identification is not independently established by the quoted local expression.\par
\par\noindent\textbf{OM-SOL-0329 -- t1 IV p.510 -- SOL}\par
Démonstration de 13.4, renvoi après le premier isomorphisme canonique\par
\emph{Reported reading}: (12.12)\par
\emph{Proposal}: \textbackslash{}hyperref[IV.prop12.14]\{12.14\}\par
This finding intersects Orgogozo's \textbackslash{}sisi material at Démonstration de 13.4, renvoi après le premier isomorphisme canonique; preserve that human entry unchanged. The proposed reference \textbackslash{}hyperref[IV.prop12.14]\{12.14\} at Démonstration de 13.4, renvoi après le premier isomorphisme canonique needs comparison with the statement actually invoked; that target has not been independently established here.\par
\par\noindent\textbf{OM-SOL-0330 -- t1 IV p.513 -- SOL}\par
Exercice 13.9 c), condition (i), définition de X\par
\emph{Reported reading}: \$X=\allowbreak f\textasciicircum{}\{-1\}(X)\$\par
\emph{Proposal}: \$X=\allowbreak f\textasciicircum{}\{-1\}(X')\$\par
The proposed reading \$X=\allowbreak f\textasciicircum{}\{-1\}(X')\$ at Exercice 13.9 c), condition (i), définition de X changes a substantive mathematical assertion or bibliographic identification; no independent derivation of this exact replacement is established in this review.\par
\par\noindent\textbf{OM-SOL-0331 -- t1 IV p.513 -- SOL}\par
Exercice 13.9 c), condition (i), égalité affichée\par
\emph{Reported reading}: X'\_\{s'\}=\allowbreak f\textasciicircum{}\{-1\}(X\_s)\par
\emph{Proposal}: X\_s=\allowbreak f\textasciicircum{}\{-1\}(X'\_\{s'\})\par
The proposed reading X\_s=\allowbreak f\textasciicircum{}\{-1\}(X'\_\{s'\}) at Exercice 13.9 c), condition (i), égalité affichée changes a substantive mathematical assertion or bibliographic identification; no independent derivation of this exact replacement is established in this review.\par
\par\noindent\textbf{OM-SOL-0332 -- t1 IV p.515 -- SOL}\par
14.1, renvoi après j\textasciicircum{}\{\textbackslash{}ast\}A\par
\emph{Reported reading}: (11.2.2)\par
\emph{Proposal}: \textbackslash{}hyperref[IV.subsubsec11.2.1]\{11.2.1\}\par
This finding intersects Orgogozo's \textbackslash{}sisi material at 14.1, renvoi après j\textasciicircum{}\{\textbackslash{}ast\}A; preserve that human entry unchanged. The proposed reading \textbackslash{}hyperref[IV.subsubsec11.2.1]\{11.2.1\} involves a printed glyph or a missing stroke at 14.1, renvoi après j\textasciicircum{}\{\textbackslash{}ast\}A; the textual record does not independently establish that glyph.\par
\par\noindent\textbf{OM-SOL-0333 -- t1 IV p.515 -- SOL}\par
14.1, morphisme associé à l'annelage de U\par
\emph{Reported reading}: et le morphisme de topos \$i\$\par
\emph{Proposal}: et le morphisme de topos \$j\$\par
The replacement et le morphisme de topos \$j\$ at 14.1, morphisme associé à l'annelage de U requires a complete check of the maps, their endpoints and the composition convention; this identification is not independently verified.\par
\par\noindent\textbf{OM-SOL-0334 -- t1 IV p.515 -- SOL}\par
14.1, morphisme associé à l'annelage de F\par
\emph{Reported reading}: \$j\$, complété de manière évidente, est alors un morphisme de topos\par
\emph{Proposal}: \$i\$, complété de manière évidente, est alors un morphisme de topos\par
The replacement \$i\$, complété de manière évidente, est alors un morphisme de topos at 14.1, morphisme associé à l'annelage de F requires a complete check of the maps, their endpoints and the composition convention; this identification is not independently verified.\par
\par\noindent\textbf{OM-SOL-0335 -- t1 IV p.518 -- SOL}\par
Proposition 14.8, parenthèse après F\par
\emph{Reported reading}: (notation de 5.9.1)\par
\emph{Proposal}: (notation de \$9.5.1\$)\par
This finding intersects Orgogozo's \textbackslash{}sisi material at Proposition 14.8, parenthèse après F; preserve that human entry unchanged. The proposed reference (notation de \$9.5.1\$) at Proposition 14.8, parenthèse après F needs comparison with the statement actually invoked; that target has not been independently established here.\par
\par\noindent\textbf{OM-SOL-0337 -- t1 IV p.518 -- SOL}\par
Démonstration de 14.10, « On tire donc de la suite exacte ... »\par
\emph{Reported reading}: 15.10.1\par
\emph{Proposal}: 14.10.1\par
This finding intersects Orgogozo's \textbackslash{}sisi material at Démonstration de 14.10, « On tire donc de la suite exacte ... »; preserve that human entry unchanged. The proposed reference 14.10.1 at Démonstration de 14.10, « On tire donc de la suite exacte ... » needs comparison with the statement actually invoked; that target has not been independently established here.\par
\par\noindent\textbf{OM-SOL-0338 -- t1 IV p.519 -- SOL}\par
Bibliographie [6], auteur des notes\par
\emph{Reported reading}: I. Coates\par
\emph{Proposal}: J. Coates\par
The proposed reading J. Coates at Bibliographie [6], auteur des notes changes a substantive mathematical assertion or bibliographic identification; no independent derivation of this exact replacement is established in this review.\par
\par\noindent\textbf{OM-SOL-0342 -- t2 V p.7 -- SOL}\par
End of the proof of c) implies a), after Z =\allowbreak  Im(u)\par
\emph{Reported reading}: Comme est pleinement fidèle\par
\emph{Proposal}: Comme \$\textbackslash{}epsilon\$ est pleinement fidèle\par
This finding intersects Orgogozo's \textbackslash{}sisi material at End of the proof of c) implies a), after Z =\allowbreak  Im(u); preserve that human entry unchanged. The proposed reading Comme \$\textbackslash{}epsilon\$ est pleinement fidèle at End of the proof of c) implies a), after Z =\allowbreak  Im(u) changes a substantive mathematical assertion or bibliographic identification; no independent derivation of this exact replacement is established in this review.\par
\par\noindent\textbf{OM-SOL-0343 -- t2 V p.7 -- SOL}\par
Same sentence, relation between u and epsilon(v)\par
\emph{Reported reading}: \$u : \textbackslash{}epsilon(v)\$\par
\emph{Proposal}: \$u =\allowbreak  \textbackslash{}epsilon(v)\$\par
This finding intersects Orgogozo's \textbackslash{}sisi material at Same sentence, relation between u and epsilon(v); preserve that human entry unchanged. The proposed reading \$u =\allowbreak  \textbackslash{}epsilon(v)\$ at Same sentence, relation between u and epsilon(v) changes a substantive mathematical assertion or bibliographic identification; no independent derivation of this exact replacement is established in this review.\par
\par\noindent\textbf{OM-SOL-0344 -- t2 V p.7 -- SOL}\par
Same sentence, Im(epsilon(V))\par
\emph{Reported reading}: \textbackslash{}epsilon(V)\par
\emph{Proposal}: \textbackslash{}epsilon(v)\par
This finding intersects Orgogozo's \textbackslash{}sisi material at Same sentence, Im(epsilon(V)); preserve that human entry unchanged. The replacement \textbackslash{}epsilon(v) at Same sentence, Im(epsilon(V)) requires a complete check of the maps, their endpoints and the composition convention; this identification is not independently verified.\par
\par\noindent\textbf{OM-SOL-0350 -- t2 V p.15 -- SOL}\par
Preuve de 1.9, premier isomorphisme après « on a donc »\par
\emph{Reported reading}: \textbackslash{}Hom\_\{C\textasciicircum{}\{\textbackslash{}sim\}\_\{\textbackslash{}sheaf V\}\}(\textbackslash{}epsilon X, \textbackslash{}Hhom\_A(M,N))\par
\emph{Proposal}: \textbackslash{}Hom\_\{C\textasciicircum{}\{\textbackslash{}sim\}\_\{\textbackslash{}sheaf V\}\}(\textbackslash{}epsilon(X), \textbackslash{}epsilon(\textbackslash{}Hhom\_A(M,N)))\par
This finding intersects Orgogozo's \textbackslash{}sisi material at Preuve de 1.9, premier isomorphisme après « on a donc »; preserve that human entry unchanged. The proposed reading \textbackslash{}Hom\_\{C\textasciicircum{}\{\textbackslash{}sim\}\_\{\textbackslash{}sheaf V\}\}(\textbackslash{}epsilon(X), \textbackslash{}epsilon(\textbackslash{}Hhom\_A(M,N))) at Preuve de 1.9, premier isomorphisme après « on a donc » changes a substantive mathematical assertion or bibliographic identification; no independent derivation of this exact replacement is established in this review.\par
\par\noindent\textbf{OM-SOL-0351 -- t2 V p.15 -- SOL}\par
1.10, indice de la somme définissant S\_n\par
\emph{Reported reading}: f : \textbackslash{}Delta \textbackslash{}to I\par
\emph{Proposal}: f : \textbackslash{}Delta\_n \textbackslash{}to I\par
This finding intersects Orgogozo's \textbackslash{}sisi material at 1.10, indice de la somme définissant S\_n; preserve that human entry unchanged. The replacement f : \textbackslash{}Delta\_n \textbackslash{}to I at 1.10, indice de la somme définissant S\_n identifies a particular indexed object or symbol; that identification is not independently established by the quoted local expression.\par
\par\noindent\textbf{OM-SOL-0354 -- t2 V p.16 -- SOL}\par
Formule 1.10.2, dernier facteur\par
\emph{Reported reading}: U\_\{f( )\}\par
\emph{Proposal}: U\_\{f(n)\}\par
This finding intersects Orgogozo's \textbackslash{}sisi material at Formule 1.10.2, dernier facteur; preserve that human entry unchanged. The proposed reading U\_\{f(n)\} at Formule 1.10.2, dernier facteur changes a substantive mathematical assertion or bibliographic identification; no independent derivation of this exact replacement is established in this review.\par
\par\noindent\textbf{OM-SOL-0358 -- t2 V p.19 -- SOL}\par
Après 2.2.2, famille des foncteurs Ext\par
\emph{Reported reading}: Le foncteur ... q > 0, forment\par
\emph{Proposal}: Les foncteurs\par
This finding intersects Orgogozo's \textbackslash{}sisi material at Après 2.2.2, famille des foncteurs Ext; preserve that human entry unchanged. The replacement Les foncteurs at Après 2.2.2, famille des foncteurs Ext identifies a particular indexed object or symbol; that identification is not independently established by the quoted local expression.\par
\par\noindent\textbf{OM-SOL-0361 -- t2 V p.21 -- SOL}\par
g), acyclicité des composants\par
\emph{Reported reading}: sont acycliques pour \textbackslash{}varprojlim\_C\par
\emph{Proposal}: \textbackslash{}varinjlim\_C\par
The replacement \textbackslash{}varinjlim\_C at g), acyclicité des composants requires a complete check of the maps, their endpoints and the composition convention; this identification is not independently verified.\par
\par\noindent\textbf{OM-SOL-0362 -- t2 V p.21 -- SOL}\par
g), définition du complexe ST homologique\par
\emph{Reported reading}: \textbackslash{}varprojlim\_C(\textbackslash{}ZZ\_\{ES(C)\}\textbackslash{}otimes M)\par
\emph{Proposal}: \textbackslash{}varinjlim\_C(\textbackslash{}ZZ\_\{ES(C)\}\textbackslash{}otimes M)\par
Changing projective to inductive limit changes the complex; its universal construction needs independent checking.\par
\par\noindent\textbf{OM-SOL-0363 -- t2 V p.24 -- SOL}\par
(2.4.6.1), parenthésage de la première égalité\par
\emph{Reported reading}: La première égalité finit par « (M) , », avec une seule parenthèse fermante après M.\par
\emph{Proposal}: \{\textbackslash{}check\{H\}\}\textasciicircum{}q(X,M)=\allowbreak \{\textbackslash{}check\{H\}\}\textasciicircum{}q(X,\textbackslash{}Hh\textasciicircum{}\{0\}(M))\par
The proposed reading \{\textbackslash{}check\{H\}\}\textasciicircum{}q(X,M)=\allowbreak \{\textbackslash{}check\{H\}\}\textasciicircum{}q(X,\textbackslash{}Hh\textasciicircum{}\{0\}(M)) involves a printed glyph or a missing stroke at (2.4.6.1), parenthésage de la première égalité; the textual record does not independently establish that glyph.\par
\par\noindent\textbf{OM-SOL-0365 -- t2 V p.25 -- SOL}\par
2.5.2, référence après le foncteur epsilon-chapeau H\textasciicircum{}0\par
\emph{Reported reading}: Parenthèses laissées vides après le but du foncteur.\par
\emph{Proposal}: (\textbackslash{}hyperref[V.subsubsec2.4.2]\{2.4.2\})\par
The proposed reference (\textbackslash{}hyperref[V.subsubsec2.4.2]\{2.4.2\}) at 2.5.2, référence après le foncteur epsilon-chapeau H\textasciicircum{}0 needs comparison with the statement actually invoked; that target has not been independently established here.\par
\par\noindent\textbf{OM-SOL-0366 -- t2 V p.26 -- SOL}\par
Proof of Proposition 3.1, final Hom functor\par
\emph{Reported reading}: N \textbackslash{}mapsto \textbackslash{}Hom\_\{C\_\{A,V\}\}(N,\textbackslash{}Hh\textasciicircum{}0(M))\par
\emph{Proposal}: N \textbackslash{}mapsto \textbackslash{}Hom\_\{C\textasciicircum{}\{\textbackslash{}wedge\}\_\{A, \ensuremath{𝒱}\}\} (N, \textbackslash{}Hh\textasciicircum{}\{0\} (M))\par
The proposed reading N \textbackslash{}mapsto \textbackslash{}Hom\_\{C\textasciicircum{}\{\textbackslash{}wedge\}\_\{A, \ensuremath{𝒱}\}\} (N, \textbackslash{}Hh\textasciicircum{}\{0\} (M)) at Proof of Proposition 3.1, final Hom functor changes a substantive mathematical assertion or bibliographic identification; no independent derivation of this exact replacement is established in this review.\par
\par\noindent\textbf{OM-SOL-0367 -- t2 V p.27 -- SOL}\par
Corollary 3.4, formula (3.4.3)\par
\emph{Reported reading}: \textbackslash{}Phi\textasciicircum{}q(M) : \textbackslash{}Hh\textasciicircum{}q(M) \textbackslash{}to \textbackslash{}Hh\textasciicircum{}q(M)\par
\emph{Proposal}: \textbackslash{}Phi\textasciicircum{}q(M) : \textbackslash{}check\{\textbackslash{}Hh\}\textasciicircum{}q(M) \textbackslash{}to \textbackslash{}Hh\textasciicircum{}q(M)\par
The replacement \textbackslash{}Phi\textasciicircum{}q(M) : \textbackslash{}check\{\textbackslash{}Hh\}\textasciicircum{}q(M) \textbackslash{}to \textbackslash{}Hh\textasciicircum{}q(M) at Corollary 3.4, formula (3.4.3) requires a complete check of the maps, their endpoints and the composition convention; this identification is not independently verified.\par
\par\noindent\textbf{OM-SOL-0368 -- t2 V p.27 -- SOL}\par
Corollary 3.4, formula (3.4.4)\par
\emph{Reported reading}: \textbackslash{}Phi\textasciicircum{}q\_X(M) : H\textasciicircum{}q(X,M) \textbackslash{}to H\textasciicircum{}q(X,M)\par
\emph{Proposal}: \textbackslash{}Phi\textasciicircum{}q\_X(M) : \textbackslash{}check\{H\}\textasciicircum{}q(X,M) \textbackslash{}to H\textasciicircum{}q(X,M)\par
The replacement \textbackslash{}Phi\textasciicircum{}q\_X(M) : \textbackslash{}check\{H\}\textasciicircum{}q(X,M) \textbackslash{}to H\textasciicircum{}q(X,M) at Corollary 3.4, formula (3.4.4) identifies a particular indexed object or symbol; that identification is not independently established by the quoted local expression.\par
\par\noindent\textbf{OM-SOL-0369 -- t2 V p.28 -- SOL}\par
Proof of Corollary 3.5, displayed direct-limit identities\par
\emph{Reported reading}: H\textasciicircum{}p(Y,M)=\allowbreak \textbackslash{}varinjlim\_\{\textbackslash{}mathcal U\}H\textasciicircum{}p(\textbackslash{}mathcal U,M), and the analogous identity for \textbackslash{}underline\{M\}\par
\emph{Proposal}: \textbackslash{}check\{H\}\textasciicircum{}p(Y,M)=\allowbreak \textbackslash{}varinjlim\_\{\textbackslash{}mathcal\{U\}\}H\textasciicircum{}p(\textbackslash{}mathcal\{U\},M)\textbackslash{}qquad \textbackslash{}check\{H\}\textasciicircum{}p(Y,\textbackslash{}underline\{M\})=\allowbreak \textbackslash{}varinjlim\_\{\textbackslash{}mathcal\{U\}\}H\textasciicircum{}p(\textbackslash{}mathcal\{U\},\textbackslash{}underline\{M\})\par
The proposed reading \textbackslash{}check\{H\}\textasciicircum{}p(Y,M)=\allowbreak \textbackslash{}varinjlim\_\{\textbackslash{}mathcal\{U\}\}H\textasciicircum{}p(\textbackslash{}mathcal\{U\},M)\textbackslash{}qquad \textbackslash{}check\{H\}\textasciicircum{}p(Y,\textbackslash{}underline\{M\})=\allowbreak \textbackslash{}varinjlim\_\{\textbackslash{}mathcal\{U\}\}H\textasciicircum{}p(\textbackslash{}mathcal\{U\},\textbackslash{}underline\{M\}) at Proof of Corollary 3.5, displayed direct-limit identities changes a substantive mathematical assertion or bibliographic identification; no independent derivation of this exact replacement is established in this review.\par
\par\noindent\textbf{OM-SOL-0370 -- t2 V p.30 -- SOL}\par
Paragraph 4.6, final parenthetical reference\par
\emph{Reported reading}: exercice 4.1.3\par
\emph{Proposal}: exercice 4.13\par
The proposed reference exercice 4.13 at Paragraph 4.6, final parenthetical reference needs comparison with the statement actually invoked; that target has not been independently established here.\par
\par\noindent\textbf{OM-SOL-0387 -- t2 V p.50 -- SOL}\par
Proposition 7.1.2, 1), inequality for coskeleta\par
\emph{Reported reading}: n \textbackslash{}leq m (resp. n \textbackslash{}geq m)\par
\emph{Proposal}: \$n \textbackslash{}leq m\$ (\textbackslash{}resp \$n \textbackslash{}leq m\$)\par
The replacement \$n \textbackslash{}leq m\$ (\textbackslash{}resp \$n \textbackslash{}leq m\$) at Proposition 7.1.2, 1), inequality for coskeleta requires a complete check of the maps, their endpoints and the composition convention; this identification is not independently verified.\par
\par\noindent\textbf{OM-SOL-0389 -- t2 V p.52 -- SOL}\par
Même nom de théorème\par
\emph{Reported reading}: théorème d'Hurewitz\par
\emph{Proposal}: théorème d'Hurewicz\par
The reported spelling needs identification of the particular named theorem; the supplied record contains two competing names.\par
\par\noindent\textbf{OM-SOL-0392 -- t2 V p.54 -- SOL}\par
Début de page, fin de 7.3.1.4\par
\emph{Reported reading}: du hyper-recouvrement de type \$p\$ de \$X\$\par
\emph{Proposal}: des hyper-recouvrements de type \$p\$ de \$X\$\par
The proposed repair changes both the article and the number of hyper-recouvrements; the intended category must be checked.\par
\par\noindent\textbf{OM-SOL-0394 -- t2 V p.55 -- SOL}\par
Preuve de 7.3.2, propriété a)\par
\emph{Reported reading}: \$K\_0\textbackslash{}to e\$\par
\emph{Proposal}: \$\{K\}\textasciicircum{}\{\textbackslash{}sim\}\_0 \textbackslash{}to e\$\par
The replacement \$\{K\}\textasciicircum{}\{\textbackslash{}sim\}\_0 \textbackslash{}to e\$ at Preuve de 7.3.2, propriété a) identifies a particular indexed object or symbol; that identification is not independently established by the quoted local expression.\par
\par\noindent\textbf{OM-SOL-0395 -- t2 V p.55 -- SOL}\par
Preuve de 7.3.2, propriété b)\par
\emph{Reported reading}: \$K\_\{n+1\}\textbackslash{}to(\textbackslash{}cosk\_n(K\_∙))\_\{n+1\}\$\par
\emph{Proposal}: \$\{K\}\textasciicircum{}\{\textbackslash{}sim\}\_\{n + 1\} \textbackslash{}to (\textbackslash{}cosk\_n(\{K\}\textasciicircum{}\{\textbackslash{}sim\}\_\{∙\}))\_\{n + 1\}\$\par
The replacement \$\{K\}\textasciicircum{}\{\textbackslash{}sim\}\_\{n + 1\} \textbackslash{}to (\textbackslash{}cosk\_n(\{K\}\textasciicircum{}\{\textbackslash{}sim\}\_\{∙\}))\_\{n + 1\}\$ at Preuve de 7.3.2, propriété b) identifies a particular indexed object or symbol; that identification is not independently established by the quoted local expression.\par
\par\noindent\textbf{OM-SOL-0396 -- t2 V p.57 -- SOL}\par
7.3.5, dernière phrase\par
\emph{Reported reading}: d'appliquer 3.4. et le sorite 7.3.0\par
\emph{Proposal}: \textbackslash{}hyperref[V.lem7.3.4]\{7.3.4\}\par
The proposed reference \textbackslash{}hyperref[V.lem7.3.4]\{7.3.4\} at 7.3.5, dernière phrase needs comparison with the statement actually invoked; that target has not been independently established here.\par
\par\noindent\textbf{OM-SOL-0397 -- t2 V p.58 -- SOL}\par
7.3.6.0, dernier argument de l'isomorphisme canonique\par
\emph{Reported reading}: \textbackslash{}Hhom\_0(M\_∙ \textbackslash{}times \textbackslash{}sheaf\{\textbackslash{}Delta\}\textasciicircum{}c\_n, N)\par
\emph{Proposal}: \textbackslash{}Hhom\_0(M\_∙ \textbackslash{}times \textbackslash{}sheaf\{\textbackslash{}Delta\}\textasciicircum{}c\_n, N\_∙)\par
The proposed reading \textbackslash{}Hhom\_0(M\_∙ \textbackslash{}times \textbackslash{}sheaf\{\textbackslash{}Delta\}\textasciicircum{}c\_n, N\_∙) at 7.3.6.0, dernier argument de l'isomorphisme canonique changes a substantive mathematical assertion or bibliographic identification; no independent derivation of this exact replacement is established in this review.\par
\par\noindent\textbf{OM-SOL-0398 -- t2 V p.61 -- SOL}\par
Fin de la preuve de 7.3.2, bornes de l'indice i\par
\emph{Reported reading}: 0 \textbackslash{}leq i \textbackslash{}leq k\par
\emph{Proposal}: 0 \textbackslash{}leq i < k\par
The proposed reading 0 \textbackslash{}leq i < k at Fin de la preuve de 7.3.2, bornes de l'indice i changes a substantive mathematical assertion or bibliographic identification; no independent derivation of this exact replacement is established in this review.\par
\par\noindent\textbf{OM-SOL-0400 -- t2 V p.65 -- SOL}\par
Définition 8.1.1, axiome (L2), base du recouvrement\par
\emph{Reported reading}: \$f\_j : V\_j \textbackslash{}to V\$ de \$V\$\par
\emph{Proposal}: \$f\_j : U\_j \textbackslash{}to U\$ de \$U\$\par
The replacement \$f\_j : U\_j \textbackslash{}to U\$ de \$U\$ at Définition 8.1.1, axiome (L2), base du recouvrement requires a complete check of the maps, their endpoints and the composition convention; this identification is not independently verified.\par
\par\noindent\textbf{OM-SOL-0401 -- t2 V p.65 -- SOL}\par
Définition 8.1.1, axiome (L2), égalisation\par
\emph{Reported reading}: \textbackslash{}phi\_0 \textbackslash{}psi\_j =\allowbreak  \textbackslash{}phi\_1, \textbackslash{}psi\_j\par
\emph{Proposal}: \textbackslash{}phi\_0 \textbackslash{}psi\_j =\allowbreak  \textbackslash{}phi\_1 \textbackslash{}psi\_j\par
The proposed reading \textbackslash{}phi\_0 \textbackslash{}psi\_j =\allowbreak  \textbackslash{}phi\_1 \textbackslash{}psi\_j at Définition 8.1.1, axiome (L2), égalisation changes a substantive mathematical assertion or bibliographic identification; no independent derivation of this exact replacement is established in this review.\par
\par\noindent\textbf{OM-SOL-0402 -- t2 V p.65 -- SOL}\par
Foncteur composé définissant c, dernier terme\par
\emph{Reported reading}: \textbackslash{}Hom\_\{\textbackslash{}Ss\}(\textbackslash{}Ll, C)\par
\emph{Proposal}: \textbackslash{}Hhom\_\{\textbackslash{}Ss\}(\textbackslash{}Ll, C)\par
The proposed reading \textbackslash{}Hhom\_\{\textbackslash{}Ss\}(\textbackslash{}Ll, C) involves a printed glyph or a missing stroke at Foncteur composé définissant c, dernier terme; the textual record does not independently establish that glyph.\par
\par\noindent\textbf{OM-SOL-0403 -- t2 V p.66 -- SOL}\par
(CF2), recouvrement\par
\emph{Reported reading}: f\_j : U\_j \textbackslash{}rightrightarrows U\par
\emph{Proposal}: f\_j : U\_j \textbackslash{}to U\par
The replacement f\_j : U\_j \textbackslash{}to U at (CF2), recouvrement requires a complete check of the maps, their endpoints and the composition convention; this identification is not independently verified.\par
\par\noindent\textbf{OM-SOL-0405 -- t2 V p.66 -- SOL}\par
Lemme 8.1.6, objets sources des flèches \ensuremath{\phi}\_i\par
\emph{Reported reading}: F(\textbackslash{}lambda'\_i)\par
\emph{Proposal}: \textbackslash{}phi\_i : F(\textbackslash{}lambda\_i) \textbackslash{}to \textbackslash{}mu \textbackslash{}quad (i=\allowbreak 1,2)\par
The replacement \textbackslash{}phi\_i : F(\textbackslash{}lambda\_i) \textbackslash{}to \textbackslash{}mu \textbackslash{}quad (i=\allowbreak 1,2) at Lemme 8.1.6, objets sources des flèches \ensuremath{\phi}\_i identifies a particular indexed object or symbol; that identification is not independently established by the quoted local expression.\par
\par\noindent\textbf{OM-SOL-0406 -- t2 V p.67 -- SOL}\par
Premier diagramme triangulaire, flèche \ensuremath{\psi}\_\{\ensuremath{\lambda}\_j\}\par
\emph{Reported reading}: X\_\{F(\textbackslash{}lambda\_j)\} \textbackslash{}longrightarrow X|U\_j\par
\emph{Proposal}: X|U\_j \textbackslash{}xrightarrow\{\textbackslash{}psi\_\{\textbackslash{}lambda\_j\}\} X\_\{F(\textbackslash{}lambda\_j)\}\par
The replacement X|U\_j \textbackslash{}xrightarrow\{\textbackslash{}psi\_\{\textbackslash{}lambda\_j\}\} X\_\{F(\textbackslash{}lambda\_j)\} at Premier diagramme triangulaire, flèche \ensuremath{\psi}\_\{\ensuremath{\lambda}\_j\} requires a complete check of the maps, their endpoints and the composition convention; this identification is not independently verified.\par
\par\noindent\textbf{OM-SOL-0407 -- t2 V p.67 -- SOL}\par
Dernière phrase avant le diagramme en losange\par
\emph{Reported reading}: Soit un diagramme commutatif dans\par
\emph{Proposal}: Soit un diagramme commutatif dans \$\textbackslash{}Ss\$\par
The replacement Soit un diagramme commutatif dans \$\textbackslash{}Ss\$ at Dernière phrase avant le diagramme en losange requires a complete check of the maps, their endpoints and the composition convention; this identification is not independently verified.\par
\par\noindent\textbf{OM-SOL-0408 -- t2 V p.68 -- SOL}\par
Dernière phrase de la démonstration\par
\emph{Reported reading}: \textbackslash{}psi =\allowbreak  F\textasciicircum{}\{\textbackslash{}ast\}\par
\emph{Proposal}: \textbackslash{}psi =\allowbreak  F\textasciicircum{}\{\textbackslash{}ast\}\textbackslash{}phi\par
The proposed reading \textbackslash{}psi =\allowbreak  F\textasciicircum{}\{\textbackslash{}ast\}\textbackslash{}phi at Dernière phrase de la démonstration changes a substantive mathematical assertion or bibliographic identification; no independent derivation of this exact replacement is established in this review.\par
\par\noindent\textbf{OM-SOL-0409 -- t2 V p.70 -- SOL}\par
Vérification (ii), choix des objets \ensuremath{\lambda}\_j\par
\emph{Reported reading}: \textbackslash{}lambda\_j \textbackslash{}in \textbackslash{}Ob \textbackslash{}Ll\_\{U\_j\}\par
\emph{Proposal}: \textbackslash{}lambda\_j \textbackslash{}in \textbackslash{}Ob \textbackslash{}Ll'\_\{U\_j\}\par
The replacement \textbackslash{}lambda\_j \textbackslash{}in \textbackslash{}Ob \textbackslash{}Ll'\_\{U\_j\} at Vérification (ii), choix des objets \ensuremath{\lambda}\_j identifies a particular indexed object or symbol; that identification is not independently established by the quoted local expression.\par
\par\noindent\textbf{OM-SOL-0410 -- t2 V p.70 -- SOL}\par
Vérification (ii), phrase commençant par « Alors »\par
\emph{Reported reading}: X\_j \textbackslash{}in E\_\{V\_j\}\par
\emph{Proposal}: X\_j \textbackslash{}in E\_\{U\_j\}\par
The replacement X\_j \textbackslash{}in E\_\{U\_j\} at Vérification (ii), phrase commençant par « Alors » identifies a particular indexed object or symbol; that identification is not independently established by the quoted local expression.\par
\par\noindent\textbf{OM-SOL-0411 -- t2 V p.70 -- SOL}\par
Vérification (v), égalisation de \ensuremath{\phi}\_0 et \ensuremath{\phi}\_1\par
\emph{Reported reading}: \textbackslash{}phi\_0,\textbackslash{}psi\_j =\allowbreak  \textbackslash{}phi\_1\textbackslash{}psi\_j\par
\emph{Proposal}: \textbackslash{}phi\_0\textbackslash{}psi\_j =\allowbreak  \textbackslash{}phi\_1\textbackslash{}psi\_j\par
The proposed reading \textbackslash{}phi\_0\textbackslash{}psi\_j =\allowbreak  \textbackslash{}phi\_1\textbackslash{}psi\_j at Vérification (v), égalisation de \ensuremath{\phi}\_0 et \ensuremath{\phi}\_1 changes a substantive mathematical assertion or bibliographic identification; no independent derivation of this exact replacement is established in this review.\par
\par\noindent\textbf{OM-SOL-0412 -- t2 V p.71 -- SOL}\par
8.1.10 (ii), source and target of f\par
\emph{Reported reading}: f : V,U,\textbackslash{}phi,\textbackslash{}lambda) \textbackslash{}longrightarrow (V',\textbackslash{}phi',U',\textbackslash{}lambda')\par
\emph{Proposal}: f:(V,U,\textbackslash{}phi,\textbackslash{}lambda)\textbackslash{}to(V',U',\textbackslash{}phi',\textbackslash{}lambda')\par
The proposed reading f:(V,U,\textbackslash{}phi,\textbackslash{}lambda)\textbackslash{}to(V',U',\textbackslash{}phi',\textbackslash{}lambda') at 8.1.10 (ii), source and target of f changes a substantive mathematical assertion or bibliographic identification; no independent derivation of this exact replacement is established in this review.\par
\par\noindent\textbf{OM-SOL-0413 -- t2 V p.71 -- SOL}\par
8.1.10 (ii), the triple defining f\par
\emph{Reported reading}: (f\_1,f\_2,F\_3)\par
\emph{Proposal}: (f\_1,f\_2,f\_3)\par
The replacement (f\_1,f\_2,f\_3) at 8.1.10 (ii), the triple defining f identifies a particular indexed object or symbol; that identification is not independently established by the quoted local expression.\par
\par\noindent\textbf{OM-SOL-0414 -- t2 V p.72 -- SOL}\par
Equalizer formula, final object\par
\emph{Reported reading}: q\textasciicircum{}\{\textbackslash{}ast\}U\_2\par
\emph{Proposal}: q\textasciicircum{}\{\textbackslash{}ast\}U\textasciicircum{}2\par
The proposed reading q\textasciicircum{}\{\textbackslash{}ast\}U\textasciicircum{}2 at Equalizer formula, final object changes a substantive mathematical assertion or bibliographic identification; no independent derivation of this exact replacement is established in this review.\par
\par\noindent\textbf{OM-SOL-0415 -- t2 V p.72 -- SOL}\par
Diagram for (L2), lower-right object\par
\emph{Reported reading}: U\textasciicircum{}2\par
\emph{Proposal}: q\textasciicircum{}\{\textbackslash{}ast\}U\textasciicircum{}\{2\}\par
The proposed reading q\textasciicircum{}\{\textbackslash{}ast\}U\textasciicircum{}\{2\} at Diagram for (L2), lower-right object changes a substantive mathematical assertion or bibliographic identification; no independent derivation of this exact replacement is established in this review.\par
\par\noindent\textbf{OM-SOL-0416 -- t2 V p.72 -- SOL}\par
Pair of arrows to which (L2) is applied\par
\emph{Reported reading}: (\textbackslash{}lambda\_1\textbackslash{}chi\_j,h\_2\textbackslash{}chi\_j)\par
\emph{Proposal}: (h\_1\textbackslash{}chi\_j,h\_2\textbackslash{}chi\_j)\par
The proposed reading (h\_1\textbackslash{}chi\_j,h\_2\textbackslash{}chi\_j) at Pair of arrows to which (L2) is applied changes a substantive mathematical assertion or bibliographic identification; no independent derivation of this exact replacement is established in this review.\par
\par\noindent\textbf{OM-SOL-0417 -- t2 V p.72 -- SOL}\par
Equality obtained after refining U\_j and V\_j\par
\emph{Reported reading}: h, \textbackslash{}chi\_j =\allowbreak  h\_2\textbackslash{}chi\_j\par
\emph{Proposal}: h\_1\textbackslash{}chi\_j=\allowbreak h\_2\textbackslash{}chi\_j\par
The replacement h\_1\textbackslash{}chi\_j=\allowbreak h\_2\textbackslash{}chi\_j at Equality obtained after refining U\_j and V\_j identifies a particular indexed object or symbol; that identification is not independently established by the quoted local expression.\par
\par\noindent\textbf{OM-SOL-0418 -- t2 V p.73 -- SOL}\par
First commutative square, lower-right vertex\par
\emph{Reported reading}: Only a raised 2 is printed at this vertex.\par
\emph{Proposal}: \textbackslash{}lambda\textasciicircum{}2\par
The proposed reading \textbackslash{}lambda\textasciicircum{}2 at First commutative square, lower-right vertex changes a substantive mathematical assertion or bibliographic identification; no independent derivation of this exact replacement is established in this review.\par
\par\noindent\textbf{OM-SOL-0419 -- t2 V p.73 -- SOL}\par
Proof (ii), site containing f:U\ensuremath{\to}V\par
\emph{Reported reading}: f:U\textbackslash{}to V dans \textbackslash{}Ss\textasciicircum{}\{\textbackslash{}ast\}\par
\emph{Proposal}: f:U\textbackslash{}to V\textbackslash{}text\{ dans \}\textbackslash{}Ss\par
The proposed reading f:U\textbackslash{}to V\textbackslash{}text\{ dans \}\textbackslash{}Ss at Proof (ii), site containing f:U\ensuremath{\to}V changes a substantive mathematical assertion or bibliographic identification; no independent derivation of this exact replacement is established in this review.\par
\par\noindent\textbf{OM-SOL-0420 -- t2 V p.73 -- SOL}\par
Proof (ii), membership condition for \ensuremath{\mu}′\par
\emph{Reported reading}: \textbackslash{}mu'\textbackslash{}in\textbackslash{}Ob\{\}\textasciicircum{}\{\textbackslash{}ast\}\_U\par
\emph{Proposal}: \textbackslash{}mu'\textbackslash{}in\textbackslash{}Ob\textbackslash{}Ss\textasciicircum{}\{\textbackslash{}ast\}\_U\par
The record documents a lacuna at Proof (ii), membership condition for \ensuremath{\mu}′; it does not uniquely establish the proposed completion \textbackslash{}mu'\textbackslash{}in\textbackslash{}Ob\textbackslash{}Ss\textasciicircum{}\{\textbackslash{}ast\}\_U.\par
\par\noindent\textbf{OM-SOL-0421 -- t2 V p.73 -- SOL}\par
Object diagram in (ii), intermediate object on the lower branch\par
\emph{Reported reading}: \textbackslash{}lambda\_j\textbackslash{}xrightarrow\{\textbackslash{}psi\_j\}\textbackslash{}mu\textbackslash{}xrightarrow\{\textbackslash{}psi\}\textbackslash{}mu\par
\emph{Proposal}: \textbackslash{}lambda\_j\textbackslash{}xrightarrow\{\textbackslash{}psi\_j\}\textbackslash{}mu'\textbackslash{}xrightarrow\{\textbackslash{}psi\}\textbackslash{}mu\par
The replacement \textbackslash{}lambda\_j\textbackslash{}xrightarrow\{\textbackslash{}psi\_j\}\textbackslash{}mu'\textbackslash{}xrightarrow\{\textbackslash{}psi\}\textbackslash{}mu at Object diagram in (ii), intermediate object on the lower branch identifies a particular indexed object or symbol; that identification is not independently established by the quoted local expression.\par
\par\noindent\textbf{OM-SOL-0422 -- t2 V p.74 -- SOL}\par
Proof (iii), fibres containing \ensuremath{\lambda}⁰, \ensuremath{\mu}⁰ and \ensuremath{\nu}⁰\par
\emph{Reported reading}: \textbackslash{}Ob\textbackslash{}Ss\_U,\textbackslash{}quad\textbackslash{}Ob\textbackslash{}Ss\_V,\textbackslash{}quad\textbackslash{}Ob\textbackslash{}Ss\_W\par
\emph{Proposal}: \textbackslash{}lambda\textasciicircum{}\{0\}\textbackslash{}in\textbackslash{}Ob\textbackslash{}Ss\textasciicircum{}\{\textbackslash{}ast\}\_U,\textbackslash{}quad\textbackslash{}mu\textasciicircum{}\{0\}\textbackslash{}in\textbackslash{}Ob\textbackslash{}Ss\textasciicircum{}\{\textbackslash{}ast\}\_V,\textbackslash{}quad\textbackslash{}nu\textasciicircum{}\{0\}\textbackslash{}in\textbackslash{}Ob\textbackslash{}Ss\textasciicircum{}\{\textbackslash{}ast\}\_W\par
The proposed reading \textbackslash{}lambda\textasciicircum{}\{0\}\textbackslash{}in\textbackslash{}Ob\textbackslash{}Ss\textasciicircum{}\{\textbackslash{}ast\}\_U,\textbackslash{}quad\textbackslash{}mu\textasciicircum{}\{0\}\textbackslash{}in\textbackslash{}Ob\textbackslash{}Ss\textasciicircum{}\{\textbackslash{}ast\}\_V,\textbackslash{}quad\textbackslash{}nu\textasciicircum{}\{0\}\textbackslash{}in\textbackslash{}Ob\textbackslash{}Ss\textasciicircum{}\{\textbackslash{}ast\}\_W at Proof (iii), fibres containing \ensuremath{\lambda}⁰, \ensuremath{\mu}⁰ and \ensuremath{\nu}⁰ changes a substantive mathematical assertion or bibliographic identification; no independent derivation of this exact replacement is established in this review.\par
\par\noindent\textbf{OM-SOL-0423 -- t2 V p.75 -- SOL}\par
Last paragraph, indexing category of the local inductive system\par
\emph{Reported reading}: indexé par \textbackslash{}lambda\textasciicircum{}\{\textbackslash{}ast\}\par
\emph{Proposal}: \textbackslash{}text\{indexé par \}\textbackslash{}Ss\textasciicircum{}\{\textbackslash{}ast\}\par
The replacement \textbackslash{}text\{indexé par \}\textbackslash{}Ss\textasciicircum{}\{\textbackslash{}ast\} at Last paragraph, indexing category of the local inductive system identifies a particular indexed object or symbol; that identification is not independently established by the quoted local expression.\par
\par\noindent\textbf{OM-SOL-0425 -- t2 V p.79 -- SOL}\par
Preuve de 8.2.7, renvoi aux hypothèses sur les sites\par
\emph{Reported reading}: les hypothèses faites en 2.6\par
\emph{Proposal}: les hypothèses faites en 8.2.6\par
The proposed reference les hypothèses faites en 8.2.6 at Preuve de 8.2.7, renvoi aux hypothèses sur les sites needs comparison with the statement actually invoked; that target has not been independently established here.\par
\par\noindent\textbf{OM-SOL-0426 -- t2 V p.80 -- SOL}\par
Diagramme de la preuve de 8.2.10, deux dernières colonnes\par
\emph{Reported reading}: \textbackslash{}Ff\_1 puis \textbackslash{}Ff\_2 dans les deux lignes du diagramme\par
\emph{Proposal}: \textbackslash{}xymatrix\{ 0 \textbackslash{}to \textbackslash{}Hhom(\textbackslash{}Ff,\textbackslash{}Gg) \textbackslash{}otimes \textbackslash{}Hh \textbackslash{}ar[r] \textbackslash{}ar[d]\& \textbackslash{}Hhom(\textbackslash{}Ff\_0, \textbackslash{}Gg) \textbackslash{}otimes \textbackslash{}Hh \textbackslash{}ar[r]\textbackslash{}ar[d]\textasciicircum{}\{S\} \& \textbackslash{}Hhom(\textbackslash{}Ff\_1,\textbackslash{}Gg) \textbackslash{}otimes \textbackslash{}Hh \textbackslash{}ar[d]\textasciicircum{}\{S\}\textbackslash{}\textbackslash{} 0 \textbackslash{}to \textbackslash{}Hhom(\textbackslash{}Ff,\textbackslash{}Gg \textbackslash{}otimes \textbackslash{}Hh) \textbackslash{}ar[r]\& \textbackslash{}Hhom(\textbackslash{}Ff\_0, \textbackslash{}Gg \textbackslash{}otimes \textbackslash{}Hh) \textbackslash{}ar[r] \& \textbackslash{}Hhom(\textbackslash{}Ff\_1,\textbackslash{}Gg \textbackslash{}otimes \textbackslash{}Hh) \}\par
The proposed reading \textbackslash{}xymatrix\{ 0 \textbackslash{}to \textbackslash{}Hhom(\textbackslash{}Ff,\textbackslash{}Gg) \textbackslash{}otimes \textbackslash{}Hh \textbackslash{}ar[r] \textbackslash{}ar[d]\& \textbackslash{}Hhom(\textbackslash{}Ff\_0, \textbackslash{}Gg) \textbackslash{}otimes \textbackslash{}Hh \textbackslash{}ar[r]\textbackslash{}ar[d]\textasciicircum{}\{S\} \& \textbackslash{}Hhom(\textbackslash{}Ff\_1,\textbackslash{}Gg) \textbackslash{}otimes \textbackslash{}Hh \textbackslash{}ar[d]\textasciicircum{}\{S\}\textbackslash{}\textbackslash{} 0 \textbackslash{}to \textbackslash{}Hhom(\textbackslash{}Ff,\textbackslash{}Gg \textbackslash{}otimes \textbackslash{}Hh) \textbackslash{}ar[r]\& \textbackslash{}Hhom(\textbackslash{}Ff\_0, \textbackslash{}Gg \textbackslash{}otimes \textbackslash{}Hh) \textbackslash{}ar[r] \& \textbackslash{}Hhom(\textbackslash{}Ff\_1,\textbackslash{}Gg \textbackslash{}otimes \textbackslash{}Hh) \} at Diagramme de la preuve de 8.2.10, deux dernières colonnes changes a substantive mathematical assertion or bibliographic identification; no independent derivation of this exact replacement is established in this review.\par
\par\noindent\textbf{OM-SOL-0427 -- t2 V p.82 -- SOL}\par
Bibliography [1], second author's surname\par
\emph{Reported reading}: B. Mazure\par
\emph{Proposal}: B. Mazur\par
This finding intersects Orgogozo's \textbackslash{}sisi material at Bibliography [1], second author's surname; preserve that human entry unchanged. The proposed reading B. Mazur at Bibliography [1], second author's surname changes a substantive mathematical assertion or bibliographic identification; no independent derivation of this exact replacement is established in this review.\par
\par\noindent\textbf{OM-SOL-0428 -- t2 V p.82 -- SOL}\par
Bibliography [1], title\par
\emph{Reported reading}: Etale Homotopy Theory\par
\emph{Proposal}: Etale Homotopy\par
This finding intersects Orgogozo's \textbackslash{}sisi material at Bibliography [1], title; preserve that human entry unchanged. The proposed reading Etale Homotopy at Bibliography [1], title changes a substantive mathematical assertion or bibliographic identification; no independent derivation of this exact replacement is established in this review.\par
\par\noindent\textbf{OM-SOL-0429 -- t2 V p.82 -- SOL}\par
Bibliography [1], Lecture Notes volume number\par
\emph{Reported reading}: Lecture Notes nᵒ , Springer Verlag.\par
\emph{Proposal}: Lecture Notes nᵒ 100\par
This finding intersects Orgogozo's \textbackslash{}sisi material at Bibliography [1], Lecture Notes volume number; preserve that human entry unchanged. The record documents a lacuna at Bibliography [1], Lecture Notes volume number; it does not uniquely establish the proposed completion Lecture Notes nᵒ 100.\par
\par\noindent\textbf{OM-SOL-0430 -- t2 V p.82 -- SOL}\par
Bibliography [2], first author's surname\par
\emph{Reported reading}: Blum et Herrera\par
\emph{Proposal}: Bloom et Herrera\par
The author-name correction requires a bibliographic authority beyond the quoted entry.\par
\par\noindent\textbf{OM-SOL-0432 -- t2 V p.82 -- SOL}\par
Bibliography [6], title, final word\par
\emph{Reported reading}: Calculus of Fraction\par
\emph{Proposal}: Calculus of Fractions\par
The title's exact plural requires comparison with the cited publication.\par
\par\noindent\textbf{OM-SOL-0433 -- t2 V p.82 -- SOL}\par
Bibliography [7], publisher\par
\emph{Reported reading}: Herman.\par
\emph{Proposal}: Hermann.\par
The publisher's exact spelling requires comparison with the cited edition.\par
\par\noindent\textbf{OM-SOL-0440 -- t2 Vbis p.94 -- SOL}\par
Définition 1.2.14, right-hand side of the adjunction\par
\emph{Reported reading}: (\textbackslash{}cdot,\textbackslash{}Phi\_\{\textbackslash{}ast\}\textbackslash{}cdot)\par
\emph{Proposal}: \textbackslash{}Hom(\textbackslash{}cdot,\textbackslash{}Phi\_\{\textbackslash{}ast\}\textbackslash{}cdot)\par
The proposed reading \textbackslash{}Hom(\textbackslash{}cdot,\textbackslash{}Phi\_\{\textbackslash{}ast\}\textbackslash{}cdot) at Définition 1.2.14, right-hand side of the adjunction changes a substantive mathematical assertion or bibliographic identification; no independent derivation of this exact replacement is established in this review.\par
\par\noindent\textbf{OM-SOL-0442 -- t2 Vbis p.95 -- SOL}\par
1.3.4, formula specifying the ring on the D-topos\par
\emph{Reported reading}: (\textbackslash{}mathcal\{O\}.X)\textasciicircum{}\{\textbackslash{}circ\}\par
\emph{Proposal}: (\textbackslash{}mathcal\{O\}. X\textasciicircum{}\{\textbackslash{}circ\})\par
The proposed reading (\textbackslash{}mathcal\{O\}. X\textasciicircum{}\{\textbackslash{}circ\}) at 1.3.4, formula specifying the ring on the D-topos changes a substantive mathematical assertion or bibliographic identification; no independent derivation of this exact replacement is established in this review.\par
\par\noindent\textbf{OM-SOL-0443 -- t2 Vbis p.96 -- SOL}\par
1.3.6, target of the inverse-image functor\par
\emph{Reported reading}: \textbackslash{}varphi\textasciicircum{}\{\textbackslash{}ast\}:\textbackslash{}underline\{\textbackslash{}Mod\}(E',A')\textbackslash{}to\textbackslash{}Mod(E,A)\par
\emph{Proposal}: \textbackslash{}varphi\textasciicircum{}\{\textbackslash{}ast\} : \textbackslash{}underline\{\textbackslash{}Mod\}(E', A') \textbackslash{}to \textbackslash{}underline\{\textbackslash{}Mod\}(E, A)\par
The replacement \textbackslash{}varphi\textasciicircum{}\{\textbackslash{}ast\} : \textbackslash{}underline\{\textbackslash{}Mod\}(E', A') \textbackslash{}to \textbackslash{}underline\{\textbackslash{}Mod\}(E, A) at 1.3.6, target of the inverse-image functor identifies a particular indexed object or symbol; that identification is not independently established by the quoted local expression.\par
\par\noindent\textbf{OM-SOL-0445 -- t2 Vbis p.97 -- SOL}\par
Proposition 1.3.10(i), hypothesis on F\par
\emph{Reported reading}: Pour tout objet F de Mod(\textbackslash{}underline\{\textbackslash{}Gamma\}(E),A)\par
\emph{Proposal}: Pour tout objet \$F\$ de \$I\_\{(E, A)\}\$ et tout objet \$i\$ de \$D\$, \$e\textasciicircum{}\{\textbackslash{}ast\}\_i(F)\$ est totalement acyclique.\par
The proposed reading Pour tout objet \$F\$ de \$I\_\{(E, A)\}\$ et tout objet \$i\$ de \$D\$, \$e\textasciicircum{}\{\textbackslash{}ast\}\_i(F)\$ est totalement acyclique. at Proposition 1.3.10(i), hypothesis on F changes a substantive mathematical assertion or bibliographic identification; no independent derivation of this exact replacement is established in this review.\par
\par\noindent\textbf{OM-SOL-0446 -- t2 Vbis p.99 -- SOL}\par
Proof of Corollaire 1.3.11, description of L\par
\emph{Reported reading}: complexes fermés d'objets\par
\emph{Proposal}: complexes formés d'objets\par
Replacing fermés by formés changes the description of the complexes and requires the construction of L.\par
\par\noindent\textbf{OM-SOL-0447 -- t2 Vbis p.99 -- SOL}\par
Proof of Corollaire 1.3.11, reference establishing that L is triangulated\par
\emph{Reported reading}: ((1.3.10), (ii))\par
\emph{Proposal}: ((\textbackslash{}hyperref[Vbis.prop1.3.10]\{1.3.10\}), (iii))\par
The proposed reference ((\textbackslash{}hyperref[Vbis.prop1.3.10]\{1.3.10\}), (iii)) at Proof of Corollaire 1.3.11, reference establishing that L is triangulated needs comparison with the statement actually invoked; that target has not been independently established here.\par
\par\noindent\textbf{OM-SOL-0448 -- t2 Vbis p.99 -- SOL}\par
Corollaire 1.3.12, description of the fibre morphism \ensuremath{\phi}\_i\par
\emph{Reported reading}: le morphisme de D-topos annelé induit\par
\emph{Proposal}: le morphisme de topos annelés induit\par
Removing the D qualification changes the objects described; the fibre construction needs independent verification.\par
\par\noindent\textbf{OM-SOL-0449 -- t2 Vbis p.101 -- SOL}\par
2.1.3, sentence immediately after the augmentation display\par
\emph{Reported reading}: désignerons par la même lettre que le morphisme fonctoriel qui lui donne naissance.\par
\emph{Proposal}: Nous la désignerons par la même lettre que le morphisme fonctoriel qui lui donne naissance.\par
The sentence is incomplete, but the exact supplied subject/object Nous la is an editorial reconstruction.\par
\par\noindent\textbf{OM-SOL-0450 -- t2 Vbis p.103 -- SOL}\par
Proof of Proposition 2.2.7, b), introduction of the truncation\par
\emph{Reported reading}: \textbackslash{}sigma\_\{\textbackslash{}leq\}(F\textasciicircum{}\{\textbackslash{}bigdot\})\par
\emph{Proposal}: \textbackslash{}sigma\_\{\textbackslash{}leq n\}(F\textasciicircum{}\{\textbackslash{}bigdot\})\par
This finding intersects Orgogozo's \textbackslash{}sisi material at Proof of Proposition 2.2.7, b), introduction of the truncation; preserve that human entry unchanged. The proposed reading \textbackslash{}sigma\_\{\textbackslash{}leq n\}(F\textasciicircum{}\{\textbackslash{}bigdot\}) at Proof of Proposition 2.2.7, b), introduction of the truncation changes a substantive mathematical assertion or bibliographic identification; no independent derivation of this exact replacement is established in this review.\par
\par\noindent\textbf{OM-SOL-0451 -- t2 Vbis p.104 -- SOL}\par
2.3.1.1, parenthetical reference\par
\emph{Reported reading}: Tohokû\par
\emph{Proposal}: Tôhoku\par
This finding intersects Orgogozo's \textbackslash{}sisi material at 2.3.1.1, parenthetical reference; preserve that human entry unchanged. The proposed reading Tôhoku at 2.3.1.1, parenthetical reference changes a substantive mathematical assertion or bibliographic identification; no independent derivation of this exact replacement is established in this review.\par
\par\noindent\textbf{OM-SOL-0452 -- t2 Vbis p.106 -- SOL}\par
First paragraph, Tôhoku reference\par
\emph{Reported reading}: Tohokû\par
\emph{Proposal}: Tôhoku\par
This finding intersects Orgogozo's \textbackslash{}sisi material at First paragraph, Tôhoku reference; preserve that human entry unchanged. The proposed reading Tôhoku at First paragraph, Tôhoku reference changes a substantive mathematical assertion or bibliographic identification; no independent derivation of this exact replacement is established in this review.\par
\par\noindent\textbf{OM-SOL-0453 -- t2 Vbis p.109 -- SOL}\par
Corollary 2.3.7, displayed isomorphism\par
\emph{Reported reading}: \textbackslash{}RR\textasciicircum{}+\textbackslash{}Gamma(\textbackslash{}theta\_\textbackslash{}ast)\textbackslash{}circ\textbackslash{}RR\textasciicircum{}+\textbackslash{}varepsilon\_\{Z\textasciicircum{}\{\textbackslash{}bigdot\}\textbackslash{}ast\}\par
\emph{Proposal}: \textbackslash{}RR\textasciicircum{}+\textbackslash{}bar\{\textbackslash{}theta\}\_\{\textbackslash{}ast\}\textbackslash{}xrightarrow\{\textbackslash{}sim\}\textbackslash{}RR\textasciicircum{}+\textbackslash{}varepsilon\_\{Z\textasciicircum{}\{\textbackslash{}bigdot\}\textbackslash{}ast\}\textbackslash{}circ\textbackslash{}RR\textasciicircum{}+\textbackslash{}Gamma(\textbackslash{}theta\_\{\textbackslash{}ast\})\par
This finding intersects Orgogozo's \textbackslash{}sisi material at Corollary 2.3.7, displayed isomorphism; preserve that human entry unchanged. The proposed reading \textbackslash{}RR\textasciicircum{}+\textbackslash{}bar\{\textbackslash{}theta\}\_\{\textbackslash{}ast\}\textbackslash{}xrightarrow\{\textbackslash{}sim\}\textbackslash{}RR\textasciicircum{}+\textbackslash{}varepsilon\_\{Z\textasciicircum{}\{\textbackslash{}bigdot\}\textbackslash{}ast\}\textbackslash{}circ\textbackslash{}RR\textasciicircum{}+\textbackslash{}Gamma(\textbackslash{}theta\_\{\textbackslash{}ast\}) at Corollary 2.3.7, displayed isomorphism changes a substantive mathematical assertion or bibliographic identification; no independent derivation of this exact replacement is established in this review.\par
\par\noindent\textbf{OM-SOL-0454 -- t2 Vbis p.109 -- SOL}\par
Remark 2.3.8, isomorphism in the first sentence\par
\emph{Reported reading}: \textbackslash{}RR\textasciicircum{}+\textbackslash{}Gamma(\textbackslash{}theta\_\textbackslash{}ast)\textbackslash{}circ\textbackslash{}RR\textasciicircum{}+\textbackslash{}varepsilon\_\textbackslash{}ast\par
\emph{Proposal}: \textbackslash{}RR\textasciicircum{}+\textbackslash{}bar\{\textbackslash{}theta\}\_\{\textbackslash{}ast\}\textbackslash{}xrightarrow\{\textbackslash{}sim\}\textbackslash{}RR\textasciicircum{}+\textbackslash{}varepsilon\_\{\textbackslash{}ast\}\textbackslash{}circ\textbackslash{}RR\textasciicircum{}+\textbackslash{}Gamma(\textbackslash{}theta\_\{\textbackslash{}ast\})\par
This finding intersects Orgogozo's \textbackslash{}sisi material at Remark 2.3.8, isomorphism in the first sentence; preserve that human entry unchanged. The proposed reading \textbackslash{}RR\textasciicircum{}+\textbackslash{}bar\{\textbackslash{}theta\}\_\{\textbackslash{}ast\}\textbackslash{}xrightarrow\{\textbackslash{}sim\}\textbackslash{}RR\textasciicircum{}+\textbackslash{}varepsilon\_\{\textbackslash{}ast\}\textbackslash{}circ\textbackslash{}RR\textasciicircum{}+\textbackslash{}Gamma(\textbackslash{}theta\_\{\textbackslash{}ast\}) at Remark 2.3.8, isomorphism in the first sentence changes a substantive mathematical assertion or bibliographic identification; no independent derivation of this exact replacement is established in this review.\par
\par\noindent\textbf{OM-SOL-0455 -- t2 Vbis p.113 -- SOL}\par
Conclusion of Lemma 2.4.4.2, immediately below the opening display\par
\emph{Reported reading}: est un morphisme.\par
\emph{Proposal}: est un isomorphisme.\par
A morphism can be mentioned without asserting invertibility; the stronger conclusion requires the lemma's proof.\par
\par\noindent\textbf{OM-SOL-0463 -- t2 Vbis p.124 -- SOL}\par
Proof of Lemma 3.0.2.4, index of the projective limit\par
\emph{Reported reading}: q \textbackslash{}geq n\par
\emph{Proposal}: q \textbackslash{}leq n\par
The proposed reading q \textbackslash{}leq n at Proof of Lemma 3.0.2.4, index of the projective limit changes a substantive mathematical assertion or bibliographic identification; no independent derivation of this exact replacement is established in this review.\par
\par\noindent\textbf{OM-SOL-0464 -- t2 Vbis p.124 -- SOL}\par
Proposition 3.1.1, last factor on the right of eta\_G\textasciicircum{}f\par
\emph{Reported reading}: \textbackslash{}mathbf\{L\}\textasciicircum{}\{+\}\textbackslash{}bar\{u\}\textasciicircum{}\{\textbackslash{}ast\}\par
\emph{Proposal}: \textbackslash{}mathbf\{L\}\_G\textasciicircum{}\{+\}\textbackslash{}bar\{u\}\textasciicircum{}\{\textbackslash{}ast\}\par
The replacement \textbackslash{}mathbf\{L\}\_G\textasciicircum{}\{+\}\textbackslash{}bar\{u\}\textasciicircum{}\{\textbackslash{}ast\} at Proposition 3.1.1, last factor on the right of eta\_G\textasciicircum{}f identifies a particular indexed object or symbol; that identification is not independently established by the quoted local expression.\par
\par\noindent\textbf{OM-SOL-0465 -- t2 Vbis p.127 -- SOL}\par
Première formule, source de la seconde flèche\par
\emph{Reported reading}: K\textasciicircum{}+(\textbackslash{}bar\{u\}\textasciicircum{}\textbackslash{}ast)(F\textasciicircum{}\{\textbackslash{}bigdot\}) \textbackslash{}to r\_0\textasciicircum{}\textbackslash{}ast(J\textasciicircum{}\{\textbackslash{}bigdot\})\par
\emph{Proposal}: K\textasciicircum{}+ (\{\textbackslash{}bar\{u\}'\}\textasciicircum{}\textbackslash{}ast)(F\textasciicircum{}\{\textbackslash{}bigdot\}) \textbackslash{}to r\textasciicircum{}\{\textbackslash{}ast\}\_0(J\textasciicircum{}\{\textbackslash{}bigdot\})\par
The replacement K\textasciicircum{}+ (\{\textbackslash{}bar\{u\}'\}\textasciicircum{}\textbackslash{}ast)(F\textasciicircum{}\{\textbackslash{}bigdot\}) \textbackslash{}to r\textasciicircum{}\{\textbackslash{}ast\}\_0(J\textasciicircum{}\{\textbackslash{}bigdot\}) at Première formule, source de la seconde flèche identifies a particular indexed object or symbol; that identification is not independently established by the quoted local expression.\par
\par\noindent\textbf{OM-SOL-0467 -- t2 Vbis p.132 -- SOL}\par
Target of the augmentation X f X'\par
\emph{Reported reading}: C\textasciicircum{}\{\textbackslash{}Delta'\}\_S\par
\emph{Proposal}: C\textasciicircum{}\{\textbackslash{}Delta'\}\_\{S'\}\par
The replacement C\textasciicircum{}\{\textbackslash{}Delta'\}\_\{S'\} at Target of the augmentation X f X' identifies a particular indexed object or symbol; that identification is not independently established by the quoted local expression.\par
\par\noindent\textbf{OM-SOL-0468 -- t2 Vbis p.137 -- SOL}\par
Lemme 3.3.3.2, introduction du morphisme f\par
\emph{Reported reading}: f : X \textbackslash{}to X'\par
\emph{Proposal}: f : X' \textbackslash{}to X\par
The replacement f : X' \textbackslash{}to X at Lemme 3.3.3.2, introduction du morphisme f requires a complete check of the maps, their endpoints and the composition convention; this identification is not independently verified.\par
\par\noindent\textbf{OM-SOL-0469 -- t2 Vbis p.138 -- SOL}\par
Diagramme en tête de page, facteurs des deux produits\par
\emph{Reported reading}: \textbackslash{}prod\_\{r\textbackslash{}in\textbackslash{}operatorname\{ob\}(\textbackslash{}Delta\textasciicircum{}+\_\{n+1[p]\})\}X'\_2 ; \textbackslash{}prod\_\{\textbackslash{}operatorname\{ob\}(\textbackslash{}Delta\textasciicircum{}+\_\{n+1[p]\})\}X\_2\par
\emph{Proposal}: \textbackslash{}prod\_\{r\textbackslash{}in\textbackslash{}operatorname\{ob\}(\textbackslash{}Delta\textasciicircum{}+\_\{n+1[p]\})\}X'\_r \textbackslash{}quad\textbackslash{}text\{et\}\textbackslash{}quad \textbackslash{}prod\_\{r\textbackslash{}in\textbackslash{}operatorname\{ob\}(\textbackslash{}Delta\textasciicircum{}+\_\{n+1[p]\})\}X\_r\par
The replacement \textbackslash{}prod\_\{r\textbackslash{}in\textbackslash{}operatorname\{ob\}(\textbackslash{}Delta\textasciicircum{}+\_\{n+1[p]\})\}X'\_r \textbackslash{}quad\textbackslash{}text\{et\}\textbackslash{}quad \textbackslash{}prod\_\{r\textbackslash{}in\textbackslash{}operatorname\{ob\}(\textbackslash{}Delta\textasciicircum{}+\_\{n+1[p]\})\}X\_r at Diagramme en tête de page, facteurs des deux produits identifies a particular indexed object or symbol; that identification is not independently established by the quoted local expression.\par
\par\noindent\textbf{OM-SOL-0470 -- t2 Vbis p.140 -- SOL}\par
Début de page, images réciproques de l'isomorphisme de descente\par
\emph{Reported reading}: de ces isomorphismes\par
\emph{Proposal}: de cet isomorphisme\par
Singularizing these isomorphisms requires identifying all isomorphisms referred to in the preceding sentence.\par
\par\noindent\textbf{OM-SOL-0471 -- t2 Vbis p.147 -- SOL}\par
Triangle de factorisation, sommet inférieur gauche\par
\emph{Reported reading}: X\par
\emph{Proposal}: X'\par
A second name for the auxiliary scheme is plausible, but the proposed prime requires comparison with the whole factorization.\par
\par\noindent\textbf{OM-SOL-0474 -- t2 Vbis p.149 -- SOL}\par
Démonstration de 5.1.3, « On pose »\par
\emph{Reported reading}: Y'\_n =\allowbreak  (\textbackslash{}cosk\_n Y)\_\{n+1\}\par
\emph{Proposal}: Y'\_\{n+1\} =\allowbreak  (\textbackslash{}cosk\_n Y)\_\{n+1\}\par
The replacement Y'\_\{n+1\} =\allowbreak  (\textbackslash{}cosk\_n Y)\_\{n+1\} at Démonstration de 5.1.3, « On pose » requires a complete check of the maps, their endpoints and the composition convention; this identification is not independently verified.\par
\par\noindent\textbf{OM-SOL-0476 -- t2 Vbis p.156 -- SOL}\par
Démonstration de 5.2.5, trois flèches horizontales du diagramme\par
\emph{Reported reading}: Les trois flèches vont des termes primés, à gauche, vers les termes non primés, à droite.\par
\emph{Proposal}: \$H\textasciicircum{}\{\textbackslash{}Zar.\}\_\{DR\}(u')\textbackslash{}longleftarrow H\_\{DR\}\textasciicircum{}\{\textbackslash{}Zar.\}(u)\$, \$H\textasciicircum{}\{\textbackslash{}bigdot\}(S'\textasciicircum{}\{\textbackslash{}an\},\textbackslash{}CC)\textbackslash{}xleftarrow\{H\textasciicircum{}\{\textbackslash{}bigdot\}(t\textasciicircum{}\{\textbackslash{}an\},\textbackslash{}CC)\}H\textasciicircum{}\{\textbackslash{}bigdot\}(S\textasciicircum{}\{\textbackslash{}an\},\textbackslash{}CC)\$, \$H\textasciicircum{}\{\textbackslash{}Zar.\}\_\{DR\}(v')\textbackslash{}longleftarrow H\_\{DR\}\textasciicircum{}\{\textbackslash{}Zar.\}(v)\$\par
The replacement \$H\textasciicircum{}\{\textbackslash{}Zar.\}\_\{DR\}(u')\textbackslash{}longleftarrow H\_\{DR\}\textasciicircum{}\{\textbackslash{}Zar.\}(u)\$, \$H\textasciicircum{}\{\textbackslash{}bigdot\}(S'\textasciicircum{}\{\textbackslash{}an\},\textbackslash{}CC)\textbackslash{}xleftarrow\{H\textasciicircum{}\{\textbackslash{}bigdot\}(t\textasciicircum{}\{\textbackslash{}an\},\textbackslash{}CC)\}H\textasciicircum{}\{\textbackslash{}bigdot\}(S\textasciicircum{}\{\textbackslash{}an\},\textbackslash{}CC)\$, \$H\textasciicircum{}\{\textbackslash{}Zar.\}\_\{DR\}(v')\textbackslash{}longleftarrow H\_\{DR\}\textasciicircum{}\{\textbackslash{}Zar.\}(v)\$ at Démonstration de 5.2.5, trois flèches horizontales du diagramme requires a complete check of the maps, their endpoints and the composition convention; this identification is not independently verified.\par
\par\noindent\textbf{OM-SOL-0477 -- t2 Vbis p.158 -- SOL}\par
Deuxième alinéa marqué +, définition des morphismes\par
\emph{Reported reading}: un morphisme de (\ensuremath{\bar{X}}, X, i) dans (\ensuremath{\bar{X}}', X', i') au-dessus de t\par
\emph{Proposal}: un morphisme de \$(\textbackslash{}bar\{X\}', X', i')\$ dans \$(\textbackslash{}bar\{X\}, X, i)\$ au-dessus de \$t\$\par
The replacement un morphisme de \$(\textbackslash{}bar\{X\}', X', i')\$ dans \$(\textbackslash{}bar\{X\}, X, i)\$ au-dessus de \$t\$ at Deuxième alinéa marqué +, définition des morphismes requires a complete check of the maps, their endpoints and the composition convention; this identification is not independently verified.\par
\par\noindent\textbf{OM-SOL-0478 -- t2 Vbis p.158 -- SOL}\par
N.B., membre droit de l'égalité\par
\emph{Reported reading}: \ensuremath{\bar{h}}⁻¹(i(X)) =\allowbreak  i(X')\par
\emph{Proposal}: \$\textbackslash{}bar\{h\}\textasciicircum{}\{-1\}(i(X)) =\allowbreak  i'(X')\$\par
The replacement \$\textbackslash{}bar\{h\}\textasciicircum{}\{-1\}(i(X)) =\allowbreak  i'(X')\$ at N.B., membre droit de l'égalité identifies a particular indexed object or symbol; that identification is not independently established by the quoted local expression.\par
\par\noindent\textbf{OM-SOL-0479 -- t2 Vbis p.160 -- SOL}\par
Premier grand diagramme, libellé de \ensuremath{\bar{Z}}' \ensuremath{\to} \ensuremath{\bar{X}}'\par
\emph{Reported reading}: h', sans barre\par
\emph{Proposal}: \$\textbackslash{}bar\{h\}'\$\par
The replacement \$\textbackslash{}bar\{h\}'\$ at Premier grand diagramme, libellé de \ensuremath{\bar{Z}}' \ensuremath{\to} \ensuremath{\bar{X}}' requires a complete check of the maps, their endpoints and the composition convention; this identification is not independently verified.\par
\par\noindent\textbf{OM-SOL-0480 -- t2 Vbis p.162 -- SOL}\par
Bibliography [1], first author's surname\par
\emph{Reported reading}: T. Blum\par
\emph{Proposal}: T. Bloom\par
The author-name correction requires a bibliographic authority beyond the quoted entry.\par
\par\noindent\textbf{OM-SOL-0482 -- t2 VI p.169 -- SOL}\par
Preuve du corollaire 1.4, les deux indices de coproducts\par
\emph{Reported reading}: \textbackslash{}amalg\_\{j=\allowbreak J\} (deux occurrences)\par
\emph{Proposal}: \textbackslash{}amalg\_\{j\textbackslash{}in J\}\par
The replacement \textbackslash{}amalg\_\{j\textbackslash{}in J\} at Preuve du corollaire 1.4, les deux indices de coproducts identifies a particular indexed object or symbol; that identification is not independently established by the quoted local expression.\par
\par\noindent\textbf{OM-SOL-0483 -- t2 VI p.170 -- SOL}\par
1.5.3, reformulation du recouvrement fini par les générateurs\par
\emph{Reported reading}: isomorphe à un objet quotient d'une somme finie d'objets de E\par
\emph{Proposal}: isomorphe à un objet quotient d'une somme finie d'objets pris parmi les \$X\_i\$\par
The proposed reading isomorphe à un objet quotient d'une somme finie d'objets pris parmi les \$X\_i\$ at 1.5.3, reformulation du recouvrement fini par les générateurs changes a substantive mathematical assertion or bibliographic identification; no independent derivation of this exact replacement is established in this review.\par
\par\noindent\textbf{OM-SOL-0484 -- t2 VI p.172 -- SOL}\par
Proposition 1.8 (ii), definition of the base-change morphism\par
\emph{Reported reading}: \$f' : X' =\allowbreak  X\textbackslash{}times\_Y Y' \textbackslash{}to\$ has no target.\par
\emph{Proposal}: \$f' : X' =\allowbreak  \textbackslash{}fprod\{X\}\{Y'\}\{Y\} \textbackslash{}to Y'\$\par
The proposed reading \$f' : X' =\allowbreak  \textbackslash{}fprod\{X\}\{Y'\}\{Y\} \textbackslash{}to Y'\$ at Proposition 1.8 (ii), definition of the base-change morphism changes a substantive mathematical assertion or bibliographic identification; no independent derivation of this exact replacement is established in this review.\par
\par\noindent\textbf{OM-SOL-0490 -- t2 VI p.190 -- SOL}\par
Omitted clause beginning « donc X\_0 est limite »\par
\emph{Reported reading}: donc \$X\_0\$ est limite inductive filtrante\par
\emph{Proposal}: donc \$X\$ est limite inductive filtrante\par
The record documents a lacuna at Omitted clause beginning « donc X\_0 est limite »; it does not uniquely establish the proposed completion donc \$X\$ est limite inductive filtrante.\par
\par\noindent\textbf{OM-SOL-0491 -- t2 VI p.190 -- SOL}\par
Definition of the quotients Coker(R\_i \ensuremath{\rightrightarrows} X\_0)\par
\emph{Reported reading}: \$\textbackslash{}Coker(R\_i \textbackslash{}rightrightarrows X\_0) =\allowbreak  X\_i\$\par
\emph{Proposal}: \$\textbackslash{}Coker (R\_i \textbackslash{}rightrightarrows X\_0) =\allowbreak  X'\_i\$\par
The replacement \$\textbackslash{}Coker (R\_i \textbackslash{}rightrightarrows X\_0) =\allowbreak  X'\_i\$ at Definition of the quotients Coker(R\_i \ensuremath{\rightrightarrows} X\_0) identifies a particular indexed object or symbol; that identification is not independently established by the quoted local expression.\par
\par\noindent\textbf{OM-SOL-0492 -- t2 VI p.190 -- SOL}\par
Conclusion of the proof of b), beginning « d'où X\_i »\par
\emph{Reported reading}: \$X\_i \textbackslash{}simeq \textbackslash{}Coker(X\_1 \textbackslash{}rightrightarrows X\_0)\$\par
\emph{Proposal}: \$X'\_i \textbackslash{}simeq \textbackslash{}Coker (X\_1 \textbackslash{}rightrightarrows X\_0)\$\par
The replacement \$X'\_i \textbackslash{}simeq \textbackslash{}Coker (X\_1 \textbackslash{}rightrightarrows X\_0)\$ at Conclusion of the proof of b), beginning « d'où X\_i » identifies a particular indexed object or symbol; that identification is not independently established by the quoted local expression.\par
\par\noindent\textbf{OM-SOL-0496 -- t2 VI p.194 -- SOL}\par
Opening line, automorphism example continued from Exercice 1.28 f)\par
\emph{Reported reading}: f(x)=\allowbreak x, f(y)=\allowbreak x+y\par
\emph{Proposal}: f(x) =\allowbreak  y,\textbackslash{}qquad f(y) =\allowbreak  x + y\textasciicircum{}2\par
Over F\_2 the printed map (x,y)↦(x,x+y) squares to the identity, so its claimed infinite order fails. The proposed quadratic replacement introduces a different example and is not uniquely recovered from this text.\par
\par\noindent\textbf{OM-SOL-0504 -- t2 VI p.214 -- SOL}\par
Démonstration de 2.10, sous-catégorie considérée\par
\emph{Reported reading}: objets noethériens quasi-séparés\par
\emph{Proposal}: objets prénoethériens quasi-séparés\par
The prefix pré changes the finiteness condition; the proof must be compared with condition (ii).\par
\par\noindent\textbf{OM-SOL-0505 -- t2 VI p.214 -- SOL}\par
Démonstration de 2.10, application de la première assertion de 1.32\par
\emph{Reported reading}: donc noethérien par la première assertion de 1.32\par
\emph{Proposal}: donc prénoethérien\par
The proposed reading donc prénoethérien at Démonstration de 2.10, application de la première assertion de 1.32 changes a substantive mathematical assertion or bibliographic identification; no independent derivation of this exact replacement is established in this review.\par
\par\noindent\textbf{OM-SOL-0507 -- t2 VI p.218 -- SOL}\par
Exercice 2.16 b), continuation, definition of T\_U\par
\emph{Reported reading}: \$T\_U =\allowbreak  G \textbackslash{}times U\$\par
\emph{Proposal}: \$T\_U =\allowbreak  U \textbackslash{}times G\$\par
The proposed reading \$T\_U =\allowbreak  U \textbackslash{}times G\$ at Exercice 2.16 b), continuation, definition of T\_U changes a substantive mathematical assertion or bibliographic identification; no independent derivation of this exact replacement is established in this review.\par
\par\noindent\textbf{OM-SOL-0508 -- t2 VI p.218 -- SOL}\par
Exercice 2.16 d), first assertion\par
\emph{Reported reading}: P \ensuremath{\to} e est quasi-compact si et seulement si G est fini.\par
\emph{Proposal}: Le morphisme structural \$P \textbackslash{}to e\$ est quasi-compact si et seulement si \$G\$ est fini ou \$X\$ vide.\par
The proposed reading Le morphisme structural \$P \textbackslash{}to e\$ est quasi-compact si et seulement si \$G\$ est fini ou \$X\$ vide. at Exercice 2.16 d), first assertion changes a substantive mathematical assertion or bibliographic identification; no independent derivation of this exact replacement is established in this review.\par
\par\noindent\textbf{OM-SOL-0510 -- t2 VI p.221 -- SOL}\par
e), target of the parallel morphisms f and g\par
\emph{Reported reading}: \$f,g : Y \textbackslash{}rightrightarrows Y\$\par
\emph{Proposal}: \$f\$, \$g : Y \textbackslash{}rightrightarrows Z\$\par
The proposed reading \$f\$, \$g : Y \textbackslash{}rightrightarrows Z\$ at e), target of the parallel morphisms f and g changes a substantive mathematical assertion or bibliographic identification; no independent derivation of this exact replacement is established in this review.\par
\par\noindent\textbf{OM-SOL-0511 -- t2 VI p.222 -- SOL}\par
h), opening factorisation hypothesis\par
\emph{Reported reading}: \$f : X \textbackslash{}to X\$\par
\emph{Proposal}: \$f: X \textbackslash{}to Y\$\par
The replacement \$f: X \textbackslash{}to Y\$ at h), opening factorisation hypothesis requires a complete check of the maps, their endpoints and the composition convention; this identification is not independently verified.\par
\par\noindent\textbf{OM-SOL-0512 -- t2 VI p.227 -- SOL}\par
Proof of 3.3, topoi generated by the epsilon images\par
\emph{Reported reading}: de \$C\$ (resp. de \$C'\$), without tildes\par
\emph{Proposal}: de \$C\textasciicircum{}\{\textbackslash{}sim\}\$ (\textbackslash{}resp de \$\{C'\}\textasciicircum{}\{\textbackslash{}sim\}\$)\par
The replacement de \$C\textasciicircum{}\{\textbackslash{}sim\}\$ (\textbackslash{}resp de \$\{C'\}\textasciicircum{}\{\textbackslash{}sim\}\$) at Proof of 3.3, topoi generated by the epsilon images identifies a particular indexed object or symbol; that identification is not independently established by the quoted local expression.\par
\par\noindent\textbf{OM-SOL-0513 -- t2 VI p.228 -- SOL}\par
3.5 (ii ter), « défini par »\par
\emph{Reported reading}: défini par \$v\_i\$\par
\emph{Proposal}: défini par \$v\_\{i'\}\$\par
The replacement défini par \$v\_\{i'\}\$ at 3.5 (ii ter), « défini par » requires a complete check of the maps, their endpoints and the composition convention; this identification is not independently verified.\par
\par\noindent\textbf{OM-SOL-0514 -- t2 VI p.228 -- SOL}\par
Proof of 3.5, final occurrence of the second projection\par
\emph{Reported reading}: \$\textbackslash{}pr\_2 : \textbackslash{}fprod\{f\textasciicircum{}\{\textbackslash{}ast\}(X)\}\{f\textasciicircum{}\{\textbackslash{}ast\}(X\_i)\}\{f\textasciicircum{}\{\textbackslash{}ast\}(Y)\}\$ est quasi-compact\par
\emph{Proposal}: \textbackslash{}pr\_2 : \textbackslash{}fprod\{f\textasciicircum{}\{\textbackslash{}ast\}(X)\}\{f\textasciicircum{}\{\textbackslash{}ast\}(X\_i)\}\{f\textasciicircum{}\{\textbackslash{}ast\}(Y)\} \textbackslash{}to f\textasciicircum{}\{\textbackslash{}ast\}(X\_i)\par
The replacement \textbackslash{}pr\_2 : \textbackslash{}fprod\{f\textasciicircum{}\{\textbackslash{}ast\}(X)\}\{f\textasciicircum{}\{\textbackslash{}ast\}(X\_i)\}\{f\textasciicircum{}\{\textbackslash{}ast\}(Y)\} \textbackslash{}to f\textasciicircum{}\{\textbackslash{}ast\}(X\_i) at Proof of 3.5, final occurrence of the second projection requires a complete check of the maps, their endpoints and the composition convention; this identification is not independently verified.\par
\par\noindent\textbf{OM-SOL-0515 -- t2 VI p.229 -- SOL}\par
Proof of 3.5, « même critère », argument of the inverse-image functor\par
\emph{Reported reading}: \$f\textasciicircum{}\{\textbackslash{}ast\}\_v(X')\$\par
\emph{Proposal}: f\textasciicircum{}\{\textbackslash{}ast\}\_v(X)\par
The replacement f\textasciicircum{}\{\textbackslash{}ast\}\_v(X) at Proof of 3.5, « même critère », argument of the inverse-image functor identifies a particular indexed object or symbol; that identification is not independently established by the quoted local expression.\par
\par\noindent\textbf{OM-SOL-0518 -- t2 VI p.237 -- SOL}\par
Proof, « identifiant (par i\_*) E' »\par
\emph{Reported reading}: identifiant (par \$i\_\{\textbackslash{}ast\}\$) \$E'\$\par
\emph{Proposal}: \$E''\$\par
The replacement \$E''\$ at Proof, « identifiant (par i\_*) E' » identifies a particular indexed object or symbol; that identification is not independently established by the quoted local expression.\par
\par\noindent\textbf{OM-SOL-0519 -- t2 VI p.238 -- SOL}\par
Opening paragraph, reference after « celle annoncée dans »\par
\emph{Reported reading}: 3.14.1\par
\emph{Proposal}: 4.2.1\par
The proposed reference 4.2.1 at Opening paragraph, reference after « celle annoncée dans » needs comparison with the statement actually invoked; that target has not been independently established here.\par
\par\noindent\textbf{OM-SOL-0520 -- t2 VI p.238 -- SOL}\par
Proof of Corollaire 4.3 a), « en vertu de 3.14 »\par
\emph{Reported reading}: 3.14\par
\emph{Proposal}: 4.2\par
The proposed reference 4.2 at Proof of Corollaire 4.3 a), « en vertu de 3.14 » needs comparison with the statement actually invoked; that target has not been independently established here.\par
\par\noindent\textbf{OM-SOL-0521 -- t2 VI p.239 -- SOL}\par
First fiber product inside i\textasciicircum{}*, following « Mais alors »\par
\emph{Reported reading}: i\textasciicircum{}\{\textbackslash{}ast\}(X\_\{\textbackslash{}alpha\}\textbackslash{}times\_\{i\_\{\textbackslash{}ast\}(Y)\}X''\_\{\textbackslash{}beta\})\par
\emph{Proposal}: \$X'\_\{\textbackslash{}alpha\}\$\par
The replacement \$X'\_\{\textbackslash{}alpha\}\$ at First fiber product inside i\textasciicircum{}*, following « Mais alors » identifies a particular indexed object or symbol; that identification is not independently established by the quoted local expression.\par
\par\noindent\textbf{OM-SOL-0528 -- t2 VI p.252 -- SOL}\par
(iv bis), catégorie contenant X''\par
\emph{Reported reading}: \$X''\textbackslash{}in\textbackslash{}Ob E'\_\{\textbackslash{}coh\}\$\par
\emph{Proposal}: \$X''\textbackslash{}in \textbackslash{}Ob E''\_\{\textbackslash{}coh\}\$\par
The replacement \$X''\textbackslash{}in \textbackslash{}Ob E''\_\{\textbackslash{}coh\}\$ at (iv bis), catégorie contenant X'' identifies a particular indexed object or symbol; that identification is not independently established by the quoted local expression.\par
\par\noindent\textbf{OM-SOL-0529 -- t2 VI p.252 -- SOL}\par
(iv bis), flèche de Y' vers X'\par
\emph{Reported reading}: \$v:Y'\textbackslash{}to X'\$\par
\emph{Proposal}: \$v' : Y'\textbackslash{}to X'\$\par
The replacement \$v' : Y'\textbackslash{}to X'\$ at (iv bis), flèche de Y' vers X' requires a complete check of the maps, their endpoints and the composition convention; this identification is not independently verified.\par
\par\noindent\textbf{OM-SOL-0530 -- t2 VI p.252 -- SOL}\par
(iv bis), condition sur Y'\par
\emph{Reported reading}: il existe une flèche v: Y' \ensuremath{\to} X' dans E' telle que u se factorise\par
\emph{Proposal}: \$Y'\textbackslash{}in \textbackslash{}Ob E'\_\{\textbackslash{}coh\}\$\par
The proposed reading \$Y'\textbackslash{}in \textbackslash{}Ob E'\_\{\textbackslash{}coh\}\$ at (iv bis), condition sur Y' changes a substantive mathematical assertion or bibliographic identification; no independent derivation of this exact replacement is established in this review.\par
\par\noindent\textbf{OM-SOL-0532 -- t2 VI p.253 -- SOL}\par
d), membre de phrase omis concernant les points\par
\emph{Reported reading}: pour tout point p de X''\par
\emph{Proposal}: pour tout point \$p\$ de \$E''\$\par
The proposed reading pour tout point \$p\$ de \$E''\$ at d), membre de phrase omis concernant les points changes a substantive mathematical assertion or bibliographic identification; no independent derivation of this exact replacement is established in this review.\par
\par\noindent\textbf{OM-SOL-0533 -- t2 VI p.253 -- SOL}\par
d), référence à l'appendice\par
\emph{Reported reading}: l'appendice (10. )\par
\emph{Proposal}: (9.0)\par
The proposed reference (9.0) at d), référence à l'appendice needs comparison with the statement actually invoked; that target has not been independently established here.\par
\par\noindent\textbf{OM-SOL-0534 -- t2 VI p.253 -- SOL}\par
d), « un topos cohérent … a assez de points »\par
\emph{Reported reading}: un topos cohérent E'\par
\emph{Proposal}: un topos cohérent \$E''\$\par
The replacement un topos cohérent \$E''\$ at d), « un topos cohérent … a assez de points » identifies a particular indexed object or symbol; that identification is not independently established by the quoted local expression.\par
\par\noindent\textbf{OM-SOL-0535 -- t2 VI p.253 -- SOL}\par
d), dernière phrase avant l'équivalence\par
\emph{Reported reading}: \$X'\textbackslash{}in\textbackslash{}Ob E\$\par
\emph{Proposal}: \$X'\textbackslash{}in \textbackslash{}Ob E'\$\par
The replacement \$X'\textbackslash{}in \textbackslash{}Ob E'\$ at d), dernière phrase avant l'équivalence identifies a particular indexed object or symbol; that identification is not independently established by the quoted local expression.\par
\par\noindent\textbf{OM-SOL-0538 -- t2 VI p.256 -- SOL}\par
Catégorie des valeurs du préfaisceau \ensuremath{\alpha} \ensuremath{\circ} F\par
\emph{Reported reading}: \textbackslash{}Pro(E'')\textasciicircum{}\{0\}\par
\emph{Proposal}: \textbackslash{}Pro(E')\textasciicircum{}\{0\}\par
The replacement \textbackslash{}Pro(E')\textasciicircum{}\{0\} at Catégorie des valeurs du préfaisceau \ensuremath{\alpha} \ensuremath{\circ} F identifies a particular indexed object or symbol; that identification is not independently established by the quoted local expression.\par
\par\noindent\textbf{OM-SOL-0539 -- t2 VI p.257 -- SOL}\par
Début de page, définition du foncteur de composition\par
\emph{Reported reading}: g \textbackslash{}mapsto \textbackslash{}alpha \textbackslash{}circ g'\par
\emph{Proposal}: \$g \textbackslash{}mapsto \textbackslash{}alpha \textbackslash{}circ g\$\par
The replacement \$g \textbackslash{}mapsto \textbackslash{}alpha \textbackslash{}circ g\$ at Début de page, définition du foncteur de composition identifies a particular indexed object or symbol; that identification is not independently established by the quoted local expression.\par
\par\noindent\textbf{OM-SOL-0540 -- t2 VI p.258 -- SOL}\par
Premier alinéa, identification de la catégorie de foncteurs\par
\emph{Reported reading}: E'\_\{\textbackslash{}coh\} =\allowbreak  \textbackslash{}Hom(E'\_\{\textbackslash{}coh\}, (\textbackslash{}Ens))\par
\emph{Proposal}: \$\textbackslash{}overset\{\textbackslash{}vee\}\{E'\}\_\{\textbackslash{}coh\} =\allowbreak  \textbackslash{}Hom(E'\_\{\textbackslash{}coh\}, (\textbackslash{}Ens))\$\par
The proposed reading \$\textbackslash{}overset\{\textbackslash{}vee\}\{E'\}\_\{\textbackslash{}coh\} =\allowbreak  \textbackslash{}Hom(E'\_\{\textbackslash{}coh\}, (\textbackslash{}Ens))\$ at Premier alinéa, identification de la catégorie de foncteurs changes a substantive mathematical assertion or bibliographic identification; no independent derivation of this exact replacement is established in this review.\par
\par\noindent\textbf{OM-SOL-0542 -- t2 VI p.259 -- SOL}\par
Condition de représentabilité des limites finies dans P\par
\emph{Reported reading}: \textbackslash{}varinjlim finies\par
\emph{Proposal}: \$\textbackslash{}varprojlim\$ finies\par
The replacement \$\textbackslash{}varprojlim\$ finies at Condition de représentabilité des limites finies dans P requires a complete check of the maps, their endpoints and the composition convention; this identification is not independently verified.\par
\par\noindent\textbf{OM-SOL-0544 -- t2 VI p.260 -- SOL}\par
Preuve de 5.1, notation de la sous-catégorie génératrice\par
\emph{Reported reading}: Soit E\_\{cohalg\} … ; puis (X \ensuremath{\in} ob C) et « tout objet X de C »\par
\emph{Proposal}: Soit \$C =\allowbreak  E\_\{\textbackslash{}text\{cohalg\}\}\$ la sous-catégorie pleine de \$E\$\par
The proposed reading Soit \$C =\allowbreak  E\_\{\textbackslash{}text\{cohalg\}\}\$ la sous-catégorie pleine de \$E\$ at Preuve de 5.1, notation de la sous-catégorie génératrice changes a substantive mathematical assertion or bibliographic identification; no independent derivation of this exact replacement is established in this review.\par
\par\noindent\textbf{OM-SOL-0545 -- t2 VI p.261 -- SOL}\par
Lemme 5.4, « par des objets de C »\par
\emph{Reported reading}: par des objets de \$C\textasciicircum{}\{\textbackslash{}bullet\}\$\par
\emph{Proposal}: \$C\$\par
The replacement \$C\$ at Lemme 5.4, « par des objets de C » identifies a particular indexed object or symbol; that identification is not independently established by the quoted local expression.\par
\par\noindent\textbf{OM-SOL-0548 -- t2 VI p.262 -- SOL}\par
Proposition 5.7, preuve, formule après « On a »\par
\emph{Reported reading}: \textbackslash{}Eext\textasciicircum{}i\_A(G,F) : \textbackslash{}Hh\textasciicircum{}i(\textbackslash{}Hhom\_A(L.,F))\par
\emph{Proposal}: \textbackslash{}Eext\textasciicircum{}\{i\}\_A(G, F) =\allowbreak  \textbackslash{}Hh\textasciicircum{}\{i\} (\textbackslash{}Hhom\_A(L., F))\par
The proposed reading \textbackslash{}Eext\textasciicircum{}\{i\}\_A(G, F) =\allowbreak  \textbackslash{}Hh\textasciicircum{}\{i\} (\textbackslash{}Hhom\_A(L., F)) at Proposition 5.7, preuve, formule après « On a » changes a substantive mathematical assertion or bibliographic identification; no independent derivation of this exact replacement is established in this review.\par
\par\noindent\textbf{OM-SOL-0549 -- t2 VI p.263 -- SOL}\par
Corollaire 5.10, borne sur i\par
\emph{Reported reading}: G de q-présentation finie ; i(i-1)/2 \textbackslash{}leq r-1\par
\emph{Proposal}: \textbackslash{}dfrac\{i(i - 1)\}\{2\} \textbackslash{}leq q - 1\par
The proposed reading \textbackslash{}dfrac\{i(i - 1)\}\{2\} \textbackslash{}leq q - 1 at Corollaire 5.10, borne sur i changes a substantive mathematical assertion or bibliographic identification; no independent derivation of this exact replacement is established in this review.\par
\par\noindent\textbf{OM-SOL-0551 -- t2 VI p.265 -- SOL}\par
Formule (6.1.2.2), premier facteur du membre de droite\par
\emph{Reported reading}: C\_\{gh,f\}\par
\emph{Proposal}: C\_\{gf, h\}\par
The replacement C\_\{gf, h\} at Formule (6.1.2.2), premier facteur du membre de droite identifies a particular indexed object or symbol; that identification is not independently established by the quoted local expression.\par
\par\noindent\textbf{OM-SOL-0552 -- t2 VI p.265 -- SOL}\par
6.1.3, isomorphisme de foncteurs C\_\{f,g\}\par
\emph{Reported reading}: C\_\{f,g\}: f\textasciicircum{}*g\textasciicircum{}* \textbackslash{}to (fg)\textasciicircum{}*\par
\emph{Proposal}: C\_\{f, g\} : f\textasciicircum{}\{\textbackslash{}ast\} g\textasciicircum{}\{\textbackslash{}ast\} \textbackslash{}to (gf)\textasciicircum{}\{\textbackslash{}ast\}\par
The proposed reading C\_\{f, g\} : f\textasciicircum{}\{\textbackslash{}ast\} g\textasciicircum{}\{\textbackslash{}ast\} \textbackslash{}to (gf)\textasciicircum{}\{\textbackslash{}ast\} at 6.1.3, isomorphisme de foncteurs C\_\{f,g\} changes a substantive mathematical assertion or bibliographic identification; no independent derivation of this exact replacement is established in this review.\par
\par\noindent\textbf{OM-SOL-0553 -- t2 VI p.267 -- SOL}\par
Description of the pseudo-functor, after the triangular diagram\par
\emph{Reported reading}: tout ξ \ensuremath{\in} ob, F\_ξ\par
\emph{Proposal}: tout \$\textbackslash{}xi \textbackslash{}in \textbackslash{}ob(\textbackslash{}Ee)\$, \$\textbackslash{}Ff\_\{\textbackslash{}xi\}\$\par
The record documents a lacuna at Description of the pseudo-functor, after the triangular diagram; it does not uniquely establish the proposed completion tout \$\textbackslash{}xi \textbackslash{}in \textbackslash{}ob(\textbackslash{}Ee)\$, \$\textbackslash{}Ff\_\{\textbackslash{}xi\}\$.\par
\par\noindent\textbf{OM-SOL-0555 -- t2 VI p.270 -- SOL}\par
Exercise 6.6, definition of S(X)\par
\emph{Reported reading}: on note S(X) la catégorie des morphismes de but X.\par
\emph{Proposal}: on note \$S(X)\$ la catégorie des morphismes de \$S\$ de but \$X\$.\par
The replacement on note \$S(X)\$ la catégorie des morphismes de \$S\$ de but \$X\$. at Exercise 6.6, definition of S(X) requires a complete check of the maps, their endpoints and the composition convention; this identification is not independently verified.\par
\par\noindent\textbf{OM-SOL-0557 -- t2 VI p.272 -- SOL}\par
Proposition 6.9, domaine du foncteur représentable\par
\emph{Reported reading}: (Cat) \textbackslash{}to (Ens)\par
\emph{Proposal}: (\textbackslash{}Cat)\textasciicircum{}\{\textbackslash{}circ\} \textbackslash{}to (\textbackslash{}Ens)\par
The replacement (\textbackslash{}Cat)\textasciicircum{}\{\textbackslash{}circ\} \textbackslash{}to (\textbackslash{}Ens) at Proposition 6.9, domaine du foncteur représentable requires a complete check of the maps, their endpoints and the composition convention; this identification is not independently verified.\par
\par\noindent\textbf{OM-SOL-0558 -- t2 VI p.275 -- SOL}\par
(7.1.6.2), facteur situé entre m\_i*(c\_\{f,g\}) et f*(b\_g)\par
\emph{Reported reading}: b\_f\textasciicircum{}\{\textbackslash{}ast\}\par
\emph{Proposal}: b\_f g\textasciicircum{}\{\textbackslash{}ast\}\par
The proposed reading b\_f g\textasciicircum{}\{\textbackslash{}ast\} at (7.1.6.2), facteur situé entre m\_i*(c\_\{f,g\}) et f*(b\_g) changes a substantive mathematical assertion or bibliographic identification; no independent derivation of this exact replacement is established in this review.\par
\par\noindent\textbf{OM-SOL-0559 -- t2 VI p.279 -- SOL}\par
7.2.2, sous-ensemble d'objets définissant la catégorie Cont\par
\emph{Reported reading}: \$\textbackslash{}underline\{\textbackslash{}mathrm\{Cont\}\}\_I(D,C)\textbackslash{}subset\textbackslash{}mathrm\{Hom\}\_I(D,C)\$\par
\emph{Proposal}: \$\textbackslash{}Cont\_I (D, C) \textbackslash{}subset \textbackslash{}Hom\_I (D, C)\$\par
The replacement \$\textbackslash{}Cont\_I (D, C) \textbackslash{}subset \textbackslash{}Hom\_I (D, C)\$ at 7.2.2, sous-ensemble d'objets définissant la catégorie Cont identifies a particular indexed object or symbol; that identification is not independently established by the quoted local expression.\par
\par\noindent\textbf{OM-SOL-0560 -- t2 VI p.279 -- SOL}\par
7.2.4, représentabilité des limites projectives\par
\emph{Reported reading}: dans les catégories fibrées\par
\emph{Proposal}: dans les catégories fibres\par
Changing fibrées to fibres changes the category referred to and needs the surrounding fibre construction.\par
\par\noindent\textbf{OM-SOL-0561 -- t2 VI p.281 -- SOL}\par
Sentence following (7.2.6.5)\par
\emph{Reported reading}: condition de cocycles de 6.2.2.2\par
\emph{Proposal}: condition de cocycles de 6.1.2.2\par
The proposed reference condition de cocycles de 6.1.2.2 at Sentence following (7.2.6.5) needs comparison with the statement actually invoked; that target has not been independently established here.\par
\par\noindent\textbf{OM-SOL-0562 -- t2 VI p.281 -- SOL}\par
Paragraph following (7.2.6.7), « que epsilon … est un morphisme »\par
\emph{Reported reading}: que \textbackslash{}epsilon\_\{c/I\} est\par
\emph{Proposal}: \textbackslash{}epsilon\_\{C/I\}\par
The replacement \textbackslash{}epsilon\_\{C/I\} at Paragraph following (7.2.6.7), « que epsilon … est un morphisme » identifies a particular indexed object or symbol; that identification is not independently established by the quoted local expression.\par
\par\noindent\textbf{OM-SOL-0563 -- t2 VI p.283 -- SOL}\par
Sentence immediately following (7.3.1.4)\par
\emph{Reported reading}: qui induise sur les topos fibrés les morphismes\par
\emph{Proposal}: qui induise sur les topos fibres les morphismes\par
Changing fibrés to fibres changes which topoi carry the induced maps and needs the localization definition.\par
\par\noindent\textbf{OM-SOL-0567 -- t2 VI p.290 -- SOL}\par
Diagram (7.4.5.1), lower-left vertex\par
\emph{Reported reading}: \$C\_j\$\par
\emph{Proposal}: \$C\textasciicircum{}\{\textbackslash{}sim\}\_j\$\par
The replacement \$C\textasciicircum{}\{\textbackslash{}sim\}\_j\$ at Diagram (7.4.5.1), lower-left vertex requires a complete check of the maps, their endpoints and the composition convention; this identification is not independently verified.\par
\par\noindent\textbf{OM-SOL-0568 -- t2 VI p.290 -- SOL}\par
Formula (7.4.5.5)\par
\emph{Reported reading}: \$m\_Y : \textbackslash{}alpha\_\{i!\} f\textasciicircum{}\{\textbackslash{}ast\} \textbackslash{}to \textbackslash{}alpha\_\{j!\}\$\par
\emph{Proposal}: \$m\_Y : \textbackslash{}alpha\_\{i!\} f\textasciicircum{}\{\textbackslash{}ast\}(Y) \textbackslash{}to \textbackslash{}alpha\_\{j!\}(Y)\$\par
The replacement \$m\_Y : \textbackslash{}alpha\_\{i!\} f\textasciicircum{}\{\textbackslash{}ast\}(Y) \textbackslash{}to \textbackslash{}alpha\_\{j!\}(Y)\$ at Formula (7.4.5.5) requires a complete check of the maps, their endpoints and the composition convention; this identification is not independently verified.\par
\par\noindent\textbf{OM-SOL-0569 -- t2 VI p.291 -- SOL}\par
Proof of Proposition 7.4.7, canonical cartesian morphism\par
\emph{Reported reading}: f(Y) \textbackslash{}to Y\par
\emph{Proposal}: f\textasciicircum{}\{\textbackslash{}ast\}(Y) \textbackslash{}to Y\par
The replacement f\textasciicircum{}\{\textbackslash{}ast\}(Y) \textbackslash{}to Y at Proof of Proposition 7.4.7, canonical cartesian morphism requires a complete check of the maps, their endpoints and the composition convention; this identification is not independently verified.\par
\par\noindent\textbf{OM-SOL-0570 -- t2 VI p.293 -- SOL}\par
Sentence immediately before (7.4.10.1)\par
\emph{Reported reading}: I\textasciicircum{}\{\textbackslash{}circ\}-foncteur cartésien\par
\emph{Proposal}: \$I\$-foncteur cartésien\par
The proposed reading \$I\$-foncteur cartésien at Sentence immediately before (7.4.10.1) changes a substantive mathematical assertion or bibliographic identification; no independent derivation of this exact replacement is established in this review.\par
\par\noindent\textbf{OM-SOL-0571 -- t2 VI p.296 -- SOL}\par
Opening paragraph, the covering sieve\par
\emph{Reported reading}: \$R\textbackslash{}hookrightarrow X\$ de \$C\_i/X\$\par
\emph{Proposal}: D\_i/X\par
The proposed reading D\_i/X at Opening paragraph, the covering sieve changes a substantive mathematical assertion or bibliographic identification; no independent derivation of this exact replacement is established in this review.\par
\par\noindent\textbf{OM-SOL-0572 -- t2 VI p.297 -- SOL}\par
7.4.13.6, immediately after the factorization\par
\emph{Reported reading}: où \$\textbackslash{}bar\{\textbackslash{}alpha\}\_i\$ est le foncteur évident\par
\emph{Proposal}: \textbackslash{}bar\{\textbackslash{}alpha\}\_\{i!\}\par
The proposed reading \textbackslash{}bar\{\textbackslash{}alpha\}\_\{i!\} at 7.4.13.6, immediately after the factorization changes a substantive mathematical assertion or bibliographic identification; no independent derivation of this exact replacement is established in this review.\par
\par\noindent\textbf{OM-SOL-0573 -- t2 VI p.298 -- SOL}\par
7.4.13.7, list of left-exact functors\par
\emph{Reported reading}: \textbackslash{}alpha\_i\textasciicircum{}\{\textbackslash{}ast\},\textbackslash{} m\textasciicircum{}\{\textbackslash{}sim\}\textbackslash{}text\{ et \}\textbackslash{}alpha\_i\textasciicircum{}\{\textbackslash{}ast\}\par
\emph{Proposal}: \textbackslash{}alpha\_i\textasciicircum{}\{\textbackslash{}ast\},\textbackslash{} m\textasciicircum{}\{\textbackslash{}sim\}\textbackslash{}text\{ et \}\textbackslash{}alpha\_\{i!\}\par
The proposed reading \textbackslash{}alpha\_i\textasciicircum{}\{\textbackslash{}ast\},\textbackslash{} m\textasciicircum{}\{\textbackslash{}sim\}\textbackslash{}text\{ et \}\textbackslash{}alpha\_\{i!\} at 7.4.13.7, list of left-exact functors changes a substantive mathematical assertion or bibliographic identification; no independent derivation of this exact replacement is established in this review.\par
\par\noindent\textbf{OM-SOL-0574 -- t2 VI p.299 -- SOL}\par
Diagram (8.1.2.1), vertical arrow\par
\emph{Reported reading}: m\textbackslash{}times\textbackslash{}id\_I\par
\emph{Proposal}: n\textbackslash{}times\textbackslash{}id\_I\par
The replacement n\textbackslash{}times\textbackslash{}id\_I at Diagram (8.1.2.1), vertical arrow requires a complete check of the maps, their endpoints and the composition convention; this identification is not independently verified.\par
\par\noindent\textbf{OM-SOL-0575 -- t2 VI p.299 -- SOL}\par
Diagram (8.1.2.1), diagonal arrow\par
\emph{Reported reading}: m\_1\par
\emph{Proposal}: m\par
The replacement m at Diagram (8.1.2.1), diagonal arrow requires a complete check of the maps, their endpoints and the composition convention; this identification is not independently verified.\par
\par\noindent\textbf{OM-SOL-0578 -- t2 VI p.302 -- SOL}\par
8.1.5, définition de \ensuremath{\mu}′\par
\emph{Reported reading}: \$\textbackslash{}mu' : \textbackslash{}underleftarrow\{F\} \textbackslash{}times I' \textbackslash{}to F\$\par
\emph{Proposal}: \textbackslash{}mu' : \textbackslash{}underleftarrow\{F\} \textbackslash{}times I' \textbackslash{}to F\_\{\textbackslash{}phi\}\par
The replacement \textbackslash{}mu' : \textbackslash{}underleftarrow\{F\} \textbackslash{}times I' \textbackslash{}to F\_\{\textbackslash{}phi\} at 8.1.5, définition de \ensuremath{\mu}′ requires a complete check of the maps, their endpoints and the composition convention; this identification is not independently verified.\par
\par\noindent\textbf{OM-SOL-0579 -- t2 VI p.303 -- SOL}\par
(8.1.6.4), but de Θ\par
\emph{Reported reading}: \$\textbackslash{}Hom(\textbackslash{}varinjlim\_\{I\textasciicircum{}\{\textbackslash{}circ\}\}F,D)\textasciicircum{}\{\textbackslash{}circ\}\$, Hom non souligné\par
\emph{Proposal}: \textbackslash{}Hhom(\textbackslash{}varinjlim\_\{I\textasciicircum{}\{\textbackslash{}circ\}\}F, D)\textasciicircum{}\{\textbackslash{}circ\}\par
The proposed reading \textbackslash{}Hhom(\textbackslash{}varinjlim\_\{I\textasciicircum{}\{\textbackslash{}circ\}\}F, D)\textasciicircum{}\{\textbackslash{}circ\} at (8.1.6.4), but de Θ changes a substantive mathematical assertion or bibliographic identification; no independent derivation of this exact replacement is established in this review.\par
\par\noindent\textbf{OM-SOL-0581 -- t2 VI p.303 -- SOL}\par
Lemme 8.2.1, but de m′\par
\emph{Reported reading}: \$m' : D \textbackslash{}times I' \textbackslash{}to F\$\par
\emph{Proposal}: \$m': D \textbackslash{}times I'\textbackslash{}to F'\$\par
The replacement \$m': D \textbackslash{}times I'\textbackslash{}to F'\$ at Lemme 8.2.1, but de m′ identifies a particular indexed object or symbol; that identification is not independently established by the quoted local expression.\par
\par\noindent\textbf{OM-SOL-0583 -- t2 VI p.305 -- SOL}\par
Même égalité, membre droit\par
\emph{Reported reading}: \$\textbackslash{}pr\_1 \textbackslash{}circ (n\textbackslash{}pi,p)\textasciicircum{}\{\textbackslash{}sim/I\}=\allowbreak n'\$\par
\emph{Proposal}: \textbackslash{}pr\_1 \textbackslash{}circ (n\textbackslash{}pi,p)\textasciicircum{}\{\textbackslash{}sim/I\}=\allowbreak n'\textbackslash{}pi'\par
The proposed reading \textbackslash{}pr\_1 \textbackslash{}circ (n\textbackslash{}pi,p)\textasciicircum{}\{\textbackslash{}sim/I\}=\allowbreak n'\textbackslash{}pi' at Même égalité, membre droit changes a substantive mathematical assertion or bibliographic identification; no independent derivation of this exact replacement is established in this review.\par
\par\noindent\textbf{OM-SOL-0587 -- t2 VI p.310 -- SOL}\par
8.2.12, équivalence citée après les trois foncteurs \ensuremath{\alpha}\par
\emph{Reported reading}: F' \textbackslash{}simeq \textbackslash{}Hhom\_\{I\textasciicircum{}\{\textbackslash{}circ\}\}(I\textasciicircum{}\{\textbackslash{}circ\},F)\par
\emph{Proposal}: F\textasciicircum{}\{\textbackslash{}sim\} \textbackslash{}simeq \textbackslash{}Hhom\_\{I\textasciicircum{}\{\textbackslash{}circ\}\}(I\textasciicircum{}\{\textbackslash{}circ\},F')\par
The proposed reading F\textasciicircum{}\{\textbackslash{}sim\} \textbackslash{}simeq \textbackslash{}Hhom\_\{I\textasciicircum{}\{\textbackslash{}circ\}\}(I\textasciicircum{}\{\textbackslash{}circ\},F') at 8.2.12, équivalence citée après les trois foncteurs \ensuremath{\alpha} changes a substantive mathematical assertion or bibliographic identification; no independent derivation of this exact replacement is established in this review.\par
\par\noindent\textbf{OM-SOL-0588 -- t2 VI p.310 -- SOL}\par
8.2.12, premier foncteur composé introduit par « correspond par le foncteur »\par
\emph{Reported reading}: \textbackslash{}xrightarrow\{\textbackslash{}pi\_\{\textbackslash{}ast\}\} F \textbackslash{}simeq \textbackslash{}Hhom\_\{I\textasciicircum{}\{\textbackslash{}circ\}\}(I\textasciicircum{}\{\textbackslash{}circ\},F')\par
\emph{Proposal}: \textbackslash{}xrightarrow\{\textbackslash{}pi\_\{\textbackslash{}ast\}\} F\textasciicircum{}\{\textbackslash{}sim\} \textbackslash{}simeq \textbackslash{}Hhom\_\{I\textasciicircum{}\{\textbackslash{}circ\}\}(I\textasciicircum{}\{\textbackslash{}circ\},F')\par
The replacement \textbackslash{}xrightarrow\{\textbackslash{}pi\_\{\textbackslash{}ast\}\} F\textasciicircum{}\{\textbackslash{}sim\} \textbackslash{}simeq \textbackslash{}Hhom\_\{I\textasciicircum{}\{\textbackslash{}circ\}\}(I\textasciicircum{}\{\textbackslash{}circ\},F') at 8.2.12, premier foncteur composé introduit par « correspond par le foncteur » identifies a particular indexed object or symbol; that identification is not independently established by the quoted local expression.\par
\par\noindent\textbf{OM-SOL-0589 -- t2 VI p.310 -- SOL}\par
8.2.12, foncteur composé introduit par « Par suite »\par
\emph{Reported reading}: \textbackslash{}xrightarrow\{\textbackslash{}pi\_\{\textbackslash{}ast\}\} F \textbackslash{}simeq \textbackslash{}Hhom\_\{I\textasciicircum{}\{\textbackslash{}circ\}\}(I\textasciicircum{}\{\textbackslash{}circ\},F')\par
\emph{Proposal}: \textbackslash{}xrightarrow\{\textbackslash{}pi\_\{\textbackslash{}ast\}\} F\textasciicircum{}\{\textbackslash{}sim\} \textbackslash{}simeq \textbackslash{}Hhom\_\{I\textasciicircum{}\{\textbackslash{}circ\}\}(I\textasciicircum{}\{\textbackslash{}circ\},F')\par
The replacement \textbackslash{}xrightarrow\{\textbackslash{}pi\_\{\textbackslash{}ast\}\} F\textasciicircum{}\{\textbackslash{}sim\} \textbackslash{}simeq \textbackslash{}Hhom\_\{I\textasciicircum{}\{\textbackslash{}circ\}\}(I\textasciicircum{}\{\textbackslash{}circ\},F') at 8.2.12, foncteur composé introduit par « Par suite » identifies a particular indexed object or symbol; that identification is not independently established by the quoted local expression.\par
\par\noindent\textbf{OM-SOL-0590 -- t2 VI p.312 -- SOL}\par
Premier diagramme, flèche verticale\par
\emph{Reported reading}: \textbackslash{}pi(u)\par
\emph{Proposal}: \textbackslash{}pi(n)\par
The replacement \textbackslash{}pi(n) at Premier diagramme, flèche verticale requires a complete check of the maps, their endpoints and the composition convention; this identification is not independently verified.\par
\par\noindent\textbf{OM-SOL-0591 -- t2 VI p.312 -- SOL}\par
Troisième diagramme, flèche verticale de gauche\par
\emph{Reported reading}: \textbackslash{}pi(u')\par
\emph{Proposal}: \textbackslash{}pi(n')\par
The replacement \textbackslash{}pi(n') at Troisième diagramme, flèche verticale de gauche requires a complete check of the maps, their endpoints and the composition convention; this identification is not independently verified.\par
\par\noindent\textbf{OM-SOL-0592 -- t2 VI p.312 -- SOL}\par
Troisième diagramme, segment de \ensuremath{\pi}(Z') à \ensuremath{\pi}(X'')\par
\emph{Reported reading}: \textbackslash{}pi(m')\par
\emph{Proposal}: \textbackslash{}pi(m'')\par
The prime on this arrow requires matching its endpoints and the preceding factorization.\par
\par\noindent\textbf{OM-SOL-0593 -- t2 VI p.313 -- SOL}\par
8.3.5, renvoi après « commutent aux produits fibrés »\par
\emph{Reported reading}: (8.3.3)\par
\emph{Proposal}: (\textbackslash{}hyperref[VI.lem8.3.4]\{8.3.4\})\par
The proposed reference (\textbackslash{}hyperref[VI.lem8.3.4]\{8.3.4\}) at 8.3.5, renvoi après « commutent aux produits fibrés » needs comparison with the statement actually invoked; that target has not been independently established here.\par
\par\noindent\textbf{OM-SOL-0595 -- t2 VI p.315 -- SOL}\par
Premier diagramme, sommet supérieur droit\par
\emph{Reported reading}: f\textasciicircum{}\{\textbackslash{}ast\}(Z\_\{\textbackslash{}alpha\textbackslash{}beta\})\par
\emph{Proposal}: f\textasciicircum{}\{\textbackslash{}ast\}\_\{\textbackslash{}alpha\}(Z\_\{\textbackslash{}alpha\textbackslash{}beta\})\par
The replacement f\textasciicircum{}\{\textbackslash{}ast\}\_\{\textbackslash{}alpha\}(Z\_\{\textbackslash{}alpha\textbackslash{}beta\}) at Premier diagramme, sommet supérieur droit identifies a particular indexed object or symbol; that identification is not independently established by the quoted local expression.\par
\par\noindent\textbf{OM-SOL-0596 -- t2 VI p.315 -- SOL}\par
Indice placé sous le deuxième diagramme\par
\emph{Reported reading}: \textbackslash{}alpha\textbackslash{}beta\textbackslash{}quad\textbackslash{}coprod\_A B\_\{\textbackslash{}alpha\}\par
\emph{Proposal}: \textbackslash{}alpha\textbackslash{}beta\textbackslash{}in\textbackslash{}coprod\_A B\_\{\textbackslash{}alpha\}\par
The proposed reading \textbackslash{}alpha\textbackslash{}beta\textbackslash{}in\textbackslash{}coprod\_A B\_\{\textbackslash{}alpha\} at Indice placé sous le deuxième diagramme changes a substantive mathematical assertion or bibliographic identification; no independent derivation of this exact replacement is established in this review.\par
\par\noindent\textbf{OM-SOL-0597 -- t2 VI p.317 -- SOL}\par
Deuxième paragraphe, source et but de q'\par
\emph{Reported reading}: q' : f'\textasciicircum{}\{\textbackslash{}ast\}(X') \textbackslash{}to f'\textasciicircum{}\{\textbackslash{}ast\}(X)\par
\emph{Proposal}: q' : f'\textasciicircum{}\{\textbackslash{}ast\}(f\textasciicircum{}\{\textbackslash{}ast\}(X')) \textbackslash{}to f'\textasciicircum{}\{\textbackslash{}ast\}(f\textasciicircum{}\{\textbackslash{}ast\}(X))\par
The proposed reading q' : f'\textasciicircum{}\{\textbackslash{}ast\}(f\textasciicircum{}\{\textbackslash{}ast\}(X')) \textbackslash{}to f'\textasciicircum{}\{\textbackslash{}ast\}(f\textasciicircum{}\{\textbackslash{}ast\}(X)) at Deuxième paragraphe, source et but de q' changes a substantive mathematical assertion or bibliographic identification; no independent derivation of this exact replacement is established in this review.\par
\par\noindent\textbf{OM-SOL-0598 -- t2 VI p.317 -- SOL}\par
8.3.12, familles couvrantes déduites par changement de base\par
\emph{Reported reading}: Z\_\{\textbackslash{}alpha\textbackslash{}beta\} \textbackslash{}to Z\par
\emph{Proposal}: Z\_\{\textbackslash{}alpha \textbackslash{}beta\} \textbackslash{}to Z\_\{\textbackslash{}alpha\}\par
The replacement Z\_\{\textbackslash{}alpha \textbackslash{}beta\} \textbackslash{}to Z\_\{\textbackslash{}alpha\} at 8.3.12, familles couvrantes déduites par changement de base identifies a particular indexed object or symbol; that identification is not independently established by the quoted local expression.\par
\par\noindent\textbf{OM-SOL-0599 -- t2 VI p.318 -- SOL}\par
Dernier paragraphe, « Soit X un objet cohérent »\par
\emph{Reported reading}: (F\_\{\textbackslash{}coh\})\textasciicircum{}\{\textbackslash{}sim\}\par
\emph{Proposal}: (\textbackslash{}underrightarrow\{F\}\_\{\textbackslash{}coh\})\textasciicircum{}\{\textbackslash{}sim\}\par
The replacement (\textbackslash{}underrightarrow\{F\}\_\{\textbackslash{}coh\})\textasciicircum{}\{\textbackslash{}sim\} at Dernier paragraphe, « Soit X un objet cohérent » requires a complete check of the maps, their endpoints and the composition convention; this identification is not independently verified.\par
\par\noindent\textbf{OM-SOL-0600 -- t2 VI p.319 -- SOL}\par
Définition de R' après X\_i =\allowbreak  coker(p\_1,p\_2)\par
\emph{Reported reading}: R' =\allowbreak  Y' \textbackslash{}times\_\{X'\_i\} Y'\par
\emph{Proposal}: R' =\allowbreak  \textbackslash{}fprod\{Y'\}\{Y'\}\{X\_i\}\par
The replacement R' =\allowbreak  \textbackslash{}fprod\{Y'\}\{Y'\}\{X\_i\} at Définition de R' après X\_i =\allowbreak  coker(p\_1,p\_2) identifies a particular indexed object or symbol; that identification is not independently established by the quoted local expression.\par
\par\noindent\textbf{OM-SOL-0601 -- t2 VI p.319 -- SOL}\par
Démonstration du corollaire 8.3.14, introduction du foncteur cartésien\par
\emph{Reported reading}: m\textasciicircum{}\{\textbackslash{}ast\}\_\{\textbackslash{}coh\} =\allowbreak  G\_\{\textbackslash{}coh\} \textbackslash{}to F\_\{\textbackslash{}coh\}\par
\emph{Proposal}: m\textasciicircum{}\{\textbackslash{}ast\}\_\{\textbackslash{}coh\} : G\_\{\textbackslash{}coh\} \textbackslash{}to F\_\{\textbackslash{}coh\}\par
The proposed reading m\textasciicircum{}\{\textbackslash{}ast\}\_\{\textbackslash{}coh\} : G\_\{\textbackslash{}coh\} \textbackslash{}to F\_\{\textbackslash{}coh\} at Démonstration du corollaire 8.3.14, introduction du foncteur cartésien changes a substantive mathematical assertion or bibliographic identification; no independent derivation of this exact replacement is established in this review.\par
\par\noindent\textbf{OM-SOL-0602 -- t2 VI p.322 -- SOL}\par
8.5.1, immediately before (8.5.1.3)\par
\emph{Reported reading}: pour tout objet \$M\$ de \$F\$\par
\emph{Proposal}: \textbackslash{}underleftarrow\{F\}\par
The replacement \textbackslash{}underleftarrow\{F\} at 8.5.1, immediately before (8.5.1.3) requires a complete check of the maps, their endpoints and the composition convention; this identification is not independently verified.\par
\par\noindent\textbf{OM-SOL-0603 -- t2 VI p.323 -- SOL}\par
(8.5.1.7), diagonal labelled beta\_i\par
\emph{Reported reading}: \textbackslash{}beta\_i:F\_i\textbackslash{}to\textbackslash{}Top(X)\par
\emph{Proposal}: \textbackslash{}beta\_i:\textbackslash{}Top(F)\textbackslash{}to F\_i\par
The replacement \textbackslash{}beta\_i:\textbackslash{}Top(F)\textbackslash{}to F\_i at (8.5.1.7), diagonal labelled beta\_i requires a complete check of the maps, their endpoints and the composition convention; this identification is not independently verified.\par
\par\noindent\textbf{OM-SOL-0604 -- t2 VI p.323 -- SOL}\par
(8.5.1.7), lower vertex\par
\emph{Reported reading}: \textbackslash{}Top(X)\par
\emph{Proposal}: \textbackslash{}Top(F)\par
The proposed reading \textbackslash{}Top(F) at (8.5.1.7), lower vertex changes a substantive mathematical assertion or bibliographic identification; no independent derivation of this exact replacement is established in this review.\par
\par\noindent\textbf{OM-SOL-0606 -- t2 VI p.327 -- SOL}\par
Premier paragraphe, description de la section beta\_i\textasciicircum{}*(M\_i)\par
\emph{Reported reading}: (f : i \textbackslash{}to j) \textbackslash{}mapsto f\textasciicircum{}\{\textbackslash{}ast\}(M\_i)\par
\emph{Proposal}: (f : j \textbackslash{}to i) \textbackslash{}mapsto f\textasciicircum{}\{\textbackslash{}ast\}(M\_i)\par
The replacement (f : j \textbackslash{}to i) \textbackslash{}mapsto f\textasciicircum{}\{\textbackslash{}ast\}(M\_i) at Premier paragraphe, description de la section beta\_i\textasciicircum{}*(M\_i) requires a complete check of the maps, their endpoints and the composition convention; this identification is not independently verified.\par
\par\noindent\textbf{OM-SOL-0607 -- t2 VI p.328 -- SOL}\par
Démonstration de 8.5.8, isomorphisme obtenu par commutativité du diagramme\par
\emph{Reported reading}: Q\textasciicircum{}\{\textbackslash{}ast\}u\textasciicircum{}\{\textbackslash{}sim\}\_\{\textbackslash{}ast\} Q\_\{\textbackslash{}ast\}(N)\par
\emph{Proposal}: Q\textasciicircum{}\{\textbackslash{}ast\}m\textasciicircum{}\{\textbackslash{}sim\}\_\{\textbackslash{}ast\} Q\_\{\textbackslash{}ast\}(N)\par
The replacement Q\textasciicircum{}\{\textbackslash{}ast\}m\textasciicircum{}\{\textbackslash{}sim\}\_\{\textbackslash{}ast\} Q\_\{\textbackslash{}ast\}(N) at Démonstration de 8.5.8, isomorphisme obtenu par commutativité du diagramme requires a complete check of the maps, their endpoints and the composition convention; this identification is not independently verified.\par
\par\noindent\textbf{OM-SOL-0608 -- t2 VI p.329 -- SOL}\par
Démonstration de 8.5.9, définition du foncteur phi\par
\emph{Reported reading}: le morphisme identique de I\par
\emph{Proposal}: le morphisme identique de \$i\$\par
The replacement le morphisme identique de \$i\$ at Démonstration de 8.5.9, définition du foncteur phi identifies a particular indexed object or symbol; that identification is not independently established by the quoted local expression.\par
\par\noindent\textbf{OM-SOL-0610 -- t2 VI p.330 -- SOL}\par
8.6.3, renvoi concernant les images inverses de faisceaux de modules\par
\emph{Reported reading}: (IV 1.1)\par
\emph{Proposal}: (IV 11.1)\par
The proposed reference (IV 11.1) at 8.6.3, renvoi concernant les images inverses de faisceaux de modules needs comparison with the statement actually invoked; that target has not been independently established here.\par
\par\noindent\textbf{OM-SOL-0615 -- t2 VI p.335 -- SOL}\par
Démonstration du corollaire 8.7.7, comparaison des termes E\_2\par
\emph{Reported reading}: il résulte de (8.7.6.2)\par
\emph{Proposal}: \textbackslash{}eqhyperref[VI.eq8.7.6.3]\{8.7.6.3\}\par
The proposed reference \textbackslash{}eqhyperref[VI.eq8.7.6.3]\{8.7.6.3\} at Démonstration du corollaire 8.7.7, comparaison des termes E\_2 needs comparison with the statement actually invoked; that target has not been independently established here.\par
\par\noindent\textbf{OM-SOL-0618 -- t2 VI p.337 -- SOL}\par
Lemme 9.3, final clause\par
\emph{Reported reading}: dans l'image de l'un des \$(V\_i)\$\par
\emph{Proposal}: dans l'image de l'un des \$Q(V\_i)\$\par
The proposed reading dans l'image de l'un des \$Q(V\_i)\$ at Lemme 9.3, final clause changes a substantive mathematical assertion or bibliographic identification; no independent derivation of this exact replacement is established in this review.\par
\par\noindent\textbf{OM-SOL-0619 -- t2 VI p.337 -- SOL}\par
Proof of Lemme 9.3, last parenthetical condition\par
\emph{Reported reading}: pour \$i > i\_\{\textbackslash{}circ\}\$\par
\emph{Proposal}: pour \$i \textbackslash{}geq i\_\{\textbackslash{}circ\}\$\par
The proposed reading pour \$i \textbackslash{}geq i\_\{\textbackslash{}circ\}\$ at Proof of Lemme 9.3, last parenthetical condition changes a substantive mathematical assertion or bibliographic identification; no independent derivation of this exact replacement is established in this review.\par
\par\noindent\textbf{OM-SOL-0620 -- t2 VI p.338 -- SOL}\par
Proof of Lemme 9.3, indexed inductive limit\par
\emph{Reported reading}: \textbackslash{}varinjlim\_\{i > i\_\textbackslash{}circ\} \textbackslash{}Ff(U\_\{i,k\})\par
\emph{Proposal}: \textbackslash{}varinjlim\_\{i \textbackslash{}geq i\_\textbackslash{}circ\} \textbackslash{}Ff(U\_\{i,k\})\par
The proposed reading \textbackslash{}varinjlim\_\{i \textbackslash{}geq i\_\textbackslash{}circ\} \textbackslash{}Ff(U\_\{i,k\}) at Proof of Lemme 9.3, indexed inductive limit changes a substantive mathematical assertion or bibliographic identification; no independent derivation of this exact replacement is established in this review.\par
\par\noindent\textbf{OM-SOL-0621 -- t2 VI p.339 -- SOL}\par
Corollaire 9.5, existence of the set of points\par
\emph{Reported reading}: il existe un ensemble très dense de points\par
\emph{Proposal}: il existe un ensemble très dense \$X\$ de points\par
The proposed reading il existe un ensemble très dense \$X\$ de points at Corollaire 9.5, existence of the set of points changes a substantive mathematical assertion or bibliographic identification; no independent derivation of this exact replacement is established in this review.\par
\par\noindent\textbf{OM-SOL-0622 -- t2 VI p.340 -- SOL}\par
Bibliography [1], title\par
\emph{Reported reading}: Calculus of Fraction\par
\emph{Proposal}: Calculus of Fractions\par
The cited book title needs verification against the publication.\par
\par\noindent\textbf{OM-SOL-0623 -- t2 VI p.340 -- SOL}\par
Bibliography [7], article title\par
\emph{Reported reading}: Sur quelques points d'algèbres homologiques\par
\emph{Proposal}: Sur quelques points d'algèbre homologique\par
The exact article title needs verification against the publication.\par
\par\noindent\textbf{OM-SOL-0626 -- t2 VII p.347 -- SOL}\par
Footnote (*), example over an algebraically closed field\par
\emph{Reported reading}: si \$S =\allowbreak  \textbackslash{}Spec k\$\par
\emph{Proposal}: si \$X =\allowbreak  \textbackslash{}Spec k\$\par
The proposed reading si \$X =\allowbreak  \textbackslash{}Spec k\$ at Footnote (*), example over an algebraically closed field changes a substantive mathematical assertion or bibliographic identification; no independent derivation of this exact replacement is established in this review.\par
\par\noindent\textbf{OM-SOL-0627 -- t2 VII p.352 -- SOL}\par
Footnote (1), reference after « les formules »\par
\emph{Reported reading}: les formules 4.1 (i)\par
\emph{Proposal}: \textbackslash{}hyperref[VII.prop4.1]\{4.1\} (ii)\par
The proposed reference \textbackslash{}hyperref[VII.prop4.1]\{4.1\} (ii) at Footnote (1), reference after « les formules » needs comparison with the statement actually invoked; that target has not been independently established here.\par
\par\noindent\textbf{OM-SOL-0629 -- t2 VII p.353 -- SOL}\par
Paragraph beginning « Ainsi »\par
\emph{Reported reading}: Ainsi, (i) montre\par
\emph{Proposal}: Ainsi, (ii) montre\par
The number of the cited subassertion requires checking both parts of Proposition 4.1.\par
\par\noindent\textbf{OM-SOL-0633 -- t2 VII p.356 -- SOL}\par
Opening proof, coefficient module in (*)\par
\emph{Reported reading}: et F un module quasi-cohérent sur X\par
\emph{Proposal}: et \$F\$ un module quasi-cohérent sur \$Y\$\par
The proposed reading et \$F\$ un module quasi-cohérent sur \$Y\$ at Opening proof, coefficient module in (*) changes a substantive mathematical assertion or bibliographic identification; no independent derivation of this exact replacement is established in this review.\par
\par\noindent\textbf{OM-SOL-0634 -- t2 VII p.356 -- SOL}\par
Remarques 4.5, replacement of the étale site\par
\emph{Reported reading}: (\textbackslash{}Sch)\_\{/X'\}\par
\emph{Proposal}: \$(\textbackslash{}Sch)\_\{/X\}\$\par
The replacement \$(\textbackslash{}Sch)\_\{/X\}\$ at Remarques 4.5, replacement of the étale site identifies a particular indexed object or symbol; that identification is not independently established by the quoted local expression.\par
\par\noindent\textbf{OM-SOL-0635 -- t2 VII p.357 -- SOL}\par
5.1, description of the algebras \ensuremath{𝒜}\_i\par
\emph{Reported reading}: où \ensuremath{𝒜}\_i est une Algèbre quasi-cohérente sur X\_i\par
\emph{Proposal}: où \$\ensuremath{𝒜}\_i\$ est une Algèbre quasi-cohérente sur \$X\_\{i\_0\}\$\par
The proposed reading où \$\ensuremath{𝒜}\_i\$ est une Algèbre quasi-cohérente sur \$X\_\{i\_0\}\$ at 5.1, description of the algebras \ensuremath{𝒜}\_i changes a substantive mathematical assertion or bibliographic identification; no independent derivation of this exact replacement is established in this review.\par
\par\noindent\textbf{OM-SOL-0636 -- t2 VII p.359 -- SOL}\par
5.3, formula (*), subscript of the inductive limit\par
\emph{Reported reading}: \textbackslash{}Ff\_X /\textbackslash{}Cat I\par
\emph{Proposal}: \textbackslash{}Ff\_\{\textbackslash{}mathfrak\{X\}\} /\textbackslash{}Cat I\par
The replacement \textbackslash{}Ff\_\{\textbackslash{}mathfrak\{X\}\} /\textbackslash{}Cat I at 5.3, formula (*), subscript of the inductive limit identifies a particular indexed object or symbol; that identification is not independently established by the quoted local expression.\par
\par\noindent\textbf{OM-SOL-0637 -- t2 VII p.360 -- SOL}\par
5.5, sentence beginning « Compte tenu des topologies »\par
\emph{Reported reading}: \textbackslash{}Gg devient un site fibré sur \textbackslash{}cat(I)\par
\emph{Proposal}: \$\textbackslash{}Gg\_\{\textbackslash{}mathfrak\{X\}\}\$ devient un \textbackslash{}emph\{site fibré\} sur \$\textbackslash{}cat(I)\$\par
The replacement \$\textbackslash{}Gg\_\{\textbackslash{}mathfrak\{X\}\}\$ devient un \textbackslash{}emph\{site fibré\} sur \$\textbackslash{}cat(I)\$ at 5.5, sentence beginning « Compte tenu des topologies » identifies a particular indexed object or symbol; that identification is not independently established by the quoted local expression.\par
\par\noindent\textbf{OM-SOL-0638 -- t2 VII p.361 -- SOL}\par
Paragraphe précédant le théorème 5.7\par
\emph{Reported reading}: les résultats de [blanc] à la situation présente\par
\emph{Proposal}: les résultats de VI à la situation présente\par
The proposed reference les résultats de VI à la situation présente at Paragraphe précédant le théorème 5.7 needs comparison with the statement actually invoked; that target has not been independently established here.\par
\par\noindent\textbf{OM-SOL-0640 -- t2 VII p.362 -- SOL}\par
Phrase introduisant le corollaire 5.8\par
\emph{Reported reading}: explicité également dans [blanc]\par
\emph{Proposal}: explicité également dans VI.\par
The proposed reference explicité également dans VI. at Phrase introduisant le corollaire 5.8 needs comparison with the statement actually invoked; that target has not been independently established here.\par
\par\noindent\textbf{OM-SOL-0641 -- t2 VII p.363 -- SOL}\par
Démonstration de 5.9, fin de la parenthèse invoquant EGA IV 8.8.2\par
\emph{Reported reading}: site étale restreint de \$G\_\{\textbackslash{}infty\}\$\par
\emph{Proposal}: site étale restreint de \$X\$\par
The proposed reading site étale restreint de \$X\$ at Démonstration de 5.9, fin de la parenthèse invoquant EGA IV 8.8.2 changes a substantive mathematical assertion or bibliographic identification; no independent derivation of this exact replacement is established in this review.\par
\par\noindent\textbf{OM-SOL-0642 -- t2 VII p.364 -- SOL}\par
Corollaire 5.11, parenthèse définissant le morphisme utilisé pour l'image inverse\par
\emph{Reported reading}: \$u\_i:X\textbackslash{}to X\_i\$\par
\emph{Proposal}: \$u\_i : Y \textbackslash{}to Y\_i\$\par
The proposed reading \$u\_i : Y \textbackslash{}to Y\_i\$ at Corollaire 5.11, parenthèse définissant le morphisme utilisé pour l'image inverse changes a substantive mathematical assertion or bibliographic identification; no independent derivation of this exact replacement is established in this review.\par
\par\noindent\textbf{OM-SOL-0644 -- t2 VII p.365 -- SOL}\par
Remarque 5.14 a), après le renvoi à la thèse de Giraud\par
\emph{Reported reading}: le contexte général de [blanc] Il est certainement possible\par
\emph{Proposal}: le contexte général de VI. Il est certainement possible\par
The proposed reference le contexte général de VI. Il est certainement possible at Remarque 5.14 a), après le renvoi à la thèse de Giraud needs comparison with the statement actually invoked; that target has not been independently established here.\par
\par\noindent\textbf{OM-SOL-0645 -- t2 VIII p.367 -- SOL}\par
Première référence dans la démonstration\par
\emph{Reported reading}: SGA IX 4.10 et 4.11\par
\emph{Proposal}: \textbackslash{}SGA 1 IX 4.10 et 4.11\par
This finding intersects Orgogozo's \textbackslash{}sisi material at Première référence dans la démonstration; preserve that human entry unchanged. The proposed reading \textbackslash{}SGA 1 IX 4.10 et 4.11 involves a printed glyph or a missing stroke at Première référence dans la démonstration; the textual record does not independently establish that glyph.\par
\par\noindent\textbf{OM-SOL-0647 -- t2 VIII p.368 -- SOL}\par
Conclusion de la démonstration, « Cela prouve… »\par
\emph{Reported reading}: la première assertion de 2.1\par
\emph{Proposal}: \textbackslash{}hyperref[VIII.thm1.1]\{1.1\}\par
This finding intersects Orgogozo's \textbackslash{}sisi material at Conclusion de la démonstration, « Cela prouve… »; preserve that human entry unchanged. The proposed reference \textbackslash{}hyperref[VIII.thm1.1]\{1.1\} at Conclusion de la démonstration, « Cela prouve… » needs comparison with the statement actually invoked; that target has not been independently established here.\par
\par\noindent\textbf{OM-SOL-0648 -- t2 VIII p.368 -- SOL}\par
Note de bas de page, « les fibres de S'\_i »\par
\emph{Reported reading}: la fibre de X\_i =\allowbreak  Spec A'\_i […] les fibres de S'\_i\par
\emph{Proposal}: les fibres de \$X\_i\$\par
The proposed reading les fibres de \$X\_i\$ at Note de bas de page, « les fibres de S'\_i » changes a substantive mathematical assertion or bibliographic identification; no independent derivation of this exact replacement is established in this review.\par
\par\noindent\textbf{OM-SOL-0649 -- t2 VIII p.369 -- SOL}\par
Fin de l'exemple 1.3 c)\par
\emph{Reported reading}: donc 2.1 s'applique\par
\emph{Proposal}: \textbackslash{}hyperref[VIII.thm1.1]\{1.1\}\par
This finding intersects Orgogozo's \textbackslash{}sisi material at Fin de l'exemple 1.3 c); preserve that human entry unchanged. The proposed reference \textbackslash{}hyperref[VIII.thm1.1]\{1.1\} at Fin de l'exemple 1.3 c) needs comparison with the statement actually invoked; that target has not been independently established here.\par
\par\noindent\textbf{OM-SOL-0650 -- t2 VIII p.369 -- SOL}\par
Remarques 1.4, « satisfait aux hypothèses de… »\par
\emph{Reported reading}: satisfait aux hypothèses de 2.1\par
\emph{Proposal}: \textbackslash{}hyperref[VIII.thm1.1]\{1.1\}\par
This finding intersects Orgogozo's \textbackslash{}sisi material at Remarques 1.4, « satisfait aux hypothèses de… »; preserve that human entry unchanged. The proposed reference \textbackslash{}hyperref[VIII.thm1.1]\{1.1\} at Remarques 1.4, « satisfait aux hypothèses de… » needs comparison with the statement actually invoked; that target has not been independently established here.\par
\par\noindent\textbf{OM-SOL-0651 -- t2 VIII p.369 -- SOL}\par
Deuxième paragraphe des remarques 1.4, référence à l'exemple\par
\emph{Reported reading}: à propos de l'exemple 2.3 c)\par
\emph{Proposal}: \textbackslash{}hyperref[VIII.exem1.3]\{1.3\} c)\par
This finding intersects Orgogozo's \textbackslash{}sisi material at Deuxième paragraphe des remarques 1.4, référence à l'exemple; preserve that human entry unchanged. The proposed reference \textbackslash{}hyperref[VIII.exem1.3]\{1.3\} c) at Deuxième paragraphe des remarques 1.4, référence à l'exemple needs comparison with the statement actually invoked; that target has not been independently established here.\par
\par\noindent\textbf{OM-SOL-0652 -- t2 VIII p.369 -- SOL}\par
Deuxième paragraphe des remarques 1.4, référence à la conclusion\par
\emph{Reported reading}: la conclusion de 2.2\par
\emph{Proposal}: \textbackslash{}hyperref[VIII.cor1.2]\{1.2\}\par
This finding intersects Orgogozo's \textbackslash{}sisi material at Deuxième paragraphe des remarques 1.4, référence à la conclusion; preserve that human entry unchanged. The proposed reference \textbackslash{}hyperref[VIII.cor1.2]\{1.2\} at Deuxième paragraphe des remarques 1.4, référence à la conclusion needs comparison with the statement actually invoked; that target has not been independently established here.\par
\par\noindent\textbf{OM-SOL-0653 -- t2 VIII p.370 -- SOL}\par
Proposition 2.1, définition du foncteur i\par
\emph{Reported reading}: i(X')(X'') =\allowbreak  Hom\_X(X', X'')\par
\emph{Proposal}: i(X')(X'') =\allowbreak  \textbackslash{}Hom\_X (X'', X')\par
The proposed reading i(X')(X'') =\allowbreak  \textbackslash{}Hom\_X (X'', X') at Proposition 2.1, définition du foncteur i changes a substantive mathematical assertion or bibliographic identification; no independent derivation of this exact replacement is established in this review.\par
\par\noindent\textbf{OM-SOL-0655 -- t2 VIII p.372 -- SOL}\par
3.1, last sentence\par
\emph{Reported reading}: à k-isomorphisme\par
\emph{Proposal}: à \$k(x)\$-isomorphisme\par
The proposed reading à \$k(x)\$-isomorphisme at 3.1, last sentence changes a substantive mathematical assertion or bibliographic identification; no independent derivation of this exact replacement is established in this review.\par
\par\noindent\textbf{OM-SOL-0656 -- t2 VIII p.373 -- SOL}\par
3.4, parenthesis about replacing the field\par
\emph{Reported reading}: la clôture algébrique séparable de \$k\$ dans \$\textbackslash{}Omega\$\par
\emph{Proposal}: la clôture algébrique séparable de \$k(x)\$ dans \$\textbackslash{}Omega\$\par
The proposed reading la clôture algébrique séparable de \$k(x)\$ dans \$\textbackslash{}Omega\$ at 3.4, parenthesis about replacing the field changes a substantive mathematical assertion or bibliographic identification; no independent derivation of this exact replacement is established in this review.\par
\par\noindent\textbf{OM-SOL-0661 -- t2 VIII p.378 -- SOL}\par
Paragraph beginning « Si \ensuremath{\bar{x}} », reference defining the functor j\par
\emph{Reported reading}: du foncteur \$j\$ de 1.1\par
\emph{Proposal}: du foncteur \$j\$ de \textbackslash{}hyperref[VIII.prop2.1]\{2.1\}\par
This finding intersects Orgogozo's \textbackslash{}sisi material at Paragraph beginning « Si \ensuremath{\bar{x}} », reference defining the functor j; preserve that human entry unchanged. The proposed reference du foncteur \$j\$ de \textbackslash{}hyperref[VIII.prop2.1]\{2.1\} at Paragraph beginning « Si \ensuremath{\bar{x}} », reference defining the functor j needs comparison with the statement actually invoked; that target has not been independently established here.\par
\par\noindent\textbf{OM-SOL-0662 -- t2 VIII p.378 -- SOL}\par
End of the paragraph discussing the two meanings of F\_x\par
\emph{Reported reading}: mais en vertu de 1.1 c'est là une contradiction inessentielle\par
\emph{Proposal}: mais en vertu de \textbackslash{}hyperref[VIII.prop2.1]\{2.1\} c'est là une contradiction inessentielle\par
This finding intersects Orgogozo's \textbackslash{}sisi material at End of the paragraph discussing the two meanings of F\_x; preserve that human entry unchanged. The proposed reference mais en vertu de \textbackslash{}hyperref[VIII.prop2.1]\{2.1\} c'est là une contradiction inessentielle at End of the paragraph discussing the two meanings of F\_x needs comparison with the statement actually invoked; that target has not been independently established here.\par
\par\noindent\textbf{OM-SOL-0663 -- t2 VIII p.378 -- SOL}\par
3.12, parenthetical reference after the displayed functor\par
\emph{Reported reading}: grâce à 1.1\par
\emph{Proposal}: grâce à \textbackslash{}hyperref[VIII.prop2.1]\{2.1\}\par
This finding intersects Orgogozo's \textbackslash{}sisi material at 3.12, parenthetical reference after the displayed functor; preserve that human entry unchanged. The proposed reference grâce à \textbackslash{}hyperref[VIII.prop2.1]\{2.1\} at 3.12, parenthetical reference after the displayed functor needs comparison with the statement actually invoked; that target has not been independently established here.\par
\par\noindent\textbf{OM-SOL-0664 -- t2 VIII p.380 -- SOL}\par
Universal property of the henselization\par
\emph{Reported reading}: dans des anneaux locaux noethériens\par
\emph{Proposal}: dans des anneaux locaux henséliens\par
This finding intersects Orgogozo's \textbackslash{}sisi material at Universal property of the henselization; preserve that human entry unchanged. The proposed reading dans des anneaux locaux henséliens at Universal property of the henselization changes a substantive mathematical assertion or bibliographic identification; no independent derivation of this exact replacement is established in this review.\par
\par\noindent\textbf{OM-SOL-0668 -- t2 VIII p.389 -- SOL}\par
Démonstration de 5.6, après le renvoi à VII 5.13\par
\emph{Reported reading}: R\textasciicircum{}q f\_\{\textbackslash{}ast\}(F\_i) =\allowbreak  0\par
\emph{Proposal}: R\textasciicircum{}\{q\} (f\_i)\_\{\textbackslash{}ast\} (F\_i) =\allowbreak  0\par
The replacement R\textasciicircum{}\{q\} (f\_i)\_\{\textbackslash{}ast\} (F\_i) =\allowbreak  0 at Démonstration de 5.6, après le renvoi à VII 5.13 requires a complete check of the maps, their endpoints and the composition convention; this identification is not independently verified.\par
\par\noindent\textbf{OM-SOL-0669 -- t2 VIII p.390 -- SOL}\par
Fin du premier paragraphe, désignation du topos étale\par
\emph{Reported reading}: étale X\_\{\textbackslash{}text\{\textbackslash{}rm ét\}\} de X\par
\emph{Proposal}: étale \$\textbackslash{}widetilde\{X\_\{\textbackslash{}text\{\textbackslash{}rm ét\}\}\}\$ de \$X\$\par
The replacement étale \$\textbackslash{}widetilde\{X\_\{\textbackslash{}text\{\textbackslash{}rm ét\}\}\}\$ de \$X\$ at Fin du premier paragraphe, désignation du topos étale identifies a particular indexed object or symbol; that identification is not independently established by the quoted local expression.\par
\par\noindent\textbf{OM-SOL-0671 -- t2 VIII p.392 -- SOL}\par
Proof of 6.3, converse: both displays and the intervening and concluding references to G\par
\emph{Reported reading}: \$G \textbackslash{}to i\_\{\textbackslash{}ast\}i\textasciicircum{}\{\textbackslash{}ast\}G\$; « l'hypothèse faite sur G »; \$i\textasciicircum{}\{\textbackslash{}ast\}(G) \textbackslash{}to i\textasciicircum{}\{\textbackslash{}ast\}(i\_\{\textbackslash{}ast\}i\textasciicircum{}\{\textbackslash{}ast\}(G))\$; « appliqué à F =\allowbreak  i\textasciicircum{}\{\textbackslash{}ast\}G »\par
\emph{Proposal}: \$F \textbackslash{}to i\_\{\textbackslash{}ast\}i\textasciicircum{}\{\textbackslash{}ast\}F\$; l'hypothèse faite sur \$F\$; \$i\textasciicircum{}\{\textbackslash{}ast\}(F) \textbackslash{}to i\textasciicircum{}\{\textbackslash{}ast\}(i\_\{\textbackslash{}ast\}i\textasciicircum{}\{\textbackslash{}ast\}(F))\$; 5.7 appliqué au faisceau \$i\textasciicircum{}\{\textbackslash{}ast\}F\$ et à \$i\$\par
The proposed reading \$F \textbackslash{}to i\_\{\textbackslash{}ast\}i\textasciicircum{}\{\textbackslash{}ast\}F\$; l'hypothèse faite sur \$F\$; \$i\textasciicircum{}\{\textbackslash{}ast\}(F) \textbackslash{}to i\textasciicircum{}\{\textbackslash{}ast\}(i\_\{\textbackslash{}ast\}i\textasciicircum{}\{\textbackslash{}ast\}(F))\$; 5.7 appliqué au faisceau \$i\textasciicircum{}\{\textbackslash{}ast\}F\$ et à \$i\$ at Proof of 6.3, converse: both displays and the intervening and concluding references to G changes a substantive mathematical assertion or bibliographic identification; no independent derivation of this exact replacement is established in this review.\par
\par\noindent\textbf{OM-SOL-0673 -- t2 VIII p.397 -- SOL}\par
7.8, renvoi après « rappelons »\par
\emph{Reported reading}: rappelons (IV )\par
\emph{Proposal}: (IV 6.1)\par
This finding intersects Orgogozo's \textbackslash{}sisi material at 7.8, renvoi après « rappelons »; preserve that human entry unchanged. The proposed reference (IV 6.1) at 7.8, renvoi après « rappelons » needs comparison with the statement actually invoked; that target has not been independently established here.\par
\par\noindent\textbf{OM-SOL-0674 -- t2 VIII p.400 -- SOL}\par
7.9.1, indice de la limite inductive dans le premier Hom\par
\emph{Reported reading}: \textbackslash{}varinjlim\_\{C\_\{\textbackslash{}varphi\}\} \textbackslash{}check\{Y\}\par
\emph{Proposal}: \textbackslash{}varinjlim\_\{C\textasciicircum{}0\_\{/\textbackslash{}varphi\}\} \textbackslash{}check\{Y\}\par
The replacement \textbackslash{}varinjlim\_\{C\textasciicircum{}0\_\{/\textbackslash{}varphi\}\} \textbackslash{}check\{Y\} at 7.9.1, indice de la limite inductive dans le premier Hom identifies a particular indexed object or symbol; that identification is not independently established by the quoted local expression.\par
\par\noindent\textbf{OM-SOL-0675 -- t2 VIII p.401 -- SOL}\par
Phrase commençant par « De plus », appartenance du couple (X′, ξ)\par
\emph{Reported reading}: \$(X',\textbackslash{}xi)\textbackslash{}in\textbackslash{}Ob X\_\{\textbackslash{}text\{ét\}\}\$\par
\emph{Proposal}: (X', \textbackslash{}xi) \textbackslash{}in \textbackslash{}Ob C\_\{/\textbackslash{}varphi\}\par
The proposed reading (X', \textbackslash{}xi) \textbackslash{}in \textbackslash{}Ob C\_\{/\textbackslash{}varphi\} at Phrase commençant par « De plus », appartenance du couple (X′, ξ) changes a substantive mathematical assertion or bibliographic identification; no independent derivation of this exact replacement is established in this review.\par
\par\noindent\textbf{OM-SOL-0676 -- t2 VIII p.402 -- SOL}\par
Proposition 8.1, aboutissement de la suite spectrale\par
\emph{Reported reading}: \$H\textasciicircum{}\{\textbackslash{}ast\}(X,F)\$\par
\emph{Proposal}: H\textasciicircum{}\{\textbackslash{}ast\}(Y,F)\par
The proposed reading H\textasciicircum{}\{\textbackslash{}ast\}(Y,F) at Proposition 8.1, aboutissement de la suite spectrale changes a substantive mathematical assertion or bibliographic identification; no independent derivation of this exact replacement is established in this review.\par
\par\noindent\textbf{OM-SOL-0678 -- t2 VIII p.409 -- SOL}\par
Premier mot de la page, suite du lemme 9.3\par
\emph{Reported reading}: X, et définissons de façon analogue\par
\emph{Proposal}: \$F\$\par
This finding intersects Orgogozo's \textbackslash{}sisi material at Premier mot de la page, suite du lemme 9.3; preserve that human entry unchanged. The proposed reading \$F\$ at Premier mot de la page, suite du lemme 9.3 changes a substantive mathematical assertion or bibliographic identification; no independent derivation of this exact replacement is established in this review.\par
\par\noindent\textbf{OM-SOL-0680 -- t2 VIII p.410 -- SOL}\par
Premier paragraphe, référence à SGA\par
\emph{Reported reading}: SGA IX 2.3\par
\emph{Proposal}: \textbackslash{}SGA 1 IX 2.3\par
This finding intersects Orgogozo's \textbackslash{}sisi material at Premier paragraphe, référence à SGA; preserve that human entry unchanged. The proposed reference \textbackslash{}SGA 1 IX 2.3 at Premier paragraphe, référence à SGA needs comparison with the statement actually invoked; that target has not been independently established here.\par
\par\noindent\textbf{OM-SOL-0683 -- t2 VIII p.412 -- SOL}\par
9.4.4, « pour un x \ensuremath{\in} X »\par
\emph{Reported reading}: \$x\textbackslash{}in X\$\par
\emph{Proposal}: x\textbackslash{}in X'\par
The replacement x\textbackslash{}in X' at 9.4.4, « pour un x \ensuremath{\in} X » identifies a particular indexed object or symbol; that identification is not independently established by the quoted local expression.\par
\par\noindent\textbf{OM-SOL-0684 -- t3 IX p.1 -- SOL}\par
Introduction, « variété abélienne sur k »\par
\emph{Reported reading}: si X est une variété abélienne sur k\par
\emph{Proposal}: si \$X\$ est une variété abélienne sur \$K\$\par
The proposed reading si \$X\$ est une variété abélienne sur \$K\$ at Introduction, « variété abélienne sur k » changes a substantive mathematical assertion or bibliographic identification; no independent derivation of this exact replacement is established in this review.\par
\par\noindent\textbf{OM-SOL-0685 -- t3 IX p.2 -- SOL}\par
Suite de l'introduction, algèbre de quaternions\par
\emph{Reported reading}: algèbre de quaternions définie sur \$\textbackslash{}mathbb\{Z\}\$\par
\emph{Proposal}: algèbre de quaternions définie sur \$\textbackslash{}mathbb\{Q\}\$\par
The replacement algèbre de quaternions définie sur \$\textbackslash{}mathbb\{Q\}\$ at Suite de l'introduction, algèbre de quaternions identifies a particular indexed object or symbol; that identification is not independently established by the quoted local expression.\par
\par\noindent\textbf{OM-SOL-0687 -- t3 IX p.3 -- SOL}\par
Proposition 1.2 (v), argument de l'image directe supérieure\par
\emph{Reported reading}: \$R\textasciicircum{}qu\_\{\textbackslash{}ast\}F'\$\par
\emph{Proposal}: R\textasciicircum{}qu\_\{\textbackslash{}ast\}F\par
The replacement R\textasciicircum{}qu\_\{\textbackslash{}ast\}F at Proposition 1.2 (v), argument de l'image directe supérieure identifies a particular indexed object or symbol; that identification is not independently established by the quoted local expression.\par
\par\noindent\textbf{OM-SOL-0688 -- t3 IX p.3 -- SOL}\par
Proposition 1.2 (v), topos portant les images directes supérieures\par
\emph{Reported reading}: sur \$T\$ pour tout \$q\$\par
\emph{Proposal}: sur \$T'\$ pour tout \$q\$\par
The replacement sur \$T'\$ pour tout \$q\$ at Proposition 1.2 (v), topos portant les images directes supérieures identifies a particular indexed object or symbol; that identification is not independently established by the quoted local expression.\par
\par\noindent\textbf{OM-SOL-0689 -- t3 IX p.4 -- SOL}\par
Démonstration de 1.2 (iii), dernière phrase\par
\emph{Reported reading}: La deuxième assertion est une conséquence de\par
\emph{Proposal}: La deuxième assertion est une conséquence de la commutation de la cohomologie aux limites inductives filtrantes.\par
The broken sentence establishes a lacuna, but it does not uniquely determine the proposed explanatory completion.\par
\par\noindent\textbf{OM-SOL-0690 -- t3 IX p.4 -- SOL}\par
Démonstration de 1.2 (iv), membre droit de l'égalité\par
\emph{Reported reading}: \$u\textasciicircum{}\{\textbackslash{}ast\}(\textbackslash{}underline\{a\}P')=\allowbreak \textbackslash{}underline\{a\}(u\textbackslash{}cdot P)\$\par
\emph{Proposal}: u\textasciicircum{}\{\textbackslash{}ast\}(\textbackslash{}underline\{a\}P')=\allowbreak \textbackslash{}underline\{a\}(u\textbackslash{}cdot P')\par
The replacement u\textasciicircum{}\{\textbackslash{}ast\}(\textbackslash{}underline\{a\}P')=\allowbreak \textbackslash{}underline\{a\}(u\textbackslash{}cdot P') at Démonstration de 1.2 (iv), membre droit de l'égalité identifies a particular indexed object or symbol; that identification is not independently established by the quoted local expression.\par
\par\noindent\textbf{OM-SOL-0692 -- t3 IX p.10 -- SOL}\par
Proposition 2.4 (v), designation of the nonempty open subset\par
\emph{Reported reading}: un ouvert non-vide de l'adhérence ... sur U\par
\emph{Proposal}: un ouvert non-vide \$U\$ de l'adhérence\par
The proposed reading un ouvert non-vide \$U\$ de l'adhérence at Proposition 2.4 (v), designation of the nonempty open subset changes a substantive mathematical assertion or bibliographic identification; no independent derivation of this exact replacement is established in this review.\par
\par\noindent\textbf{OM-SOL-0694 -- t3 IX p.12 -- SOL}\par
Démonstration de 2.7, produit contenant H\par
\emph{Reported reading}: C'est un sous-faisceau de G \ensuremath{\times}\_S G\par
\emph{Proposal}: sous-faisceau de \$\textbackslash{}fprod\{G\}\{G\}\{X\}\$\par
The proposed reading sous-faisceau de \$\textbackslash{}fprod\{G\}\{G\}\{X\}\$ at Démonstration de 2.7, produit contenant H changes a substantive mathematical assertion or bibliographic identification; no independent derivation of this exact replacement is established in this review.\par
\par\noindent\textbf{OM-SOL-0695 -- t3 IX p.13 -- SOL}\par
Démonstration de 2.7.2, indice suivant la définition des F\_\ensuremath{\ell}\par
\emph{Reported reading}: \ensuremath{\ell} =\allowbreak  (H', J) \ensuremath{\in} L\par
\emph{Proposal}: \$\textbackslash{}ell=\allowbreak (J,H')\textbackslash{}in L\$\par
The proposed reading \$\textbackslash{}ell=\allowbreak (J,H')\textbackslash{}in L\$ at Démonstration de 2.7.2, indice suivant la définition des F\_\ensuremath{\ell} changes a substantive mathematical assertion or bibliographic identification; no independent derivation of this exact replacement is established in this review.\par
\par\noindent\textbf{OM-SOL-0696 -- t3 IX p.14 -- SOL}\par
Limite inductive à J fixé\par
\emph{Reported reading}: pour \ensuremath{\ell} =\allowbreak  (N', J) avec J fixé\par
\emph{Proposal}: \$\textbackslash{}ell=\allowbreak (J,N')\$\par
The replacement \$\textbackslash{}ell=\allowbreak (J,N')\$ at Limite inductive à J fixé identifies a particular indexed object or symbol; that identification is not independently established by the quoted local expression.\par
\par\noindent\textbf{OM-SOL-0697 -- t3 IX p.14 -- SOL}\par
Sous-faisceau en groupes engendré par le faisceau d'ensembles F\_J\par
\emph{Reported reading}: le sous-faisceau en groupes F\_J de F\par
\emph{Proposal}: sous-faisceau en groupes \$\textbackslash{}bar\{F\_J\}\$ de \$F\$\par
The replacement sous-faisceau en groupes \$\textbackslash{}bar\{F\_J\}\$ de \$F\$ at Sous-faisceau en groupes engendré par le faisceau d'ensembles F\_J identifies a particular indexed object or symbol; that identification is not independently established by the quoted local expression.\par
\par\noindent\textbf{OM-SOL-0699 -- t3 IX p.15 -- SOL}\par
Corollaire 2.7.4, après « catégories fibrées, définie dans »\par
\emph{Reported reading}: définie dans [blanc] En particulier\par
\emph{Proposal}: définie dans [référence à restituer]. En particulier\par
The proposed reference définie dans [référence à restituer]. En particulier at Corollaire 2.7.4, après « catégories fibrées, définie dans » needs comparison with the statement actually invoked; that target has not been independently established here.\par
\par\noindent\textbf{OM-SOL-0700 -- t3 IX p.16 -- SOL}\par
Suite du lemme 2.8.1, parenthèse « en particulier »\par
\emph{Reported reading}: chaque \$Y\_i\$ est constructible dans \$X\$\par
\emph{Proposal}: chaque \$Y\_i\$ est constructible dans \$Y\$\par
The proposed reading chaque \$Y\_i\$ est constructible dans \$Y\$ at Suite du lemme 2.8.1, parenthèse « en particulier » changes a substantive mathematical assertion or bibliographic identification; no independent derivation of this exact replacement is established in this review.\par
\par\noindent\textbf{OM-SOL-0701 -- t3 IX p.16 -- SOL}\par
Démonstration du lemme, réunion des images T\_i\par
\emph{Reported reading}: la réunion des \$T\_i\$ est \$X\$ puisque \$f\$ est surjectif\par
\emph{Proposal}: la réunion des \$T\_i\$ est \$Y\$ puisque \$f\$ est surjectif\par
The proposed reading la réunion des \$T\_i\$ est \$Y\$ puisque \$f\$ est surjectif at Démonstration du lemme, réunion des images T\_i changes a substantive mathematical assertion or bibliographic identification; no independent derivation of this exact replacement is established in this review.\par
\par\noindent\textbf{OM-SOL-0702 -- t3 IX p.16 -- SOL}\par
Démonstration du lemme, extraction d'une réunion finie\par
\emph{Reported reading}: que \$X\$ est réunion d'un nombre fini des \$T\_i\$\par
\emph{Proposal}: que \$Y\$ est réunion d'un nombre fini des \$T\_i\$\par
The proposed reading que \$Y\$ est réunion d'un nombre fini des \$T\_i\$ at Démonstration du lemme, extraction d'une réunion finie changes a substantive mathematical assertion or bibliographic identification; no independent derivation of this exact replacement is established in this review.\par
\par\noindent\textbf{OM-SOL-0704 -- t3 IX p.16 -- SOL}\par
Démonstration de 2.8, factorisation obtenue par le lemme\par
\emph{Reported reading}: \$f\$ est de la forme \$gh\$\par
\emph{Proposal}: est de la forme \$hg\$\par
The replacement est de la forme \$hg\$ at Démonstration de 2.8, factorisation obtenue par le lemme requires a complete check of the maps, their endpoints and the composition convention; this identification is not independently verified.\par
\par\noindent\textbf{OM-SOL-0706 -- t3 IX p.22 -- SOL}\par
Démonstration de 2.13 (iii), phrase suivant le diagramme\par
\emph{Reported reading}: Comme c(F\_Q) =\allowbreak  c(F\_P)\par
\emph{Proposal}: \textbackslash{}operatorname\{card\}(F\_Q)=\allowbreak \textbackslash{}operatorname\{card\}(F\_P)\par
The proposed reading \textbackslash{}operatorname\{card\}(F\_Q)=\allowbreak \textbackslash{}operatorname\{card\}(F\_P) at Démonstration de 2.13 (iii), phrase suivant le diagramme changes a substantive mathematical assertion or bibliographic identification; no independent derivation of this exact replacement is established in this review.\par
\par\noindent\textbf{OM-SOL-0707 -- t3 IX p.23 -- SOL}\par
Proposition 2.14 (i), désignation du faisceau\par
\emph{Reported reading}: et tout faisceau d'ensembles […] constructible sur X', f\_*(F) est constructible\par
\emph{Proposal}: et tout faisceau \$F\$ d'ensembles\par
The proposed reading et tout faisceau \$F\$ d'ensembles at Proposition 2.14 (i), désignation du faisceau changes a substantive mathematical assertion or bibliographic identification; no independent derivation of this exact replacement is established in this review.\par
\par\noindent\textbf{OM-SOL-0708 -- t3 IX p.24 -- SOL}\par
Démonstration de 2.14 (ii), produit des images directes par p\_i r\_i\par
\emph{Reported reading}: p\_\{i*\}(r\_\{i*\}(C\_\{\textbackslash{}bar Z\_i\})) =\allowbreak  t\_\{i*\}(C\_\{i\textbackslash{}bar Z\_i\})\par
\emph{Proposal}: p\_\{i\textbackslash{}ast\}(r\_\{i\textbackslash{}ast\}(C\_\{i\textbackslash{}bar\{Z\}\_i\}))=\allowbreak t\_\{i\textbackslash{}ast\}(C\_\{i\textbackslash{}bar\{Z\}\_i\})\par
The replacement p\_\{i\textbackslash{}ast\}(r\_\{i\textbackslash{}ast\}(C\_\{i\textbackslash{}bar\{Z\}\_i\}))=\allowbreak t\_\{i\textbackslash{}ast\}(C\_\{i\textbackslash{}bar\{Z\}\_i\}) at Démonstration de 2.14 (ii), produit des images directes par p\_i r\_i identifies a particular indexed object or symbol; that identification is not independently established by the quoted local expression.\par
\par\noindent\textbf{OM-SOL-0712 -- t3 IX p.35 -- SOL}\par
Théorème 4.6, suite exacte, terme H\textasciicircum{}1\par
\emph{Reported reading}: H\textasciicircum{}1(X,(\ensuremath{𝛍}\_m)\_X)\par
\emph{Proposal}: H\textasciicircum{}1(X,(\ensuremath{𝛍}\_n)\_X)\par
The replacement H\textasciicircum{}1(X,(\ensuremath{𝛍}\_n)\_X) at Théorème 4.6, suite exacte, terme H\textasciicircum{}1 identifies a particular indexed object or symbol; that identification is not independently established by the quoted local expression.\par
\par\noindent\textbf{OM-SOL-0713 -- t3 IX p.38 -- SOL}\par
Diagramme (5.3'), étiquette de la flèche composée.\par
\emph{Reported reading}: \$d'\$\par
\emph{Proposal}: \$d\$\par
The replacement \$d\$ at Diagramme (5.3'), étiquette de la flèche composée. identifies a particular indexed object or symbol; that identification is not independently established by the quoted local expression.\par
\par\noindent\textbf{OM-SOL-0714 -- t3 IX p.39 -- SOL}\par
Démonstration de 5.5, choix du revêtement principal U''.\par
\emph{Reported reading}: tel que \$g\$ devienne constant sur \$U''\$\par
\emph{Proposal}: tel que \$G\$ devienne constant sur \$U''\$\par
The proposed reading tel que \$G\$ devienne constant sur \$U''\$ at Démonstration de 5.5, choix du revêtement principal U''. changes a substantive mathematical assertion or bibliographic identification; no independent derivation of this exact replacement is established in this review.\par
\par\noindent\textbf{OM-SOL-0715 -- t3 IX p.42 -- SOL}\par
Bibliography [2], Ogg's title\par
\emph{Reported reading}: over functions fields\par
\emph{Proposal}: over function fields\par
The exact title wording requires comparison with the cited article.\par
\par\noindent\textbf{OM-SOL-0718 -- t3 X p.45 -- SOL}\par
Cas a), définition du groupe de Galois H\par
\emph{Reported reading}: \$H=\allowbreak G(\textbackslash{}widetilde\{k\}/k')\$\par
\emph{Proposal}: \$H=\allowbreak G(\textbackslash{}bar\{k\}/k')\$\par
The replacement \$H=\allowbreak G(\textbackslash{}bar\{k\}/k')\$ at Cas a), définition du groupe de Galois H identifies a particular indexed object or symbol; that identification is not independently established by the quoted local expression.\par
\par\noindent\textbf{OM-SOL-0719 -- t3 X p.47 -- SOL}\par
Fin de la démonstration du lemme 2.2.1\par
\emph{Reported reading}: et comme \$K'\$ en est l'anneau total des fractions\par
\emph{Proposal}: \textbackslash{}widehat\{K\}'\par
The replacement \textbackslash{}widehat\{K\}' at Fin de la démonstration du lemme 2.2.1 identifies a particular indexed object or symbol; that identification is not independently established by the quoted local expression.\par
\par\noindent\textbf{OM-SOL-0720 -- t3 X p.55 -- SOL}\par
Démonstration de 4.1, paragraphe « Pour chaque F\_i »\par
\emph{Reported reading}: tels que \$F\_x\textbackslash{}neq 0\$\par
\emph{Proposal}: \$(F\_i)\_x\textbackslash{}neq 0\$\par
The replacement \$(F\_i)\_x\textbackslash{}neq 0\$ at Démonstration de 4.1, paragraphe « Pour chaque F\_i » identifies a particular indexed object or symbol; that identification is not independently established by the quoted local expression.\par
\par\noindent\textbf{OM-SOL-0722 -- t3 X p.58 -- SOL}\par
Démonstration de 5.1, première formule d'annulation\par
\emph{Reported reading}: H\textasciicircum{}q(X,i\_1(\textbackslash{}ZZ/p)\_U)=\allowbreak 0\par
\emph{Proposal}: H\textasciicircum{}q(X,i\_!(\textbackslash{}ZZ/p)\_U)=\allowbreak 0\par
This finding intersects Orgogozo's \textbackslash{}sisi material at Démonstration de 5.1, première formule d'annulation; preserve that human entry unchanged. The replacement H\textasciicircum{}q(X,i\_!(\textbackslash{}ZZ/p)\_U)=\allowbreak 0 at Démonstration de 5.1, première formule d'annulation identifies a particular indexed object or symbol; that identification is not independently established by the quoted local expression.\par
\par\noindent\textbf{OM-SOL-0723 -- t3 X p.58 -- SOL}\par
Démonstration de 5.1, fin du paragraphe commençant « Tenant compte »\par
\emph{Reported reading}: pour \$q>\textbackslash{}cdqc X\$\par
\emph{Proposal}: pour \$q>\textbackslash{}cdqc X+1\$\par
The proposed reading pour \$q>\textbackslash{}cdqc X+1\$ at Démonstration de 5.1, fin du paragraphe commençant « Tenant compte » changes a substantive mathematical assertion or bibliographic identification; no independent derivation of this exact replacement is established in this review.\par
\par\noindent\textbf{OM-SOL-0733 -- t3 XI p.72 -- SOL}\par
4.0, référence précédant (4.1)\par
\emph{Reported reading}: morphisme (inclusion) de sites (IV )\par
\emph{Proposal}: (IV\textasciitilde{}4.9.1)\par
The proposed reference (IV\textasciitilde{}4.9.1) at 4.0, référence précédant (4.1) needs comparison with the statement actually invoked; that target has not been independently established here.\par
\par\noindent\textbf{OM-SOL-0734 -- t3 XI p.72 -- SOL}\par
4.0, justification de l'équivalence des deux topos\par
\emph{Reported reading}: et on déduit de ( ) que les topos\par
\emph{Proposal}: (III\textasciitilde{}4.1)\par
The proposed reference (III\textasciitilde{}4.1) at 4.0, justification de l'équivalence des deux topos needs comparison with the statement actually invoked; that target has not been independently established here.\par
\par\noindent\textbf{OM-SOL-0735 -- t3 XI p.75 -- SOL}\par
Théorème 4.4 (iii), renvoi initial\par
\emph{Reported reading}: Avec les notations de (i)\par
\emph{Proposal}: Avec les notations de \{\textbackslash{}rm(ii)\}\par
The proposed reference Avec les notations de \{\textbackslash{}rm(ii)\} at Théorème 4.4 (iii), renvoi initial needs comparison with the statement actually invoked; that target has not been independently established here.\par
\par\noindent\textbf{OM-SOL-0737 -- t3 XI p.76 -- SOL}\par
Dernière phrase, domaine d'application de la suite exacte (*)\par
\emph{Reported reading}: Lorsque q \ensuremath{\geq} 1\par
\emph{Proposal}: Lorsque \$q>1\$\par
The proposed reading Lorsque \$q>1\$ at Dernière phrase, domaine d'application de la suite exacte (*) changes a substantive mathematical assertion or bibliographic identification; no independent derivation of this exact replacement is established in this review.\par
\par\noindent\textbf{OM-SOL-0738 -- t3 XI p.78 -- SOL}\par
Bibliographie [4], pagination de l'article de Serre\par
\emph{Reported reading}: p. 1-12 (1956)\par
\emph{Proposal}: p. 1-42 (1956)\par
The proposed reading p. 1-42 (1956) at Bibliographie [4], pagination de l'article de Serre changes a substantive mathematical assertion or bibliographic identification; no independent derivation of this exact replacement is established in this review.\par
\par\noindent\textbf{OM-SOL-0739 -- t3 XII p.83 -- SOL}\par
Proof of Proposition 3.3, finite product indexed by 1\ensuremath{\leq}i\ensuremath{\leq}n\par
\emph{Reported reading}: F\textbackslash{}longrightarrow\textbackslash{}prod\_\{1\textbackslash{}leq i\textbackslash{}leq n\}i\_\{x\textbackslash{}ast\}i\_x\textasciicircum{}\{\textbackslash{}ast\}F\par
\emph{Proposal}: F\textbackslash{}longrightarrow \textbackslash{}prod\_\{1\textbackslash{}leq i\textbackslash{}leq n\}(i\_\{x\_i\})\_\{\textbackslash{}ast\}i\_\{x\_i\}\textasciicircum{}\{\textbackslash{}ast\}F\par
The replacement F\textbackslash{}longrightarrow \textbackslash{}prod\_\{1\textbackslash{}leq i\textbackslash{}leq n\}(i\_\{x\_i\})\_\{\textbackslash{}ast\}i\_\{x\_i\}\textasciicircum{}\{\textbackslash{}ast\}F at Proof of Proposition 3.3, finite product indexed by 1\ensuremath{\leq}i\ensuremath{\leq}n identifies a particular indexed object or symbol; that identification is not independently established by the quoted local expression.\par
\par\noindent\textbf{OM-SOL-0740 -- t3 XII p.85 -- SOL}\par
Notation for the derived base-change morphism, before (4.2)\par
\emph{Reported reading}: \textbackslash{}varphi\textasciicircum{}q=\allowbreak \textbackslash{}varphi\textasciicircum{}g g\par
\emph{Proposal}: \textbackslash{}varphi\textasciicircum{}q=\allowbreak \textbackslash{}varphi\_g\textasciicircum{}q\par
The replacement \textbackslash{}varphi\textasciicircum{}q=\allowbreak \textbackslash{}varphi\_g\textasciicircum{}q at Notation for the derived base-change morphism, before (4.2) requires a complete check of the maps, their endpoints and the composition convention; this identification is not independently verified.\par
\par\noindent\textbf{OM-SOL-0743 -- t3 XII p.88 -- SOL}\par
Paragraph preceding Corollary 5.2, after the cohomology equality\par
\emph{Reported reading}: pour chaque faisceau G sur X\par
\emph{Proposal}: pour chaque faisceau \$G\$ sur \$X'\$\par
The replacement pour chaque faisceau \$G\$ sur \$X'\$ at Paragraph preceding Corollary 5.2, after the cohomology equality identifies a particular indexed object or symbol; that identification is not independently established by the quoted local expression.\par
\par\noindent\textbf{OM-SOL-0745 -- t3 XII p.91 -- SOL}\par
Same paragraph, “d'un élément de…”\par
\emph{Reported reading}: H\textasciicircum{}1(G,H\textasciicircum{}0(\textbackslash{}bar\{X\},\textbackslash{}bar\{F\}))\par
\emph{Proposal}: H\textasciicircum{}1(G,H\textasciicircum{}\{\textbackslash{}circ\}(\textbackslash{}bar\{X\},\textbackslash{}bar\{F\}\textasciicircum{}\{\textbackslash{}alpha\}))\par
The proposed reading H\textasciicircum{}1(G,H\textasciicircum{}\{\textbackslash{}circ\}(\textbackslash{}bar\{X\},\textbackslash{}bar\{F\}\textasciicircum{}\{\textbackslash{}alpha\})) at Same paragraph, “d'un élément de…” changes a substantive mathematical assertion or bibliographic identification; no independent derivation of this exact replacement is established in this review.\par
\par\noindent\textbf{OM-SOL-0746 -- t3 XII p.92 -- SOL}\par
Godement argument: parenthesis below the square and lower row of the following diagram\par
\emph{Reported reading}: H\textasciicircum{}1(\textbackslash{}bar\{X\},\textbackslash{}bar\{G\}); H\textasciicircum{}1(X\_0,G\_0); H\textasciicircum{}1(\textbackslash{}bar\{X\}\_0,\textbackslash{}bar\{G\}\_0)\textasciicircum{}G\par
\emph{Proposal}: H\textasciicircum{}1(\textbackslash{}bar\{X\},\textbackslash{}bar\{C\})=\allowbreak \textbackslash{}varinjlim H\textasciicircum{}1(X',C)=\allowbreak 0;\textbackslash{}qquad H\textasciicircum{}\{1\}(X\_0,C\_0)\textbackslash{}longrightarrow H\textasciicircum{}\{1\}(\textbackslash{}bar\{X\}\_0,\textbackslash{}bar\{C\}\_0)\textasciicircum{}G\par
The proposed reading H\textasciicircum{}1(\textbackslash{}bar\{X\},\textbackslash{}bar\{C\})=\allowbreak \textbackslash{}varinjlim H\textasciicircum{}1(X',C)=\allowbreak 0;\textbackslash{}qquad H\textasciicircum{}\{1\}(X\_0,C\_0)\textbackslash{}longrightarrow H\textasciicircum{}\{1\}(\textbackslash{}bar\{X\}\_0,\textbackslash{}bar\{C\}\_0)\textasciicircum{}G at Godement argument: parenthesis below the square and lower row of the following diagram changes a substantive mathematical assertion or bibliographic identification; no independent derivation of this exact replacement is established in this review.\par
\par\noindent\textbf{OM-SOL-0747 -- t3 XII p.92 -- SOL}\par
Immediately after the commutative square\par
\emph{Reported reading}: donc est nul.\par
\emph{Proposal}: donc \$\textbackslash{}alpha\$ est nul.\par
The proposed reading donc \$\textbackslash{}alpha\$ est nul. at Immediately after the commutative square changes a substantive mathematical assertion or bibliographic identification; no independent derivation of this exact replacement is established in this review.\par
\par\noindent\textbf{OM-SOL-0750 -- t3 XII p.93 -- SOL}\par
5.6, definition of the constant sheaf\par
\emph{Reported reading}: F=\allowbreak S\_X, followed by “de valeur T”\par
\emph{Proposal}: F=\allowbreak T\_X\par
The proposed reading F=\allowbreak T\_X at 5.6, definition of the constant sheaf changes a substantive mathematical assertion or bibliographic identification; no independent derivation of this exact replacement is established in this review.\par
\par\noindent\textbf{OM-SOL-0751 -- t3 XII p.93 -- SOL}\par
Corollary 5.7, map on connected components\par
\emph{Reported reading}: \textbackslash{}pi\_0(X)\textbackslash{}longrightarrow\textbackslash{}pi\_0(X\_0), induced by X\_0\textbackslash{}longrightarrow X\par
\emph{Proposal}: \textbackslash{}pi\_\{\textbackslash{}circ\}(X\_\{\textbackslash{}circ\})\textbackslash{}longrightarrow\textbackslash{}pi\_\{\textbackslash{}circ\}(X)\par
The replacement \textbackslash{}pi\_\{\textbackslash{}circ\}(X\_\{\textbackslash{}circ\})\textbackslash{}longrightarrow\textbackslash{}pi\_\{\textbackslash{}circ\}(X) at Corollary 5.7, map on connected components requires a complete check of the maps, their endpoints and the composition convention; this identification is not independently verified.\par
\par\noindent\textbf{OM-SOL-0753 -- t3 XII p.96 -- SOL}\par
Classe de la gerbe, après l'équivalence des catégories de sections\par
\emph{Reported reading}: \$\textbackslash{}xi\textbackslash{}in H\textasciicircum{}2(X,\textbackslash{}mathcal\{G\})\$\par
\emph{Proposal}: \$\textbackslash{}xi\textbackslash{}in H\textasciicircum{}2(X,L)\$\par
This finding intersects Orgogozo's \textbackslash{}sisi material at Classe de la gerbe, après l'équivalence des catégories de sections; preserve that human entry unchanged. The proposed reading \$\textbackslash{}xi\textbackslash{}in H\textasciicircum{}2(X,L)\$ at Classe de la gerbe, après l'équivalence des catégories de sections changes a substantive mathematical assertion or bibliographic identification; no independent derivation of this exact replacement is established in this review.\par
\par\noindent\textbf{OM-SOL-0754 -- t3 XII p.97 -- SOL}\par
Introduction des localisés stricts, avant le carré commutatif\par
\emph{Reported reading}: \$\textbackslash{}bar\{S\}\$ et \$\textbackslash{}bar\{S\}'\$ de \$S\$ en \$\textbackslash{}bar\{s\}\$ et en \$\textbackslash{}bar\{s\}'\$ respectivement\par
\emph{Proposal}: \$\textbackslash{}bar\{S\}\$ et \$\textbackslash{}bar\{S\}'\$ de \$S\$ et de \$S'\$ en \$\textbackslash{}bar\{s\}\$ et en \$\textbackslash{}bar\{s\}'\$ respectivement\par
The proposed reading \$\textbackslash{}bar\{S\}\$ et \$\textbackslash{}bar\{S\}'\$ de \$S\$ et de \$S'\$ en \$\textbackslash{}bar\{s\}\$ et en \$\textbackslash{}bar\{s\}'\$ respectivement at Introduction des localisés stricts, avant le carré commutatif changes a substantive mathematical assertion or bibliographic identification; no independent derivation of this exact replacement is established in this review.\par
\par\noindent\textbf{OM-SOL-0755 -- t3 XII p.98 -- SOL}\par
Début de 6.1.1\par
\emph{Reported reading}: \$(F,f,-)\$\par
\emph{Proposal}: \$(f,F,-)\$\par
The proposed reading \$(f,F,-)\$ at Début de 6.1.1 changes a substantive mathematical assertion or bibliographic identification; no independent derivation of this exact replacement is established in this review.\par
\par\noindent\textbf{OM-SOL-0756 -- t3 XII p.99 -- SOL}\par
Démonstration de 6.2, cas initial de la récurrence\par
\emph{Reported reading}: est triviale si \$n=\allowbreak 1\$\par
\emph{Proposal}: est triviale si \$n=\allowbreak -1\$\par
The proposed reading est triviale si \$n=\allowbreak -1\$ at Démonstration de 6.2, cas initial de la récurrence changes a substantive mathematical assertion or bibliographic identification; no independent derivation of this exact replacement is established in this review.\par
\par\noindent\textbf{OM-SOL-0757 -- t3 XII p.100 -- SOL}\par
Diagramme de suites exactes, troisième terme de la ligne supérieure\par
\emph{Reported reading}: \$R\textasciicircum{}\{n+1\}(A)\$\par
\emph{Proposal}: \$T\textasciicircum{}\{n+1\}(A)\$\par
The proposed reading \$T\textasciicircum{}\{n+1\}(A)\$ at Diagramme de suites exactes, troisième terme de la ligne supérieure changes a substantive mathematical assertion or bibliographic identification; no independent derivation of this exact replacement is established in this review.\par
\par\noindent\textbf{OM-SOL-0758 -- t3 XII p.101 -- SOL}\par
Opening sentence, criterion for surjectivity\par
\emph{Reported reading}: l'image de \$T\textasciicircum{}\{n+1\}(L\textasciicircum{}\{\textbackslash{}bigdot\})\$ dans \$T\textasciicircum{}\{n+1\}(K\textasciicircum{}\{\textbackslash{}bigdot\})\$\par
\emph{Proposal}: l'image de \$T'\textasciicircum{}\{n+1\}(L\textasciicircum{}\{\textbackslash{}bigdot\})\$ dans \$T'\textasciicircum{}\{n+1\}(K\textasciicircum{}\{\textbackslash{}bigdot\})\$\par
The replacement l'image de \$T'\textasciicircum{}\{n+1\}(L\textasciicircum{}\{\textbackslash{}bigdot\})\$ dans \$T'\textasciicircum{}\{n+1\}(K\textasciicircum{}\{\textbackslash{}bigdot\})\$ at Opening sentence, criterion for surjectivity identifies a particular indexed object or symbol; that identification is not independently established by the quoted local expression.\par
\par\noindent\textbf{OM-SOL-0759 -- t3 XII p.108 -- SOL}\par
Description d'un F-torseur par réduction du groupe structural\par
\emph{Reported reading}: d'un torseur sous G … muni d'une section de Q/F\par
\emph{Proposal}: d'un torseur \$Q\$ sous \$G\$\par
The proposed reading d'un torseur \$Q\$ sous \$G\$ at Description d'un F-torseur par réduction du groupe structural changes a substantive mathematical assertion or bibliographic identification; no independent derivation of this exact replacement is established in this review.\par
\par\noindent\textbf{OM-SOL-0761 -- t3 XII p.108 -- SOL}\par
Application de 6.2 dans la démonstration de 6.5 (iii), catégorie C'\par
\emph{Reported reading}: C' la catégorie des faisceaux de groupes abéliens\par
\emph{Proposal}: \$C'\$ la catégorie des groupes abéliens\par
The proposed reading \$C'\$ la catégorie des groupes abéliens at Application de 6.2 dans la démonstration de 6.5 (iii), catégorie C' changes a substantive mathematical assertion or bibliographic identification; no independent derivation of this exact replacement is established in this review.\par
\par\noindent\textbf{OM-SOL-0762 -- t3 XII p.109 -- SOL}\par
Dernier paragraphe, notation des faisceaux constants\par
\emph{Reported reading}: G\_\{X'\_i\}\par
\emph{Proposal}: G\_\{i X'\_i\}\par
The replacement G\_\{i X'\_i\} at Dernier paragraphe, notation des faisceaux constants identifies a particular indexed object or symbol; that identification is not independently established by the quoted local expression.\par
\par\noindent\textbf{OM-SOL-0764 -- t3 XII p.114 -- SOL}\par
Diagramme commutatif, sommet supérieur droit\par
\emph{Reported reading}: H\textasciicircum{}\{\textbackslash{}circ\}(Y,h\textasciicircum{}\{\textbackslash{}ast\}(\textbackslash{}bar\{F\}))\par
\emph{Proposal}: H\textasciicircum{}\{\textbackslash{}circ\}(Y,h\textasciicircum{}\{\textbackslash{}ast\}(p\_\{\textbackslash{}ast\}(\textbackslash{}bar\{F\})))\par
The replacement H\textasciicircum{}\{\textbackslash{}circ\}(Y,h\textasciicircum{}\{\textbackslash{}ast\}(p\_\{\textbackslash{}ast\}(\textbackslash{}bar\{F\}))) at Diagramme commutatif, sommet supérieur droit identifies a particular indexed object or symbol; that identification is not independently established by the quoted local expression.\par
\par\noindent\textbf{OM-SOL-0768 -- t3 XII p.116 -- SOL}\par
Premier paragraphe, homomorphisme sur R\textasciicircum{}1\par
\emph{Reported reading}: H\textasciicircum{}\{\textbackslash{}ast\}(R\textasciicircum{}1(p\_\{\textbackslash{}ast\}(\textbackslash{}bar\{F\}))\textbackslash{}longrightarrow R\textasciicircum{}1p\_\{\textbackslash{}ast\}(\textbackslash{}bar\{h\}\textasciicircum{}\{\textbackslash{}ast\}(\textbackslash{}bar\{F\}))\par
\emph{Proposal}: h\textasciicircum{}\{\textbackslash{}ast\}(R\textasciicircum{}1p\_\{\textbackslash{}ast\}(\textbackslash{}bar\{F\}))\textbackslash{}longrightarrow R\textasciicircum{}1p\_\{Y\textbackslash{}ast\}(\textbackslash{}bar\{h\}\textasciicircum{}\{\textbackslash{}ast\}(\textbackslash{}bar\{F\}))\par
The replacement h\textasciicircum{}\{\textbackslash{}ast\}(R\textasciicircum{}1p\_\{\textbackslash{}ast\}(\textbackslash{}bar\{F\}))\textbackslash{}longrightarrow R\textasciicircum{}1p\_\{Y\textbackslash{}ast\}(\textbackslash{}bar\{h\}\textasciicircum{}\{\textbackslash{}ast\}(\textbackslash{}bar\{F\})) at Premier paragraphe, homomorphisme sur R\textasciicircum{}1 identifies a particular indexed object or symbol; that identification is not independently established by the quoted local expression.\par
\par\noindent\textbf{OM-SOL-0769 -- t3 XII p.116 -- SOL}\par
Formule (2), premier membre\par
\emph{Reported reading}: H\textasciicircum{}\{\textbackslash{}ast\}(Rp\_\{\textbackslash{}ast\}(\textbackslash{}bar\{F\}))\par
\emph{Proposal}: h\textasciicircum{}\{\textbackslash{}ast\}(Rp\_\{\textbackslash{}ast\}(\textbackslash{}bar\{F\}))\par
The replacement h\textasciicircum{}\{\textbackslash{}ast\}(Rp\_\{\textbackslash{}ast\}(\textbackslash{}bar\{F\})) at Formule (2), premier membre identifies a particular indexed object or symbol; that identification is not independently established by the quoted local expression.\par
\par\noindent\textbf{OM-SOL-0772 -- t3 XII p.118 -- SOL}\par
Fin du corollaire 6.9, égalité des groupes de cohomologie\par
\emph{Reported reading}: H\textasciicircum{}\{n+1\}(\textbackslash{}bar\{Y\},p\_\{Y\textbackslash{}ast\}(h\textasciicircum{}\{\textbackslash{}ast\}(F))=\allowbreak H\textasciicircum{}\{n+1\}(\textbackslash{}bar\{Y\},\textbackslash{}bar\{h\}\textasciicircum{}\{\textbackslash{}ast\}(\textbackslash{}bar\{F\}))\par
\emph{Proposal}: H\textasciicircum{}\{n+1\}(\textbackslash{}bar\{Y\},p\_Y\textasciicircum{}\{\textbackslash{}ast\}(h\textasciicircum{}\{\textbackslash{}ast\}(F)))=\allowbreak H\textasciicircum{}\{n+1\}(\textbackslash{}bar\{Y\},\textbackslash{}bar\{h\}\textasciicircum{}\{\textbackslash{}ast\}(\textbackslash{}bar\{F\}))\par
The replacement H\textasciicircum{}\{n+1\}(\textbackslash{}bar\{Y\},p\_Y\textasciicircum{}\{\textbackslash{}ast\}(h\textasciicircum{}\{\textbackslash{}ast\}(F)))=\allowbreak H\textasciicircum{}\{n+1\}(\textbackslash{}bar\{Y\},\textbackslash{}bar\{h\}\textasciicircum{}\{\textbackslash{}ast\}(\textbackslash{}bar\{F\})) at Fin du corollaire 6.9, égalité des groupes de cohomologie identifies a particular indexed object or symbol; that identification is not independently established by the quoted local expression.\par
\par\noindent\textbf{OM-SOL-0773 -- t3 XII p.118 -- SOL}\par
Corollaire 6.10 (i), condition 2°, variante des ind-LL-groupes\par
\emph{Reported reading}: tout faisceau de ind-\textbackslash{}LL-groupes \textbackslash{}bar\{F\} sur X\par
\emph{Proposal}: tout faisceau de ind-\$\textbackslash{}LL\$-groupes \$\textbackslash{}bar\{F\}\$ sur \$\textbackslash{}bar\{X\}\$\par
The replacement tout faisceau de ind-\$\textbackslash{}LL\$-groupes \$\textbackslash{}bar\{F\}\$ sur \$\textbackslash{}bar\{X\}\$ at Corollaire 6.10 (i), condition 2°, variante des ind-LL-groupes identifies a particular indexed object or symbol; that identification is not independently established by the quoted local expression.\par
\par\noindent\textbf{OM-SOL-0776 -- t3 XII p.121 -- SOL}\par
Première ligne, exposants de R\par
\emph{Reported reading}: h\textasciicircum{}\{\textbackslash{}ast\}(R\textasciicircum{}1p\_\{\textbackslash{}ast\}(\textbackslash{}bar\{F\}))\textbackslash{}longrightarrow R\textasciicircum{}1p\_\{Y\textasciicircum{}\{\textbackslash{}ast\}\}(\textbackslash{}bar\{h\}\textasciicircum{}\{\textbackslash{}ast\}(\textbackslash{}bar\{F\})), puis « pour i=\allowbreak 0,1 »\par
\emph{Proposal}: h\textasciicircum{}\{\textbackslash{}ast\}(R\textasciicircum{}\{i\}p\_\{\textbackslash{}ast\}(\textbackslash{}bar\{F\}))\textbackslash{}longrightarrow R\textasciicircum{}\{i\}p\_\{Y\textasciicircum{}\{\textbackslash{}ast\}\}(\textbackslash{}bar\{h\}\textasciicircum{}\{\textbackslash{}ast\}(\textbackslash{}bar\{F\}))\par
The replacement h\textasciicircum{}\{\textbackslash{}ast\}(R\textasciicircum{}\{i\}p\_\{\textbackslash{}ast\}(\textbackslash{}bar\{F\}))\textbackslash{}longrightarrow R\textasciicircum{}\{i\}p\_\{Y\textasciicircum{}\{\textbackslash{}ast\}\}(\textbackslash{}bar\{h\}\textasciicircum{}\{\textbackslash{}ast\}(\textbackslash{}bar\{F\})) at Première ligne, exposants de R identifies a particular indexed object or symbol; that identification is not independently established by the quoted local expression.\par
\par\noindent\textbf{OM-SOL-0781 -- t3 XII p.122 -- SOL}\par
Parenthèse « qui est une résolution de »\par
\emph{Reported reading}: h\textasciicircum{}\{\textbackslash{}ast\}(\textbackslash{}bar\{F\})\par
\emph{Proposal}: \textbackslash{}bar\{h\}\textasciicircum{}\{\textbackslash{}ast\}(\textbackslash{}bar\{F\})\par
The replacement \textbackslash{}bar\{h\}\textasciicircum{}\{\textbackslash{}ast\}(\textbackslash{}bar\{F\}) at Parenthèse « qui est une résolution de » identifies a particular indexed object or symbol; that identification is not independently established by the quoted local expression.\par
\par\noindent\textbf{OM-SOL-0785 -- t3 XII p.123 -- SOL}\par
Remarque 6.12, renvoi à l'argument d'effaçabilité\par
\emph{Reported reading}: 6.10 (ii)\par
\emph{Proposal}: \textbackslash{}hyperref[XII.prop6.11]\{6.11\} (ii)\par
The proposed reference \textbackslash{}hyperref[XII.prop6.11]\{6.11\} (ii) at Remarque 6.12, renvoi à l'argument d'effaçabilité needs comparison with the statement actually invoked; that target has not been independently established here.\par
\par\noindent\textbf{OM-SOL-0787 -- t3 XII p.124 -- SOL}\par
Propriété (*), application canonique sur les sections\par
\emph{Reported reading}: \textbackslash{}Gamma(Y'/Y)\textbackslash{}rightarrow \textbackslash{}Gamma(X'/X)\par
\emph{Proposal}: \textbackslash{}Gamma(X'/X)\textbackslash{}rightarrow \textbackslash{}Gamma(Y'/Y)\par
The proposed reading \textbackslash{}Gamma(X'/X)\textbackslash{}rightarrow \textbackslash{}Gamma(Y'/Y) at Propriété (*), application canonique sur les sections changes a substantive mathematical assertion or bibliographic identification; no independent derivation of this exact replacement is established in this review.\par
\par\noindent\textbf{OM-SOL-0789 -- t3 XII p.126 -- SOL}\par
Dernier paragraphe, restriction du recouvrement après remplacement de X par X'\par
\emph{Reported reading}: les U par les U\_k\textbackslash{}cap V non-vides\par
\emph{Proposal}: U\_k\textbackslash{}cap X'\par
The proposed reading U\_k\textbackslash{}cap X' at Dernier paragraphe, restriction du recouvrement après remplacement de X par X' changes a substantive mathematical assertion or bibliographic identification; no independent derivation of this exact replacement is established in this review.\par
\par\noindent\textbf{OM-SOL-0790 -- t3 XII p.127 -- SOL}\par
Premier paragraphe de preuve, « Enfin, comme X est projectif »\par
\emph{Reported reading}: Enfin, comme X est projectif sur S\par
\emph{Proposal}: Enfin, comme \$\textbackslash{}bar\{X\}\$ est projectif sur \$S\$\par
The replacement Enfin, comme \$\textbackslash{}bar\{X\}\$ est projectif sur \$S\$ at Premier paragraphe de preuve, « Enfin, comme X est projectif » identifies a particular indexed object or symbol; that identification is not independently established by the quoted local expression.\par
\par\noindent\textbf{OM-SOL-0791 -- t3 XII p.128 -- SOL}\par
Premier paragraphe, corps résiduels des A\_i\par
\emph{Reported reading}: la limite inductive des corps k des A\_i\par
\emph{Proposal}: des corps \$k\_i\$ des \$A\_i\$\par
The replacement des corps \$k\_i\$ des \$A\_i\$ at Premier paragraphe, corps résiduels des A\_i identifies a particular indexed object or symbol; that identification is not independently established by the quoted local expression.\par
\par\noindent\textbf{OM-SOL-0792 -- t3 XII p.131 -- SOL}\par
Kummer diagram, upper middle term\par
\emph{Reported reading}: H\textasciicircum{}\{2\}(X, \{\}\_n), with a blank before the subscript n\par
\emph{Proposal}: H\textasciicircum{}\{2\}(X, \ensuremath{𝛍}\_n)\par
The proposed reading H\textasciicircum{}\{2\}(X, \ensuremath{𝛍}\_n) at Kummer diagram, upper middle term changes a substantive mathematical assertion or bibliographic identification; no independent derivation of this exact replacement is established in this review.\par
\par\noindent\textbf{OM-SOL-0793 -- t3 XII p.131 -- SOL}\par
Kummer diagram, lower middle term\par
\emph{Reported reading}: H\textasciicircum{}\{2\}(X\_0, \{\}\_n), with a blank before the subscript n\par
\emph{Proposal}: H\textasciicircum{}2(X\_0, \ensuremath{𝛍}\_n)\par
The proposed reading H\textasciicircum{}2(X\_0, \ensuremath{𝛍}\_n) at Kummer diagram, lower middle term changes a substantive mathematical assertion or bibliographic identification; no independent derivation of this exact replacement is established in this review.\par
\par\noindent\textbf{OM-SOL-0795 -- t3 XIII p.134 -- SOL}\par
Démonstration du lemme 1.3, première mention de la fibre Z\_z\par
\emph{Reported reading}: la fibre \$Z\_z\$ au point \$z\$ image de \$s'\$\par
\emph{Proposal}: la fibre \$(X\_Z)\_z\$ au point \$z\$ image de \$s'\$\par
The proposed reading la fibre \$(X\_Z)\_z\$ au point \$z\$ image de \$s'\$ at Démonstration du lemme 1.3, première mention de la fibre Z\_z changes a substantive mathematical assertion or bibliographic identification; no independent derivation of this exact replacement is established in this review.\par
\par\noindent\textbf{OM-SOL-0796 -- t3 XIII p.134 -- SOL}\par
Démonstration du lemme 1.3, définition de l'ouvert U\par
\emph{Reported reading}: \$(M\_u)\_z\$ soit plat sur \$Z\_z\$\par
\emph{Proposal}: \$(M\_u)\_z\$ soit plat sur \$(X\_Z)\_z\$\par
The proposed reading \$(M\_u)\_z\$ soit plat sur \$(X\_Z)\_z\$ at Démonstration du lemme 1.3, définition de l'ouvert U changes a substantive mathematical assertion or bibliographic identification; no independent derivation of this exact replacement is established in this review.\par
\par\noindent\textbf{OM-SOL-0797 -- t3 XIII p.134 -- SOL}\par
Démonstration du lemme 1.3, reformulation après « i.e. »\par
\emph{Reported reading}: \$M\_u\$ soit plat sur \$Z\$ en les points de \$Z\_s\$\par
\emph{Proposal}: \$M\_u\$ soit plat sur \$X\_Z\$ en les points de \$(X\_Z)\_z\$\par
The proposed reading \$M\_u\$ soit plat sur \$X\_Z\$ en les points de \$(X\_Z)\_z\$ at Démonstration du lemme 1.3, reformulation après « i.e. » changes a substantive mathematical assertion or bibliographic identification; no independent derivation of this exact replacement is established in this review.\par
\par\noindent\textbf{OM-SOL-0798 -- t3 XIII p.134 -- SOL}\par
Dernière ligne, hypothèse de platitude avant l'application d'EGA III 7.7.6\par
\emph{Reported reading}: \$\textbackslash{}Hhom(M'\textbackslash{}otimes M',M')\$ est \$S\$-plat\par
\emph{Proposal}: \$\textbackslash{}Hhom(M\textbackslash{}otimes M,M)\$ est \$S\$-plat\par
The replacement \$\textbackslash{}Hhom(M\textbackslash{}otimes M,M)\$ est \$S\$-plat at Dernière ligne, hypothèse de platitude avant l'application d'EGA III 7.7.6 identifies a particular indexed object or symbol; that identification is not independently established by the quoted local expression.\par
\par\noindent\textbf{OM-SOL-0801 -- t3 XIII p.139 -- SOL}\par
Démonstration du lemme 2.2, définition de J₀, première variable\par
\emph{Reported reading}: k[X\_\{\textbackslash{}circ\},\textbackslash{}ldots,x\_N]\par
\emph{Proposal}: k[x\_\{\textbackslash{}circ\},\textbackslash{}ldots,x\_N]\par
The proposed reading k[x\_\{\textbackslash{}circ\},\textbackslash{}ldots,x\_N] at Démonstration du lemme 2.2, définition de J₀, première variable changes a substantive mathematical assertion or bibliographic identification; no independent derivation of this exact replacement is established in this review.\par
\par\noindent\textbf{OM-SOL-0802 -- t3 XIII p.140 -- SOL}\par
Phrase précédant le diagramme\par
\emph{Reported reading}: correspondant de schémas\par
\emph{Proposal}: correspondant de catégories\par
The diagram may be induced from schemes; replacing its descriptive noun requires the intended level of the construction.\par
\par\noindent\textbf{OM-SOL-0803 -- t3 XIII p.141 -- SOL}\par
Remarque 2.4, determination of quasi-coherent modules\par
\emph{Reported reading}: Modules quasi-cohérents sur X\par
\emph{Proposal}: Modules quasi-cohérents sur \$Z\$\par
The proposed reading Modules quasi-cohérents sur \$Z\$ at Remarque 2.4, determination of quasi-coherent modules changes a substantive mathematical assertion or bibliographic identification; no independent derivation of this exact replacement is established in this review.\par
\par\noindent\textbf{OM-SOL-0804 -- t3 XIII p.142 -- SOL}\par
Proof of Proposition 3.1(i), local equation f\_x\par
\emph{Reported reading}: \$f\_x\textbackslash{}in\textbackslash{}rad(\textbackslash{}Oo\_\{X\_\{\textbackslash{}circ\}\},x)\$\par
\emph{Proposal}: \$f\_x\textbackslash{}in\textbackslash{}rad(\textbackslash{}Oo\_\{X\_\{\textbackslash{}circ\},x\})\$\par
The replacement \$f\_x\textbackslash{}in\textbackslash{}rad(\textbackslash{}Oo\_\{X\_\{\textbackslash{}circ\},x\})\$ at Proof of Proposition 3.1(i), local equation f\_x identifies a particular indexed object or symbol; that identification is not independently established by the quoted local expression.\par
\par\noindent\textbf{OM-SOL-0805 -- t3 XIII p.144 -- SOL}\par
Corollaire 3.2, target of the canonical Picard map\par
\emph{Reported reading}: \$\textbackslash{}Pic(X)\textbackslash{}rightarrow\textbackslash{}Pic(X\_\{\textbackslash{}circ\})\$\par
\emph{Proposal}: \$\textbackslash{}Pic(X)\textbackslash{}rightarrow\textbackslash{}Pic(X'\_\{\textbackslash{}circ\})\$\par
The replacement \$\textbackslash{}Pic(X)\textbackslash{}rightarrow\textbackslash{}Pic(X'\_\{\textbackslash{}circ\})\$ at Corollaire 3.2, target of the canonical Picard map identifies a particular indexed object or symbol; that identification is not independently established by the quoted local expression.\par
\par\noindent\textbf{OM-SOL-0806 -- t3 XIV p.146 -- SOL}\par
Démonstration du lemme 1.5, après l'inclusion f\_*F \ensuremath{\subset} f\_*G\par
\emph{Reported reading}: (IX 2.9 (iii))\par
\emph{Proposal}: IX 2.9 (ii)\par
This finding intersects Orgogozo's \textbackslash{}sisi material at Démonstration du lemme 1.5, après l'inclusion f\_*F \ensuremath{\subset} f\_*G; preserve that human entry unchanged. The proposed reference IX 2.9 (ii) at Démonstration du lemme 1.5, après l'inclusion f\_*F \ensuremath{\subset} f\_*G needs comparison with the statement actually invoked; that target has not been independently established here.\par
\par\noindent\textbf{OM-SOL-0807 -- t3 XIV p.146 -- SOL}\par
Démonstration du lemme 1.5, constructibilité d'un produit fini\par
\emph{Reported reading}: (IX 2.14)\par
\emph{Proposal}: IX 2.6 (ibis)\par
This finding intersects Orgogozo's \textbackslash{}sisi material at Démonstration du lemme 1.5, constructibilité d'un produit fini; preserve that human entry unchanged. The proposed reference IX 2.6 (ibis) at Démonstration du lemme 1.5, constructibilité d'un produit fini needs comparison with the statement actually invoked; that target has not been independently established here.\par
\par\noindent\textbf{OM-SOL-0808 -- t3 XIV p.147 -- SOL}\par
Première formule\par
\emph{Reported reading}: \$C(S)=\allowbreak \textbackslash{}pi\_0(X\_\{\textbackslash{}bar\{S\}\})\$\par
\emph{Proposal}: \$C(s)=\allowbreak \textbackslash{}pi\_0(X\_\{\textbackslash{}bar\{s\}\})\$\par
The proposed reading \$C(s)=\allowbreak \textbackslash{}pi\_0(X\_\{\textbackslash{}bar\{s\}\})\$ involves a printed glyph or a missing stroke at Première formule; the textual record does not independently establish that glyph.\par
\par\noindent\textbf{OM-SOL-0809 -- t3 XIV p.147 -- SOL}\par
Fin du lemme 1.7\par
\emph{Reported reading}: Alors il est également vrai pour\par
\emph{Proposal}: Alors il est également vrai pour \$f\$.\par
This finding intersects Orgogozo's \textbackslash{}sisi material at Fin du lemme 1.7; preserve that human entry unchanged. The proposed reading Alors il est également vrai pour \$f\$. at Fin du lemme 1.7 changes a substantive mathematical assertion or bibliographic identification; no independent derivation of this exact replacement is established in this review.\par
\par\noindent\textbf{OM-SOL-0810 -- t3 XIV p.148 -- SOL}\par
Premier paragraphe de la démonstration du lemme 1.7, constructibilité de i\_*i\textasciicircum{}*F\par
\emph{Reported reading}: IX 2.15 (i)\par
\emph{Proposal}: IX 2.14 (i)\par
This finding intersects Orgogozo's \textbackslash{}sisi material at Premier paragraphe de la démonstration du lemme 1.7, constructibilité de i\_*i\textasciicircum{}*F; preserve that human entry unchanged. The proposed reference IX 2.14 (i) at Premier paragraphe de la démonstration du lemme 1.7, constructibilité de i\_*i\textasciicircum{}*F needs comparison with the statement actually invoked; that target has not been independently established here.\par
\par\noindent\textbf{OM-SOL-0812 -- t3 XIV p.148 -- SOL}\par
Dernier paragraphe, « même raisonnement que pour… »\par
\emph{Reported reading}: même raisonnement que pour \$\textbackslash{}bar\{j\}\_\{!\}\textbackslash{}bar\{j\}\textasciicircum{}\{\textbackslash{}ast\}\textbackslash{}bar\{F\}\$\par
\emph{Proposal}: même raisonnement que pour \$j\_\{!\}j\textasciicircum{}\{\textbackslash{}ast\}F\$\par
This finding intersects Orgogozo's \textbackslash{}sisi material at Dernier paragraphe, « même raisonnement que pour… »; preserve that human entry unchanged. The replacement même raisonnement que pour \$j\_\{!\}j\textasciicircum{}\{\textbackslash{}ast\}F\$ at Dernier paragraphe, « même raisonnement que pour… » identifies a particular indexed object or symbol; that identification is not independently established by the quoted local expression.\par
\par\noindent\textbf{OM-SOL-0813 -- t3 XIV p.148 -- SOL}\par
Dernier paragraphe, membre de l'isomorphisme auquel s'applique l'hypothèse\par
\emph{Reported reading}: le membre de droite\par
\emph{Proposal}: le membre de gauche\par
This finding intersects Orgogozo's \textbackslash{}sisi material at Dernier paragraphe, membre de l'isomorphisme auquel s'applique l'hypothèse; preserve that human entry unchanged. The replacement le membre de gauche at Dernier paragraphe, membre de l'isomorphisme auquel s'applique l'hypothèse identifies a particular indexed object or symbol; that identification is not independently established by the quoted local expression.\par
\par\noindent\textbf{OM-SOL-0814 -- t3 XIV p.150 -- SOL}\par
Démonstration du lemme 1.8, section de D\_i induite par \ensuremath{\alpha}\par
\emph{Reported reading}: une section \$\textbackslash{}alpha\$ de \$D\_i\$\par
\emph{Proposal}: une section \$\textbackslash{}bar\{\textbackslash{}alpha\}\$ de \$D\_i\$\par
This finding intersects Orgogozo's \textbackslash{}sisi material at Démonstration du lemme 1.8, section de D\_i induite par \ensuremath{\alpha}; preserve that human entry unchanged. The replacement une section \$\textbackslash{}bar\{\textbackslash{}alpha\}\$ de \$D\_i\$ at Démonstration du lemme 1.8, section de D\_i induite par \ensuremath{\alpha} identifies a particular indexed object or symbol; that identification is not independently established by the quoted local expression.\par
\par\noindent\textbf{OM-SOL-0815 -- t3 XIV p.150 -- SOL}\par
Phrase suivante, application du lemme 1.5\par
\emph{Reported reading}: Puisque \$f\_\{\textbackslash{}ast\}G\$ est constructible\par
\emph{Proposal}: Puisque \$f\_\{\textbackslash{}ast\}C\textasciicircum{}\{\textbackslash{}alpha\}\$ est constructible\par
The replacement Puisque \$f\_\{\textbackslash{}ast\}C\textasciicircum{}\{\textbackslash{}alpha\}\$ est constructible at Phrase suivante, application du lemme 1.5 requires a complete check of the maps, their endpoints and the composition convention; this identification is not independently verified.\par
\par\noindent\textbf{OM-SOL-0816 -- t3 XIV p.151 -- SOL}\par
Reference after «même constant et constructible»\par
\emph{Reported reading}: IX 2.8\par
\emph{Proposal}: \textbackslash{}hyperrefext[IX][IX.prop2.4]\{2.4\}\textasciitilde{}(i)\par
This finding intersects Orgogozo's \textbackslash{}sisi material at Reference after «même constant et constructible»; preserve that human entry unchanged. The proposed reference \textbackslash{}hyperrefext[IX][IX.prop2.4]\{2.4\}\textasciitilde{}(i) at Reference after «même constant et constructible» needs comparison with the statement actually invoked; that target has not been independently established here.\par
\par\noindent\textbf{OM-SOL-0817 -- t3 XIV p.151 -- SOL}\par
Displayed exact sequence after «Soit», obtained by twisting\par
\emph{Reported reading}: 0\textbackslash{}longrightarrow (R\textasciicircum{}\{1\}f\_\{\textbackslash{}ast\})\textbackslash{}pi\_\{\textbackslash{}ast\}\textbackslash{}bar\{F\}\textbackslash{}longrightarrow (R\textasciicircum{}\{1\}\textbackslash{}bar\{f\}\_\{\textbackslash{}ast\})\textbackslash{}bar\{F\}\textbackslash{}longrightarrow f\_\{\textbackslash{}ast\}(R\textasciicircum{}\{1\}\textbackslash{}pi\_\{\textbackslash{}ast\})\textbackslash{}bar\{F\}\par
\emph{Proposal}: 0\textbackslash{}longrightarrow (R\textasciicircum{}\{1\}f\_\{\textbackslash{}ast\})\textbackslash{}pi\_\{\textbackslash{}ast\}\textbackslash{}bar\{F\}\textasciicircum{}\{\textbackslash{}alpha\}\textbackslash{}longrightarrow (R\textasciicircum{}\{1\}\textbackslash{}bar\{f\}\_\{\textbackslash{}ast\})\textbackslash{}bar\{F\}\textasciicircum{}\{\textbackslash{}alpha\}\textbackslash{}longrightarrow f\_\{\textbackslash{}ast\}(R\textasciicircum{}\{1\}\textbackslash{}pi\_\{\textbackslash{}ast\})\textbackslash{}bar\{F\}\textasciicircum{}\{\textbackslash{}alpha\}\par
The replacement 0\textbackslash{}longrightarrow (R\textasciicircum{}\{1\}f\_\{\textbackslash{}ast\})\textbackslash{}pi\_\{\textbackslash{}ast\}\textbackslash{}bar\{F\}\textasciicircum{}\{\textbackslash{}alpha\}\textbackslash{}longrightarrow (R\textasciicircum{}\{1\}\textbackslash{}bar\{f\}\_\{\textbackslash{}ast\})\textbackslash{}bar\{F\}\textasciicircum{}\{\textbackslash{}alpha\}\textbackslash{}longrightarrow f\_\{\textbackslash{}ast\}(R\textasciicircum{}\{1\}\textbackslash{}pi\_\{\textbackslash{}ast\})\textbackslash{}bar\{F\}\textasciicircum{}\{\textbackslash{}alpha\} at Displayed exact sequence after «Soit», obtained by twisting identifies a particular indexed object or symbol; that identification is not independently established by the quoted local expression.\par
\par\noindent\textbf{OM-SOL-0818 -- t3 XIV p.152 -- SOL}\par
Proof of Lemma 1.10, reduction to projective space\par
\emph{Reported reading}: IX 2.15 (i)\par
\emph{Proposal}: \textbackslash{}hyperrefext[IX][IX.prop2.14]\{2.14\} (i)\par
This finding intersects Orgogozo's \textbackslash{}sisi material at Proof of Lemma 1.10, reduction to projective space; preserve that human entry unchanged. The proposed reading \textbackslash{}hyperrefext[IX][IX.prop2.14]\{2.14\} (i) at Proof of Lemma 1.10, reduction to projective space changes a substantive mathematical assertion or bibliographic identification; no independent derivation of this exact replacement is established in this review.\par
\par\noindent\textbf{OM-SOL-0820 -- t3 XIV p.153 -- SOL}\par
§1.11, reference after «C\_i constant et constructible»\par
\emph{Reported reading}: IX 2.15\par
\emph{Proposal}: \textbackslash{}hyperrefext[IX][IX.prop2.14]\{2.14\}\par
This finding intersects Orgogozo's \textbackslash{}sisi material at §1.11, reference after «C\_i constant et constructible»; preserve that human entry unchanged. The proposed reading \textbackslash{}hyperrefext[IX][IX.prop2.14]\{2.14\} at §1.11, reference after «C\_i constant et constructible» changes a substantive mathematical assertion or bibliographic identification; no independent derivation of this exact replacement is established in this review.\par
\par\noindent\textbf{OM-SOL-0821 -- t3 XIV p.153 -- SOL}\par
§1.11, reference used to construct the resolution of an A-module sheaf\par
\emph{Reported reading}: en appliquant IX 2.15\par
\emph{Proposal}: \textbackslash{}hyperrefext[IX][IX.prop2.14]\{2.14\}\par
This finding intersects Orgogozo's \textbackslash{}sisi material at §1.11, reference used to construct the resolution of an A-module sheaf; preserve that human entry unchanged. The proposed reading \textbackslash{}hyperrefext[IX][IX.prop2.14]\{2.14\} at §1.11, reference used to construct the resolution of an A-module sheaf changes a substantive mathematical assertion or bibliographic identification; no independent derivation of this exact replacement is established in this review.\par
\par\noindent\textbf{OM-SOL-0822 -- t3 XIV p.153 -- SOL}\par
Spectral sequence of the resolution\par
\emph{Reported reading}: E\textasciicircum{}\{p,q\}\_2=\allowbreak H\textasciicircum{}\{q\}((R\textasciicircum{}\{q\}f\_\{\textbackslash{}ast\})G\textasciicircum{}\{\textbackslash{}bigdot\})\par
\emph{Proposal}: E\textasciicircum{}\{p,q\}\_2=\allowbreak H\textasciicircum{}\{p\}((R\textasciicircum{}\{q\}f\_\{\textbackslash{}ast\})G\textasciicircum{}\{\textbackslash{}bigdot\})\par
This finding intersects Orgogozo's \textbackslash{}sisi material at Spectral sequence of the resolution; preserve that human entry unchanged. The replacement E\textasciicircum{}\{p,q\}\_2=\allowbreak H\textasciicircum{}\{p\}((R\textasciicircum{}\{q\}f\_\{\textbackslash{}ast\})G\textasciicircum{}\{\textbackslash{}bigdot\}) at Spectral sequence of the resolution identifies a particular indexed object or symbol; that identification is not independently established by the quoted local expression.\par
\par\noindent\textbf{OM-SOL-0824 -- t3 XIV p.154 -- SOL}\par
Opening reference concerning inductive limits\par
\emph{Reported reading}: VI 1.6\par
\emph{Proposal}: \textbackslash{}hyperrefext[VI][VI.thm5.1]\{5.1\} et \textbackslash{}hyperrefext[VII][VII.prop3.3]\{3.3\}\par
This finding intersects Orgogozo's \textbackslash{}sisi material at Opening reference concerning inductive limits; preserve that human entry unchanged. The proposed reference \textbackslash{}hyperrefext[VI][VI.thm5.1]\{5.1\} et \textbackslash{}hyperrefext[VII][VII.prop3.3]\{3.3\} at Opening reference concerning inductive limits needs comparison with the statement actually invoked; that target has not been independently established here.\par
\par\noindent\textbf{OM-SOL-0825 -- t3 XIV p.154 -- SOL}\par
Construction of the smooth projective morphism over S′\par
\emph{Reported reading}: et un morphisme fini surjectif, un morphisme lisse et projectif\par
\emph{Proposal}: un morphisme lisse et projectif\par
The duplicated phrase may signal an unfinished construction; deleting it requires the complete sentence and its morphisms.\par
\par\noindent\textbf{OM-SOL-0826 -- t3 XIV p.155 -- SOL}\par
Beginning of the treatment of a constant sheaf of groups\par
\emph{Reported reading}: f=\allowbreak G\_X\par
\emph{Proposal}: F=\allowbreak G\_X\par
This finding intersects Orgogozo's \textbackslash{}sisi material at Beginning of the treatment of a constant sheaf of groups; preserve that human entry unchanged. The replacement F=\allowbreak G\_X at Beginning of the treatment of a constant sheaf of groups requires a complete check of the maps, their endpoints and the composition convention; this identification is not independently verified.\par
\par\noindent\textbf{OM-SOL-0827 -- t3 XIV p.156 -- SOL}\par
Reference following « Alors, d'après »\par
\emph{Reported reading}: XII 5.2 (ii)\par
\emph{Proposal}: \textbackslash{}hyperrefext[XII][XII.cor5.2]\{5.2\} (iii)\par
This finding intersects Orgogozo's \textbackslash{}sisi material at Reference following « Alors, d'après »; preserve that human entry unchanged. The proposed reading \textbackslash{}hyperrefext[XII][XII.cor5.2]\{5.2\} (iii) at Reference following « Alors, d'après » changes a substantive mathematical assertion or bibliographic identification; no independent derivation of this exact replacement is established in this review.\par
\par\noindent\textbf{OM-SOL-0828 -- t3 XIV p.156 -- SOL}\par
Reference immediately after XII 5.2 (ii)\par
\emph{Reported reading}: X 4.7\par
\emph{Proposal}: \textbackslash{}hyperrefext[X][X.cor5.2]\{5.2\}\par
This finding intersects Orgogozo's \textbackslash{}sisi material at Reference immediately after XII 5.2 (ii); preserve that human entry unchanged. The proposed reference \textbackslash{}hyperrefext[X][X.cor5.2]\{5.2\} at Reference immediately after XII 5.2 (ii) needs comparison with the statement actually invoked; that target has not been independently established here.\par
\par\noindent\textbf{OM-SOL-0829 -- t3 XIV p.156 -- SOL}\par
Remarque 1.11, « on se donne sur … un faisceau »\par
\emph{Reported reading}: sur Y un faisceau\par
\emph{Proposal}: sur \$S\$ un faisceau\par
This finding intersects Orgogozo's \textbackslash{}sisi material at Remarque 1.11, « on se donne sur … un faisceau »; preserve that human entry unchanged. The proposed reading sur \$S\$ un faisceau at Remarque 1.11, « on se donne sur … un faisceau » changes a substantive mathematical assertion or bibliographic identification; no independent derivation of this exact replacement is established in this review.\par
\par\noindent\textbf{OM-SOL-0830 -- t3 XIV p.158 -- SOL}\par
Proof of Proposition 2.3, introduction of x′ and y′\par
\emph{Reported reading}: Soient x′ \ensuremath{\in} X′ et y′ =\allowbreak  f(x′).\par
\emph{Proposal}: \$x'\textbackslash{}in X\$ et \$y'=\allowbreak f(x')\$\par
The replacement \$x'\textbackslash{}in X\$ et \$y'=\allowbreak f(x')\$ at Proof of Proposition 2.3, introduction of x′ and y′ identifies a particular indexed object or symbol; that identification is not independently established by the quoted local expression.\par
\par\noindent\textbf{OM-SOL-0831 -- t3 XIV p.158 -- SOL}\par
Last equality in the proof of Proposition 2.3\par
\emph{Reported reading}: =\allowbreak  \ensuremath{\delta}(x) =\allowbreak  1\par
\emph{Proposal}: =\allowbreak \textbackslash{}delta(x)-1\par
This finding intersects Orgogozo's \textbackslash{}sisi material at Last equality in the proof of Proposition 2.3; preserve that human entry unchanged. The proposed reading =\allowbreak \textbackslash{}delta(x)-1 at Last equality in the proof of Proposition 2.3 changes a substantive mathematical assertion or bibliographic identification; no independent derivation of this exact replacement is established in this review.\par
\par\noindent\textbf{OM-SOL-0832 -- t3 XIV p.160 -- SOL}\par
Proof of Corollary 3.3, « trivial à partir de … »\par
\emph{Reported reading}: à partir de 2.2\par
\emph{Proposal}: \textbackslash{}hyperref[XIV.cor3.2]\{3.2\}\par
This finding intersects Orgogozo's \textbackslash{}sisi material at Proof of Corollary 3.3, « trivial à partir de … »; preserve that human entry unchanged. The proposed reference \textbackslash{}hyperref[XIV.cor3.2]\{3.2\} at Proof of Corollary 3.3, « trivial à partir de … » needs comparison with the statement actually invoked; that target has not been independently established here.\par
\par\noindent\textbf{OM-SOL-0833 -- t3 XIV p.160 -- SOL}\par
Reference for the exact sequence associated with a closed subset\par
\emph{Reported reading}: (IV 3.7.1)\par
\emph{Proposal}: (\textbackslash{}hyperrefext[V][V.6.5.4]\{6.5.4\})\par
This finding intersects Orgogozo's \textbackslash{}sisi material at Reference for the exact sequence associated with a closed subset; preserve that human entry unchanged. The proposed reference (\textbackslash{}hyperrefext[V][V.6.5.4]\{6.5.4\}) at Reference for the exact sequence associated with a closed subset needs comparison with the statement actually invoked; that target has not been independently established here.\par
\par\noindent\textbf{OM-SOL-0834 -- t3 XIV p.160 -- SOL}\par
Reference after « et la suite spectrale »\par
\emph{Reported reading}: et la suite spectrale ( ),\par
\emph{Proposal}: (\textbackslash{}hyperrefext[V][V.6.4.3]\{6.4.3\})\par
This finding intersects Orgogozo's \textbackslash{}sisi material at Reference after « et la suite spectrale »; preserve that human entry unchanged. The proposed reference (\textbackslash{}hyperrefext[V][V.6.4.3]\{6.4.3\}) at Reference after « et la suite spectrale » needs comparison with the statement actually invoked; that target has not been independently established here.\par
\par\noindent\textbf{OM-SOL-0835 -- t3 XIV p.161 -- SOL}\par
The scheme on which F' is induced\par
\emph{Reported reading}: F' le faisceau induit de F sur X\par
\emph{Proposal}: \$F'\$ le faisceau induit de \$F\$ sur \$X'\$\par
The replacement \$F'\$ le faisceau induit de \$F\$ sur \$X'\$ at The scheme on which F' is induced identifies a particular indexed object or symbol; that identification is not independently established by the quoted local expression.\par
\par\noindent\textbf{OM-SOL-0836 -- t3 XIV p.162 -- SOL}\par
Last words of Remarque 3.6\par
\emph{Reported reading}: lorsqu'on suppose de plus Y de caractéristique nulle\par
\emph{Proposal}: lorsqu'on suppose de plus \$Y'\$ de caractéristique nulle.\par
The replacement lorsqu'on suppose de plus \$Y'\$ de caractéristique nulle. at Last words of Remarque 3.6 identifies a particular indexed object or symbol; that identification is not independently established by the quoted local expression.\par
\par\noindent\textbf{OM-SOL-0837 -- t3 XIV p.163 -- SOL}\par
Opening reference concerning constructibility and the support\par
\emph{Reported reading}: (IX 2.3 (ii))\par
\emph{Proposal}: (\textbackslash{}hyperrefext[IX][IX.def2.3]\{2.3\})\par
This finding intersects Orgogozo's \textbackslash{}sisi material at Opening reference concerning constructibility and the support; preserve that human entry unchanged. The proposed reference (\textbackslash{}hyperrefext[IX][IX.def2.3]\{2.3\}) at Opening reference concerning constructibility and the support needs comparison with the statement actually invoked; that target has not been independently established here.\par
\par\noindent\textbf{OM-SOL-0838 -- t3 XIV p.163 -- SOL}\par
Conclusion of the proof of Lemme 4.1\par
\emph{Reported reading}: d'où d(R\textasciicircum{}1 f\_*F) \ensuremath{\leq} 1\par
\emph{Proposal}: d'où \$d(R\textasciicircum{}\{1\}f\_\{\textbackslash{}ast\}F)\textbackslash{}leq 0\$\par
The proposed reading d'où \$d(R\textasciicircum{}\{1\}f\_\{\textbackslash{}ast\}F)\textbackslash{}leq 0\$ at Conclusion of the proof of Lemme 4.1 changes a substantive mathematical assertion or bibliographic identification; no independent derivation of this exact replacement is established in this review.\par
\par\noindent\textbf{OM-SOL-0839 -- t3 XIV p.163 -- SOL}\par
End of the statement of Lemme 4.2\par
\emph{Reported reading}: et où f est le morphisme structural\par
\emph{Proposal}: et où \$f'\$ est le morphisme structural.\par
The replacement et où \$f'\$ est le morphisme structural. at End of the statement of Lemme 4.2 requires a complete check of the maps, their endpoints and the composition convention; this identification is not independently verified.\par
\par\noindent\textbf{OM-SOL-0840 -- t3 XIV p.165 -- SOL}\par
Last formula, definition of the open subscheme of the strict localization\par
\emph{Reported reading}: \textbackslash{}widetilde\{U\}=\allowbreak \textbackslash{}widetilde\{X\}'-\textbackslash{}widetilde\{X\}\textbackslash{}times\_\{\textbackslash{}mathbf\{P\}\textasciicircum{}\{1\}\}Y'\_\{\textbackslash{}infty\}\par
\emph{Proposal}: \textbackslash{}widetilde\{U\}=\allowbreak \textbackslash{}widetilde\{X\}'-\textbackslash{}fprod\{\textbackslash{}widetilde\{X\}'\}\{Y'\_\{\textbackslash{}infty\}\}\{\textbackslash{}PP\textasciicircum{}\{1\}\}\par
The replacement \textbackslash{}widetilde\{U\}=\allowbreak \textbackslash{}widetilde\{X\}'-\textbackslash{}fprod\{\textbackslash{}widetilde\{X\}'\}\{Y'\_\{\textbackslash{}infty\}\}\{\textbackslash{}PP\textasciicircum{}\{1\}\} at Last formula, definition of the open subscheme of the strict localization identifies a particular indexed object or symbol; that identification is not independently established by the quoted local expression.\par
\par\noindent\textbf{OM-SOL-0841 -- t3 XIV p.166 -- SOL}\par
Proof of Lemme 4.5, introduction of the data\par
\emph{Reported reading}: Soient les données \$X'\$, \$f\$, \$U\$\par
\emph{Proposal}: Soient les données \$X\$, \$f\$, \$U\$\par
The replacement Soient les données \$X\$, \$f\$, \$U\$ at Proof of Lemme 4.5, introduction of the data identifies a particular indexed object or symbol; that identification is not independently established by the quoted local expression.\par
\par\noindent\textbf{OM-SOL-0842 -- t3 XIV p.167 -- SOL}\par
Hochschild–Serre spectral sequence, abutment\par
\emph{Reported reading}: \$H\textasciicircum{}\{n\}(U,F)\$\par
\emph{Proposal}: \$H\textasciicircum{}\{n\}(U'',F)\$\par
The replacement \$H\textasciicircum{}\{n\}(U'',F)\$ at Hochschild–Serre spectral sequence, abutment identifies a particular indexed object or symbol; that identification is not independently established by the quoted local expression.\par
\par\noindent\textbf{OM-SOL-0843 -- t3 XV p.168 -- SOL}\par
Introduction, dernière proposition du premier paragraphe\par
\emph{Reported reading}: implique que f est acyclique.\par
\emph{Proposal}: implique que \$g\$ est acyclique.\par
The replacement implique que \$g\$ est acyclique. at Introduction, dernière proposition du premier paragraphe requires a complete check of the maps, their endpoints and the composition convention; this identification is not independently verified.\par
\par\noindent\textbf{OM-SOL-0844 -- t3 XV p.171 -- SOL}\par
Adjunction map immediately after « D'après (i), on a »\par
\emph{Reported reading}: \$f\_\{\textbackslash{}ast\}(F)\textbackslash{}rightarrow g\_\{\textbackslash{}ast\}g\textasciicircum{}\{\textbackslash{}ast\}f\textasciicircum{}\{\textbackslash{}ast\}F\$).\par
\emph{Proposal}: on a \$f\_\{\textbackslash{}ast\}(F)\textbackslash{}rightarrow g\_\{\textbackslash{}ast\}g\textasciicircum{}\{\textbackslash{}ast\}f\_\{\textbackslash{}ast\}F\$.\par
This finding intersects Orgogozo's \textbackslash{}sisi material at Adjunction map immediately after « D'après (i), on a »; preserve that human entry unchanged. The replacement on a \$f\_\{\textbackslash{}ast\}(F)\textbackslash{}rightarrow g\_\{\textbackslash{}ast\}g\textasciicircum{}\{\textbackslash{}ast\}f\_\{\textbackslash{}ast\}F\$. at Adjunction map immediately after « D'après (i), on a » identifies a particular indexed object or symbol; that identification is not independently established by the quoted local expression.\par
\par\noindent\textbf{OM-SOL-0845 -- t3 XV p.174 -- SOL}\par
Nonabelian cohomology exact sequence, final term\par
\emph{Reported reading}: \$H\textasciicircum{}\{0\}(Y',(R\textasciicircum{}\{1\}\{g'\}\_\{\textbackslash{}ast\})g\textasciicircum{}\{\textbackslash{}ast\}F)\$\par
\emph{Proposal}: \$H\textasciicircum{}\{0\}(Y',(R\textasciicircum{}\{1\}\{g'\}\_\{\textbackslash{}ast\})\{g'\}\textasciicircum{}\{\textbackslash{}ast\}F)\$\par
The replacement \$H\textasciicircum{}\{0\}(Y',(R\textasciicircum{}\{1\}\{g'\}\_\{\textbackslash{}ast\})\{g'\}\textasciicircum{}\{\textbackslash{}ast\}F)\$ at Nonabelian cohomology exact sequence, final term identifies a particular indexed object or symbol; that identification is not independently established by the quoted local expression.\par
\par\noindent\textbf{OM-SOL-0846 -- t3 XV p.179 -- SOL}\par
Définition du morphisme canonique f′\ensuremath{\alpha}\par
\emph{Reported reading}: f'\_\{\textbackslash{}alpha\}:\textbackslash{}widetilde\{Y\}'\_\{\textbackslash{}alpha\}\textbackslash{}longrightarrow Y'\_\{\textbackslash{}alpha\}\par
\emph{Proposal}: f'\_\{\textbackslash{}alpha\}:\textbackslash{}widetilde\{Y\}'\textbackslash{}rightarrow Y'\_\{\textbackslash{}alpha\}\par
The replacement f'\_\{\textbackslash{}alpha\}:\textbackslash{}widetilde\{Y\}'\textbackslash{}rightarrow Y'\_\{\textbackslash{}alpha\} at Définition du morphisme canonique f′\ensuremath{\alpha} requires a complete check of the maps, their endpoints and the composition convention; this identification is not independently verified.\par
\par\noindent\textbf{OM-SOL-0847 -- t3 XV p.179 -- SOL}\par
Renvoi après « D'après », près du bas de la page\par
\emph{Reported reading}: XII 1.2 (v) et 1.6 (iii)\par
\emph{Proposal}: \textbackslash{}hyperrefext[IX][IX.prop1.2]\{1.2\} et \textbackslash{}hyperrefext[IX][IX.prop1.6]\{1.6\}\par
This finding intersects Orgogozo's \textbackslash{}sisi material at Renvoi après « D'après », près du bas de la page; preserve that human entry unchanged. The proposed reference \textbackslash{}hyperrefext[IX][IX.prop1.2]\{1.2\} et \textbackslash{}hyperrefext[IX][IX.prop1.6]\{1.6\} at Renvoi après « D'après », près du bas de la page needs comparison with the statement actually invoked; that target has not been independently established here.\par
\par\noindent\textbf{OM-SOL-0848 -- t3 XV p.179 -- SOL}\par
Dernière phrase, « les faisceaux f′\ensuremath{\alpha}\ensuremath{\ast}F »\par
\emph{Reported reading}: f'\_\{\textbackslash{}alpha\}\{\}\_\{\textbackslash{}ast\}F\par
\emph{Proposal}: f'\_\{\textbackslash{}alpha\}\{\}\_\{\textbackslash{}ast\}\textbackslash{}widetilde\{F\}\par
The replacement f'\_\{\textbackslash{}alpha\}\{\}\_\{\textbackslash{}ast\}\textbackslash{}widetilde\{F\} at Dernière phrase, « les faisceaux f′\ensuremath{\alpha}\ensuremath{\ast}F » identifies a particular indexed object or symbol; that identification is not independently established by the quoted local expression.\par
\par\noindent\textbf{OM-SOL-0849 -- t3 XV p.181 -- SOL}\par
Paragraphe après la définition 1.11, première phrase\par
\emph{Reported reading}: les énoncés (i) à (iii) de 1.10\par
\emph{Proposal}: les énoncés (i) à (iv) de \textbackslash{}hyperref[XV.prop1.10]\{1.10\}\par
The proposed reading les énoncés (i) à (iv) de \textbackslash{}hyperref[XV.prop1.10]\{1.10\} at Paragraphe après la définition 1.11, première phrase changes a substantive mathematical assertion or bibliographic identification; no independent derivation of this exact replacement is established in this review.\par
\par\noindent\textbf{OM-SOL-0851 -- t3 XV p.184 -- SOL}\par
Corollaire 1.17, clause concernant les fibres géométriques\par
\emph{Reported reading}: chaque fibre géométrique X\_\{\textbackslash{}bar\{z\}\} de g, relativement à un point géométrique \textbackslash{}bar\{z\} de Y algébrique sur un point z de Y\par
\emph{Proposal}: chaque fibre géométrique \$\textbackslash{}widetilde\{X\}\_\{\textbackslash{}bar\{z\}\}\$ de \$\textbackslash{}widetilde\{g\}\$, relativement à un point géométrique \$\textbackslash{}bar\{z\}\$ de \$\textbackslash{}widetilde\{Y\}\$ algébrique sur un point \$z\$ de \$\textbackslash{}widetilde\{Y\}\$\par
The replacement chaque fibre géométrique \$\textbackslash{}widetilde\{X\}\_\{\textbackslash{}bar\{z\}\}\$ de \$\textbackslash{}widetilde\{g\}\$, relativement à un point géométrique \$\textbackslash{}bar\{z\}\$ de \$\textbackslash{}widetilde\{Y\}\$ algébrique sur un point \$z\$ de \$\textbackslash{}widetilde\{Y\}\$ at Corollaire 1.17, clause concernant les fibres géométriques requires a complete check of the maps, their endpoints and the composition convention; this identification is not independently verified.\par
\par\noindent\textbf{OM-SOL-0853 -- t3 XV p.187 -- SOL}\par
Premier paragraphe, définition de Z dans la parenthèse\par
\emph{Reported reading}: Z=\allowbreak X-U\par
\emph{Proposal}: Z=\allowbreak Y-U\par
The proposed reading Z=\allowbreak Y-U at Premier paragraphe, définition de Z dans la parenthèse changes a substantive mathematical assertion or bibliographic identification; no independent derivation of this exact replacement is established in this review.\par
\par\noindent\textbf{OM-SOL-0854 -- t3 XV p.187 -- SOL}\par
Diagramme des inclusions, flèches j' et j\par
\emph{Reported reading}: Y'=\allowbreak X\textbackslash{}xrightarrow\{j'\}Z',\textbackslash{}qquad Y\textbackslash{}xrightarrow\{j\}Z\par
\emph{Proposal}: Y'=\allowbreak X\textbackslash{}xleftarrow\{j'\}Z',\textbackslash{}qquad Y\textbackslash{}xleftarrow\{j\}Z\par
The replacement Y'=\allowbreak X\textbackslash{}xleftarrow\{j'\}Z',\textbackslash{}qquad Y\textbackslash{}xleftarrow\{j\}Z at Diagramme des inclusions, flèches j' et j requires a complete check of the maps, their endpoints and the composition convention; this identification is not independently verified.\par
\par\noindent\textbf{OM-SOL-0855 -- t3 XV p.189 -- SOL}\par
Dernier paragraphe, après « en vertu de »\par
\emph{Reported reading}: « en vertu de » suivi d'un blanc avant la virgule\par
\emph{Proposal}: en vertu de \textbackslash{}text\{IV \}9.5.4,\par
This finding intersects Orgogozo's \textbackslash{}sisi material at Dernier paragraphe, après « en vertu de »; preserve that human entry unchanged. The proposed reference en vertu de \textbackslash{}text\{IV \}9.5.4, at Dernier paragraphe, après « en vertu de » needs comparison with the statement actually invoked; that target has not been independently established here.\par
\par\noindent\textbf{OM-SOL-0856 -- t3 XV p.195 -- SOL}\par
Définition de la fibre géométrique X\_alpha en y\_alpha\par
\emph{Reported reading}: la fibre géométrique de \$\textbackslash{}Spec A\_\{\textbackslash{}alpha\}\textbackslash{}\{t\textbackslash{}\}\$ sur \$\textbackslash{}Spec A\$ en \$\textbackslash{}bar\{y\}\_\{\textbackslash{}alpha\}\$\par
\emph{Proposal}: la fibre géométrique de \$\textbackslash{}Spec A\_\{\textbackslash{}alpha\}\textbackslash{}\{t\textbackslash{}\}\$ sur \$\textbackslash{}Spec A\_\{\textbackslash{}alpha\}\$ en\textasciitilde{}\$\textbackslash{}bar\{y\}\_\{\textbackslash{}alpha\}\$\par
The replacement la fibre géométrique de \$\textbackslash{}Spec A\_\{\textbackslash{}alpha\}\textbackslash{}\{t\textbackslash{}\}\$ sur \$\textbackslash{}Spec A\_\{\textbackslash{}alpha\}\$ en\textasciitilde{}\$\textbackslash{}bar\{y\}\_\{\textbackslash{}alpha\}\$ at Définition de la fibre géométrique X\_alpha en y\_alpha identifies a particular indexed object or symbol; that identification is not independently established by the quoted local expression.\par
\par\noindent\textbf{OM-SOL-0858 -- t3 XV p.198 -- SOL}\par
Définition de s, après « Soient donc ... »\par
\emph{Reported reading}: s le p.g.c.d. des ordres des groupes d'inertie\par
\emph{Proposal}: soit \$s\$ le p.p.c.m. des ordres des groupes d'inertie\par
The proposed reading soit \$s\$ le p.p.c.m. des ordres des groupes d'inertie at Définition de s, après « Soient donc ... » changes a substantive mathematical assertion or bibliographic identification; no independent derivation of this exact replacement is established in this review.\par
\par\noindent\textbf{OM-SOL-0859 -- t3 XV p.201 -- SOL}\par
Premier paragraphe, définition de \textbackslash{}mathcal\{C\}\par
\emph{Reported reading}: \textbackslash{}Oo\_\{U'\}-algèbres définissant U'\par
\emph{Proposal}: \textbackslash{}Oo\_U\par
The replacement \textbackslash{}Oo\_U at Premier paragraphe, définition de \textbackslash{}mathcal\{C\} identifies a particular indexed object or symbol; that identification is not independently established by the quoted local expression.\par
\par\noindent\textbf{OM-SOL-0860 -- t3 XVI p.208 -- SOL}\par
Same proof, immediately after the base-change formula\par
\emph{Reported reading}: est encore universellement 0-acyclique\par
\emph{Proposal}: encore universellement localement \$0\$-acyclique\par
The proposed reading encore universellement localement \$0\$-acyclique at Same proof, immediately after the base-change formula changes a substantive mathematical assertion or bibliographic identification; no independent derivation of this exact replacement is established in this review.\par
\par\noindent\textbf{OM-SOL-0861 -- t3 XVI p.210 -- SOL}\par
End of the paragraph completing the proof of 1.1\par
\emph{Reported reading}: (VI 5)\par
\emph{Proposal}: \textbackslash{}hyperrefext[VII][VII.sec5]\{5\}\par
The proposed reference \textbackslash{}hyperrefext[VII][VII.sec5]\{5\} at End of the paragraph completing the proof of 1.1 needs comparison with the statement actually invoked; that target has not been independently established here.\par
\par\noindent\textbf{OM-SOL-0863 -- t3 XVI p.214 -- SOL}\par
Corollaire 2.2, hypothèse sur F\par
\emph{Reported reading}: constructible sur X\par
\emph{Proposal}: constructible et localement constant sur \$X\$\par
The replacement constructible et localement constant sur \$X\$ at Corollaire 2.2, hypothèse sur F requires a complete check of the maps, their endpoints and the composition convention; this identification is not independently verified.\par
\par\noindent\textbf{OM-SOL-0864 -- t3 XVI p.216 -- SOL}\par
Lemme 2.4, dernière phrase\par
\emph{Reported reading}: \$n=\allowbreak 0\$ i.e. \$g\$ localement acyclique\par
\emph{Proposal}: \$n=\allowbreak 0\$ \textbackslash{}ie \$f\$ localement acyclique\par
The replacement \$n=\allowbreak 0\$ \textbackslash{}ie \$f\$ localement acyclique at Lemme 2.4, dernière phrase requires a complete check of the maps, their endpoints and the composition convention; this identification is not independently verified.\par
\par\noindent\textbf{OM-SOL-0865 -- t3 XVI p.217 -- SOL}\par
Premier paragraphe, hypothèse de normalité\par
\emph{Reported reading}: l'hypothèse que \$X\$ est normal\par
\emph{Proposal}: l'hypothèse que \$Y\$ est normal\par
This finding intersects Orgogozo's \textbackslash{}sisi material at Premier paragraphe, hypothèse de normalité; preserve that human entry unchanged. The proposed reading l'hypothèse que \$Y\$ est normal at Premier paragraphe, hypothèse de normalité changes a substantive mathematical assertion or bibliographic identification; no independent derivation of this exact replacement is established in this review.\par
\par\noindent\textbf{OM-SOL-0866 -- t3 XVI p.217 -- SOL}\par
Premier paragraphe, « D'autre part, l'hypothèse… »\par
\emph{Reported reading}: l'hypothèse que \$g\$ est localement \$0\$-acyclique\par
\emph{Proposal}: l'hypothèse que \$f\$ est localement \$0\$-acyclique\par
The proposed reading l'hypothèse que \$f\$ est localement \$0\$-acyclique at Premier paragraphe, « D'autre part, l'hypothèse… » changes a substantive mathematical assertion or bibliographic identification; no independent derivation of this exact replacement is established in this review.\par
\par\noindent\textbf{OM-SOL-0867 -- t3 XVI p.217 -- SOL}\par
Même phrase, images directes dont la formation commute au changement de base\par
\emph{Reported reading}: R\textasciicircum{}\{q\}g\_\{\textbackslash{}ast\}(G\_y)\par
\emph{Proposal}: R\textasciicircum{}\{q\}g\_\{\textbackslash{}ast\}(G\_\{y\_1\})\par
The replacement R\textasciicircum{}\{q\}g\_\{\textbackslash{}ast\}(G\_\{y\_1\}) at Même phrase, images directes dont la formation commute au changement de base identifies a particular indexed object or symbol; that identification is not independently established by the quoted local expression.\par
\par\noindent\textbf{OM-SOL-0868 -- t3 XVI p.219 -- SOL}\par
Paragraphe introduisant les faisceaux de cohomologie locale\par
\emph{Reported reading}: V 4\par
\emph{Proposal}: \textbackslash{}hyperrefext[V][V.sec6]\{6\}\par
This finding intersects Orgogozo's \textbackslash{}sisi material at Paragraphe introduisant les faisceaux de cohomologie locale; preserve that human entry unchanged. The proposed reference \textbackslash{}hyperrefext[V][V.sec6]\{6\} at Paragraphe introduisant les faisceaux de cohomologie locale needs comparison with the statement actually invoked; that target has not been independently established here.\par
\par\noindent\textbf{OM-SOL-0869 -- t3 XVI p.220 -- SOL}\par
Démonstration, choix des points géométriques\par
\emph{Reported reading}: \$\textbackslash{}bar\{x\}\$ un point géométrique de \$X\$ au-dessus de \$S\$\par
\emph{Proposal}: \$\textbackslash{}bar\{x\}\$ un point géométrique de \$X\$ au-dessus de \$\textbackslash{}bar\{s\}\$\par
The replacement \$\textbackslash{}bar\{x\}\$ un point géométrique de \$X\$ au-dessus de \$\textbackslash{}bar\{s\}\$ at Démonstration, choix des points géométriques requires a complete check of the maps, their endpoints and the composition convention; this identification is not independently verified.\par
\par\noindent\textbf{OM-SOL-0872 -- t3 XVI p.224 -- SOL}\par
Démonstration du corollaire 3.4, formule des fibres\par
\emph{Reported reading}: (R\textasciicircum{}\{q\}j\_*g\textasciicircum{}*F)\_\{\textbackslash{}bar\{x\}\}=\allowbreak H\textasciicircum{}\{q\}(\textbackslash{}widetilde\{U\},G\_\{\textbackslash{}widetilde\{U\}\})\textbackslash{}simeq H\textasciicircum{}\{q\}(\textbackslash{}widetilde\{X\},G\_\{\textbackslash{}widetilde\{X\}\})\par
\emph{Proposal}: (R\textasciicircum{}\{1\}j\_\{\textbackslash{}ast\}g\textasciicircum{}\{\textbackslash{}ast\}F)\_\{\textbackslash{}bar\{x\}\}=\allowbreak H\textasciicircum{}\{1\}(\textbackslash{}widetilde\{U\},G\_\{\textbackslash{}widetilde\{U\}\})\textbackslash{}simeq H\textasciicircum{}\{1\}(\textbackslash{}widetilde\{X\},G\_\{\textbackslash{}widetilde\{X\}\})\par
The replacement (R\textasciicircum{}\{1\}j\_\{\textbackslash{}ast\}g\textasciicircum{}\{\textbackslash{}ast\}F)\_\{\textbackslash{}bar\{x\}\}=\allowbreak H\textasciicircum{}\{1\}(\textbackslash{}widetilde\{U\},G\_\{\textbackslash{}widetilde\{U\}\})\textbackslash{}simeq H\textasciicircum{}\{1\}(\textbackslash{}widetilde\{X\},G\_\{\textbackslash{}widetilde\{X\}\}) at Démonstration du corollaire 3.4, formule des fibres requires a complete check of the maps, their endpoints and the composition convention; this identification is not independently verified.\par
\par\noindent\textbf{OM-SOL-0874 -- t3 XVI p.225 -- SOL}\par
Réduction au cas m=\allowbreak 1, après « c'est-à-dire »\par
\emph{Reported reading}: \textbackslash{}bar\{U\}\_m=\allowbreak \textbackslash{}widetilde\{U\}\par
\emph{Proposal}: \$\textbackslash{}widetilde\{U\}\_m=\allowbreak \textbackslash{}widetilde\{U\}\$\par
The replacement \$\textbackslash{}widetilde\{U\}\_m=\allowbreak \textbackslash{}widetilde\{U\}\$ at Réduction au cas m=\allowbreak 1, après « c'est-à-dire » identifies a particular indexed object or symbol; that identification is not independently established by the quoted local expression.\par
\par\noindent\textbf{OM-SOL-0875 -- t3 XVI p.228 -- SOL}\par
Dernière égalité, après « enfin »\par
\emph{Reported reading}: \textbackslash{}RR\textasciicircum{}\{p\}f\_\{\textbackslash{}ast\}j\_\{\textbackslash{}ast\}(F|U)=\allowbreak \textbackslash{}RR\textasciicircum{}\{p\}f\_\{\textbackslash{}ast\}(F)\par
\emph{Proposal}: R\textasciicircum{}\{p\}f\_\{\textbackslash{}ast\}j\_\{\textbackslash{}ast\}(F|U)=\allowbreak R\textasciicircum{}\{p\}f\_\{\textbackslash{}ast\}(F)\par
This finding intersects Orgogozo's \textbackslash{}sisi material at Dernière égalité, après « enfin »; preserve that human entry unchanged. The replacement R\textasciicircum{}\{p\}f\_\{\textbackslash{}ast\}j\_\{\textbackslash{}ast\}(F|U)=\allowbreak R\textasciicircum{}\{p\}f\_\{\textbackslash{}ast\}(F) at Dernière égalité, après « enfin » identifies a particular indexed object or symbol; that identification is not independently established by the quoted local expression.\par
\par\noindent\textbf{OM-SOL-0876 -- t3 XVI p.233 -- SOL}\par
Théorème 4.1, reference for the base-change morphisms\par
\emph{Reported reading}: XII 4.1.2\par
\emph{Proposal}: \textbackslash{}hyperrefext[XII][XII.eq4.2]\{4.2\}\par
This finding intersects Orgogozo's \textbackslash{}sisi material at Théorème 4.1, reference for the base-change morphisms; preserve that human entry unchanged. The proposed reference \textbackslash{}hyperrefext[XII][XII.eq4.2]\{4.2\} at Théorème 4.1, reference for the base-change morphisms needs comparison with the statement actually invoked; that target has not been independently established here.\par
\par\noindent\textbf{OM-SOL-0877 -- t3 XVI p.239 -- SOL}\par
Paragraphe « Pour traiter ce dernier cas », renvoi sur la dimension cohomologique\par
\emph{Reported reading}: IX 7.7\par
\emph{Proposal}: \textbackslash{}hyperrefext[IX][IX.cor5.7]\{5.7\}\par
This finding intersects Orgogozo's \textbackslash{}sisi material at Paragraphe « Pour traiter ce dernier cas », renvoi sur la dimension cohomologique; preserve that human entry unchanged. The proposed reference \textbackslash{}hyperrefext[IX][IX.cor5.7]\{5.7\} at Paragraphe « Pour traiter ce dernier cas », renvoi sur la dimension cohomologique needs comparison with the statement actually invoked; that target has not been independently established here.\par
\par\noindent\textbf{OM-SOL-0878 -- t3 XVI p.247 -- SOL}\par
Proof of 5.1, reference after « Ensuite on applique »\par
\emph{Reported reading}: IX 2.15\par
\emph{Proposal}: \textbackslash{}hyperrefext[IX][IX.prop2.14]\{2.14\}\par
This finding intersects Orgogozo's \textbackslash{}sisi material at Proof of 5.1, reference after « Ensuite on applique »; preserve that human entry unchanged. The proposed reference \textbackslash{}hyperrefext[IX][IX.prop2.14]\{2.14\} at Proof of 5.1, reference after « Ensuite on applique » needs comparison with the statement actually invoked; that target has not been independently established here.\par
\par\noindent\textbf{OM-SOL-0879 -- t3 XVI p.249 -- SOL}\par
Bibliography [2], Godement's initial\par
\emph{Reported reading}: Godement, P.\par
\emph{Proposal}: Godement, R.\par
The proposed reading Godement, R. involves a printed glyph or a missing stroke at Bibliography [2], Godement's initial; the textual record does not independently establish that glyph.\par
\par\noindent\textbf{OM-SOL-0883 -- t3 XVII p.254 -- SOL}\par
0.3, référence après les foncteurs de translation\par
\emph{Reported reading}: (CD I 1.10)\par
\emph{Proposal}: \textbackslash{}cite\{SGA4-17:CD\} I 1.1.0\par
This finding intersects Orgogozo's \textbackslash{}sisi material at 0.3, référence après les foncteurs de translation; preserve that human entry unchanged. The proposed reference \textbackslash{}cite\{SGA4-17:CD\} I 1.1.0 at 0.3, référence après les foncteurs de translation needs comparison with the statement actually invoked; that target has not been independently established here.\par
\par\noindent\textbf{OM-SOL-0884 -- t3 XVII p.255 -- SOL}\par
Première ligne, référence après « complétant »\par
\emph{Reported reading}: ([CD] I 1.1)\par
\emph{Proposal}: \textbackslash{}cite\{SGA4-17:CD\} II 2.1.1\par
This finding intersects Orgogozo's \textbackslash{}sisi material at Première ligne, référence après « complétant »; preserve that human entry unchanged. The proposed reference \textbackslash{}cite\{SGA4-17:CD\} II 2.1.1 at Première ligne, référence après « complétant » needs comparison with the statement actually invoked; that target has not been independently established here.\par
\par\noindent\textbf{OM-SOL-0885 -- t3 XVII p.256 -- SOL}\par
0.10, final parenthesis\par
\emph{Reported reading}: l'univers fini\par
\emph{Proposal}: l'univers fixé\par
This finding intersects Orgogozo's \textbackslash{}sisi material at 0.10, final parenthesis; preserve that human entry unchanged. The proposed reading l'univers fixé at 0.10, final parenthesis changes a substantive mathematical assertion or bibliographic identification; no independent derivation of this exact replacement is established in this review.\par
\par\noindent\textbf{OM-SOL-0886 -- t3 XVII p.258 -- SOL}\par
1.1.1, upper row of the commutative diagram, between the second 0 and the first X\par
\emph{Reported reading}: 0=\allowbreak X\par
\emph{Proposal}: 0\textbackslash{}longrightarrow X\par
This finding intersects Orgogozo's \textbackslash{}sisi material at 1.1.1, upper row of the commutative diagram, between the second 0 and the first X; preserve that human entry unchanged. The replacement 0\textbackslash{}longrightarrow X at 1.1.1, upper row of the commutative diagram, between the second 0 and the first X requires a complete check of the maps, their endpoints and the composition convention; this identification is not independently verified.\par
\par\noindent\textbf{OM-SOL-0888 -- t3 XVII p.260 -- SOL}\par
(1.1.4.2), second condition beneath the summation sign\par
\emph{Reported reading}: i<\_2 i\par
\emph{Proposal}: j<\_2 i\par
This finding intersects Orgogozo's \textbackslash{}sisi material at (1.1.4.2), second condition beneath the summation sign; preserve that human entry unchanged. The proposed reading j<\_2 i at (1.1.4.2), second condition beneath the summation sign changes a substantive mathematical assertion or bibliographic identification; no independent derivation of this exact replacement is established in this review.\par
\par\noindent\textbf{OM-SOL-0890 -- t3 XVII p.261 -- SOL}\par
(1.1.4.4), second row\par
\emph{Reported reading}: d\_j[1\_i]\textasciicircum{}\{\textbackslash{}underline\{k\}\}=\allowbreak -d\_j[1\_i]\textasciicircum{}\{\textbackslash{}underline\{k\}+1\_i\}\par
\emph{Proposal}: d\_j[1\_i]\textasciicircum{}\{\textbackslash{}sheafk \} =\allowbreak  -d\_j\textasciicircum{}\{\textbackslash{}sheafk +1\_i\}\par
The proposed reading d\_j[1\_i]\textasciicircum{}\{\textbackslash{}sheafk \} =\allowbreak  -d\_j\textasciicircum{}\{\textbackslash{}sheafk +1\_i\} at (1.1.4.4), second row changes a substantive mathematical assertion or bibliographic identification; no independent derivation of this exact replacement is established in this review.\par
\par\noindent\textbf{OM-SOL-0891 -- t3 XVII p.262 -- SOL}\par
(1.1.5.2), exponent of sigma\par
\emph{Reported reading}: \textbackslash{}sigma\textasciicircum{}\{\textbackslash{}underline\{k\}\}=\allowbreak (-1)\textasciicircum{}\{k+1\}\textbackslash{}id\_\{F(K\textasciicircum{}\{-k-1\})\}\par
\emph{Proposal}: \textbackslash{}sigma\textasciicircum{}\{k\}=\allowbreak (-1)\textasciicircum{}\{k+1\}\textbackslash{}id\_\{F(K\textasciicircum{}\{-k-1\})\}\par
This finding intersects Orgogozo's \textbackslash{}sisi material at (1.1.5.2), exponent of sigma; preserve that human entry unchanged. The replacement \textbackslash{}sigma\textasciicircum{}\{k\}=\allowbreak (-1)\textasciicircum{}\{k+1\}\textbackslash{}id\_\{F(K\textasciicircum{}\{-k-1\})\} at (1.1.5.2), exponent of sigma identifies a particular indexed object or symbol; that identification is not independently established by the quoted local expression.\par
\par\noindent\textbf{OM-SOL-0892 -- t3 XVII p.263 -- SOL}\par
Sentence following the dual short exact sequence\par
\emph{Reported reading}: f\textasciicircum{}\{\textbackslash{}ast\} de F(K) dans F(L)\par
\emph{Proposal}: f\textasciicircum{}\{\textbackslash{}ast\}\textbackslash{}text\{ de \}F(K)\textbackslash{}text\{ dans \}F(M)\par
This finding intersects Orgogozo's \textbackslash{}sisi material at Sentence following the dual short exact sequence; preserve that human entry unchanged. The replacement f\textasciicircum{}\{\textbackslash{}ast\}\textbackslash{}text\{ de \}F(K)\textbackslash{}text\{ dans \}F(M) at Sentence following the dual short exact sequence requires a complete check of the maps, their endpoints and the composition convention; this identification is not independently verified.\par
\par\noindent\textbf{OM-SOL-0893 -- t3 XVII p.263 -- SOL}\par
Tuple naming the triangle, in the paragraph following the dual exact sequence\par
\emph{Reported reading}: (k,L,M,u,v,f)\par
\emph{Proposal}: (K,L,M,u,v,f)\par
The replacement (K,L,M,u,v,f) at Tuple naming the triangle, in the paragraph following the dual exact sequence identifies a particular indexed object or symbol; that identification is not independently established by the quoted local expression.\par
\par\noindent\textbf{OM-SOL-0897 -- t3 XVII p.266 -- SOL}\par
Formule (1.1.8.1) et définition du signe qui la suit\par
\emph{Reported reading}: \textbackslash{}rho\textasciicircum{}\{\textbackslash{}sheafk\}=\allowbreak \textbackslash{}prod\_i\textbackslash{}eta\_i\textasciicircum{}\{k\_i\}, avec l'exposant k\_i\textbackslash{}epsilon(i)\textbackslash{}epsilon(\textbackslash{}Psi(i))\par
\emph{Proposal}: \textbackslash{}rho\textasciicircum{}\{\textbackslash{}sheafk\}=\allowbreak (-1)\textasciicircum{}\{A(\textbackslash{}sheafk)\}:\textbackslash{};\textbackslash{}mbox\{\textbackslash{}rm automorphisme de \}F\textbackslash{}circ(G\_j)(K\_i\textasciicircum{}\{k\_i\textbackslash{}epsilon\_I(i)\textbackslash{}epsilon\_J(\textbackslash{}Psi(i))\}),\textbackslash{}qquad A(\textbackslash{}sheafk)=\allowbreak \textbackslash{}sum\_\{\textbackslash{}epsilon\_I(i)=\allowbreak \textbackslash{}epsilon\_J(\textbackslash{}Psi(i))=\allowbreak -\}k\_i+\textbackslash{}sum\_\{\textbackslash{}epsilon\_J(j)=\allowbreak -\}\textbackslash{}sum\_\{\textbackslash{}substack\{\textbackslash{}Psi(a)=\allowbreak \textbackslash{}Psi(b)=\allowbreak j\textbackslash{}\textbackslash{}a<b\}\}k\_a k\_b\par
This finding intersects Orgogozo's \textbackslash{}sisi material at Formule (1.1.8.1) et définition du signe qui la suit; preserve that human entry unchanged. The replacement \textbackslash{}rho\textasciicircum{}\{\textbackslash{}sheafk\}=\allowbreak (-1)\textasciicircum{}\{A(\textbackslash{}sheafk)\}:\textbackslash{};\textbackslash{}mbox\{\textbackslash{}rm automorphisme de \}F\textbackslash{}circ(G\_j)(K\_i\textasciicircum{}\{k\_i\textbackslash{}epsilon\_I(i)\textbackslash{}epsilon\_J(\textbackslash{}Psi(i))\}),\textbackslash{}qquad A(\textbackslash{}sheafk)=\allowbreak \textbackslash{}sum\_\{\textbackslash{}epsilon\_I(i)=\allowbreak \textbackslash{}epsilon\_J(\textbackslash{}Psi(i))=\allowbreak -\}k\_i+\textbackslash{}sum\_\{\textbackslash{}epsilon\_J(j)=\allowbreak -\}\textbackslash{}sum\_\{\textbackslash{}substack\{\textbackslash{}Psi(a)=\allowbreak \textbackslash{}Psi(b)=\allowbreak j\textbackslash{}\textbackslash{}a<b\}\}k\_a k\_b at Formule (1.1.8.1) et définition du signe qui la suit identifies a particular indexed object or symbol; that identification is not independently established by the quoted local expression.\par
\par\noindent\textbf{OM-SOL-0898 -- t3 XVII p.267 -- SOL}\par
1.1.9, définition d'un morphisme de contraction\par
\emph{Reported reading}: « pour tout X\textbackslash{}in\textbackslash{}Aa »\par
\emph{Proposal}: X\textbackslash{}in\textbackslash{}Aa\_i\par
The replacement X\textbackslash{}in\textbackslash{}Aa\_i at 1.1.9, définition d'un morphisme de contraction identifies a particular indexed object or symbol; that identification is not independently established by the quoted local expression.\par
\par\noindent\textbf{OM-SOL-0900 -- t3 XVII p.270 -- SOL}\par
Première phrase, famille des objets de (D)\par
\emph{Reported reading}: (X\_n)\_\{n\textbackslash{}in\textbackslash{}ZZ\}\par
\emph{Proposal}: (X\textasciicircum{}\{n\})\_\{n\textbackslash{}in\textbackslash{}ZZ\}\par
This finding intersects Orgogozo's \textbackslash{}sisi material at Première phrase, famille des objets de (D); preserve that human entry unchanged. The replacement (X\textasciicircum{}\{n\})\_\{n\textbackslash{}in\textbackslash{}ZZ\} at Première phrase, famille des objets de (D) identifies a particular indexed object or symbol; that identification is not independently established by the quoted local expression.\par
\par\noindent\textbf{OM-SOL-0901 -- t3 XVII p.270 -- SOL}\par
Définition du complexe D après (1.1.11.1)\par
\emph{Reported reading}: d\textasciicircum{}n=\allowbreak d\textasciicircum{}n\par
\emph{Proposal}: d\textasciicircum{}\{n\}=\allowbreak d\_n\par
The replacement d\textasciicircum{}\{n\}=\allowbreak d\_n at Définition du complexe D après (1.1.11.1) requires a complete check of the maps, their endpoints and the composition convention; this identification is not independently verified.\par
\par\noindent\textbf{OM-SOL-0902 -- t3 XVII p.271 -- SOL}\par
(1.1.12.3), right-hand Hom\par
\emph{Reported reading}: \textbackslash{}Hom(K,F\_\{\textbackslash{}bigdot\}(L))\par
\emph{Proposal}: \textbackslash{}Hom\textasciicircum{}\{\textbackslash{}bigdot\}(K,F\_\{\textbackslash{}bigdot\}(L))\par
This finding intersects Orgogozo's \textbackslash{}sisi material at (1.1.12.3), right-hand Hom; preserve that human entry unchanged. The proposed reading \textbackslash{}Hom\textasciicircum{}\{\textbackslash{}bigdot\}(K,F\_\{\textbackslash{}bigdot\}(L)) at (1.1.12.3), right-hand Hom changes a substantive mathematical assertion or bibliographic identification; no independent derivation of this exact replacement is established in this review.\par
\par\noindent\textbf{OM-SOL-0903 -- t3 XVII p.272 -- SOL}\par
Paragraph after Proposition 1.1.14, « il résulte de »\par
\emph{Reported reading}: 1.1.4\par
\emph{Proposal}: 1.1.14\par
This finding intersects Orgogozo's \textbackslash{}sisi material at Paragraph after Proposition 1.1.14, « il résulte de »; preserve that human entry unchanged. The proposed reference 1.1.14 at Paragraph after Proposition 1.1.14, « il résulte de » needs comparison with the statement actually invoked; that target has not been independently established here.\par
\par\noindent\textbf{OM-SOL-0907 -- t3 XVII p.276 -- SOL}\par
Première phrase, morphisme de triangles\par
\emph{Reported reading}: \$\textbackslash{}Delta'\textbackslash{}xrightarrow\{s\_1\}\textbackslash{}Delta''\_1\$ dans \$S\$, admettant \$s\$ pour composante.\par
\emph{Proposal}: \$\textbackslash{}Delta'\textbackslash{}xrightarrow\{s\}\textbackslash{}Delta''\_0\$ dans \$S\$, admettant \$s\_X\$ pour composante.\par
This finding intersects Orgogozo's \textbackslash{}sisi material at Première phrase, morphisme de triangles; preserve that human entry unchanged. The replacement \$\textbackslash{}Delta'\textbackslash{}xrightarrow\{s\}\textbackslash{}Delta''\_0\$ dans \$S\$, admettant \$s\_X\$ pour composante. at Première phrase, morphisme de triangles requires a complete check of the maps, their endpoints and the composition convention; this identification is not independently verified.\par
\par\noindent\textbf{OM-SOL-0908 -- t3 XVII p.276 -- SOL}\par
Diagramme (3), étiquettes des flèches\par
\emph{Reported reading}: \$Z\textbackslash{}xrightarrow\{f\}Z''\_0\$, \$X\textbackslash{}xrightarrow\{f\}X''\_0\$, verticale de droite \$w'\$.\par
\emph{Proposal}: \textbackslash{}xymatrix\{Z\textbackslash{}ar[r]\textasciicircum{}\{f\_1\}\textbackslash{}ar[d]\_\{w\}\&Z''\_0\textbackslash{}ar[d]\_\{w''\}\textbackslash{}\textbackslash{}X\textbackslash{}ar[r]\textasciicircum{}\{f\_1\}\&X''\_0\}\par
This finding intersects Orgogozo's \textbackslash{}sisi material at Diagramme (3), étiquettes des flèches; preserve that human entry unchanged. The proposed reading \textbackslash{}xymatrix\{Z\textbackslash{}ar[r]\textasciicircum{}\{f\_1\}\textbackslash{}ar[d]\_\{w\}\&Z''\_0\textbackslash{}ar[d]\_\{w''\}\textbackslash{}\textbackslash{}X\textbackslash{}ar[r]\textasciicircum{}\{f\_1\}\&X''\_0\} at Diagramme (3), étiquettes des flèches changes a substantive mathematical assertion or bibliographic identification; no independent derivation of this exact replacement is established in this review.\par
\par\noindent\textbf{OM-SOL-0909 -- t3 XVII p.276 -- SOL}\par
Égalité de commutation après le diagramme (3)\par
\emph{Reported reading}: \$tw''f=\allowbreak tfw\$\par
\emph{Proposal}: \$t\_0w''f\_1=\allowbreak t\_0f\_1w\$\par
This finding intersects Orgogozo's \textbackslash{}sisi material at Égalité de commutation après le diagramme (3); preserve that human entry unchanged. The replacement \$t\_0w''f\_1=\allowbreak t\_0f\_1w\$ at Égalité de commutation après le diagramme (3) identifies a particular indexed object or symbol; that identification is not independently established by the quoted local expression.\par
\par\noindent\textbf{OM-SOL-0910 -- t3 XVII p.276 -- SOL}\par
Fin de l'alinéa (b)\par
\emph{Reported reading}: \$Sf=\allowbreak Sg\$\par
\emph{Proposal}: \$sf=\allowbreak sg\$\par
This finding intersects Orgogozo's \textbackslash{}sisi material at Fin de l'alinéa (b); preserve that human entry unchanged. The replacement \$sf=\allowbreak sg\$ at Fin de l'alinéa (b) requires a complete check of the maps, their endpoints and the composition convention; this identification is not independently verified.\par
\par\noindent\textbf{OM-SOL-0911 -- t3 XVII p.277 -- SOL}\par
Premier paragraphe, remplacement des objets par des buts de flèches\par
\emph{Reported reading}: une flèche de \$S\$ de source \$A\_i\$, \$A'\_1\$, \$A\_1\$ ou \$A''\_1\$\par
\emph{Proposal}: \$S\_i\$ ou \$S\_1\$\par
This finding intersects Orgogozo's \textbackslash{}sisi material at Premier paragraphe, remplacement des objets par des buts de flèches; preserve that human entry unchanged. The proposed reading \$S\_i\$ ou \$S\_1\$ at Premier paragraphe, remplacement des objets par des buts de flèches changes a substantive mathematical assertion or bibliographic identification; no independent derivation of this exact replacement is established in this review.\par
\par\noindent\textbf{OM-SOL-0912 -- t3 XVII p.277 -- SOL}\par
Grand diagramme final, case supérieure gauche\par
\emph{Reported reading}: \$\textbackslash{}Hom'Y,X''[1])\$\par
\emph{Proposal}: \textbackslash{}Hom(Y,X''[-1])\par
This finding intersects Orgogozo's \textbackslash{}sisi material at Grand diagramme final, case supérieure gauche; preserve that human entry unchanged. The proposed reading \textbackslash{}Hom(Y,X''[-1]) at Grand diagramme final, case supérieure gauche changes a substantive mathematical assertion or bibliographic identification; no independent derivation of this exact replacement is established in this review.\par
\par\noindent\textbf{OM-SOL-0913 -- t3 XVII p.277 -- SOL}\par
Grand diagramme final, case inférieure gauche\par
\emph{Reported reading}: \$\textbackslash{}Hom(Y,RF(A'',A\_i)[1])\$\par
\emph{Proposal}: \textbackslash{}Hom(Y,RF(A''\_1,A\_i)[-1])\par
This finding intersects Orgogozo's \textbackslash{}sisi material at Grand diagramme final, case inférieure gauche; preserve that human entry unchanged. The replacement \textbackslash{}Hom(Y,RF(A''\_1,A\_i)[-1]) at Grand diagramme final, case inférieure gauche identifies a particular indexed object or symbol; that identification is not independently established by the quoted local expression.\par
\par\noindent\textbf{OM-SOL-0916 -- t3 XVII p.279 -- SOL}\par
Paragraphe « En particulier », domaine de la restriction\par
\emph{Reported reading}: \$D\textasciicircum{}\{+\}(\textbackslash{}Bb)\$\par
\emph{Proposal}: \$D\textasciicircum{}\{+\}(\textbackslash{}Aa)\$\par
This finding intersects Orgogozo's \textbackslash{}sisi material at Paragraphe « En particulier », domaine de la restriction; preserve that human entry unchanged. The proposed reading \$D\textasciicircum{}\{+\}(\textbackslash{}Aa)\$ at Paragraphe « En particulier », domaine de la restriction changes a substantive mathematical assertion or bibliographic identification; no independent derivation of this exact replacement is established in this review.\par
\par\noindent\textbf{OM-SOL-0917 -- t3 XVII p.280 -- SOL}\par
Premier paragraphe de démonstration, objet résolu\par
\emph{Reported reading}: résolutions gauches de \$A\_0\$\par
\emph{Proposal}: résolutions gauches de \$A\$\par
This finding intersects Orgogozo's \textbackslash{}sisi material at Premier paragraphe de démonstration, objet résolu; preserve that human entry unchanged. The replacement résolutions gauches de \$A\$ at Premier paragraphe de démonstration, objet résolu identifies a particular indexed object or symbol; that identification is not independently established by the quoted local expression.\par
\par\noindent\textbf{OM-SOL-0919 -- t3 XVII p.280 -- SOL}\par
Proposition 1.2.7, première phrase\par
\emph{Reported reading}: Soit \$F\$ un foncteur additif d'une catégorie abélienne dans une catégorie abélienne \$\textbackslash{}Bb\$.\par
\emph{Proposal}: Soit \$F\$ un foncteur additif d'une catégorie abélienne \$\textbackslash{}Aa\$ dans une catégorie abélienne \$\textbackslash{}Bb\$.\par
The proposed reading Soit \$F\$ un foncteur additif d'une catégorie abélienne \$\textbackslash{}Aa\$ dans une catégorie abélienne \$\textbackslash{}Bb\$. at Proposition 1.2.7, première phrase changes a substantive mathematical assertion or bibliographic identification; no independent derivation of this exact replacement is established in this review.\par
\par\noindent\textbf{OM-SOL-0921 -- t3 XVII p.281 -- SOL}\par
Deuxième paragraphe, après « Puisque »\par
\emph{Reported reading}: H\textasciicircum{}\{j\}(\textbackslash{}tau\_\{\textbackslash{}geq i\}(L))=\allowbreak 0\par
\emph{Proposal}: H\textasciicircum{}\{j\}(\textbackslash{}tau\_\{\textbackslash{}geq i\}(M))=\allowbreak 0\par
This finding intersects Orgogozo's \textbackslash{}sisi material at Deuxième paragraphe, après « Puisque »; preserve that human entry unchanged. The proposed reading H\textasciicircum{}\{j\}(\textbackslash{}tau\_\{\textbackslash{}geq i\}(M))=\allowbreak 0 at Deuxième paragraphe, après « Puisque » changes a substantive mathematical assertion or bibliographic identification; no independent derivation of this exact replacement is established in this review.\par
\par\noindent\textbf{OM-SOL-0925 -- t3 XVII p.284 -- SOL}\par
Proposition 2.1.3, nom initial de la catégorie\par
\emph{Reported reading}: Soient \$\textbackslash{}epsilon\$ une catégorie sur \$\textbackslash{}Bb\$\par
\emph{Proposal}: Soient \$\textbackslash{}Ee\$ une catégorie sur \$\textbackslash{}Bb\$\par
The proposed reading Soient \$\textbackslash{}Ee\$ une catégorie sur \$\textbackslash{}Bb\$ at Proposition 2.1.3, nom initial de la catégorie changes a substantive mathematical assertion or bibliographic identification; no independent derivation of this exact replacement is established in this review.\par
\par\noindent\textbf{OM-SOL-0926 -- t3 XVII p.284 -- SOL}\par
Diagramme (2.1.3.1), flèche horizontale inférieure\par
\emph{Reported reading}: x'\textbackslash{}longrightarrow x, sans étiquette\par
\emph{Proposal}: x'\textbackslash{}xrightarrow\{g\}x\par
The proposed reading x'\textbackslash{}xrightarrow\{g\}x at Diagramme (2.1.3.1), flèche horizontale inférieure changes a substantive mathematical assertion or bibliographic identification; no independent derivation of this exact replacement is established in this review.\par
\par\noindent\textbf{OM-SOL-0927 -- t3 XVII p.285 -- SOL}\par
Première phrase, seconde image directe supplémentaire\par
\emph{Reported reading}: g'\_\{\textbackslash{}ast\}f'\_\{\textbackslash{}ast\}\{g'\}\textasciicircum{}\{\textbackslash{}ast\}F\par
\emph{Proposal}: g\_\{\textbackslash{}ast\}f'\_\{\textbackslash{}ast\}\{g'\}\textasciicircum{}\{\textbackslash{}ast\}F\par
The replacement g\_\{\textbackslash{}ast\}f'\_\{\textbackslash{}ast\}\{g'\}\textasciicircum{}\{\textbackslash{}ast\}F at Première phrase, seconde image directe supplémentaire identifies a particular indexed object or symbol; that identification is not independently established by the quoted local expression.\par
\par\noindent\textbf{OM-SOL-0928 -- t3 XVII p.285 -- SOL}\par
Deuxième phrase, seconde image réciproque supplémentaire\par
\emph{Reported reading}: \{f'\}\textasciicircum{}\{\textbackslash{}ast\}f'\_\{\textbackslash{}ast\}F\par
\emph{Proposal}: f\textasciicircum{}\{\textbackslash{}ast\}f\_\{\textbackslash{}ast\}F\par
The replacement f\textasciicircum{}\{\textbackslash{}ast\}f\_\{\textbackslash{}ast\}F at Deuxième phrase, seconde image réciproque supplémentaire requires a complete check of the maps, their endpoints and the composition convention; this identification is not independently verified.\par
\par\noindent\textbf{OM-SOL-0929 -- t3 XVII p.285 -- SOL}\par
Preuve de 2.1.3, première ligne des prétendues bijections\par
\emph{Reported reading}: \textbackslash{}Hom\_y(F,F)=\allowbreak \textbackslash{}Hom\_\{g'\}(\{g'\}\textasciicircum{}\{\textbackslash{}ast\}F,F)=\allowbreak \textbackslash{}Hom\_\{fg'\}(\{g'\}\textasciicircum{}\{\textbackslash{}ast\}F,f\_\{\textbackslash{}ast\}F)\par
\emph{Proposal}: \textbackslash{}[ \textbackslash{}Hom\_\{x'\}(f'\_\{\textbackslash{}ast\}\{g'\}\textasciicircum{}\{\textbackslash{}ast\}F,g\textasciicircum{}\{\textbackslash{}ast\}f\_\{\textbackslash{}ast\}F)=\allowbreak \textbackslash{}Hom\_\{f'\}(\{g'\}\textasciicircum{}\{\textbackslash{}ast\}F,g\textasciicircum{}\{\textbackslash{}ast\}f\_\{\textbackslash{}ast\}F)=\allowbreak \textbackslash{}Hom\_\{gf'\}(\{g'\}\textasciicircum{}\{\textbackslash{}ast\}F,f\_\{\textbackslash{}ast\}F). \textbackslash{}]\par
The replacement \textbackslash{}[ \textbackslash{}Hom\_\{x'\}(f'\_\{\textbackslash{}ast\}\{g'\}\textasciicircum{}\{\textbackslash{}ast\}F,g\textasciicircum{}\{\textbackslash{}ast\}f\_\{\textbackslash{}ast\}F)=\allowbreak \textbackslash{}Hom\_\{f'\}(\{g'\}\textasciicircum{}\{\textbackslash{}ast\}F,g\textasciicircum{}\{\textbackslash{}ast\}f\_\{\textbackslash{}ast\}F)=\allowbreak \textbackslash{}Hom\_\{gf'\}(\{g'\}\textasciicircum{}\{\textbackslash{}ast\}F,f\_\{\textbackslash{}ast\}F). \textbackslash{}] at Preuve de 2.1.3, première ligne des prétendues bijections identifies a particular indexed object or symbol; that identification is not independently established by the quoted local expression.\par
\par\noindent\textbf{OM-SOL-0930 -- t3 XVII p.285 -- SOL}\par
Preuve, référence suivant « les flèches du diagramme »\par
\emph{Reported reading}: (2.1.2.2)\par
\emph{Proposal}: \textbackslash{}eqhyperref[XVII.eq2.1.3.2]\{2.1.3.2\}\par
This finding intersects Orgogozo's \textbackslash{}sisi material at Preuve, référence suivant « les flèches du diagramme »; preserve that human entry unchanged. The proposed reference \textbackslash{}eqhyperref[XVII.eq2.1.3.2]\{2.1.3.2\} at Preuve, référence suivant « les flèches du diagramme » needs comparison with the statement actually invoked; that target has not been independently established here.\par
\par\noindent\textbf{OM-SOL-0931 -- t3 XVII p.285 -- SOL}\par
Preuve, dernière chaîne de Hom avant « et l'image de 1\_F »\par
\emph{Reported reading}: induisent des bijections ; \textbackslash{}Hom\_y(F,F)=\allowbreak \textbackslash{}Hom\_\{fg'\}(\{g'\}\textasciicircum{}\{\textbackslash{}ast\}F,f\_\{\textbackslash{}ast\}F)=\allowbreak \textbackslash{}Hom\_\{x'\}(f'\_\{\textbackslash{}ast\}\{g'\}\textasciicircum{}\{\textbackslash{}ast\}F,g\textasciicircum{}\{\textbackslash{}ast\}f\_\{\textbackslash{}ast\}F)\par
\emph{Proposal}: induisent des applications \textbackslash{}[ \textbackslash{}Hom\_y(F,F)\textbackslash{}longrightarrow\textbackslash{}Hom\_\{fg'\}(\{g'\}\textasciicircum{}\{\textbackslash{}ast\}F,f\_\{\textbackslash{}ast\}F)=\allowbreak \textbackslash{}Hom\_\{x'\}(f'\_\{\textbackslash{}ast\}\{g'\}\textasciicircum{}\{\textbackslash{}ast\}F,g\textasciicircum{}\{\textbackslash{}ast\}f\_\{\textbackslash{}ast\}F), \textbackslash{}]\par
This finding intersects Orgogozo's \textbackslash{}sisi material at Preuve, dernière chaîne de Hom avant « et l'image de 1\_F »; preserve that human entry unchanged. The proposed reading induisent des applications \textbackslash{}[ \textbackslash{}Hom\_y(F,F)\textbackslash{}longrightarrow\textbackslash{}Hom\_\{fg'\}(\{g'\}\textasciicircum{}\{\textbackslash{}ast\}F,f\_\{\textbackslash{}ast\}F)=\allowbreak \textbackslash{}Hom\_\{x'\}(f'\_\{\textbackslash{}ast\}\{g'\}\textasciicircum{}\{\textbackslash{}ast\}F,g\textasciicircum{}\{\textbackslash{}ast\}f\_\{\textbackslash{}ast\}F), \textbackslash{}] at Preuve, dernière chaîne de Hom avant « et l'image de 1\_F » changes a substantive mathematical assertion or bibliographic identification; no independent derivation of this exact replacement is established in this review.\par
\par\noindent\textbf{OM-SOL-0932 -- t3 XVII p.285 -- SOL}\par
Dernier paragraphe, définition de u\_3\par
\emph{Reported reading}: u\_3=\allowbreak g\textasciicircum{}\{1\textbackslash{}ast\}(u\_2)\par
\emph{Proposal}: u\_3=\allowbreak \{g'\}\textasciicircum{}\{\textbackslash{}ast\}(u\_2)\par
This finding intersects Orgogozo's \textbackslash{}sisi material at Dernier paragraphe, définition de u\_3; preserve that human entry unchanged. The replacement u\_3=\allowbreak \{g'\}\textasciicircum{}\{\textbackslash{}ast\}(u\_2) at Dernier paragraphe, définition de u\_3 identifies a particular indexed object or symbol; that identification is not independently established by the quoted local expression.\par
\par\noindent\textbf{OM-SOL-0933 -- t3 XVII p.286 -- SOL}\par
Last displayed formula, target of the homomorphism induced by F\par
\emph{Reported reading}: \textbackslash{}Hom\_\{\textbackslash{}Aa\_2,f\}(X,Y)\par
\emph{Proposal}: \textbackslash{}Hom\_\{\textbackslash{}Aa\_\{2\},f\}(F(X),F(Y))\par
The proposed reading \textbackslash{}Hom\_\{\textbackslash{}Aa\_\{2\},f\}(F(X),F(Y)) at Last displayed formula, target of the homomorphism induced by F changes a substantive mathematical assertion or bibliographic identification; no independent derivation of this exact replacement is established in this review.\par
\par\noindent\textbf{OM-SOL-0934 -- t3 XVII p.287 -- SOL}\par
Definition 2.2.3(ii), citation after « catégorie triangulée »\par
\emph{Reported reading}: (CD 1.1)\par
\emph{Proposal}: (\textbackslash{}cite\{SGA4-17:CD\} I\textasciitilde{}1.1.1)\par
This finding intersects Orgogozo's \textbackslash{}sisi material at Definition 2.2.3(ii), citation after « catégorie triangulée »; preserve that human entry unchanged. The proposed reference (\textbackslash{}cite\{SGA4-17:CD\} I\textasciitilde{}1.1.1) at Definition 2.2.3(ii), citation after « catégorie triangulée » needs comparison with the statement actually invoked; that target has not been independently established here.\par
\par\noindent\textbf{OM-SOL-0937 -- t3 XVII p.292 -- SOL}\par
Cas triangulé, choix du second triangle\par
\emph{Reported reading}: et soit \$(f\textasciicircum{}\{\textbackslash{}ast\}Y',f\textasciicircum{}\{\textbackslash{}ast\}Y)\$.\par
\emph{Proposal}: \$(f\textasciicircum{}\{\textbackslash{}ast\}Y',f\textasciicircum{}\{\textbackslash{}ast\}Y,X'')\$ un triangle distingué dans \$\textbackslash{}Aa\_x\$\par
This finding intersects Orgogozo's \textbackslash{}sisi material at Cas triangulé, choix du second triangle; preserve that human entry unchanged. The proposed reading \$(f\textasciicircum{}\{\textbackslash{}ast\}Y',f\textasciicircum{}\{\textbackslash{}ast\}Y,X'')\$ un triangle distingué dans \$\textbackslash{}Aa\_x\$ at Cas triangulé, choix du second triangle changes a substantive mathematical assertion or bibliographic identification; no independent derivation of this exact replacement is established in this review.\par
\par\noindent\textbf{OM-SOL-0938 -- t3 XVII p.292 -- SOL}\par
Diagramme des triangles, deux côtés du triangle central\par
\emph{Reported reading}: Les flèches de f\textasciicircum{}\{\textbackslash{}ast\}Y' vers f\textasciicircum{}\{\textbackslash{}ast\}Y et de f\textasciicircum{}\{\textbackslash{}ast\}Y vers f\textasciicircum{}\{\textbackslash{}ast\}Y'' ne sont pas tracées.\par
\emph{Proposal}: f\textasciicircum{}\{\textbackslash{}ast\}Y'\textbackslash{}longrightarrow f\textasciicircum{}\{\textbackslash{}ast\}Y,\textbackslash{}qquad f\textasciicircum{}\{\textbackslash{}ast\}Y\textbackslash{}longrightarrow f\textasciicircum{}\{\textbackslash{}ast\}Y''\par
This finding intersects Orgogozo's \textbackslash{}sisi material at Diagramme des triangles, deux côtés du triangle central; preserve that human entry unchanged. The proposed reading f\textasciicircum{}\{\textbackslash{}ast\}Y'\textbackslash{}longrightarrow f\textasciicircum{}\{\textbackslash{}ast\}Y,\textbackslash{}qquad f\textasciicircum{}\{\textbackslash{}ast\}Y\textbackslash{}longrightarrow f\textasciicircum{}\{\textbackslash{}ast\}Y'' at Diagramme des triangles, deux côtés du triangle central changes a substantive mathematical assertion or bibliographic identification; no independent derivation of this exact replacement is established in this review.\par
\par\noindent\textbf{OM-SOL-0941 -- t3 XVII p.293 -- SOL}\par
(2.3.1), membre de gauche\par
\emph{Reported reading}: \textbackslash{}Aa'(S\textasciicircum{}\{-1\})\par
\emph{Proposal}: \textbackslash{}Aa'(\{S'\}\textasciicircum{}\{-1\})\par
The replacement \textbackslash{}Aa'(\{S'\}\textasciicircum{}\{-1\}) at (2.3.1), membre de gauche identifies a particular indexed object or symbol; that identification is not independently established by the quoted local expression.\par
\par\noindent\textbf{OM-SOL-0942 -- t3 XVII p.294 -- SOL}\par
Démonstration de 2.3.2, « on vérifie que… »\par
\emph{Reported reading}: \textbackslash{}Aa'(S\textasciicircum{}\{-1\})\par
\emph{Proposal}: \textbackslash{}Aa'(\{S'\}\textasciicircum{}\{-1\})\par
The replacement \textbackslash{}Aa'(\{S'\}\textasciicircum{}\{-1\}) at Démonstration de 2.3.2, « on vérifie que… » identifies a particular indexed object or symbol; that identification is not independently established by the quoted local expression.\par
\par\noindent\textbf{OM-SOL-0943 -- t3 XVII p.298 -- SOL}\par
End of 2.4.0, reference for a saturated multiplicative system\par
\emph{Reported reading}: [CD 1.2.1]\par
\emph{Proposal}: \textbackslash{}cite[I 1.2.1]\{SGA4-17:CD\}\par
This finding intersects Orgogozo's \textbackslash{}sisi material at End of 2.4.0, reference for a saturated multiplicative system; preserve that human entry unchanged. The proposed reference \textbackslash{}cite[I 1.2.1]\{SGA4-17:CD\} at End of 2.4.0, reference for a saturated multiplicative system needs comparison with the statement actually invoked; that target has not been independently established here.\par
\par\noindent\textbf{OM-SOL-0946 -- t3 XVII p.299 -- SOL}\par
Diagram (2.4.1.2), lower horizontal arrow\par
\emph{Reported reading}: A\textbackslash{}longrightarrow B, without a label\par
\emph{Proposal}: A\textbackslash{}xrightarrow\{f\}B\par
This finding intersects Orgogozo's \textbackslash{}sisi material at Diagram (2.4.1.2), lower horizontal arrow; preserve that human entry unchanged. The replacement A\textbackslash{}xrightarrow\{f\}B at Diagram (2.4.1.2), lower horizontal arrow requires a complete check of the maps, their endpoints and the composition convention; this identification is not independently verified.\par
\par\noindent\textbf{OM-SOL-0947 -- t3 XVII p.300 -- SOL}\par
First diagram, left vertical arrow\par
\emph{Reported reading}: s\par
\emph{Proposal}: F\_\{\textbackslash{}ast\}(s)\par
This finding intersects Orgogozo's \textbackslash{}sisi material at First diagram, left vertical arrow; preserve that human entry unchanged. The replacement F\_\{\textbackslash{}ast\}(s) at First diagram, left vertical arrow requires a complete check of the maps, their endpoints and the composition convention; this identification is not independently verified.\par
\par\noindent\textbf{OM-SOL-0948 -- t3 XVII p.300 -- SOL}\par
Construction of the quasi-inverse, source and target of f\par
\emph{Reported reading}: de but A, de source B\par
\emph{Proposal}: de but \$B\$, de source \$A\$\par
The replacement de but \$B\$, de source \$A\$ at Construction of the quasi-inverse, source and target of f requires a complete check of the maps, their endpoints and the composition convention; this identification is not independently verified.\par
\par\noindent\textbf{OM-SOL-0949 -- t3 XVII p.300 -- SOL}\par
Immediately before diagram (2.4.1.4), auxiliary object to be found\par
\emph{Reported reading}: B\textasciicircum{}1\par
\emph{Proposal}: B'\par
This finding intersects Orgogozo's \textbackslash{}sisi material at Immediately before diagram (2.4.1.4), auxiliary object to be found; preserve that human entry unchanged. The proposed reading B' at Immediately before diagram (2.4.1.4), auxiliary object to be found changes a substantive mathematical assertion or bibliographic identification; no independent derivation of this exact replacement is established in this review.\par
\par\noindent\textbf{OM-SOL-0950 -- t3 XVII p.300 -- SOL}\par
Immediately before diagram (2.4.1.4), name of the auxiliary morphism\par
\emph{Reported reading}: v et s\textbackslash{}in S\textasciicircum{}1\par
\emph{Proposal}: \$u\$ et \$s\textbackslash{}in S\textasciicircum{}1\$\par
This finding intersects Orgogozo's \textbackslash{}sisi material at Immediately before diagram (2.4.1.4), name of the auxiliary morphism; preserve that human entry unchanged. The replacement \$u\$ et \$s\textbackslash{}in S\textasciicircum{}1\$ at Immediately before diagram (2.4.1.4), name of the auxiliary morphism requires a complete check of the maps, their endpoints and the composition convention; this identification is not independently verified.\par
\par\noindent\textbf{OM-SOL-0951 -- t3 XVII p.300 -- SOL}\par
Diagram (2.4.1.4), the object labels A\_1, F\_*A\_1 and B\_1\par
\emph{Reported reading}: A\_1,\textbackslash{}quad F\_\{\textbackslash{}ast\}A\_1,\textbackslash{}quad B\_1\par
\emph{Proposal}: A\textasciicircum{}1,\textbackslash{}quad F\_\{\textbackslash{}ast\}A\textasciicircum{}1,\textbackslash{}quad B\textasciicircum{}1\par
This finding intersects Orgogozo's \textbackslash{}sisi material at Diagram (2.4.1.4), the object labels A\_1, F\_*A\_1 and B\_1; preserve that human entry unchanged. The replacement A\textasciicircum{}1,\textbackslash{}quad F\_\{\textbackslash{}ast\}A\textasciicircum{}1,\textbackslash{}quad B\textasciicircum{}1 at Diagram (2.4.1.4), the object labels A\_1, F\_*A\_1 and B\_1 identifies a particular indexed object or symbol; that identification is not independently established by the quoted local expression.\par
\par\noindent\textbf{OM-SOL-0952 -- t3 XVII p.300 -- SOL}\par
Diagram (2.4.1.4), upper arrow labelled s\_B\par
\emph{Reported reading}: B\textbackslash{}xrightarrow\{s\_B\}B\_1\par
\emph{Proposal}: B\textbackslash{}xleftarrow\{s\_B\}B\textasciicircum{}1\par
The replacement B\textbackslash{}xleftarrow\{s\_B\}B\textasciicircum{}1 at Diagram (2.4.1.4), upper arrow labelled s\_B requires a complete check of the maps, their endpoints and the composition convention; this identification is not independently verified.\par
\par\noindent\textbf{OM-SOL-0954 -- t3 XVII p.303 -- SOL}\par
Introduction, renvoi à la définition de l'image directe à support propre\par
\emph{Reported reading}: cf. 5.3 dans le cas des schémas\par
\emph{Proposal}: \textbackslash{}hyperref[XVII.subsec6.1]\{6.1\}\par
This finding intersects Orgogozo's \textbackslash{}sisi material at Introduction, renvoi à la définition de l'image directe à support propre; preserve that human entry unchanged. The proposed reference \textbackslash{}hyperref[XVII.subsec6.1]\{6.1\} at Introduction, renvoi à la définition de l'image directe à support propre needs comparison with the statement actually invoked; that target has not been independently established here.\par
\par\noindent\textbf{OM-SOL-0956 -- t3 XVII p.306 -- SOL}\par
Diagramme initial, sommet supérieur\par
\emph{Reported reading}: Q\textbackslash{}times P\textbackslash{}times Z\par
\emph{Proposal}: \textbackslash{}fprod\{\textbackslash{}fprod\{Q\}\{P\}\{S\}\}\{Z\}\{S\}\par
This finding intersects Orgogozo's \textbackslash{}sisi material at Diagramme initial, sommet supérieur; preserve that human entry unchanged. The proposed reading \textbackslash{}fprod\{\textbackslash{}fprod\{Q\}\{P\}\{S\}\}\{Z\}\{S\} at Diagramme initial, sommet supérieur changes a substantive mathematical assertion or bibliographic identification; no independent derivation of this exact replacement is established in this review.\par
\par\noindent\textbf{OM-SOL-0957 -- t3 XVII p.307 -- SOL}\par
Démonstration de 3.2.3, « où p est propre et j quasi-fini séparé »\par
\emph{Reported reading}: j quasi-fini séparé\par
\emph{Proposal}: j'\par
This finding intersects Orgogozo's \textbackslash{}sisi material at Démonstration de 3.2.3, « où p est propre et j quasi-fini séparé »; preserve that human entry unchanged. The replacement j' at Démonstration de 3.2.3, « où p est propre et j quasi-fini séparé » identifies a particular indexed object or symbol; that identification is not independently established by the quoted local expression.\par
\par\noindent\textbf{OM-SOL-0958 -- t3 XVII p.308 -- SOL}\par
Diagramme (3.2.5.2), flèche verticale de Z vers Y\par
\emph{Reported reading}: Flèche Z\ensuremath{\to}Y sans étiquette\par
\emph{Proposal}: \textbackslash{}ar[d]\_g\par
This finding intersects Orgogozo's \textbackslash{}sisi material at Diagramme (3.2.5.2), flèche verticale de Z vers Y; preserve that human entry unchanged. The replacement \textbackslash{}ar[d]\_g at Diagramme (3.2.5.2), flèche verticale de Z vers Y requires a complete check of the maps, their endpoints and the composition convention; this identification is not independently verified.\par
\par\noindent\textbf{OM-SOL-0959 -- t3 XVII p.308 -- SOL}\par
Diagramme (3.2.5.2), trait oblique de \textbackslash{}bar\{Z\} vers Y\par
\emph{Reported reading}: Trait étiqueté \textbackslash{}bar\{g\}, sans pointe de flèche\par
\emph{Proposal}: \textbackslash{}ar[ld]\textasciicircum{}\{\textbackslash{}bar\{g\}\}\par
This finding intersects Orgogozo's \textbackslash{}sisi material at Diagramme (3.2.5.2), trait oblique de \textbackslash{}bar\{Z\} vers Y; preserve that human entry unchanged. The proposed reading \textbackslash{}ar[ld]\textasciicircum{}\{\textbackslash{}bar\{g\}\} involves a printed glyph or a missing stroke at Diagramme (3.2.5.2), trait oblique de \textbackslash{}bar\{Z\} vers Y; the textual record does not independently establish that glyph.\par
\par\noindent\textbf{OM-SOL-0960 -- t3 XVII p.308 -- SOL}\par
Diagramme (3.2.5.3), trois immersions horizontales de la ligne supérieure\par
\emph{Reported reading}: h, h', h''\par
\emph{Proposal}: k,\textbackslash{} k',\textbackslash{} k''\par
This finding intersects Orgogozo's \textbackslash{}sisi material at Diagramme (3.2.5.3), trois immersions horizontales de la ligne supérieure; preserve that human entry unchanged. The replacement k,\textbackslash{} k',\textbackslash{} k'' at Diagramme (3.2.5.3), trois immersions horizontales de la ligne supérieure requires a complete check of the maps, their endpoints and the composition convention; this identification is not independently verified.\par
\par\noindent\textbf{OM-SOL-0961 -- t3 XVII p.308 -- SOL}\par
Proposition 3.2.6 (iii), après « Si des diagrammes »\par
\emph{Reported reading}: (3.2.4.2)\_i\par
\emph{Proposal}: (3.2.5.2)\_i\par
This finding intersects Orgogozo's \textbackslash{}sisi material at Proposition 3.2.6 (iii), après « Si des diagrammes »; preserve that human entry unchanged. The proposed reference (3.2.5.2)\_i at Proposition 3.2.6 (iii), après « Si des diagrammes » needs comparison with the statement actually invoked; that target has not been independently established here.\par
\par\noindent\textbf{OM-SOL-0962 -- t3 XVII p.308 -- SOL}\par
Proposition 3.2.6 (iii), après « il existe une compactification »\par
\emph{Reported reading}: (3.2.4.2)\_3\par
\emph{Proposal}: (3.2.5.2)\_3\par
This finding intersects Orgogozo's \textbackslash{}sisi material at Proposition 3.2.6 (iii), après « il existe une compactification »; preserve that human entry unchanged. The proposed reference (3.2.5.2)\_3 at Proposition 3.2.6 (iii), après « il existe une compactification » needs comparison with the statement actually invoked; that target has not been independently established here.\par
\par\noindent\textbf{OM-SOL-0963 -- t3 XVII p.308 -- SOL}\par
Proposition 3.2.6 (iii), après « des morphismes de compactifications de »\par
\emph{Reported reading}: (3.2.4.2)\_3\par
\emph{Proposal}: (3.2.5.2)\_3\par
This finding intersects Orgogozo's \textbackslash{}sisi material at Proposition 3.2.6 (iii), après « des morphismes de compactifications de »; preserve that human entry unchanged. The proposed reference (3.2.5.2)\_3 at Proposition 3.2.6 (iii), après « des morphismes de compactifications de » needs comparison with the statement actually invoked; that target has not been independently established here.\par
\par\noindent\textbf{OM-SOL-0964 -- t3 XVII p.308 -- SOL}\par
Proposition 3.2.6 (iii), dernier renvoi, après « dans les compactifications »\par
\emph{Reported reading}: (3.2.4.2)\_i\par
\emph{Proposal}: (3.2.5.2)\_i\par
This finding intersects Orgogozo's \textbackslash{}sisi material at Proposition 3.2.6 (iii), dernier renvoi, après « dans les compactifications »; preserve that human entry unchanged. The proposed reference (3.2.5.2)\_i at Proposition 3.2.6 (iii), dernier renvoi, après « dans les compactifications » needs comparison with the statement actually invoked; that target has not been independently established here.\par
\par\noindent\textbf{OM-SOL-0965 -- t3 XVII p.309 -- SOL}\par
Démonstration de 3.2.6 (ii), nom de la projection\par
\emph{Reported reading}: \textbackslash{}rho\par
\emph{Proposal}: p\par
This finding intersects Orgogozo's \textbackslash{}sisi material at Démonstration de 3.2.6 (ii), nom de la projection; preserve that human entry unchanged. The proposed reading p at Démonstration de 3.2.6 (ii), nom de la projection changes a substantive mathematical assertion or bibliographic identification; no independent derivation of this exact replacement is established in this review.\par
\par\noindent\textbf{OM-SOL-0966 -- t3 XVII p.309 -- SOL}\par
Démonstration de 3.2.6 (ii), produit des deux compactifications\par
\emph{Reported reading}: \textbackslash{}fprod\{\textbackslash{}bar\{Y\}\_1\}\{\textbackslash{}bar\{Y\}\_2\}\{S\}\par
\emph{Proposal}: \textbackslash{}fprod\{\textbackslash{}bar\{Y\}\_1\}\{\textbackslash{}bar\{Y\}\_2\}\{X\}\par
The proposed reading \textbackslash{}fprod\{\textbackslash{}bar\{Y\}\_1\}\{\textbackslash{}bar\{Y\}\_2\}\{X\} at Démonstration de 3.2.6 (ii), produit des deux compactifications changes a substantive mathematical assertion or bibliographic identification; no independent derivation of this exact replacement is established in this review.\par
\par\noindent\textbf{OM-SOL-0967 -- t3 XVII p.309 -- SOL}\par
Démonstration de 3.2.6 (ii), premier renvoi après « domine »\par
\emph{Reported reading}: (3.2.4.1)\_1\par
\emph{Proposal}: (3.2.5.1)\_1\par
This finding intersects Orgogozo's \textbackslash{}sisi material at Démonstration de 3.2.6 (ii), premier renvoi après « domine »; preserve that human entry unchanged. The proposed reading (3.2.5.1)\_1 at Démonstration de 3.2.6 (ii), premier renvoi après « domine » changes a substantive mathematical assertion or bibliographic identification; no independent derivation of this exact replacement is established in this review.\par
\par\noindent\textbf{OM-SOL-0968 -- t3 XVII p.309 -- SOL}\par
Démonstration de 3.2.6 (ii), second renvoi après « domine »\par
\emph{Reported reading}: (3.2.4.1)\_2\par
\emph{Proposal}: (3.2.5.1)\_2\par
This finding intersects Orgogozo's \textbackslash{}sisi material at Démonstration de 3.2.6 (ii), second renvoi après « domine »; preserve that human entry unchanged. The proposed reference (3.2.5.1)\_2 at Démonstration de 3.2.6 (ii), second renvoi après « domine » needs comparison with the statement actually invoked; that target has not been independently established here.\par
\par\noindent\textbf{OM-SOL-0969 -- t3 XVII p.309 -- SOL}\par
Démonstration de 3.2.6 (ii), « deux morphismes de … »\par
\emph{Reported reading}: (3.2.4.1)\_1\par
\emph{Proposal}: (3.2.5.1)\_1\par
This finding intersects Orgogozo's \textbackslash{}sisi material at Démonstration de 3.2.6 (ii), « deux morphismes de … »; preserve that human entry unchanged. The proposed reference (3.2.5.1)\_1 at Démonstration de 3.2.6 (ii), « deux morphismes de … » needs comparison with the statement actually invoked; that target has not been independently established here.\par
\par\noindent\textbf{OM-SOL-0970 -- t3 XVII p.309 -- SOL}\par
Démonstration de 3.2.6 (ii), « deux morphismes de … dans … »\par
\emph{Reported reading}: (3.2.4.1)\_2\par
\emph{Proposal}: (3.2.5.1)\_2\par
This finding intersects Orgogozo's \textbackslash{}sisi material at Démonstration de 3.2.6 (ii), « deux morphismes de … dans … »; preserve that human entry unchanged. The proposed reference (3.2.5.1)\_2 at Démonstration de 3.2.6 (ii), « deux morphismes de … dans … » needs comparison with the statement actually invoked; that target has not been independently established here.\par
\par\noindent\textbf{OM-SOL-0971 -- t3 XVII p.309 -- SOL}\par
Démonstration de 3.2.6 (iii), « compactifications de f et g déduites des … »\par
\emph{Reported reading}: (3.2.4.2)\_i\par
\emph{Proposal}: (3.2.5.2)\_i\par
This finding intersects Orgogozo's \textbackslash{}sisi material at Démonstration de 3.2.6 (iii), « compactifications de f et g déduites des … »; preserve that human entry unchanged. The proposed reference (3.2.5.2)\_i at Démonstration de 3.2.6 (iii), « compactifications de f et g déduites des … » needs comparison with the statement actually invoked; that target has not been independently established here.\par
\par\noindent\textbf{OM-SOL-0972 -- t3 XVII p.311 -- SOL}\par
Données analogues pour \$(S,i)\$, premier intitulé\par
\emph{Reported reading}: (i) et (ii)'\par
\emph{Proposal}: (i')\par
This finding intersects Orgogozo's \textbackslash{}sisi material at Données analogues pour \$(S,i)\$, premier intitulé; preserve that human entry unchanged. The proposed reading (i') involves a printed glyph or a missing stroke at Données analogues pour \$(S,i)\$, premier intitulé; the textual record does not independently establish that glyph.\par
\par\noindent\textbf{OM-SOL-0973 -- t3 XVII p.311 -- SOL}\par
Conditions (a') et (b')\par
\emph{Reported reading}: conditions analogues pour \$(S,p)\$.\par
\emph{Proposal}: conditions analogues pour \$(S,i)\$.\par
The proposed reading conditions analogues pour \$(S,i)\$. at Conditions (a') et (b') changes a substantive mathematical assertion or bibliographic identification; no independent derivation of this exact replacement is established in this review.\par
\par\noindent\textbf{OM-SOL-0975 -- t3 XVII p.314 -- SOL}\par
Référence aux deux compactifications\par
\emph{Reported reading}: \$(3.2.4.2)\_i\$\par
\emph{Proposal}: \$(3.2.5.2)\_i\$\par
This finding intersects Orgogozo's \textbackslash{}sisi material at Référence aux deux compactifications; preserve that human entry unchanged. The proposed reference \$(3.2.5.2)\_i\$ at Référence aux deux compactifications needs comparison with the statement actually invoked; that target has not been independently established here.\par
\par\noindent\textbf{OM-SOL-0976 -- t3 XVII p.314 -- SOL}\par
Grand diagramme, deux triangles supérieurs\par
\emph{Reported reading}: Les deux triangles supérieurs ne portent aucune étiquette ; seules (3), (4) et (5) sont imprimées.\par
\emph{Proposal}: \$(1)\textbackslash{}qquad(2)\$\par
This finding intersects Orgogozo's \textbackslash{}sisi material at Grand diagramme, deux triangles supérieurs; preserve that human entry unchanged. The proposed reading \$(1)\textbackslash{}qquad(2)\$ involves a printed glyph or a missing stroke at Grand diagramme, deux triangles supérieurs; the textual record does not independently establish that glyph.\par
\par\noindent\textbf{OM-SOL-0978 -- t3 XVII p.315 -- SOL}\par
Renvoi au premier triangle du diagramme de la page 314\par
\emph{Reported reading}: la commutativité de (i)\par
\emph{Proposal}: la commutativité de (1)\par
This finding intersects Orgogozo's \textbackslash{}sisi material at Renvoi au premier triangle du diagramme de la page 314; preserve that human entry unchanged. The proposed reading la commutativité de (1) involves a printed glyph or a missing stroke at Renvoi au premier triangle du diagramme de la page 314; the textual record does not independently establish that glyph.\par
\par\noindent\textbf{OM-SOL-0980 -- t3 XVII p.316 -- SOL}\par
Hexagone, sommet intérieur gauche\par
\emph{Reported reading}: i\_2(g\_2p)\par
\emph{Proposal}: i\_2(\textbackslash{}bar\{g\}\_2p)\par
The replacement i\_2(\textbackslash{}bar\{g\}\_2p) at Hexagone, sommet intérieur gauche requires a complete check of the maps, their endpoints and the composition convention; this identification is not independently verified.\par
\par\noindent\textbf{OM-SOL-0981 -- t3 XVII p.317 -- SOL}\par
3.3.2.4, référence au diagramme d'une compactification du triple\par
\emph{Reported reading}: (3.2.4.3)\par
\emph{Proposal}: \textbackslash{}eqhyperref[XVII.eq3.2.5.3]\{3.2.5.3\}\par
This finding intersects Orgogozo's \textbackslash{}sisi material at 3.3.2.4, référence au diagramme d'une compactification du triple; preserve that human entry unchanged. The proposed reference \textbackslash{}eqhyperref[XVII.eq3.2.5.3]\{3.2.5.3\} at 3.3.2.4, référence au diagramme d'une compactification du triple needs comparison with the statement actually invoked; that target has not been independently established here.\par
\par\noindent\textbf{OM-SOL-0982 -- t3 XVII p.317 -- SOL}\par
Diagramme de 3.3.2.4, sommet inférieur droit\par
\emph{Reported reading}: \textbackslash{}bar\{f\}\textbackslash{}bar\{\textbackslash{}bar\{g\}\}\textbackslash{}bar\{\textbackslash{}bar\{h\}\}k''k'k\par
\emph{Proposal}: \textbackslash{}bar\{f\}\textbackslash{}bar\{\textbackslash{}bar\{g\}\}\textbackslash{}bar\{\textbackslash{}bar\{\textbackslash{}bar\{h\}\}\}k''k'k\par
The replacement \textbackslash{}bar\{f\}\textbackslash{}bar\{\textbackslash{}bar\{g\}\}\textbackslash{}bar\{\textbackslash{}bar\{\textbackslash{}bar\{h\}\}\}k''k'k at Diagramme de 3.3.2.4, sommet inférieur droit identifies a particular indexed object or symbol; that identification is not independently established by the quoted local expression.\par
\par\noindent\textbf{OM-SOL-0984 -- t3 XVII p.318 -- SOL}\par
Proposition 3.3.4, nature de la catégorie F\par
\emph{Reported reading}: une catégorie F cofibrée sur S\par
\emph{Proposal}: une catégorie \$F\$ fibrée sur \$S\$\par
The replacement une catégorie \$F\$ fibrée sur \$S\$ at Proposition 3.3.4, nature de la catégorie F requires a complete check of the maps, their endpoints and the composition convention; this identification is not independently verified.\par
\par\noindent\textbf{OM-SOL-0985 -- t3 XVII p.319 -- SOL}\par
Preuve de 4.1.1, référence concernant la conservation des Modules plats\par
\emph{Reported reading}: (IV 15)\par
\emph{Proposal}: \textbackslash{}hyperrefext[V][V.sec1]\{1\}\par
This finding intersects Orgogozo's \textbackslash{}sisi material at Preuve de 4.1.1, référence concernant la conservation des Modules plats; preserve that human entry unchanged. The proposed reference \textbackslash{}hyperrefext[V][V.sec1]\{1\} at Preuve de 4.1.1, référence concernant la conservation des Modules plats needs comparison with the statement actually invoked; that target has not been independently established here.\par
\par\noindent\textbf{OM-SOL-0986 -- t3 XVII p.319 -- SOL}\par
Dernier paragraphe, après « avec P plat »\par
\emph{Reported reading}: avec P plat (4.1.2).\par
\emph{Proposal}: avec \$P\$ plat.\par
This finding intersects Orgogozo's \textbackslash{}sisi material at Dernier paragraphe, après « avec P plat »; preserve that human entry unchanged. The proposed reference avec \$P\$ plat. at Dernier paragraphe, après « avec P plat » needs comparison with the statement actually invoked; that target has not been independently established here.\par
\par\noindent\textbf{OM-SOL-0988 -- t3 XVII p.321 -- SOL}\par
Formula (4.1.2.2)\par
\emph{Reported reading}: Lf\textasciicircum{}\{\textbackslash{}ast\}:D\textasciicircum{}\{-\}(S)\textbackslash{}longrightarrow D\textasciicircum{}\{-\}(S')\par
\emph{Proposal}: Lf\textasciicircum{}\{\textbackslash{}ast\}:D\textasciicircum{}\{-\}(S')\textbackslash{}longrightarrow D\textasciicircum{}\{-\}(S)\par
This finding intersects Orgogozo's \textbackslash{}sisi material at Formula (4.1.2.2); preserve that human entry unchanged. The proposed reading Lf\textasciicircum{}\{\textbackslash{}ast\}:D\textasciicircum{}\{-\}(S')\textbackslash{}longrightarrow D\textasciicircum{}\{-\}(S) at Formula (4.1.2.2) changes a substantive mathematical assertion or bibliographic identification; no independent derivation of this exact replacement is established in this review.\par
\par\noindent\textbf{OM-SOL-0990 -- t3 XVII p.322 -- SOL}\par
Proof of Scholie 4.1.3, « Dans le cas (a), on utilisera »\par
\emph{Reported reading}: 2.4.3\par
\emph{Proposal}: \textbackslash{}hyperref[XVII.prop2.4.1]\{2.4.1\} (\textbackslash{}resp \textbackslash{}hyperref[XVII.prop2.4.3]\{2.4.3\})\par
This finding intersects Orgogozo's \textbackslash{}sisi material at Proof of Scholie 4.1.3, « Dans le cas (a), on utilisera »; preserve that human entry unchanged. The proposed reference \textbackslash{}hyperref[XVII.prop2.4.1]\{2.4.1\} (\textbackslash{}resp \textbackslash{}hyperref[XVII.prop2.4.3]\{2.4.3\}) at Proof of Scholie 4.1.3, « Dans le cas (a), on utilisera » needs comparison with the statement actually invoked; that target has not been independently established here.\par
\par\noindent\textbf{OM-SOL-0991 -- t3 XVII p.323 -- SOL}\par
First paragraph, right-hand side of the equality involving f\_*\par
\emph{Reported reading}: \textbackslash{}tau''\_\{<\}f\_\{\textbackslash{}ast\}(K\textasciicircum{}\{\textbackslash{}ast\})\par
\emph{Proposal}: \textbackslash{}tau''\_\{<n\}f\_\{\textbackslash{}ast\}(K\textasciicircum{}\{\textbackslash{}ast\})\par
This finding intersects Orgogozo's \textbackslash{}sisi material at First paragraph, right-hand side of the equality involving f\_*; preserve that human entry unchanged. The proposed reading \textbackslash{}tau''\_\{<n\}f\_\{\textbackslash{}ast\}(K\textasciicircum{}\{\textbackslash{}ast\}) at First paragraph, right-hand side of the equality involving f\_* changes a substantive mathematical assertion or bibliographic identification; no independent derivation of this exact replacement is established in this review.\par
\par\noindent\textbf{OM-SOL-0992 -- t3 XVII p.323 -- SOL}\par
First paragraph, « où … désigne un tronqué »\par
\emph{Reported reading}: \textbackslash{}tau\textasciicircum{}\{\textbackslash{}ast\}\par
\emph{Proposal}: \textbackslash{}tau''\par
This finding intersects Orgogozo's \textbackslash{}sisi material at First paragraph, « où … désigne un tronqué »; preserve that human entry unchanged. The replacement \textbackslash{}tau'' at First paragraph, « où … désigne un tronqué » identifies a particular indexed object or symbol; that identification is not independently established by the quoted local expression.\par
\par\noindent\textbf{OM-SOL-0993 -- t3 XVII p.323 -- SOL}\par
First paragraph, « vérifie l'hypothèse de »\par
\emph{Reported reading}: 2.4.3\par
\emph{Proposal}: \textbackslash{}hyperref[XVII.prop2.4.1]\{2.4.1\} (\textbackslash{}resp \textbackslash{}hyperref[XVII.prop2.4.3]\{2.4.3\})\par
This finding intersects Orgogozo's \textbackslash{}sisi material at First paragraph, « vérifie l'hypothèse de »; preserve that human entry unchanged. The proposed reference \textbackslash{}hyperref[XVII.prop2.4.1]\{2.4.1\} (\textbackslash{}resp \textbackslash{}hyperref[XVII.prop2.4.3]\{2.4.3\}) at First paragraph, « vérifie l'hypothèse de » needs comparison with the statement actually invoked; that target has not been independently established here.\par
\par\noindent\textbf{OM-SOL-0994 -- t3 XVII p.323 -- SOL}\par
4.1.4, paragraph following diagram (4.1.4.1), « on utilise alors »\par
\emph{Reported reading}: 2.1.2\par
\emph{Proposal}: \textbackslash{}hyperref[XVII.prop2.1.3]\{2.1.3\}\par
This finding intersects Orgogozo's \textbackslash{}sisi material at 4.1.4, paragraph following diagram (4.1.4.1), « on utilise alors »; preserve that human entry unchanged. The proposed reference \textbackslash{}hyperref[XVII.prop2.1.3]\{2.1.3\} at 4.1.4, paragraph following diagram (4.1.4.1), « on utilise alors » needs comparison with the statement actually invoked; that target has not been independently established here.\par
\par\noindent\textbf{OM-SOL-0995 -- t3 XVII p.323 -- SOL}\par
Proposition 4.1.5, opening sentence, « La flèche »\par
\emph{Reported reading}: 2.1.2\par
\emph{Proposal}: \textbackslash{}hyperref[XVII.prop2.1.3]\{2.1.3\}\par
This finding intersects Orgogozo's \textbackslash{}sisi material at Proposition 4.1.5, opening sentence, « La flèche »; preserve that human entry unchanged. The proposed reference \textbackslash{}hyperref[XVII.prop2.1.3]\{2.1.3\} at Proposition 4.1.5, opening sentence, « La flèche » needs comparison with the statement actually invoked; that target has not been independently established here.\par
\par\noindent\textbf{OM-SOL-0996 -- t3 XVII p.324 -- SOL}\par
Opening paragraph, « à partir de la description »\par
\emph{Reported reading}: 2.1.2\par
\emph{Proposal}: \textbackslash{}hyperref[XVII.prop2.1.3]\{2.1.3\}\par
This finding intersects Orgogozo's \textbackslash{}sisi material at Opening paragraph, « à partir de la description »; preserve that human entry unchanged. The proposed reference \textbackslash{}hyperref[XVII.prop2.1.3]\{2.1.3\} at Opening paragraph, « à partir de la description » needs comparison with the statement actually invoked; that target has not been independently established here.\par
\par\noindent\textbf{OM-SOL-0997 -- t3 XVII p.324 -- SOL}\par
4.1.6, first factor in the source of the tensor-product functor\par
\emph{Reported reading}: \textbackslash{}Mod\_d(S,\textbackslash{}Aa\textasciicircum{}\{\textbackslash{}circ\})\par
\emph{Proposal}: \textbackslash{}Mod\_g(S,\textbackslash{}Aa\textasciicircum{}\{\textbackslash{}circ\})\par
This finding intersects Orgogozo's \textbackslash{}sisi material at 4.1.6, first factor in the source of the tensor-product functor; preserve that human entry unchanged. The proposed reading \textbackslash{}Mod\_g(S,\textbackslash{}Aa\textasciicircum{}\{\textbackslash{}circ\}) at 4.1.6, first factor in the source of the tensor-product functor changes a substantive mathematical assertion or bibliographic identification; no independent derivation of this exact replacement is established in this review.\par
\par\noindent\textbf{OM-SOL-0998 -- t3 XVII p.325 -- SOL}\par
Définition 4.1.9(iii), qualification of K'\par
\emph{Reported reading}: un complexe \$K'\$, à composantes plates\par
\emph{Proposal}: un complexe borné \$K'\$, à composantes plates et nulles pour \$i<-n\$\par
This finding intersects Orgogozo's \textbackslash{}sisi material at Définition 4.1.9(iii), qualification of K'; preserve that human entry unchanged. The proposed reading un complexe borné \$K'\$, à composantes plates et nulles pour \$i<-n\$ at Définition 4.1.9(iii), qualification of K' changes a substantive mathematical assertion or bibliographic identification; no independent derivation of this exact replacement is established in this review.\par
\par\noindent\textbf{OM-SOL-1000 -- t3 XVII p.331 -- SOL}\par
Condition (b), tensor functor\par
\emph{Reported reading}: M\_x\textbackslash{}otimes\_\{\textbackslash{}Aa\_x\}\textbackslash{}ast\par
\emph{Proposal}: Q\textbackslash{}otimes\_\{\textbackslash{}Aa\_x\}\textbackslash{}ast\par
This finding intersects Orgogozo's \textbackslash{}sisi material at Condition (b), tensor functor; preserve that human entry unchanged. The proposed reading Q\textbackslash{}otimes\_\{\textbackslash{}Aa\_x\}\textbackslash{}ast at Condition (b), tensor functor changes a substantive mathematical assertion or bibliographic identification; no independent derivation of this exact replacement is established in this review.\par
\par\noindent\textbf{OM-SOL-1001 -- t3 XVII p.331 -- SOL}\par
Condition (c), domain of f\par
\emph{Reported reading}: f:X\textasciicircum{}\{\textbackslash{}sim\}\textbackslash{}rightarrow E\par
\emph{Proposal}: f:S\textbackslash{}rightarrow E\par
This finding intersects Orgogozo's \textbackslash{}sisi material at Condition (c), domain of f; preserve that human entry unchanged. The proposed reading f:S\textbackslash{}rightarrow E at Condition (c), domain of f changes a substantive mathematical assertion or bibliographic identification; no independent derivation of this exact replacement is established in this review.\par
\par\noindent\textbf{OM-SOL-1002 -- t3 XVII p.331 -- SOL}\par
First proof paragraph, canonical map a\par
\emph{Reported reading}: a:K\_1=\allowbreak \textbackslash{}sigma\_\{\textbackslash{}leq i\}K\textbackslash{}rightarrow K\par
\emph{Proposal}: a:K\_1=\allowbreak \textbackslash{}tau\_\{\textbackslash{}leq i\}K\textbackslash{}rightarrow K\par
This finding intersects Orgogozo's \textbackslash{}sisi material at First proof paragraph, canonical map a; preserve that human entry unchanged. The proposed reading a:K\_1=\allowbreak \textbackslash{}tau\_\{\textbackslash{}leq i\}K\textbackslash{}rightarrow K at First proof paragraph, canonical map a changes a substantive mathematical assertion or bibliographic identification; no independent derivation of this exact replacement is established in this review.\par
\par\noindent\textbf{OM-SOL-1003 -- t3 XVII p.331 -- SOL}\par
First proof paragraph, definition of K\_2\par
\emph{Reported reading}: K\_2=\allowbreak \textbackslash{}tau''\_\{\textbackslash{}geq N\}K\_1\par
\emph{Proposal}: K\_2=\allowbreak \textbackslash{}tau''\_\{\textbackslash{}geq -N\}K\_1\par
This finding intersects Orgogozo's \textbackslash{}sisi material at First proof paragraph, definition of K\_2; preserve that human entry unchanged. The proposed reading K\_2=\allowbreak \textbackslash{}tau''\_\{\textbackslash{}geq -N\}K\_1 at First proof paragraph, definition of K\_2 changes a substantive mathematical assertion or bibliographic identification; no independent derivation of this exact replacement is established in this review.\par
\par\noindent\textbf{OM-SOL-1004 -- t3 XVII p.331 -- SOL}\par
Second proof paragraph, choice of i before the map b\par
\emph{Reported reading}: pour i assez grand, b\par
\emph{Proposal}: pour \$i\$ assez petit\par
The replacement pour \$i\$ assez petit at Second proof paragraph, choice of i before the map b requires a complete check of the maps, their endpoints and the composition convention; this identification is not independently verified.\par
\par\noindent\textbf{OM-SOL-1005 -- t3 XVII p.331 -- SOL}\par
Second proof paragraph, map b\par
\emph{Reported reading}: b:K\textbackslash{}rightarrow\textbackslash{}sigma\_\{\textbackslash{}geq i\}K\_1=\allowbreak K\_3\par
\emph{Proposal}: b:K\textbackslash{}rightarrow\textbackslash{}tau\_\{\textbackslash{}geq i\}K\_1=\allowbreak K\_3\par
This finding intersects Orgogozo's \textbackslash{}sisi material at Second proof paragraph, map b; preserve that human entry unchanged. The proposed reading b:K\textbackslash{}rightarrow\textbackslash{}tau\_\{\textbackslash{}geq i\}K\_1=\allowbreak K\_3 at Second proof paragraph, map b changes a substantive mathematical assertion or bibliographic identification; no independent derivation of this exact replacement is established in this review.\par
\par\noindent\textbf{OM-SOL-1007 -- t3 XVII p.331 -- SOL}\par
Second proof paragraph, definition of K\_4\par
\emph{Reported reading}: K\_4=\allowbreak \textbackslash{}sigma''\_\{\textbackslash{}leq M\}C\textasciicircum{}\{\textbackslash{}ast\}K\_3\par
\emph{Proposal}: K\_4=\allowbreak \textbackslash{}tau''\_\{\textbackslash{}leq M\}C\textasciicircum{}\{\textbackslash{}ast\}K\_3\par
This finding intersects Orgogozo's \textbackslash{}sisi material at Second proof paragraph, definition of K\_4; preserve that human entry unchanged. The proposed reading K\_4=\allowbreak \textbackslash{}tau''\_\{\textbackslash{}leq M\}C\textasciicircum{}\{\textbackslash{}ast\}K\_3 at Second proof paragraph, definition of K\_4 changes a substantive mathematical assertion or bibliographic identification; no independent derivation of this exact replacement is established in this review.\par
\par\noindent\textbf{OM-SOL-1008 -- t3 XVII p.332 -- SOL}\par
Diagram (4.2.12.2), vertical arrows k and k'\par
\emph{Reported reading}: X',X are above Q',Q, with k',k pointing downward\par
\emph{Proposal}: \textbackslash{}xymatrix\{Q'\textbackslash{}ar[r]\textasciicircum{}\{g'\_0\}\textbackslash{}ar[d]\textasciicircum{}\{k'\}\&Q\textbackslash{}ar[d]\textasciicircum{}\{k\}\textbackslash{}\textbackslash{}X'\textbackslash{}ar[r]\textasciicircum{}\{g'\}\&X\}\par
This finding intersects Orgogozo's \textbackslash{}sisi material at Diagram (4.2.12.2), vertical arrows k and k'; preserve that human entry unchanged. The replacement \textbackslash{}xymatrix\{Q'\textbackslash{}ar[r]\textasciicircum{}\{g'\_0\}\textbackslash{}ar[d]\textasciicircum{}\{k'\}\&Q\textbackslash{}ar[d]\textasciicircum{}\{k\}\textbackslash{}\textbackslash{}X'\textbackslash{}ar[r]\textasciicircum{}\{g'\}\&X\} at Diagram (4.2.12.2), vertical arrows k and k' requires a complete check of the maps, their endpoints and the composition convention; this identification is not independently verified.\par
\par\noindent\textbf{OM-SOL-1009 -- t3 XVII p.333 -- SOL}\par
Bound on N, after « sur »\par
\emph{Reported reading}: N\textbackslash{}geq\textbackslash{}sup\_i\textbackslash{}text\{(tor dim de \$L\textasciicircum{}i\$ sur \$g\textasciicircum{}\{\textbackslash{}ast\}\$)\}\par
\emph{Proposal}: N\textbackslash{}geq\textbackslash{}sup\_i\textbackslash{}text\{(tor-dim de \$L\textasciicircum{}i\$ sur \$g\textasciicircum{}\{\textbackslash{}ast\}\textbackslash{}Aa\$)\}\par
This finding intersects Orgogozo's \textbackslash{}sisi material at Bound on N, after « sur »; preserve that human entry unchanged. The proposed reading N\textbackslash{}geq\textbackslash{}sup\_i\textbackslash{}text\{(tor-dim de \$L\textasciicircum{}i\$ sur \$g\textasciicircum{}\{\textbackslash{}ast\}\textbackslash{}Aa\$)\} at Bound on N, after « sur » changes a substantive mathematical assertion or bibliographic identification; no independent derivation of this exact replacement is established in this review.\par
\par\noindent\textbf{OM-SOL-1010 -- t3 XVII p.333 -- SOL}\par
Functor immediately before reference (4.2.8), tensor subscript\par
\emph{Reported reading}: \{f'\}\textasciicircum{}\{\textbackslash{}ast\}L\textbackslash{}overset\{\textbackslash{}LL\}\{\textbackslash{}otimes\}\_\{\textbackslash{}Aa'\}\{g'\}\textasciicircum{}\{\textbackslash{}ast\}\par
\emph{Proposal}: \{f'\}\textasciicircum{}\{\textbackslash{}ast\}L\textbackslash{}overset\{\textbackslash{}LL\}\{\textbackslash{}otimes\}\_\{\{f'\}\textasciicircum{}\{\textbackslash{}ast\}g\textasciicircum{}\{\textbackslash{}ast\}\textbackslash{}Aa\}\textbackslash{}LL\{g'\}\textasciicircum{}\{\textbackslash{}ast\}\par
This finding intersects Orgogozo's \textbackslash{}sisi material at Functor immediately before reference (4.2.8), tensor subscript; preserve that human entry unchanged. The proposed reading \{f'\}\textasciicircum{}\{\textbackslash{}ast\}L\textbackslash{}overset\{\textbackslash{}LL\}\{\textbackslash{}otimes\}\_\{\{f'\}\textasciicircum{}\{\textbackslash{}ast\}g\textasciicircum{}\{\textbackslash{}ast\}\textbackslash{}Aa\}\textbackslash{}LL\{g'\}\textasciicircum{}\{\textbackslash{}ast\} at Functor immediately before reference (4.2.8), tensor subscript changes a substantive mathematical assertion or bibliographic identification; no independent derivation of this exact replacement is established in this review.\par
\par\noindent\textbf{OM-SOL-1011 -- t3 XVII p.333 -- SOL}\par
Displayed canonical map \ensuremath{\beta}\par
\emph{Reported reading}: \textbackslash{}beta:\{f'\}\textasciicircum{}\{\textbackslash{}ast\}L\textbackslash{}overset\{\textbackslash{}LL\}\{\textbackslash{}otimes\}\textbackslash{}LL g\textasciicircum{}\{\textbackslash{}ast\}K\textbackslash{}longrightarrow L\textbackslash{}otimes g\textasciicircum{}\{\textbackslash{}ast\}K\par
\emph{Proposal}: \textbackslash{}beta:\{f'\}\textasciicircum{}\{\textbackslash{}ast\}L\textbackslash{}overset\{\textbackslash{}LL\}\{\textbackslash{}otimes\}\textbackslash{}LL\{g'\}\textasciicircum{}\{\textbackslash{}ast\}K\textbackslash{}longrightarrow \{f'\}\textasciicircum{}\{\textbackslash{}ast\}L\textbackslash{}otimes \{g'\}\textasciicircum{}\{\textbackslash{}ast\}K\par
This finding intersects Orgogozo's \textbackslash{}sisi material at Displayed canonical map \ensuremath{\beta}; preserve that human entry unchanged. The replacement \textbackslash{}beta:\{f'\}\textasciicircum{}\{\textbackslash{}ast\}L\textbackslash{}overset\{\textbackslash{}LL\}\{\textbackslash{}otimes\}\textbackslash{}LL\{g'\}\textasciicircum{}\{\textbackslash{}ast\}K\textbackslash{}longrightarrow \{f'\}\textasciicircum{}\{\textbackslash{}ast\}L\textbackslash{}otimes \{g'\}\textasciicircum{}\{\textbackslash{}ast\}K at Displayed canonical map \ensuremath{\beta} identifies a particular indexed object or symbol; that identification is not independently established by the quoted local expression.\par
\par\noindent\textbf{OM-SOL-1012 -- t3 XVII p.333 -- SOL}\par
Last displayed identification of I\par
\emph{Reported reading}: I\textbackslash{}simeq L\textbackslash{}otimes\textbackslash{}LL g\textasciicircum{}\{\textbackslash{}ast\}K\_1\par
\emph{Proposal}: I\textbackslash{}simeq L\textbackslash{}otimes g\textasciicircum{}\{\textbackslash{}ast\}K\_1\par
This finding intersects Orgogozo's \textbackslash{}sisi material at Last displayed identification of I; preserve that human entry unchanged. The proposed reading I\textbackslash{}simeq L\textbackslash{}otimes g\textasciicircum{}\{\textbackslash{}ast\}K\_1 at Last displayed identification of I changes a substantive mathematical assertion or bibliographic identification; no independent derivation of this exact replacement is established in this review.\par
\par\noindent\textbf{OM-SOL-1013 -- t3 XVII p.334 -- SOL}\par
Composite defining base change, arrow into f'\_*K\_2\par
\emph{Reported reading}: \textbackslash{}xrightarrow\{\textbackslash{}alpha\}f'\_\{\textbackslash{}ast\}K\_2\par
\emph{Proposal}: \textbackslash{}xrightarrow\{\textbackslash{}delta\}f'\_\{\textbackslash{}ast\}K\_2\par
This finding intersects Orgogozo's \textbackslash{}sisi material at Composite defining base change, arrow into f'\_*K\_2; preserve that human entry unchanged. The replacement \textbackslash{}xrightarrow\{\textbackslash{}delta\}f'\_\{\textbackslash{}ast\}K\_2 at Composite defining base change, arrow into f'\_*K\_2 requires a complete check of the maps, their endpoints and the composition convention; this identification is not independently verified.\par
\par\noindent\textbf{OM-SOL-1014 -- t3 XVII p.334 -- SOL}\par
Proposition 4.2.13, category containing K\par
\emph{Reported reading}: K\textbackslash{}in\textbackslash{}Ob C\textasciicircum{}s(X',\textbackslash{}Aa)\par
\emph{Proposal}: K\textbackslash{}in\textbackslash{}Ob C\textasciicircum{}s(X,\textbackslash{}Aa)\par
This finding intersects Orgogozo's \textbackslash{}sisi material at Proposition 4.2.13, category containing K; preserve that human entry unchanged. The replacement K\textbackslash{}in\textbackslash{}Ob C\textasciicircum{}s(X,\textbackslash{}Aa) at Proposition 4.2.13, category containing K identifies a particular indexed object or symbol; that identification is not independently established by the quoted local expression.\par
\par\noindent\textbf{OM-SOL-1015 -- t3 XVII p.334 -- SOL}\par
Proposition 4.2.13, category containing K\_1\par
\emph{Reported reading}: K\_1\textbackslash{}in\textbackslash{}Ob C\textasciicircum{}s(X,\textbackslash{}Aa)\par
\emph{Proposal}: K\_1\textbackslash{}in\textbackslash{}Ob C\textasciicircum{}s(S,\textbackslash{}Aa)\par
This finding intersects Orgogozo's \textbackslash{}sisi material at Proposition 4.2.13, category containing K\_1; preserve that human entry unchanged. The proposed reading K\_1\textbackslash{}in\textbackslash{}Ob C\textasciicircum{}s(S,\textbackslash{}Aa) at Proposition 4.2.13, category containing K\_1 changes a substantive mathematical assertion or bibliographic identification; no independent derivation of this exact replacement is established in this review.\par
\par\noindent\textbf{OM-SOL-1016 -- t3 XVII p.334 -- SOL}\par
Proposition 4.2.13, category containing K\_2\par
\emph{Reported reading}: K\_2\textbackslash{}in\textbackslash{}Ob C\textasciicircum{}s(S,\textbackslash{}Aa)\par
\emph{Proposal}: K\_2\textbackslash{}in\textbackslash{}Ob C\textasciicircum{}s(X',\textbackslash{}Aa')\par
This finding intersects Orgogozo's \textbackslash{}sisi material at Proposition 4.2.13, category containing K\_2; preserve that human entry unchanged. The replacement K\_2\textbackslash{}in\textbackslash{}Ob C\textasciicircum{}s(X',\textbackslash{}Aa') at Proposition 4.2.13, category containing K\_2 identifies a particular indexed object or symbol; that identification is not independently established by the quoted local expression.\par
\par\noindent\textbf{OM-SOL-1017 -- t3 XVII p.334 -- SOL}\par
Proposition 4.2.13, target of c\_\{u,v\}\par
\emph{Reported reading}: f'\_\{\textbackslash{}ast\}k\_2\par
\emph{Proposal}: f'\_\{\textbackslash{}ast\}K\_2\par
This finding intersects Orgogozo's \textbackslash{}sisi material at Proposition 4.2.13, target of c\_\{u,v\}; preserve that human entry unchanged. The replacement f'\_\{\textbackslash{}ast\}K\_2 at Proposition 4.2.13, target of c\_\{u,v\} identifies a particular indexed object or symbol; that identification is not independently established by the quoted local expression.\par
\par\noindent\textbf{OM-SOL-1018 -- t3 XVII p.334 -- SOL}\par
Bottom display, definition of v'\par
\emph{Reported reading}: v':L\textbackslash{}overset\{\textbackslash{}LL\}\{\textbackslash{}otimes\}\textbackslash{}LL\{g'\}\textasciicircum{}\{\textbackslash{}ast\}K\textbackslash{}rightarrow L\textbackslash{}otimes \{g'\}\textasciicircum{}\{\textbackslash{}ast\}K\textbackslash{}xrightarrow\{v\}K\_2\par
\emph{Proposal}: v':\{f'\}\textasciicircum{}\{\textbackslash{}ast\}L\textbackslash{}overset\{\textbackslash{}LL\}\{\textbackslash{}otimes\}\textbackslash{}LL\{g'\}\textasciicircum{}\{\textbackslash{}ast\}K\textbackslash{}rightarrow \{f'\}\textasciicircum{}\{\textbackslash{}ast\}L\textbackslash{}otimes \{g'\}\textasciicircum{}\{\textbackslash{}ast\}K\textbackslash{}xrightarrow\{v\}K\_2\par
The proposed reading v':\{f'\}\textasciicircum{}\{\textbackslash{}ast\}L\textbackslash{}overset\{\textbackslash{}LL\}\{\textbackslash{}otimes\}\textbackslash{}LL\{g'\}\textasciicircum{}\{\textbackslash{}ast\}K\textbackslash{}rightarrow \{f'\}\textasciicircum{}\{\textbackslash{}ast\}L\textbackslash{}otimes \{g'\}\textasciicircum{}\{\textbackslash{}ast\}K\textbackslash{}xrightarrow\{v\}K\_2 at Bottom display, definition of v' changes a substantive mathematical assertion or bibliographic identification; no independent derivation of this exact replacement is established in this review.\par
\par\noindent\textbf{OM-SOL-1019 -- t3 XVII p.334 -- SOL}\par
Bottom display, first two terms defining c'\_\{u,v\}\par
\emph{Reported reading}: L\textbackslash{}overset\{\textbackslash{}LL\}\{\textbackslash{}otimes\}\textbackslash{}LL\{g'\}\textasciicircum{}\{\textbackslash{}ast\}K\_1\textbackslash{}rightarrow L\textbackslash{}otimes \{g'\}\textasciicircum{}\{\textbackslash{}ast\}K\_1\par
\emph{Proposal}: L\textbackslash{}overset\{\textbackslash{}LL\}\{\textbackslash{}otimes\}\textbackslash{}LL g\textasciicircum{}\{\textbackslash{}ast\}K\_1\textbackslash{}rightarrow L\textbackslash{}otimes g\textasciicircum{}\{\textbackslash{}ast\}K\_1\par
The replacement L\textbackslash{}overset\{\textbackslash{}LL\}\{\textbackslash{}otimes\}\textbackslash{}LL g\textasciicircum{}\{\textbackslash{}ast\}K\_1\textbackslash{}rightarrow L\textbackslash{}otimes g\textasciicircum{}\{\textbackslash{}ast\}K\_1 at Bottom display, first two terms defining c'\_\{u,v\} identifies a particular indexed object or symbol; that identification is not independently established by the quoted local expression.\par
\par\noindent\textbf{OM-SOL-1020 -- t3 XVII p.335 -- SOL}\par
Diagram (4.2.13.1), four derived-functor positions\par
\emph{Reported reading}: A plain L precedes g\textasciicircum{}\{\textbackslash{}ast\} in both left vertices and the left arrow label, and g'\textasciicircum{}\{\textbackslash{}ast\} in the upper-right vertex.\par
\emph{Proposal}: L\textbackslash{}overset\{\textbackslash{}LL\}\{\textbackslash{}otimes\}\textbackslash{}LL g\textasciicircum{}\{\textbackslash{}ast\}Rf\_\{\textbackslash{}ast\}K;\textbackslash{}quad Rf'\_\{\textbackslash{}ast\}(f'\textasciicircum{}\{\textbackslash{}ast\}L\textbackslash{}overset\{\textbackslash{}LL\}\{\textbackslash{}otimes\}\textbackslash{}LL g'\textasciicircum{}\{\textbackslash{}ast\}K);\textbackslash{}quad L\textbackslash{}overset\{\textbackslash{}LL\}\{\textbackslash{}otimes\}\textbackslash{}LL g\textasciicircum{}\{\textbackslash{}ast\}K\_1;\textbackslash{}quad L\textbackslash{}overset\{\textbackslash{}LL\}\{\textbackslash{}otimes\}\textbackslash{}LL g\textasciicircum{}\{\textbackslash{}ast\}(u')\par
This finding intersects Orgogozo's \textbackslash{}sisi material at Diagram (4.2.13.1), four derived-functor positions; preserve that human entry unchanged. The proposed reading L\textbackslash{}overset\{\textbackslash{}LL\}\{\textbackslash{}otimes\}\textbackslash{}LL g\textasciicircum{}\{\textbackslash{}ast\}Rf\_\{\textbackslash{}ast\}K;\textbackslash{}quad Rf'\_\{\textbackslash{}ast\}(f'\textasciicircum{}\{\textbackslash{}ast\}L\textbackslash{}overset\{\textbackslash{}LL\}\{\textbackslash{}otimes\}\textbackslash{}LL g'\textasciicircum{}\{\textbackslash{}ast\}K);\textbackslash{}quad L\textbackslash{}overset\{\textbackslash{}LL\}\{\textbackslash{}otimes\}\textbackslash{}LL g\textasciicircum{}\{\textbackslash{}ast\}K\_1;\textbackslash{}quad L\textbackslash{}overset\{\textbackslash{}LL\}\{\textbackslash{}otimes\}\textbackslash{}LL g\textasciicircum{}\{\textbackslash{}ast\}(u') at Diagram (4.2.13.1), four derived-functor positions changes a substantive mathematical assertion or bibliographic identification; no independent derivation of this exact replacement is established in this review.\par
\par\noindent\textbf{OM-SOL-1021 -- t3 XVII p.335 -- SOL}\par
Functor for which K\textasciicircum{}a is deployed\par
\emph{Reported reading}: L\{f'\}\textasciicircum{}\{\textbackslash{}ast\}L\textbackslash{}overset\{\textbackslash{}LL\}\{\textbackslash{}otimes\}L\{g'\}\textasciicircum{}\{\textbackslash{}ast\}\par
\emph{Proposal}: \{f'\}\textasciicircum{}\{\textbackslash{}ast\}L\textbackslash{}overset\{\textbackslash{}LL\}\{\textbackslash{}otimes\}\textbackslash{}LL\{g'\}\textasciicircum{}\{\textbackslash{}ast\}\par
This finding intersects Orgogozo's \textbackslash{}sisi material at Functor for which K\textasciicircum{}a is deployed; preserve that human entry unchanged. The proposed reading \{f'\}\textasciicircum{}\{\textbackslash{}ast\}L\textbackslash{}overset\{\textbackslash{}LL\}\{\textbackslash{}otimes\}\textbackslash{}LL\{g'\}\textasciicircum{}\{\textbackslash{}ast\} at Functor for which K\textasciicircum{}a is deployed changes a substantive mathematical assertion or bibliographic identification; no independent derivation of this exact replacement is established in this review.\par
\par\noindent\textbf{OM-SOL-1022 -- t3 XVII p.335 -- SOL}\par
Last paragraph, parenthetical functor applied to the diagram\par
\emph{Reported reading}: L\textbackslash{}otimes \{g'\}\textasciicircum{}\{\textbackslash{}ast\}\par
\emph{Proposal}: \{f'\}\textasciicircum{}\{\textbackslash{}ast\}L\textbackslash{}otimes \{g'\}\textasciicircum{}\{\textbackslash{}ast\}\par
The proposed reading \{f'\}\textasciicircum{}\{\textbackslash{}ast\}L\textbackslash{}otimes \{g'\}\textasciicircum{}\{\textbackslash{}ast\} at Last paragraph, parenthetical functor applied to the diagram changes a substantive mathematical assertion or bibliographic identification; no independent derivation of this exact replacement is established in this review.\par
\par\noindent\textbf{OM-SOL-1023 -- t3 XVII p.336 -- SOL}\par
Sentence immediately before the large prism\par
\emph{Reported reading}: de complexes sur S\par
\emph{Proposal}: S'\par
The base of all complexes in the prism requires an endpoint check of its functors.\par
\par\noindent\textbf{OM-SOL-1024 -- t3 XVII p.336 -- SOL}\par
Large prism, second row, second object\par
\emph{Reported reading}: L\textbackslash{}otimes g\textasciicircum{}\{\textbackslash{}ast\}K\_1\par
\emph{Proposal}: L\textbackslash{}otimes g\textasciicircum{}\{\textbackslash{}ast\}K\_1\textasciicircum{}a\par
This finding intersects Orgogozo's \textbackslash{}sisi material at Large prism, second row, second object; preserve that human entry unchanged. The replacement L\textbackslash{}otimes g\textasciicircum{}\{\textbackslash{}ast\}K\_1\textasciicircum{}a at Large prism, second row, second object identifies a particular indexed object or symbol; that identification is not independently established by the quoted local expression.\par
\par\noindent\textbf{OM-SOL-1025 -- t3 XVII p.336 -- SOL}\par
Large prism, third row, leftmost object\par
\emph{Reported reading}: L\textbackslash{}otimes g\textasciicircum{}\{\textbackslash{}ast\}K\_1\textasciicircum{}\{c'\}\par
\emph{Proposal}: L\textbackslash{}otimes g\textasciicircum{}\{\textbackslash{}ast\}K\_1\textasciicircum{}\{b'\}\par
This finding intersects Orgogozo's \textbackslash{}sisi material at Large prism, third row, leftmost object; preserve that human entry unchanged. The replacement L\textbackslash{}otimes g\textasciicircum{}\{\textbackslash{}ast\}K\_1\textasciicircum{}\{b'\} at Large prism, third row, leftmost object identifies a particular indexed object or symbol; that identification is not independently established by the quoted local expression.\par
\par\noindent\textbf{OM-SOL-1026 -- t3 XVII p.336 -- SOL}\par
Large prism, third row, second object\par
\emph{Reported reading}: L\textbackslash{}otimes g\textasciicircum{}\{\textbackslash{}ast\}K\textasciicircum{}\{b\}\_2\par
\emph{Proposal}: L\textbackslash{}otimes g\textasciicircum{}\{\textbackslash{}ast\}K\textasciicircum{}\{b\}\_1\par
The proposed reading L\textbackslash{}otimes g\textasciicircum{}\{\textbackslash{}ast\}K\textasciicircum{}\{b\}\_1 at Large prism, third row, second object changes a substantive mathematical assertion or bibliographic identification; no independent derivation of this exact replacement is established in this review.\par
\par\noindent\textbf{OM-SOL-1027 -- t3 XVII p.336 -- SOL}\par
Large prism, bottom row, second object\par
\emph{Reported reading}: L\textbackslash{}otimes g\textasciicircum{}\{\textbackslash{}ast\}K\textasciicircum{}\{c\}\_2\par
\emph{Proposal}: L\textbackslash{}otimes g\textasciicircum{}\{\textbackslash{}ast\}K\textasciicircum{}\{c\}\_1\par
The proposed reading L\textbackslash{}otimes g\textasciicircum{}\{\textbackslash{}ast\}K\textasciicircum{}\{c\}\_1 at Large prism, bottom row, second object changes a substantive mathematical assertion or bibliographic identification; no independent derivation of this exact replacement is established in this review.\par
\par\noindent\textbf{OM-SOL-1029 -- t3 XVII p.337 -- SOL}\par
Paragraph beginning « Dans le cas général »\par
\emph{Reported reading}: K'\_\{2\}=\allowbreak \{f'\}\textasciicircum{}\{\textbackslash{}ast\}L\textbackslash{}otimes g\textasciicircum{}\{\textbackslash{}ast\}K\par
\emph{Proposal}: K'\_\{2\}=\allowbreak \{f'\}\textasciicircum{}\{\textbackslash{}ast\}L\textbackslash{}otimes \{g'\}\textasciicircum{}\{\textbackslash{}ast\}K\par
The replacement K'\_\{2\}=\allowbreak \{f'\}\textasciicircum{}\{\textbackslash{}ast\}L\textbackslash{}otimes \{g'\}\textasciicircum{}\{\textbackslash{}ast\}K at Paragraph beginning « Dans le cas général » identifies a particular indexed object or symbol; that identification is not independently established by the quoted local expression.\par
\par\noindent\textbf{OM-SOL-1030 -- t3 XVII p.338 -- SOL}\par
Theorem 4.3.1, displayed base-change morphism\par
\emph{Reported reading}: \textbackslash{}varphi\_K:L\textbackslash{}overset\{\textbackslash{}LL\}\{\textbackslash{}otimes\}Lg\textasciicircum{}\{\textbackslash{}ast\}K\textbackslash{}longrightarrow\textbackslash{}cdots\par
\emph{Proposal}: \textbackslash{}varphi\_K:L\textbackslash{}overset\{\textbackslash{}LL\}\{\textbackslash{}otimes\}\textbackslash{}LL g\textasciicircum{}\{\textbackslash{}ast\}Rf\_\{\textbackslash{}ast\}K\par
This finding intersects Orgogozo's \textbackslash{}sisi material at Theorem 4.3.1, displayed base-change morphism; preserve that human entry unchanged. The replacement \textbackslash{}varphi\_K:L\textbackslash{}overset\{\textbackslash{}LL\}\{\textbackslash{}otimes\}\textbackslash{}LL g\textasciicircum{}\{\textbackslash{}ast\}Rf\_\{\textbackslash{}ast\}K at Theorem 4.3.1, displayed base-change morphism requires a complete check of the maps, their endpoints and the composition convention; this identification is not independently verified.\par
\par\noindent\textbf{OM-SOL-1031 -- t3 XVII p.339 -- SOL}\par
First displayed formula, target of \textbackslash{}varphi\textasciicircum{}q\par
\emph{Reported reading}: R\textasciicircum{}\{q\}f'\_\{\textbackslash{}ast\}(\{f'\}\textasciicircum{}\{\textbackslash{}ast\}L\textbackslash{}otimes g\textasciicircum{}\{\textbackslash{}ast\}K)\par
\emph{Proposal}: R\textasciicircum{}\{q\}f'\_\{\textbackslash{}ast\}(\{f'\}\textasciicircum{}\{\textbackslash{}ast\}L\textbackslash{}otimes \{g'\}\textasciicircum{}\{\textbackslash{}ast\}K)\par
The replacement R\textasciicircum{}\{q\}f'\_\{\textbackslash{}ast\}(\{f'\}\textasciicircum{}\{\textbackslash{}ast\}L\textbackslash{}otimes \{g'\}\textasciicircum{}\{\textbackslash{}ast\}K) at First displayed formula, target of \textbackslash{}varphi\textasciicircum{}q identifies a particular indexed object or symbol; that identification is not independently established by the quoted local expression.\par
\par\noindent\textbf{OM-SOL-1033 -- t3 XVII p.339 -- SOL}\par
Sentence introducing the filtered colimit of the maps \textbackslash{}varphi\_i\textasciicircum{}q\par
\emph{Reported reading}: La flèche \textbackslash{}varphi\textasciicircum{}\{q\} est donc limite inductive\par
\emph{Proposal}: La flèche \$\textbackslash{}varphi\textasciicircum{}\{q\}\_s\$ est donc limite\par
The replacement La flèche \$\textbackslash{}varphi\textasciicircum{}\{q\}\_s\$ est donc limite at Sentence introducing the filtered colimit of the maps \textbackslash{}varphi\_i\textasciicircum{}q requires a complete check of the maps, their endpoints and the composition convention; this identification is not independently verified.\par
\par\noindent\textbf{OM-SOL-1034 -- t3 XVII p.339 -- SOL}\par
Final conclusion\par
\emph{Reported reading}: et \textbackslash{}varphi est un isomorphisme\par
\emph{Proposal}: \textbackslash{}varphi\textasciicircum{}\{q\}\par
This finding intersects Orgogozo's \textbackslash{}sisi material at Final conclusion; preserve that human entry unchanged. The replacement \textbackslash{}varphi\textasciicircum{}\{q\} at Final conclusion requires a complete check of the maps, their endpoints and the composition convention; this identification is not independently verified.\par
\par\noindent\textbf{OM-SOL-1036 -- t3 XVII p.341 -- SOL}\par
Opening paragraph, reference after the restriction functor\par
\emph{Reported reading}: (IV followed by an unfilled blank)\par
\emph{Proposal}: \textbackslash{}hyperrefext[IV][IV.prop11.3.1]\{11.3.1\}\par
This finding intersects Orgogozo's \textbackslash{}sisi material at Opening paragraph, reference after the restriction functor; preserve that human entry unchanged. The proposed reference \textbackslash{}hyperrefext[IV][IV.prop11.3.1]\{11.3.1\} at Opening paragraph, reference after the restriction functor needs comparison with the statement actually invoked; that target has not been independently established here.\par
\par\noindent\textbf{OM-SOL-1037 -- t3 XVII p.344 -- SOL}\par
Conclusion of the proof of Lemma 5.1.6\par
\emph{Reported reading}: nuls en dehors de U\par
\emph{Proposal}: nuls en dehors de \$V\$\par
The proposed reading nuls en dehors de \$V\$ at Conclusion of the proof of Lemma 5.1.6 changes a substantive mathematical assertion or bibliographic identification; no independent derivation of this exact replacement is established in this review.\par
\par\noindent\textbf{OM-SOL-1038 -- t3 XVII p.346 -- SOL}\par
Théorème 5.1.8, (a)(iii), formule de l’isomorphisme composé\par
\emph{Reported reading}: c\_\{f,j'\} \textbackslash{}circ c\textasciicircum{}\{-1\}\_\{j,f'\}\par
\emph{Proposal}: c\_\{f,j\_1\} \textbackslash{}circ c\textasciicircum{}\{-1\}\_\{j\_2,g\}\par
This finding intersects Orgogozo's \textbackslash{}sisi material at Théorème 5.1.8, (a)(iii), formule de l’isomorphisme composé; preserve that human entry unchanged. The replacement c\_\{f,j\_1\} \textbackslash{}circ c\textasciicircum{}\{-1\}\_\{j\_2,g\} at Théorème 5.1.8, (a)(iii), formule de l’isomorphisme composé requires a complete check of the maps, their endpoints and the composition convention; this identification is not independently verified.\par
\par\noindent\textbf{OM-SOL-1040 -- t3 XVII p.347 -- SOL}\par
Variante 5.1.11, référence finale\par
\emph{Reported reading}: analogues à celles de 5.1.4\par
\emph{Proposal}: 5.1.8\par
This finding intersects Orgogozo's \textbackslash{}sisi material at Variante 5.1.11, référence finale; preserve that human entry unchanged. The proposed reference 5.1.8 at Variante 5.1.11, référence finale needs comparison with the statement actually invoked; that target has not been independently established here.\par
\par\noindent\textbf{OM-SOL-1042 -- t3 XVII p.348 -- SOL}\par
Suite exacte non numérotée, les deux termes entourant la flèche \ensuremath{\alpha}\par
\emph{Reported reading}: R\textasciicircum{}\{1\}f\_\{2\textbackslash{}ast\}(p\_\{\textbackslash{}ast\}i\_\{1!\}G)\textbackslash{}xrightarrow\{\textbackslash{}alpha\}R\textasciicircum{}\{1\}f\_\{1\textbackslash{}ast\}(i\_\{1!\}G)\par
\emph{Proposal}: R\textasciicircum{}\{1\}\textbackslash{}bar\{f\}\_\{2\textbackslash{}ast\}(p\_\{\textbackslash{}ast\}i\_\{1!\}G)\textbackslash{}xrightarrow\{\textbackslash{}alpha\}R\textasciicircum{}\{1\}\textbackslash{}bar\{f\}\_\{1\textbackslash{}ast\}(i\_\{1!\}G)\par
This finding intersects Orgogozo's \textbackslash{}sisi material at Suite exacte non numérotée, les deux termes entourant la flèche \ensuremath{\alpha}; preserve that human entry unchanged. The replacement R\textasciicircum{}\{1\}\textbackslash{}bar\{f\}\_\{2\textbackslash{}ast\}(p\_\{\textbackslash{}ast\}i\_\{1!\}G)\textbackslash{}xrightarrow\{\textbackslash{}alpha\}R\textasciicircum{}\{1\}\textbackslash{}bar\{f\}\_\{1\textbackslash{}ast\}(i\_\{1!\}G) at Suite exacte non numérotée, les deux termes entourant la flèche \ensuremath{\alpha} requires a complete check of the maps, their endpoints and the composition convention; this identification is not independently verified.\par
\par\noindent\textbf{OM-SOL-1043 -- t3 XVII p.348 -- SOL}\par
Phrase définissant la flèche \ensuremath{\beta}\par
\emph{Reported reading}: de R\textasciicircum{}\{1\}f\_\{2\textbackslash{}ast\}i\_\{2!\}G dans R\textasciicircum{}\{1\}f\_\{1\textbackslash{}ast\}i\_\{1!\}G\par
\emph{Proposal}: R\textasciicircum{}\{1\}\textbackslash{}bar\{f\}\_\{2\textbackslash{}ast\}i\_\{2!\}G\textbackslash{}longrightarrow R\textasciicircum{}\{1\}\textbackslash{}bar\{f\}\_\{1\textbackslash{}ast\}i\_\{1!\}G\par
This finding intersects Orgogozo's \textbackslash{}sisi material at Phrase définissant la flèche \ensuremath{\beta}; preserve that human entry unchanged. The replacement R\textasciicircum{}\{1\}\textbackslash{}bar\{f\}\_\{2\textbackslash{}ast\}i\_\{2!\}G\textbackslash{}longrightarrow R\textasciicircum{}\{1\}\textbackslash{}bar\{f\}\_\{1\textbackslash{}ast\}i\_\{1!\}G at Phrase définissant la flèche \ensuremath{\beta} requires a complete check of the maps, their endpoints and the composition convention; this identification is not independently verified.\par
\par\noindent\textbf{OM-SOL-1044 -- t3 XVII p.349 -- SOL}\par
Suite exacte en tête de page, les deux termes entourant \ensuremath{\beta}\par
\emph{Reported reading}: R\textasciicircum{}\{1\}f\_\{2\textbackslash{}ast\}(i\_\{2!\}G)\textbackslash{}xrightarrow\{\textbackslash{}beta\}R\textasciicircum{}\{1\}f\_\{1\textbackslash{}ast\}(i\_\{1!\}G)\par
\emph{Proposal}: R\textasciicircum{}\{1\}\textbackslash{}bar\{f\}\_\{2\textbackslash{}ast\}(i\_\{2!\}G)\textbackslash{}xrightarrow\{\textbackslash{}beta\}R\textasciicircum{}\{1\}\textbackslash{}bar\{f\}\_\{1\textbackslash{}ast\}(i\_\{1!\}G)\par
This finding intersects Orgogozo's \textbackslash{}sisi material at Suite exacte en tête de page, les deux termes entourant \ensuremath{\beta}; preserve that human entry unchanged. The replacement R\textasciicircum{}\{1\}\textbackslash{}bar\{f\}\_\{2\textbackslash{}ast\}(i\_\{2!\}G)\textbackslash{}xrightarrow\{\textbackslash{}beta\}R\textasciicircum{}\{1\}\textbackslash{}bar\{f\}\_\{1\textbackslash{}ast\}(i\_\{1!\}G) at Suite exacte en tête de page, les deux termes entourant \ensuremath{\beta} requires a complete check of the maps, their endpoints and the composition convention; this identification is not independently verified.\par
\par\noindent\textbf{OM-SOL-1045 -- t3 XVII p.349 -- SOL}\par
Variante 5.1.14, référence finale\par
\emph{Reported reading}: analogues à celles de 5.1.4\par
\emph{Proposal}: 5.1.8\par
The proposed reference 5.1.8 at Variante 5.1.14, référence finale needs comparison with the statement actually invoked; that target has not been independently established here.\par
\par\noindent\textbf{OM-SOL-1046 -- t3 XVII p.349 -- SOL}\par
Titre de la dernière variante\par
\emph{Reported reading}: Variante 5.1.5.\par
\emph{Proposal}: Variante 5.1.15\par
This finding intersects Orgogozo's \textbackslash{}sisi material at Titre de la dernière variante; preserve that human entry unchanged. The proposed reading Variante 5.1.15 at Titre de la dernière variante changes a substantive mathematical assertion or bibliographic identification; no independent derivation of this exact replacement is established in this review.\par
\par\noindent\textbf{OM-SOL-1048 -- t3 XVII p.353 -- SOL}\par
Premier paragraphe, second foncteur associé à (5.2.1.2)\par
\emph{Reported reading}: j'\_\{!\}(\{f'\}\textasciicircum{}\{\textbackslash{}ast\}L\textbackslash{}otimes \{g'\}\textasciicircum{}\{\textbackslash{}ast\}K)\par
\emph{Proposal}: j'\_\{!\}(\{f'\}\textasciicircum{}\{\textbackslash{}ast\}L\textbackslash{}overset\{\textbackslash{}LL\}\{\textbackslash{}otimes\}\{g'\}\textasciicircum{}\{\textbackslash{}ast\}K)\par
This finding intersects Orgogozo's \textbackslash{}sisi material at Premier paragraphe, second foncteur associé à (5.2.1.2); preserve that human entry unchanged. The proposed reading j'\_\{!\}(\{f'\}\textasciicircum{}\{\textbackslash{}ast\}L\textbackslash{}overset\{\textbackslash{}LL\}\{\textbackslash{}otimes\}\{g'\}\textasciicircum{}\{\textbackslash{}ast\}K) involves a printed glyph or a missing stroke at Premier paragraphe, second foncteur associé à (5.2.1.2); the textual record does not independently establish that glyph.\par
\par\noindent\textbf{OM-SOL-1049 -- t3 XVII p.354 -- SOL}\par
Phrase suivant la formule définissant c'\_1 et c'\_2\par
\emph{Reported reading}: après s'être restreint à X'\_2\par
\emph{Proposal}: X'\par
This finding intersects Orgogozo's \textbackslash{}sisi material at Phrase suivant la formule définissant c'\_1 et c'\_2; preserve that human entry unchanged. The proposed reading X' at Phrase suivant la formule définissant c'\_1 et c'\_2 changes a substantive mathematical assertion or bibliographic identification; no independent derivation of this exact replacement is established in this review.\par
\par\noindent\textbf{OM-SOL-1050 -- t3 XVII p.355 -- SOL}\par
Diagramme initial, sommet inférieur central\par
\emph{Reported reading}: Ru'\_\{!\}(u'\textasciicircum{}\{\textbackslash{}ast\}L\textbackslash{}overset\{\textbackslash{}LL\}\{\textbackslash{}otimes\}g'\textasciicircum{}\{\textbackslash{}ast\}Rv\_\{!\})\par
\emph{Proposal}: Ru'\_\{!\}(u'\textasciicircum{}\{\textbackslash{}ast\}L\textbackslash{}overset\{\textbackslash{}LL\}\{\textbackslash{}otimes\}g'\textasciicircum{}\{\textbackslash{}ast\}Rv\_\{!\}K)\par
This finding intersects Orgogozo's \textbackslash{}sisi material at Diagramme initial, sommet inférieur central; preserve that human entry unchanged. The proposed reading Ru'\_\{!\}(u'\textasciicircum{}\{\textbackslash{}ast\}L\textbackslash{}overset\{\textbackslash{}LL\}\{\textbackslash{}otimes\}g'\textasciicircum{}\{\textbackslash{}ast\}Rv\_\{!\}K) at Diagramme initial, sommet inférieur central changes a substantive mathematical assertion or bibliographic identification; no independent derivation of this exact replacement is established in this review.\par
\par\noindent\textbf{OM-SOL-1051 -- t3 XVII p.355 -- SOL}\par
Diagramme (5.2.4.1), étiquettes des deux colonnes de flèches verticales\par
\emph{Reported reading}: X'\textbackslash{}xrightarrow\{u'\}Y'\textbackslash{}xrightarrow\{v'\}S',\textbackslash{}qquad X\textbackslash{}xrightarrow\{u\}Y\textbackslash{}xrightarrow\{v\}S\par
\emph{Proposal}: X'\textbackslash{}xrightarrow\{v'\}Y'\textbackslash{}xrightarrow\{u'\}S',\textbackslash{}qquad X\textbackslash{}xrightarrow\{v\}Y\textbackslash{}xrightarrow\{u\}S\par
This finding intersects Orgogozo's \textbackslash{}sisi material at Diagramme (5.2.4.1), étiquettes des deux colonnes de flèches verticales; preserve that human entry unchanged. The proposed reading X'\textbackslash{}xrightarrow\{v'\}Y'\textbackslash{}xrightarrow\{u'\}S',\textbackslash{}qquad X\textbackslash{}xrightarrow\{v\}Y\textbackslash{}xrightarrow\{u\}S involves a printed glyph or a missing stroke at Diagramme (5.2.4.1), étiquettes des deux colonnes de flèches verticales; the textual record does not independently establish that glyph.\par
\par\noindent\textbf{OM-SOL-1052 -- t3 XVII p.357 -- SOL}\par
Lemma 5.2.5, lower-left entry of the second diagram\par
\emph{Reported reading}: Lg\_1\textasciicircum{}\{\textbackslash{}ast\}Lg\_2\textasciicircum{}\{\textbackslash{}ast\}Rf\_\{!\}\par
\emph{Proposal}: Lg\_2\textasciicircum{}\{\textbackslash{}ast\}Lg\_1\textasciicircum{}\{\textbackslash{}ast\}Rf\_\{!\}\par
This finding intersects Orgogozo's \textbackslash{}sisi material at Lemma 5.2.5, lower-left entry of the second diagram; preserve that human entry unchanged. The replacement Lg\_2\textasciicircum{}\{\textbackslash{}ast\}Lg\_1\textasciicircum{}\{\textbackslash{}ast\}Rf\_\{!\} at Lemma 5.2.5, lower-left entry of the second diagram identifies a particular indexed object or symbol; that identification is not independently established by the quoted local expression.\par
\par\noindent\textbf{OM-SOL-1053 -- t3 XVII p.357 -- SOL}\par
Lemma 5.2.5, lower-middle entry of the second diagram\par
\emph{Reported reading}: Lg\_1\textasciicircum{}\{\textbackslash{}ast\}Rf'\_\{!\}Lg\_2\textasciicircum{}\{\textbackslash{}ast\}\par
\emph{Proposal}: Lg\_2\textasciicircum{}\{\textbackslash{}ast\}Rf'\_\{!\}Lg'\textasciicircum{}\{\textbackslash{}ast\}\_1\par
This finding intersects Orgogozo's \textbackslash{}sisi material at Lemma 5.2.5, lower-middle entry of the second diagram; preserve that human entry unchanged. The proposed reading Lg\_2\textasciicircum{}\{\textbackslash{}ast\}Rf'\_\{!\}Lg'\textasciicircum{}\{\textbackslash{}ast\}\_1 at Lemma 5.2.5, lower-middle entry of the second diagram changes a substantive mathematical assertion or bibliographic identification; no independent derivation of this exact replacement is established in this review.\par
\par\noindent\textbf{OM-SOL-1054 -- t3 XVII p.357 -- SOL}\par
Lemma 5.2.5, lower-right entry of the second diagram\par
\emph{Reported reading}: Rf''\_\{!\}Lg'\textasciicircum{}\{\textbackslash{}ast\}\_1Lg'\textasciicircum{}\{\textbackslash{}ast\}\_2\par
\emph{Proposal}: Rf''\_\{!\}Lg'\textasciicircum{}\{\textbackslash{}ast\}\_2Lg'\textasciicircum{}\{\textbackslash{}ast\}\_1\par
This finding intersects Orgogozo's \textbackslash{}sisi material at Lemma 5.2.5, lower-right entry of the second diagram; preserve that human entry unchanged. The replacement Rf''\_\{!\}Lg'\textasciicircum{}\{\textbackslash{}ast\}\_2Lg'\textasciicircum{}\{\textbackslash{}ast\}\_1 at Lemma 5.2.5, lower-right entry of the second diagram identifies a particular indexed object or symbol; that identification is not independently established by the quoted local expression.\par
\par\noindent\textbf{OM-SOL-1055 -- t3 XVII p.358 -- SOL}\par
Proof of Corollary 5.2.8.1, dimension bound\par
\emph{Reported reading}: \textbackslash{}dim(\textbackslash{}bar\{X\})=\allowbreak \textbackslash{}dim(X)\textbackslash{}leq D\par
\emph{Proposal}: \textbackslash{}dim(\textbackslash{}bar\{X\})=\allowbreak \textbackslash{}dim(X)\textbackslash{}leq d\par
This finding intersects Orgogozo's \textbackslash{}sisi material at Proof of Corollary 5.2.8.1, dimension bound; preserve that human entry unchanged. The replacement \textbackslash{}dim(\textbackslash{}bar\{X\})=\allowbreak \textbackslash{}dim(X)\textbackslash{}leq d at Proof of Corollary 5.2.8.1, dimension bound identifies a particular indexed object or symbol; that identification is not independently established by the quoted local expression.\par
\par\noindent\textbf{OM-SOL-1056 -- t3 XVII p.360 -- SOL}\par
Last paragraph, reference to a component of the resolution\par
\emph{Reported reading}: une quelconque composante de L\_\{\textbackslash{}ast\}\par
\emph{Proposal}: une quelconque composante de \$L\textasciicircum{}\{\textbackslash{}ast\}\$\par
This finding intersects Orgogozo's \textbackslash{}sisi material at Last paragraph, reference to a component of the resolution; preserve that human entry unchanged. The proposed reading une quelconque composante de \$L\textasciicircum{}\{\textbackslash{}ast\}\$ at Last paragraph, reference to a component of the resolution changes a substantive mathematical assertion or bibliographic identification; no independent derivation of this exact replacement is established in this review.\par
\par\noindent\textbf{OM-SOL-1058 -- t3 XVII p.367 -- SOL}\par
Sentence following the diagram, morphism from L\_1\par
\emph{Reported reading}: L\_\{1\}\textbackslash{}longrightarrow g\_\{\textbackslash{}ast\}K\par
\emph{Proposal}: L\_\{1\}\textbackslash{}longrightarrow g\_\{\textbackslash{}ast\}L\par
The proposed reading L\_\{1\}\textbackslash{}longrightarrow g\_\{\textbackslash{}ast\}L at Sentence following the diagram, morphism from L\_1 changes a substantive mathematical assertion or bibliographic identification; no independent derivation of this exact replacement is established in this review.\par
\par\noindent\textbf{OM-SOL-1059 -- t3 XVII p.367 -- SOL}\par
5.4.2, final diagram, arrow labelled p'\par
\emph{Reported reading}: X'\textbackslash{}xrightarrow\{p'\}Z'\par
\emph{Proposal}: X'\textbackslash{}xleftarrow\{p'\}Z'\par
This finding intersects Orgogozo's \textbackslash{}sisi material at 5.4.2, final diagram, arrow labelled p'; preserve that human entry unchanged. The replacement X'\textbackslash{}xleftarrow\{p'\}Z' at 5.4.2, final diagram, arrow labelled p' requires a complete check of the maps, their endpoints and the composition convention; this identification is not independently verified.\par
\par\noindent\textbf{OM-SOL-1060 -- t3 XVII p.367 -- SOL}\par
5.4.2, final diagram, arrow labelled q'\par
\emph{Reported reading}: Z'\textbackslash{}xleftarrow\{q'\}Y'\par
\emph{Proposal}: Z'\textbackslash{}xrightarrow\{q'\}Y'\par
This finding intersects Orgogozo's \textbackslash{}sisi material at 5.4.2, final diagram, arrow labelled q'; preserve that human entry unchanged. The replacement Z'\textbackslash{}xrightarrow\{q'\}Y' at 5.4.2, final diagram, arrow labelled q' requires a complete check of the maps, their endpoints and the composition convention; this identification is not independently verified.\par
\par\noindent\textbf{OM-SOL-1061 -- t3 XVII p.367 -- SOL}\par
5.4.2, final diagram, arrow labelled bar p\par
\emph{Reported reading}: \textbackslash{}bar\{X\}'\textbackslash{}xrightarrow\{\textbackslash{}bar\{p\}\}\textbackslash{}bar\{Z\}'\par
\emph{Proposal}: \textbackslash{}bar\{X\}'\textbackslash{}xleftarrow\{\textbackslash{}bar\{p\}\}\textbackslash{}bar\{Z\}'\par
This finding intersects Orgogozo's \textbackslash{}sisi material at 5.4.2, final diagram, arrow labelled bar p; preserve that human entry unchanged. The replacement \textbackslash{}bar\{X\}'\textbackslash{}xleftarrow\{\textbackslash{}bar\{p\}\}\textbackslash{}bar\{Z\}' at 5.4.2, final diagram, arrow labelled bar p requires a complete check of the maps, their endpoints and the composition convention; this identification is not independently verified.\par
\par\noindent\textbf{OM-SOL-1062 -- t3 XVII p.367 -- SOL}\par
5.4.2, final diagram, arrow labelled bar q\par
\emph{Reported reading}: \textbackslash{}bar\{Z\}'\textbackslash{}xleftarrow\{\textbackslash{}bar\{q\}\}\textbackslash{}bar\{Y\}'\par
\emph{Proposal}: \textbackslash{}bar\{Z\}'\textbackslash{}xrightarrow\{\textbackslash{}bar\{q\}\}\textbackslash{}bar\{Y\}'\par
This finding intersects Orgogozo's \textbackslash{}sisi material at 5.4.2, final diagram, arrow labelled bar q; preserve that human entry unchanged. The replacement \textbackslash{}bar\{Z\}'\textbackslash{}xrightarrow\{\textbackslash{}bar\{q\}\}\textbackslash{}bar\{Y\}' at 5.4.2, final diagram, arrow labelled bar q requires a complete check of the maps, their endpoints and the composition convention; this identification is not independently verified.\par
\par\noindent\textbf{OM-SOL-1063 -- t3 XVII p.368 -- SOL}\par
(5.4.3.1), upper arrow labelled q\_1\par
\emph{Reported reading}: S\_1\textbackslash{}xrightarrow\{q\_1\}Y\_1\par
\emph{Proposal}: S\_1\textbackslash{}xleftarrow\{q\_1\}Y\_1\par
This finding intersects Orgogozo's \textbackslash{}sisi material at (5.4.3.1), upper arrow labelled q\_1; preserve that human entry unchanged. The replacement S\_1\textbackslash{}xleftarrow\{q\_1\}Y\_1 at (5.4.3.1), upper arrow labelled q\_1 requires a complete check of the maps, their endpoints and the composition convention; this identification is not independently verified.\par
\par\noindent\textbf{OM-SOL-1064 -- t3 XVII p.368 -- SOL}\par
(5.4.3.1), lower arrow labelled q\par
\emph{Reported reading}: S\textbackslash{}xrightarrow\{q\}Y\par
\emph{Proposal}: S\textbackslash{}xleftarrow\{q\}Y\par
This finding intersects Orgogozo's \textbackslash{}sisi material at (5.4.3.1), lower arrow labelled q; preserve that human entry unchanged. The replacement S\textbackslash{}xleftarrow\{q\}Y at (5.4.3.1), lower arrow labelled q requires a complete check of the maps, their endpoints and the composition convention; this identification is not independently verified.\par
\par\noindent\textbf{OM-SOL-1065 -- t3 XVII p.369 -- SOL}\par
Lemme 5.4.3.2, lower-right vertex\par
\emph{Reported reading}: Rh\_\{!\}(C\textasciicircum{}\{\textbackslash{}ast\}K\textbackslash{}overset\{L\}\{\textbackslash{}boxtimes\}\_\{\textbackslash{}Aa\}C\textasciicircum{}\{\textbackslash{}ast\}L)\par
\emph{Proposal}: Rh\_\{1!\}(C\textasciicircum{}\{\textbackslash{}ast\}K\textbackslash{}overset\{L\}\{\textbackslash{}boxtimes\}\_\{\textbackslash{}Aa\}C\textasciicircum{}\{\textbackslash{}ast\}L)\par
This finding intersects Orgogozo's \textbackslash{}sisi material at Lemme 5.4.3.2, lower-right vertex; preserve that human entry unchanged. The replacement Rh\_\{1!\}(C\textasciicircum{}\{\textbackslash{}ast\}K\textbackslash{}overset\{L\}\{\textbackslash{}boxtimes\}\_\{\textbackslash{}Aa\}C\textasciicircum{}\{\textbackslash{}ast\}L) at Lemme 5.4.3.2, lower-right vertex identifies a particular indexed object or symbol; that identification is not independently established by the quoted local expression.\par
\par\noindent\textbf{OM-SOL-1066 -- t3 XVII p.370 -- SOL}\par
Lemme 5.4.3.5, lower-right vertex\par
\emph{Reported reading}: Rf\_\{!\}Rp\_\{!\}(p\textasciicircum{}\{\textbackslash{}ast\}K\textbackslash{}otimes q\textasciicircum{}\{\textbackslash{}ast\}L)\par
\emph{Proposal}: Rf\_\{!\}Rp\_\{!\}(p\textasciicircum{}\{\textbackslash{}ast\}K\textbackslash{}overset\{\textbackslash{}LL\}\{\textbackslash{}otimes\}\_\{\textbackslash{}Aa\}q\textasciicircum{}\{\textbackslash{}ast\}L)\par
This finding intersects Orgogozo's \textbackslash{}sisi material at Lemme 5.4.3.5, lower-right vertex; preserve that human entry unchanged. The replacement Rf\_\{!\}Rp\_\{!\}(p\textasciicircum{}\{\textbackslash{}ast\}K\textbackslash{}overset\{\textbackslash{}LL\}\{\textbackslash{}otimes\}\_\{\textbackslash{}Aa\}q\textasciicircum{}\{\textbackslash{}ast\}L) at Lemme 5.4.3.5, lower-right vertex identifies a particular indexed object or symbol; that identification is not independently established by the quoted local expression.\par
\par\noindent\textbf{OM-SOL-1067 -- t3 XVII p.370 -- SOL}\par
Explanation of Lemme 5.4.3.5, “les flèches de Künneth”\par
\emph{Reported reading}: les flèches de Künneth\par
\emph{Proposal}: la flèche de Künneth\par
This finding intersects Orgogozo's \textbackslash{}sisi material at Explanation of Lemme 5.4.3.5, “les flèches de Künneth”; preserve that human entry unchanged. The replacement la flèche de Künneth at Explanation of Lemme 5.4.3.5, “les flèches de Künneth” requires a complete check of the maps, their endpoints and the composition convention; this identification is not independently verified.\par
\par\noindent\textbf{OM-SOL-1068 -- t3 XVII p.371 -- SOL}\par
5.5.1, annonce des propriétés dues à Roby\par
\emph{Reported reading}: \$TS\textasciicircum{}n\$ et \$T\textasciicircum{}n\$\par
\emph{Proposal}: TS\textasciicircum{}\{n\}\textbackslash{}text\{ et \}\textbackslash{}Gamma\textasciicircum{}\{n\}\par
The proposed reading TS\textasciicircum{}\{n\}\textbackslash{}text\{ et \}\textbackslash{}Gamma\textasciicircum{}\{n\} at 5.5.1, annonce des propriétés dues à Roby changes a substantive mathematical assertion or bibliographic identification; no independent derivation of this exact replacement is established in this review.\par
\par\noindent\textbf{OM-SOL-1069 -- t3 XVII p.372 -- SOL}\par
5.5.2.2, référence pour la définition d'une loi polynôme\par
\emph{Reported reading}: ROBY [1, I 1 pg.226]\par
\emph{Proposal}: \textbackslash{}cite[I 2 p.\textasciitilde{}219]\{SGA4-17:Ro\}\par
This finding intersects Orgogozo's \textbackslash{}sisi material at 5.5.2.2, référence pour la définition d'une loi polynôme; preserve that human entry unchanged. The proposed reference \textbackslash{}cite[I 2 p.\textasciitilde{}219]\{SGA4-17:Ro\} at 5.5.2.2, référence pour la définition d'une loi polynôme needs comparison with the statement actually invoked; that target has not been independently established here.\par
\par\noindent\textbf{OM-SOL-1070 -- t3 XVII p.373 -- SOL}\par
Définition d'une loi polynôme homogène, référence à Roby\par
\emph{Reported reading}: ROBY [1, I 1 pg.226]\par
\emph{Proposal}: \textbackslash{}cite[I 8 p.\textasciitilde{}226]\{SGA4-17:Ro\}\par
This finding intersects Orgogozo's \textbackslash{}sisi material at Définition d'une loi polynôme homogène, référence à Roby; preserve that human entry unchanged. The proposed reference \textbackslash{}cite[I 8 p.\textasciitilde{}226]\{SGA4-17:Ro\} at Définition d'une loi polynôme homogène, référence à Roby needs comparison with the statement actually invoked; that target has not been independently established here.\par
\par\noindent\textbf{OM-SOL-1071 -- t3 XVII p.374 -- SOL}\par
Début de page, référence pour l'application universelle \ensuremath{\gamma}\textasciicircum{}n\par
\emph{Reported reading}: ROBY [1, IV 1 pg.226]\par
\emph{Proposal}: \textbackslash{}cite[IV 1 p.\textasciitilde{}265]\{SGA4-17:Ro\}\par
This finding intersects Orgogozo's \textbackslash{}sisi material at Début de page, référence pour l'application universelle \ensuremath{\gamma}\textasciicircum{}n; preserve that human entry unchanged. The proposed reference \textbackslash{}cite[IV 1 p.\textasciitilde{}265]\{SGA4-17:Ro\} at Début de page, référence pour l'application universelle \ensuremath{\gamma}\textasciicircum{}n needs comparison with the statement actually invoked; that target has not been independently established here.\par
\par\noindent\textbf{OM-SOL-1072 -- t3 XVII p.375 -- SOL}\par
5.5.2.6, application précédant l'isomorphisme (5.5.2.6.1)\par
\emph{Reported reading}: m\textbackslash{}longmapsto (\textbackslash{}gamma\textasciicircum{}\{n'\}m\textbackslash{}otimes\textbackslash{}gamma\textasciicircum{}\{n''\}m)\_\{n'+n''=\allowbreak n\}\par
\emph{Proposal}: m=\allowbreak m'+m''\textbackslash{}longmapsto (\textbackslash{}gamma\textasciicircum{}\{n'\}m'\textbackslash{}otimes\textbackslash{}gamma\textasciicircum{}\{n''\}m'')\_\{n'+n''=\allowbreak n\}\par
The proposed reading m=\allowbreak m'+m''\textbackslash{}longmapsto (\textbackslash{}gamma\textasciicircum{}\{n'\}m'\textbackslash{}otimes\textbackslash{}gamma\textasciicircum{}\{n''\}m'')\_\{n'+n''=\allowbreak n\} at 5.5.2.6, application précédant l'isomorphisme (5.5.2.6.1) changes a substantive mathematical assertion or bibliographic identification; no independent derivation of this exact replacement is established in this review.\par
\par\noindent\textbf{OM-SOL-1073 -- t3 XVII p.375 -- SOL}\par
5.5.2.6, référence pour l'isomorphisme canonique\par
\emph{Reported reading}: ROBY [1, III 4]\par
\emph{Proposal}: \textbackslash{}cite[III 9 p.\textasciitilde{}262]\{SGA4-17:Ro\}\par
This finding intersects Orgogozo's \textbackslash{}sisi material at 5.5.2.6, référence pour l'isomorphisme canonique; preserve that human entry unchanged. The proposed reference \textbackslash{}cite[III 9 p.\textasciitilde{}262]\{SGA4-17:Ro\} at 5.5.2.6, référence pour l'isomorphisme canonique needs comparison with the statement actually invoked; that target has not been independently established here.\par
\par\noindent\textbf{OM-SOL-1074 -- t3 XVII p.375 -- SOL}\par
5.5.2.6, conclusion après les morphismes (5.5.2.6.2)\par
\emph{Reported reading}: donc aussi pour M' plat\par
\emph{Proposal}: \textbackslash{}text\{donc aussi pour \}M''\textbackslash{}text\{ plat\}\par
The proposed reading \textbackslash{}text\{donc aussi pour \}M''\textbackslash{}text\{ plat\} at 5.5.2.6, conclusion après les morphismes (5.5.2.6.2) changes a substantive mathematical assertion or bibliographic identification; no independent derivation of this exact replacement is established in this review.\par
\par\noindent\textbf{OM-SOL-1075 -- t3 XVII p.375 -- SOL}\par
5.5.2.7, référence pour l'extension des scalaires\par
\emph{Reported reading}: ROBY [1, III 3 pg.262]\par
\emph{Proposal}: \textbackslash{}cite[III 8 p.\textasciitilde{}262]\{SGA4-17:Ro\}\par
This finding intersects Orgogozo's \textbackslash{}sisi material at 5.5.2.7, référence pour l'extension des scalaires; preserve that human entry unchanged. The proposed reference \textbackslash{}cite[III 8 p.\textasciitilde{}262]\{SGA4-17:Ro\} at 5.5.2.7, référence pour l'extension des scalaires needs comparison with the statement actually invoked; that target has not been independently established here.\par
\par\noindent\textbf{OM-SOL-1076 -- t3 XVII p.377 -- SOL}\par
5.5.3.4, final Dold–Puppe reference\par
\emph{Reported reading}: 4.2.3 pg.233\par
\emph{Proposal}: 4.23\par
This finding intersects Orgogozo's \textbackslash{}sisi material at 5.5.3.4, final Dold–Puppe reference; preserve that human entry unchanged. The proposed reference 4.23 at 5.5.3.4, final Dold–Puppe reference needs comparison with the statement actually invoked; that target has not been independently established here.\par
\par\noindent\textbf{OM-SOL-1077 -- t3 XVII p.378 -- SOL}\par
First paragraph, target of the inclusion\par
\emph{Reported reading}: l'inclusion de K\textasciicircum{}\{\textbackslash{}geq 0\}(\textbackslash{}Aa) dans \textbackslash{}Aa\par
\emph{Proposal}: K(\textbackslash{}Aa)\par
This finding intersects Orgogozo's \textbackslash{}sisi material at First paragraph, target of the inclusion; preserve that human entry unchanged. The proposed reading K(\textbackslash{}Aa) at First paragraph, target of the inclusion changes a substantive mathematical assertion or bibliographic identification; no independent derivation of this exact replacement is established in this review.\par
\par\noindent\textbf{OM-SOL-1079 -- t3 XVII p.380 -- SOL}\par
Opening paragraph, category containing N\par
\emph{Reported reading}: N\textasciicircum{}\{\textbackslash{}bigdot\}\textbackslash{}in B\textasciicircum{}\{b,\textbackslash{}geq 0\}(A)\par
\emph{Proposal}: N\textasciicircum{}\{\textbackslash{}bigdot\}\textbackslash{}in C\textasciicircum{}\{b,\textbackslash{}geq 0\}(A)\par
This finding intersects Orgogozo's \textbackslash{}sisi material at Opening paragraph, category containing N; preserve that human entry unchanged. The proposed reading N\textasciicircum{}\{\textbackslash{}bigdot\}\textbackslash{}in C\textasciicircum{}\{b,\textbackslash{}geq 0\}(A) at Opening paragraph, category containing N changes a substantive mathematical assertion or bibliographic identification; no independent derivation of this exact replacement is established in this review.\par
\par\noindent\textbf{OM-SOL-1081 -- t3 XVII p.381 -- SOL}\par
C), but de l'application H\textasciicircum{}i\_\ensuremath{\alpha}\par
\emph{Reported reading}: L\textasciicircum{}\{i\}\_\{\textbackslash{}beta(\textbackslash{}alpha)\}\par
\emph{Proposal}: L\textasciicircum{}\{i-1\}\_\{\textbackslash{}beta(\textbackslash{}alpha)\}\par
This finding intersects Orgogozo's \textbackslash{}sisi material at C), but de l'application H\textasciicircum{}i\_\ensuremath{\alpha}; preserve that human entry unchanged. The replacement L\textasciicircum{}\{i-1\}\_\{\textbackslash{}beta(\textbackslash{}alpha)\} at C), but de l'application H\textasciicircum{}i\_\ensuremath{\alpha} identifies a particular indexed object or symbol; that identification is not independently established by the quoted local expression.\par
\par\noindent\textbf{OM-SOL-1082 -- t3 XVII p.383 -- SOL}\par
5.5.7.1, « Un point x de… »\par
\emph{Reported reading}: \textbackslash{}Sym\_\{S\}(X)\par
\emph{Proposal}: \textbackslash{}Sym\textasciicircum{}\{k\}\_\{S\}(X)\par
The replacement \textbackslash{}Sym\textasciicircum{}\{k\}\_\{S\}(X) at 5.5.7.1, « Un point x de… » identifies a particular indexed object or symbol; that identification is not independently established by the quoted local expression.\par
\par\noindent\textbf{OM-SOL-1083 -- t3 XVII p.383 -- SOL}\par
5.5.7.1, « définit un point géométrique de… »\par
\emph{Reported reading}: \textbackslash{}Sym\_\{S\}(X)\par
\emph{Proposal}: \textbackslash{}Sym\textasciicircum{}\{k\}\_\{S\}(X)\par
The replacement \textbackslash{}Sym\textasciicircum{}\{k\}\_\{S\}(X) at 5.5.7.1, « définit un point géométrique de… » identifies a particular indexed object or symbol; that identification is not independently established by the quoted local expression.\par
\par\noindent\textbf{OM-SOL-1084 -- t3 XVII p.383 -- SOL}\par
Après (5.5.7.2), anneau sur lequel sont pris les produits tensoriels\par
\emph{Reported reading}: \textbackslash{}Aa\_\{S\}\par
\emph{Proposal}: \textbackslash{}Aa\_\{s\}\par
The proposed reading \textbackslash{}Aa\_\{s\} involves a printed glyph or a missing stroke at Après (5.5.7.2), anneau sur lequel sont pris les produits tensoriels; the textual record does not independently establish that glyph.\par
\par\noindent\textbf{OM-SOL-1085 -- t3 XVII p.386 -- SOL}\par
5.5.11, unnumbered stalkwise morphism, second tensor product\par
\emph{Reported reading}: \textbackslash{}Gamma\textasciicircum{}\{n\_i-i\_\textbackslash{}alpha\}(F''\_\{y\_\textbackslash{}alpha\})\par
\emph{Proposal}: \textbackslash{}Gamma\textasciicircum{}\{n\_\{\textbackslash{}alpha\}-i\_\{\textbackslash{}alpha\}\}(F''\_\{y\_\{\textbackslash{}alpha\}\})\par
The replacement \textbackslash{}Gamma\textasciicircum{}\{n\_\{\textbackslash{}alpha\}-i\_\{\textbackslash{}alpha\}\}(F''\_\{y\_\{\textbackslash{}alpha\}\}) at 5.5.11, unnumbered stalkwise morphism, second tensor product identifies a particular indexed object or symbol; that identification is not independently established by the quoted local expression.\par
\par\noindent\textbf{OM-SOL-1086 -- t3 XVII p.386 -- SOL}\par
5.5.11, reference immediately before (5.5.11.2)\par
\emph{Reported reading}: 5.5.2.6.2\par
\emph{Proposal}: \textbackslash{}eqhyperref[XVII.eq5.5.2.6.1]\{5.5.2.6.1\}\par
This finding intersects Orgogozo's \textbackslash{}sisi material at 5.5.11, reference immediately before (5.5.11.2); preserve that human entry unchanged. The proposed reference \textbackslash{}eqhyperref[XVII.eq5.5.2.6.1]\{5.5.2.6.1\} at 5.5.11, reference immediately before (5.5.11.2) needs comparison with the statement actually invoked; that target has not been independently established here.\par
\par\noindent\textbf{OM-SOL-1087 -- t3 XVII p.386 -- SOL}\par
5.5.12, adjunction isomorphism between the Hom groups\par
\emph{Reported reading}: \textbackslash{}underset\{n\}\{\textbackslash{}simeq\}\par
\emph{Proposal}: \textbackslash{}Hom(u\textasciicircum{}\{\textbackslash{}ast\}G,F)\textbackslash{}simeq\textbackslash{}Hom(G,u\_\{\textbackslash{}ast\}F)\par
This finding intersects Orgogozo's \textbackslash{}sisi material at 5.5.12, adjunction isomorphism between the Hom groups; preserve that human entry unchanged. The replacement \textbackslash{}Hom(u\textasciicircum{}\{\textbackslash{}ast\}G,F)\textbackslash{}simeq\textbackslash{}Hom(G,u\_\{\textbackslash{}ast\}F) at 5.5.12, adjunction isomorphism between the Hom groups requires a complete check of the maps, their endpoints and the composition convention; this identification is not independently verified.\par
\par\noindent\textbf{OM-SOL-1091 -- t3 XVII p.391 -- SOL}\par
Condition (ii), morphisme de topos f\par
\emph{Reported reading}: relatif à un isomorphisme de topos\par
\emph{Proposal}: relatif à un morphisme de topos\par
Changing isomorphisme to morphisme enlarges the hypothesis materially and requires the proof of the Künneth criterion.\par
\par\noindent\textbf{OM-SOL-1092 -- t3 XVII p.392 -- SOL}\par
(5.5.17.1), terme immédiatement après la flèche \ensuremath{\gamma}\par
\emph{Reported reading}: L\textbackslash{}Gamma\textasciicircum{}\{n\}u\_\{\textbackslash{}ast\}K\par
\emph{Proposal}: L\textbackslash{}Gamma\textasciicircum{}\{n\}\_\{\textbackslash{}ext\}u\_\{\textbackslash{}ast\}K\par
The replacement L\textbackslash{}Gamma\textasciicircum{}\{n\}\_\{\textbackslash{}ext\}u\_\{\textbackslash{}ast\}K at (5.5.17.1), terme immédiatement après la flèche \ensuremath{\gamma} identifies a particular indexed object or symbol; that identification is not independently established by the quoted local expression.\par
\par\noindent\textbf{OM-SOL-1093 -- t3 XVII p.393 -- SOL}\par
Proposition 5.5.18, catégorie contenant le diagramme\par
\emph{Reported reading}: D\textasciicircum{}\{+\}(\textbackslash{}Sym\_\{S\}(Y),\textbackslash{}Aa)\par
\emph{Proposal}: D\textasciicircum{}\{+\}(\textbackslash{}Sym\textasciicircum{}\{n\}\_\{S\}(Y),\textbackslash{}Aa)\par
The replacement D\textasciicircum{}\{+\}(\textbackslash{}Sym\textasciicircum{}\{n\}\_\{S\}(Y),\textbackslash{}Aa) at Proposition 5.5.18, catégorie contenant le diagramme identifies a particular indexed object or symbol; that identification is not independently established by the quoted local expression.\par
\par\noindent\textbf{OM-SOL-1094 -- t3 XVII p.394 -- SOL}\par
(5.5.19.2), membre de droite\par
\emph{Reported reading}: R\textbackslash{}Sym\_\{S\}(u)\_\{!\}L\textbackslash{}Gamma\textasciicircum{}\{n\}\_\{\textbackslash{}ext\}\par
\emph{Proposal}: R\textbackslash{}Sym\textasciicircum{}\{n\}\_\{S\}(u)\_\{!\}L\textbackslash{}Gamma\textasciicircum{}\{n\}\_\{\textbackslash{}ext\}\par
The replacement R\textbackslash{}Sym\textasciicircum{}\{n\}\_\{S\}(u)\_\{!\}L\textbackslash{}Gamma\textasciicircum{}\{n\}\_\{\textbackslash{}ext\} at (5.5.19.2), membre de droite identifies a particular indexed object or symbol; that identification is not independently established by the quoted local expression.\par
\par\noindent\textbf{OM-SOL-1095 -- t3 XVII p.394 -- SOL}\par
Renvoi placé sur la dernière flèche du composé définissant (5.5.19.2)\par
\emph{Reported reading}: (5.5.13.5 (vii''))\par
\emph{Proposal}: R\textbackslash{}Sym\textasciicircum{}\{n\}\_\{S\}(\textbackslash{}bar\{u\})\_\{\textbackslash{}ast\}L\textbackslash{}Gamma\textasciicircum{}\{n\}\_\{\textbackslash{}ext\}j\_\{!\}\textbackslash{}xrightarrow[\textbackslash{}text\{(5.5.14.4 (vii'))\}]\{\textbackslash{}sim\}R\textbackslash{}Sym\textasciicircum{}\{n\}\_\{S\}(\textbackslash{}bar\{u\})\_\{\textbackslash{}ast\}\textbackslash{}Sym\textasciicircum{}\{n\}\_\{S\}(j)\_\{!\}L\textbackslash{}Gamma\textasciicircum{}\{n\}\_\{\textbackslash{}ext\}\textbackslash{}longrightarrow R\textbackslash{}Sym\textasciicircum{}\{n\}\_\{S\}(u)\_\{!\}L\textbackslash{}Gamma\textasciicircum{}\{n\}\_\{\textbackslash{}ext\}\par
This finding intersects Orgogozo's \textbackslash{}sisi material at Renvoi placé sur la dernière flèche du composé définissant (5.5.19.2); preserve that human entry unchanged. The proposed reference R\textbackslash{}Sym\textasciicircum{}\{n\}\_\{S\}(\textbackslash{}bar\{u\})\_\{\textbackslash{}ast\}L\textbackslash{}Gamma\textasciicircum{}\{n\}\_\{\textbackslash{}ext\}j\_\{!\}\textbackslash{}xrightarrow[\textbackslash{}text\{(5.5.14.4 (vii'))\}]\{\textbackslash{}sim\}R\textbackslash{}Sym\textasciicircum{}\{n\}\_\{S\}(\textbackslash{}bar\{u\})\_\{\textbackslash{}ast\}\textbackslash{}Sym\textasciicircum{}\{n\}\_\{S\}(j)\_\{!\}L\textbackslash{}Gamma\textasciicircum{}\{n\}\_\{\textbackslash{}ext\}\textbackslash{}longrightarrow R\textbackslash{}Sym\textasciicircum{}\{n\}\_\{S\}(u)\_\{!\}L\textbackslash{}Gamma\textasciicircum{}\{n\}\_\{\textbackslash{}ext\} at Renvoi placé sur la dernière flèche du composé définissant (5.5.19.2) needs comparison with the statement actually invoked; that target has not been independently established here.\par
\par\noindent\textbf{OM-SOL-1096 -- t3 XVII p.395 -- SOL}\par
5.5.20.3, renvoi aux morphismes (iv'), (v'), (vi')\par
\emph{Reported reading}: 5.5.13.5 (iv') (v') (vi')\par
\emph{Proposal}: \textbackslash{}hyperref[XVII.par5.5.14.4]\{5.5.14.4\}\par
This finding intersects Orgogozo's \textbackslash{}sisi material at 5.5.20.3, renvoi aux morphismes (iv'), (v'), (vi'); preserve that human entry unchanged. The proposed reference \textbackslash{}hyperref[XVII.par5.5.14.4]\{5.5.14.4\} at 5.5.20.3, renvoi aux morphismes (iv'), (v'), (vi') needs comparison with the statement actually invoked; that target has not been independently established here.\par
\par\noindent\textbf{OM-SOL-1097 -- t3 XVII p.396 -- SOL}\par
(v''), diagramme de foncteurs, les deux termes inférieurs de droite\par
\emph{Reported reading}: Les deux images directes sont écrites R\textbackslash{}Sym\textasciicircum{}\{n\}\_\{S\}(u')\_\{!\}.\par
\emph{Proposal}: R\textbackslash{}Sym\textasciicircum{}\{n\}\_\{S'\}(u')\_\{!\}\par
The replacement R\textbackslash{}Sym\textasciicircum{}\{n\}\_\{S'\}(u')\_\{!\} at (v''), diagramme de foncteurs, les deux termes inférieurs de droite identifies a particular indexed object or symbol; that identification is not independently established by the quoted local expression.\par
\par\noindent\textbf{OM-SOL-1098 -- t3 XVII p.397 -- SOL}\par
(vi''), diagramme supérieur, terme inférieur gauche\par
\emph{Reported reading}: \textbackslash{}Sym\textasciicircum{}\{n\}\_\{S\}(g)\textasciicircum{}\{\textbackslash{}ast\}R\textbackslash{}Sym\textasciicircum{}\{n\}\_\{S\}(u\_!)\par
\emph{Proposal}: \textbackslash{}Sym\textasciicircum{}\{n\}\_\{S\}(g)\textasciicircum{}\{\textbackslash{}ast\}R\textbackslash{}Sym\textasciicircum{}\{n\}\_\{S\}(u)\_\{!\}L\textbackslash{}Gamma\textasciicircum{}\{n\}\_\{\textbackslash{}ext\}\par
This finding intersects Orgogozo's \textbackslash{}sisi material at (vi''), diagramme supérieur, terme inférieur gauche; preserve that human entry unchanged. The proposed reading \textbackslash{}Sym\textasciicircum{}\{n\}\_\{S\}(g)\textasciicircum{}\{\textbackslash{}ast\}R\textbackslash{}Sym\textasciicircum{}\{n\}\_\{S\}(u)\_\{!\}L\textbackslash{}Gamma\textasciicircum{}\{n\}\_\{\textbackslash{}ext\} at (vi''), diagramme supérieur, terme inférieur gauche changes a substantive mathematical assertion or bibliographic identification; no independent derivation of this exact replacement is established in this review.\par
\par\noindent\textbf{OM-SOL-1099 -- t3 XVII p.397 -- SOL}\par
(vi''), diagramme supérieur, terme inférieur central\par
\emph{Reported reading}: R\textbackslash{}Sym\textasciicircum{}\{n\}\_\{S\}(u')\_\{!\}g'\textasciicircum{}\{\textbackslash{}ast\}L\textbackslash{}Gamma\textasciicircum{}\{n\}\_\{\textbackslash{}ext\}\par
\emph{Proposal}: R\textbackslash{}Sym\textasciicircum{}\{n\}\_\{S\}(u')\_\{!\}\textbackslash{}Sym\textasciicircum{}\{n\}\_\{S\}(g')\textasciicircum{}\{\textbackslash{}ast\}L\textbackslash{}Gamma\textasciicircum{}\{n\}\_\{\textbackslash{}ext\}\par
The proposed reading R\textbackslash{}Sym\textasciicircum{}\{n\}\_\{S\}(u')\_\{!\}\textbackslash{}Sym\textasciicircum{}\{n\}\_\{S\}(g')\textasciicircum{}\{\textbackslash{}ast\}L\textbackslash{}Gamma\textasciicircum{}\{n\}\_\{\textbackslash{}ext\} at (vi''), diagramme supérieur, terme inférieur central changes a substantive mathematical assertion or bibliographic identification; no independent derivation of this exact replacement is established in this review.\par
\par\noindent\textbf{OM-SOL-1102 -- t3 XVII p.399 -- SOL}\par
Premier paragraphe, choix de N\par
\emph{Reported reading}: au double de la dimension des flèches de u\par
\emph{Proposal}: dimension des fibres de \$u\$\par
The replacement dimension des fibres de \$u\$ at Premier paragraphe, choix de N requires a complete check of the maps, their endpoints and the composition convention; this identification is not independently verified.\par
\par\noindent\textbf{OM-SOL-1103 -- t3 XVII p.399 -- SOL}\par
« Pour le voir, on note que grâce à … »\par
\emph{Reported reading}: grâce à 5.5.14\par
\emph{Proposal}: \textbackslash{}hyperref[XVII.subsubsec5.5.15]\{5.5.15\}\par
This finding intersects Orgogozo's \textbackslash{}sisi material at « Pour le voir, on note que grâce à … »; preserve that human entry unchanged. The proposed reference \textbackslash{}hyperref[XVII.subsubsec5.5.15]\{5.5.15\} at « Pour le voir, on note que grâce à … » needs comparison with the statement actually invoked; that target has not been independently established here.\par
\par\noindent\textbf{OM-SOL-1104 -- t3 XVII p.399 -- SOL}\par
Formule (b\_i), membre de droite, signe entre les deux puissances divisées\par
\emph{Reported reading}: L\textbackslash{}Gamma\textasciicircum{}\{i\}\_\{\textbackslash{}ext\}K'\textbackslash{}vee L\textbackslash{}Gamma\textasciicircum{}\{n-i\}\_\{\textbackslash{}ext\}K''\par
\emph{Proposal}: R\textbackslash{}Sym\textasciicircum{}\{n\}(u)\_\{\textbackslash{}ast\}(L\textbackslash{}Gamma\textasciicircum{}\{i\}\_\{\textbackslash{}ext\}K'\textbackslash{}overset\{\textbackslash{}LL\}\{\textbackslash{}vee\}L\textbackslash{}Gamma\textasciicircum{}\{n-i\}\_\{\textbackslash{}ext\}K'')\par
This finding intersects Orgogozo's \textbackslash{}sisi material at Formule (b\_i), membre de droite, signe entre les deux puissances divisées; preserve that human entry unchanged. The proposed reading R\textbackslash{}Sym\textasciicircum{}\{n\}(u)\_\{\textbackslash{}ast\}(L\textbackslash{}Gamma\textasciicircum{}\{i\}\_\{\textbackslash{}ext\}K'\textbackslash{}overset\{\textbackslash{}LL\}\{\textbackslash{}vee\}L\textbackslash{}Gamma\textasciicircum{}\{n-i\}\_\{\textbackslash{}ext\}K'') at Formule (b\_i), membre de droite, signe entre les deux puissances divisées changes a substantive mathematical assertion or bibliographic identification; no independent derivation of this exact replacement is established in this review.\par
\par\noindent\textbf{OM-SOL-1105 -- t3 XVII p.399 -- SOL}\par
Développement des flèches (b\_i), terme obtenu par 5.4.3\par
\emph{Reported reading}: (R\textbackslash{}Sym\textasciicircum{}\{i\}(u)\_\{\textbackslash{}ast\}\textbackslash{}times R\textbackslash{}Sym\textasciicircum{}\{n-i\}(u)\_\{\textbackslash{}ast\})\par
\emph{Proposal}: R(\textbackslash{}Sym\textasciicircum{}\{i\}(u)\textbackslash{}times\textbackslash{}Sym\textasciicircum{}\{n-i\}(u))\_\{\textbackslash{}ast\}\par
The replacement R(\textbackslash{}Sym\textasciicircum{}\{i\}(u)\textbackslash{}times\textbackslash{}Sym\textasciicircum{}\{n-i\}(u))\_\{\textbackslash{}ast\} at Développement des flèches (b\_i), terme obtenu par 5.4.3 requires a complete check of the maps, their endpoints and the composition convention; this identification is not independently verified.\par
\par\noindent\textbf{OM-SOL-1106 -- t3 XVII p.400 -- SOL}\par
5.5.22.4, première phrase\par
\emph{Reported reading}: ses restrictions aux flèches géométriques de \textbackslash{}Sym\textasciicircum{}\{n\}\_\{S\}(Y)\par
\emph{Proposal}: aux fibres géométriques de \$\textbackslash{}Sym\textasciicircum{}\{n\}\_\{S\}(Y)\$\par
The replacement aux fibres géométriques de \$\textbackslash{}Sym\textasciicircum{}\{n\}\_\{S\}(Y)\$ at 5.5.22.4, première phrase requires a complete check of the maps, their endpoints and the composition convention; this identification is not independently verified.\par
\par\noindent\textbf{OM-SOL-1109 -- t3 XVII p.406 -- SOL}\par
First paragraph, geometric generic fibre\par
\emph{Reported reading}: \$X\_\{\textbackslash{}bar\{\textbackslash{}eta\}\}\$ de \$X\$ sur \$S\$\par
\emph{Proposal}: \$X'\_\{\textbackslash{}bar\{\textbackslash{}eta\}\}\$ de \$X'\$ sur \$S\$\par
The replacement \$X'\_\{\textbackslash{}bar\{\textbackslash{}eta\}\}\$ de \$X'\$ sur \$S\$ at First paragraph, geometric generic fibre identifies a particular indexed object or symbol; that identification is not independently established by the quoted local expression.\par
\par\noindent\textbf{OM-SOL-1110 -- t3 XVII p.407 -- SOL}\par
Diagram (5.5.27.1), upper horizontal arrow\par
\emph{Reported reading}: k\textasciicircum{}\{k\}\par
\emph{Proposal}: k\textasciicircum{}\{n\}\par
This finding intersects Orgogozo's \textbackslash{}sisi material at Diagram (5.5.27.1), upper horizontal arrow; preserve that human entry unchanged. The replacement k\textasciicircum{}\{n\} at Diagram (5.5.27.1), upper horizontal arrow identifies a particular indexed object or symbol; that identification is not independently established by the quoted local expression.\par
\par\noindent\textbf{OM-SOL-1112 -- t3 XVII p.408 -- SOL}\par
Conclusion of 5.5.28, vanishing of cohomology\par
\emph{Reported reading}: Les \$H\textasciicircum{}\{i\}\$ des deux membres de 5.5.28.1 sont donc nuls\par
\emph{Proposal}: Les \$H\textasciicircum{}\{i\}\$ (\$i\textbackslash{}ne 0\$) des deux membres de \textbackslash{}eqhyperref[XVII.eq5.5.28.1]\{5.5.28.1\} sont donc nuls\par
The proposed reading Les \$H\textasciicircum{}\{i\}\$ (\$i\textbackslash{}ne 0\$) des deux membres de \textbackslash{}eqhyperref[XVII.eq5.5.28.1]\{5.5.28.1\} sont donc nuls at Conclusion of 5.5.28, vanishing of cohomology changes a substantive mathematical assertion or bibliographic identification; no independent derivation of this exact replacement is established in this review.\par
\par\noindent\textbf{OM-SOL-1113 -- t3 XVII p.408 -- SOL}\par
5.5.29, definition of p\par
\emph{Reported reading}: p:X\textbackslash{}longrightarrow X\par
\emph{Proposal}: p:X\textbackslash{}rightarrow S\par
This finding intersects Orgogozo's \textbackslash{}sisi material at 5.5.29, definition of p; preserve that human entry unchanged. The proposed reading p:X\textbackslash{}rightarrow S at 5.5.29, definition of p changes a substantive mathematical assertion or bibliographic identification; no independent derivation of this exact replacement is established in this review.\par
\par\noindent\textbf{OM-SOL-1114 -- t3 XVII p.409 -- SOL}\par
Formula (5.5.29.1), index on the symmetric-group sign\par
\emph{Reported reading}: (\textbackslash{}sigma\_\{\textbackslash{}ast\}\textbackslash{}bigotimes\textasciicircum{}\{n\}\textbackslash{}underline\{F\})\textasciicircum{}\{\textbackslash{}Sf\textasciicircum{}\{n\}\}\par
\emph{Proposal}: (\textbackslash{}sigma\_\{\textbackslash{}ast\}\textbackslash{}bigotimes\textasciicircum{}\{n\}\ensuremath{\mathcal{F}})\textasciicircum{}\{\textbackslash{}Sf\_\{n\}\}\par
The replacement (\textbackslash{}sigma\_\{\textbackslash{}ast\}\textbackslash{}bigotimes\textasciicircum{}\{n\}\ensuremath{\mathcal{F}})\textasciicircum{}\{\textbackslash{}Sf\_\{n\}\} at Formula (5.5.29.1), index on the symmetric-group sign identifies a particular indexed object or symbol; that identification is not independently established by the quoted local expression.\par
\par\noindent\textbf{OM-SOL-1115 -- t3 XVII p.409 -- SOL}\par
Exact sequence following the flat presentation, last term\par
\emph{Reported reading}: \textbackslash{}Gamma\textasciicircum{}\{n\}\_\{\textbackslash{}ext\}(F)\par
\emph{Proposal}: \textbackslash{}Gamma\textasciicircum{}\{n\}\_\{\textbackslash{}ext\}(\ensuremath{\mathcal{F}})\par
This finding intersects Orgogozo's \textbackslash{}sisi material at Exact sequence following the flat presentation, last term; preserve that human entry unchanged. The replacement \textbackslash{}Gamma\textasciicircum{}\{n\}\_\{\textbackslash{}ext\}(\ensuremath{\mathcal{F}}) at Exact sequence following the flat presentation, last term identifies a particular indexed object or symbol; that identification is not independently established by the quoted local expression.\par
\par\noindent\textbf{OM-SOL-1116 -- t3 XVII p.409 -- SOL}\par
Formula (5.5.29.2), argument of the divided power on the right\par
\emph{Reported reading}: \textbackslash{}Gamma\_\{\textbackslash{}Oo\_\{S\}\}\textasciicircum{}\{n\_\{i\}\}(F\_\{y\_\{i\}\})\par
\emph{Proposal}: \textbackslash{}Gamma\_\{\textbackslash{}Oo\_\{S\}\}\textasciicircum{}\{n\_\{i\}\}(\ensuremath{\mathcal{F}}\_\{y\_\{i\}\})\par
This finding intersects Orgogozo's \textbackslash{}sisi material at Formula (5.5.29.2), argument of the divided power on the right; preserve that human entry unchanged. The replacement \textbackslash{}Gamma\_\{\textbackslash{}Oo\_\{S\}\}\textasciicircum{}\{n\_\{i\}\}(\ensuremath{\mathcal{F}}\_\{y\_\{i\}\}) at Formula (5.5.29.2), argument of the divided power on the right identifies a particular indexed object or symbol; that identification is not independently established by the quoted local expression.\par
\par\noindent\textbf{OM-SOL-1117 -- t3 XVII p.409 -- SOL}\par
Sentence following (5.5.29.2), free stalks\par
\emph{Reported reading}: les fibres de \$F\$ sont libres\par
\emph{Proposal}: les fibres de \$\ensuremath{\mathcal{F}}\$ sont libres\par
This finding intersects Orgogozo's \textbackslash{}sisi material at Sentence following (5.5.29.2), free stalks; preserve that human entry unchanged. The replacement les fibres de \$\ensuremath{\mathcal{F}}\$ sont libres at Sentence following (5.5.29.2), free stalks identifies a particular indexed object or symbol; that identification is not independently established by the quoted local expression.\par
\par\noindent\textbf{OM-SOL-1119 -- t3 XVII p.410 -- SOL}\par
Condition (5.5.31.1), coefficient sheaf\par
\emph{Reported reading}: C\textasciicircum{}\{b,\textbackslash{}geq 0\}(X,O\_\{X\})\par
\emph{Proposal}: C\textasciicircum{}\{b,\textbackslash{}geq 0\}(X,\textbackslash{}Oo\_\{X\})\par
This finding intersects Orgogozo's \textbackslash{}sisi material at Condition (5.5.31.1), coefficient sheaf; preserve that human entry unchanged. The proposed reading C\textasciicircum{}\{b,\textbackslash{}geq 0\}(X,\textbackslash{}Oo\_\{X\}) at Condition (5.5.31.1), coefficient sheaf changes a substantive mathematical assertion or bibliographic identification; no independent derivation of this exact replacement is established in this review.\par
\par\noindent\textbf{OM-SOL-1120 -- t3 XVII p.412 -- SOL}\par
Preuve de 5.5.34, c), première phrase\par
\emph{Reported reading}: modules quasi-cohérents sur S\par
\emph{Proposal}: modules quasi-cohérents sur \$X\$\par
The proposed reading modules quasi-cohérents sur \$X\$ at Preuve de 5.5.34, c), première phrase changes a substantive mathematical assertion or bibliographic identification; no independent derivation of this exact replacement is established in this review.\par
\par\noindent\textbf{OM-SOL-1121 -- t3 XVII p.412 -- SOL}\par
Preuve de 5.5.34, d), composantes de K₂\par
\emph{Reported reading}: à composantes plates et quasi-cohérentes\par
\emph{Proposal}: à composantes \$S\$-plates et quasi-cohérentes\par
This finding intersects Orgogozo's \textbackslash{}sisi material at Preuve de 5.5.34, d), composantes de K₂; preserve that human entry unchanged. The proposed reading à composantes \$S\$-plates et quasi-cohérentes at Preuve de 5.5.34, d), composantes de K₂ changes a substantive mathematical assertion or bibliographic identification; no independent derivation of this exact replacement is established in this review.\par
\par\noindent\textbf{OM-SOL-1122 -- t3 XVII p.412 -- SOL}\par
Preuve de 5.5.34, d), renvoi après « quasi-cohérentes »\par
\emph{Reported reading}: Parenthèses laissées vides.\par
\emph{Proposal}: (\textbackslash{}hyperref[XVII.def4.1.9]\{4.1.9\})\par
This finding intersects Orgogozo's \textbackslash{}sisi material at Preuve de 5.5.34, d), renvoi après « quasi-cohérentes »; preserve that human entry unchanged. The proposed reference (\textbackslash{}hyperref[XVII.def4.1.9]\{4.1.9\}) at Preuve de 5.5.34, d), renvoi après « quasi-cohérentes » needs comparison with the statement actually invoked; that target has not been independently established here.\par
\par\noindent\textbf{OM-SOL-1123 -- t3 XVII p.413 -- SOL}\par
Suite exacte d'Artin-Schreier, dernier terme\par
\emph{Reported reading}: \textbackslash{}Oo\par
\emph{Proposal}: 0\par
This finding intersects Orgogozo's \textbackslash{}sisi material at Suite exacte d'Artin-Schreier, dernier terme; preserve that human entry unchanged. The proposed reading 0 at Suite exacte d'Artin-Schreier, dernier terme changes a substantive mathematical assertion or bibliographic identification; no independent derivation of this exact replacement is established in this review.\par
\par\noindent\textbf{OM-SOL-1124 -- t3 XVII p.415 -- SOL}\par
Proposition 6.1.1 (iii), dernière occurrence du foncteur Φ\par
\emph{Reported reading}: \textbackslash{}Phi\par
\emph{Proposal}: \textbackslash{}Phi\_X\par
This finding intersects Orgogozo's \textbackslash{}sisi material at Proposition 6.1.1 (iii), dernière occurrence du foncteur Φ; preserve that human entry unchanged. The replacement \textbackslash{}Phi\_X at Proposition 6.1.1 (iii), dernière occurrence du foncteur Φ identifies a particular indexed object or symbol; that identification is not independently established by the quoted local expression.\par
\par\noindent\textbf{OM-SOL-1125 -- t3 XVII p.418 -- SOL}\par
First displayed formula in the proof\par
\emph{Reported reading}: \textbackslash{}Gamma\_\{p\}(\textbackslash{}Ff)\par
\emph{Proposal}: \textbackslash{}Gamma\_\{P\}(\textbackslash{}Ff)\par
This finding intersects Orgogozo's \textbackslash{}sisi material at First displayed formula in the proof; preserve that human entry unchanged. The replacement \textbackslash{}Gamma\_\{P\}(\textbackslash{}Ff) at First displayed formula in the proof identifies a particular indexed object or symbol; that identification is not independently established by the quoted local expression.\par
\par\noindent\textbf{OM-SOL-1126 -- t3 XVII p.418 -- SOL}\par
Paragraph beginning “L'assertion (iii)”, after “s'envoie dans”\par
\emph{Reported reading}: f\_\{!\}(\textbackslash{}Ff|U)\par
\emph{Proposal}: f\_\{!\}(\textbackslash{}Ff)\par
This finding intersects Orgogozo's \textbackslash{}sisi material at Paragraph beginning “L'assertion (iii)”, after “s'envoie dans”; preserve that human entry unchanged. The proposed reading f\_\{!\}(\textbackslash{}Ff) at Paragraph beginning “L'assertion (iii)”, after “s'envoie dans” changes a substantive mathematical assertion or bibliographic identification; no independent derivation of this exact replacement is established in this review.\par
\par\noindent\textbf{OM-SOL-1128 -- t3 XVII p.418 -- SOL}\par
Application of Chow's lemma, after the definition of p⁰\par
\emph{Reported reading}: fp\_\{0\}\par
\emph{Proposal}: fp\textasciicircum{}\{0\}\par
This finding intersects Orgogozo's \textbackslash{}sisi material at Application of Chow's lemma, after the definition of p⁰; preserve that human entry unchanged. The replacement fp\textasciicircum{}\{0\} at Application of Chow's lemma, after the definition of p⁰ requires a complete check of the maps, their endpoints and the composition convention; this identification is not independently verified.\par
\par\noindent\textbf{OM-SOL-1129 -- t3 XVII p.418 -- SOL}\par
Reference following the exact sequence\par
\emph{Reported reading}: VII 9.2\par
\emph{Proposal}: VIII 9.2\par
This finding intersects Orgogozo's \textbackslash{}sisi material at Reference following the exact sequence; preserve that human entry unchanged. The proposed reference VIII 9.2 at Reference following the exact sequence needs comparison with the statement actually invoked; that target has not been independently established here.\par
\par\noindent\textbf{OM-SOL-1130 -- t3 XVII p.419 -- SOL}\par
Opening paragraph, the two pairs after “pour les couples”\par
\emph{Reported reading}: \$(p\textasciicircum{}\{0\}\_\{\textbackslash{}ast\}\textbackslash{}mathcal\{G\})\$ et \$(fp\textasciicircum{}\{0\}\_\{\textbackslash{}ast\},\textbackslash{}mathcal\{G\})\$\par
\emph{Proposal}: \$(p\textasciicircum{}\{0\},\textbackslash{}mathcal\{G\})\$ et \$(fp\textasciicircum{}\{0\},\textbackslash{}mathcal\{G\})\$\par
This finding intersects Orgogozo's \textbackslash{}sisi material at Opening paragraph, the two pairs after “pour les couples”; preserve that human entry unchanged. The replacement \$(p\textasciicircum{}\{0\},\textbackslash{}mathcal\{G\})\$ et \$(fp\textasciicircum{}\{0\},\textbackslash{}mathcal\{G\})\$ at Opening paragraph, the two pairs after “pour les couples” requires a complete check of the maps, their endpoints and the composition convention; this identification is not independently verified.\par
\par\noindent\textbf{OM-SOL-1132 -- t3 XVII p.420 -- SOL}\par
Unnumbered decomposition formula following “On a”\par
\emph{Reported reading}: \textbackslash{}bigoplus\_\{i\textbackslash{}neq 0\} f\_\{1\textbackslash{}ast\}f\_\{i\}\textasciicircum{}\{\textbackslash{}ast\}(\textbackslash{}Ff|U)\par
\emph{Proposal}: \textbackslash{}bigoplus\_\{i\textbackslash{}neq 0\} f\_\{i\textbackslash{}ast\}f\_\{i\}\textasciicircum{}\{\textbackslash{}ast\}(\textbackslash{}Ff|U)\par
The proposed reading \textbackslash{}bigoplus\_\{i\textbackslash{}neq 0\} f\_\{i\textbackslash{}ast\}f\_\{i\}\textasciicircum{}\{\textbackslash{}ast\}(\textbackslash{}Ff|U) involves a printed glyph or a missing stroke at Unnumbered decomposition formula following “On a”; the textual record does not independently establish that glyph.\par
\par\noindent\textbf{OM-SOL-1134 -- t3 XVII p.421 -- SOL}\par
Proposition 6.2.2 (ii), codomain of \ensuremath{\tau}\par
\emph{Reported reading}: f\_\{!\}f\textasciicircum{}\{\textbackslash{}ast\}\textbackslash{}Ff\par
\emph{Proposal}: f\_\{!\}f\textasciicircum{}\{\textbackslash{}ast\}\textbackslash{}Ff|U\par
This finding intersects Orgogozo's \textbackslash{}sisi material at Proposition 6.2.2 (ii), codomain of \ensuremath{\tau}; preserve that human entry unchanged. The proposed reading f\_\{!\}f\textasciicircum{}\{\textbackslash{}ast\}\textbackslash{}Ff|U at Proposition 6.2.2 (ii), codomain of \ensuremath{\tau} changes a substantive mathematical assertion or bibliographic identification; no independent derivation of this exact replacement is established in this review.\par
\par\noindent\textbf{OM-SOL-1140 -- t3 XVII p.432 -- SOL}\par
Lemme 6.3.2, nom de la suite exacte\par
\emph{Reported reading}: F:0\textbackslash{}longrightarrow L'\textbackslash{}longrightarrow L\textbackslash{}longrightarrow L''\textbackslash{}longrightarrow 0\par
\emph{Proposal}: E:0\textbackslash{}longrightarrow L'\textbackslash{}longrightarrow L\textbackslash{}longrightarrow L''\textbackslash{}longrightarrow 0\par
This finding intersects Orgogozo's \textbackslash{}sisi material at Lemme 6.3.2, nom de la suite exacte; preserve that human entry unchanged. The proposed reading E:0\textbackslash{}longrightarrow L'\textbackslash{}longrightarrow L\textbackslash{}longrightarrow L''\textbackslash{}longrightarrow 0 at Lemme 6.3.2, nom de la suite exacte changes a substantive mathematical assertion or bibliographic identification; no independent derivation of this exact replacement is established in this review.\par
\par\noindent\textbf{OM-SOL-1141 -- t3 XVII p.433 -- SOL}\par
Lemme 6.3.2, décomposition de TS\textasciicircum{}n(X\textbackslash{}oplus Y)\par
\emph{Reported reading}: TS\textasciicircum{}\{n'\_\{1\}\}(X)\par
\emph{Proposal}: TS\textasciicircum{}\{n'\}(X)\par
This finding intersects Orgogozo's \textbackslash{}sisi material at Lemme 6.3.2, décomposition de TS\textasciicircum{}n(X\textbackslash{}oplus Y); preserve that human entry unchanged. The replacement TS\textasciicircum{}\{n'\}(X) at Lemme 6.3.2, décomposition de TS\textasciicircum{}n(X\textbackslash{}oplus Y) identifies a particular indexed object or symbol; that identification is not independently established by the quoted local expression.\par
\par\noindent\textbf{OM-SOL-1142 -- t3 XVII p.433 -- SOL}\par
Lemme 6.3.2, diagramme, sommet supérieur central\par
\emph{Reported reading}: \textbackslash{}Ens(L'')\par
\emph{Proposal}: \textbackslash{}End(L'')\par
The proposed reading \textbackslash{}End(L'') at Lemme 6.3.2, diagramme, sommet supérieur central changes a substantive mathematical assertion or bibliographic identification; no independent derivation of this exact replacement is established in this review.\par
\par\noindent\textbf{OM-SOL-1143 -- t3 XVII p.433 -- SOL}\par
Lemme 6.3.2, diagramme, sommet inférieur gauche\par
\emph{Reported reading}: TS\textasciicircum{}\{n\}(L)\par
\emph{Proposal}: TS\textasciicircum{}\{n\}(\textbackslash{}End(L))\par
This finding intersects Orgogozo's \textbackslash{}sisi material at Lemme 6.3.2, diagramme, sommet inférieur gauche; preserve that human entry unchanged. The proposed reading TS\textasciicircum{}\{n\}(\textbackslash{}End(L)) at Lemme 6.3.2, diagramme, sommet inférieur gauche changes a substantive mathematical assertion or bibliographic identification; no independent derivation of this exact replacement is established in this review.\par
\par\noindent\textbf{OM-SOL-1144 -- t3 XVII p.433 -- SOL}\par
Lemme 6.3.3, but de la flèche sigma\par
\emph{Reported reading}: TS\textasciicircum{}\{n\}\_\{A\}TS\textasciicircum{}\{n\}\_\{B\}(C)\par
\emph{Proposal}: TS\textasciicircum{}\{n\}\_\{A\}TS\textasciicircum{}\{m\}\_\{B\}(C)\par
This finding intersects Orgogozo's \textbackslash{}sisi material at Lemme 6.3.3, but de la flèche sigma; preserve that human entry unchanged. The replacement TS\textasciicircum{}\{n\}\_\{A\}TS\textasciicircum{}\{m\}\_\{B\}(C) at Lemme 6.3.3, but de la flèche sigma identifies a particular indexed object or symbol; that identification is not independently established by the quoted local expression.\par
\par\noindent\textbf{OM-SOL-1145 -- t3 XVII p.434 -- SOL}\par
6.3.4, référence finale à SGA 1\par
\emph{Reported reading}: (SGA 1 IV 3)\par
\emph{Proposal}: (\textbackslash{}SGA 1 V 1)\par
This finding intersects Orgogozo's \textbackslash{}sisi material at 6.3.4, référence finale à SGA 1; preserve that human entry unchanged. The proposed reference (\textbackslash{}SGA 1 V 1) at 6.3.4, référence finale à SGA 1 needs comparison with the statement actually invoked; that target has not been independently established here.\par
\par\noindent\textbf{OM-SOL-1148 -- t3 XVII p.436 -- SOL}\par
Proposition 6.3.7, formule finale\par
\emph{Reported reading}: u(\textbackslash{}sigma(\textbackslash{}Ff ))=\allowbreak \textbackslash{}sigma(u\_\{\textbackslash{}ast\}\textbackslash{}Ff )\par
\emph{Proposal}: (\textbackslash{}Sym\textasciicircum{}n\_S u)(\textbackslash{}sigma(\textbackslash{}Ff ))=\allowbreak \textbackslash{}sigma(u\_\{\textbackslash{}ast\}\textbackslash{}Ff )\par
This finding intersects Orgogozo's \textbackslash{}sisi material at Proposition 6.3.7, formule finale; preserve that human entry unchanged. The replacement (\textbackslash{}Sym\textasciicircum{}n\_S u)(\textbackslash{}sigma(\textbackslash{}Ff ))=\allowbreak \textbackslash{}sigma(u\_\{\textbackslash{}ast\}\textbackslash{}Ff ) at Proposition 6.3.7, formule finale requires a complete check of the maps, their endpoints and the composition convention; this identification is not independently verified.\par
\par\noindent\textbf{OM-SOL-1149 -- t3 XVII p.437 -- SOL}\par
6.3.8, définition du faisceau des diviseurs\par
\emph{Reported reading}: qui sont finis sur S et de degré n\par
\emph{Proposal}: qui sont \textbackslash{}emph\{finis\} sur \$S'\$ et de degré \$n\$\par
The replacement qui sont \textbackslash{}emph\{finis\} sur \$S'\$ et de degré \$n\$ at 6.3.8, définition du faisceau des diviseurs identifies a particular indexed object or symbol; that identification is not independently established by the quoted local expression.\par
\par\noindent\textbf{OM-SOL-1150 -- t3 XVII p.440 -- SOL}\par
Paragraphe « Pour vérifier 6.3.11 »\par
\emph{Reported reading}: Pour vérifier 6.3.11\par
\emph{Proposal}: Pour vérifier \textbackslash{}hyperref[XVII.prop6.3.11.1]\{6.3.11.1\}\par
This finding intersects Orgogozo's \textbackslash{}sisi material at Paragraphe « Pour vérifier 6.3.11 »; preserve that human entry unchanged. The proposed reference Pour vérifier \textbackslash{}hyperref[XVII.prop6.3.11.1]\{6.3.11.1\} at Paragraphe « Pour vérifier 6.3.11 » needs comparison with the statement actually invoked; that target has not been independently established here.\par
\par\noindent\textbf{OM-SOL-1151 -- t3 XVII p.440 -- SOL}\par
Corollaire 6.3.11.2, parenthèse sur les conoyaux\par
\emph{Reported reading}: sur S et à fibres quasi-finies\par
\emph{Proposal}: sur \$S\$ et à supports quasi-finis\par
This finding intersects Orgogozo's \textbackslash{}sisi material at Corollaire 6.3.11.2, parenthèse sur les conoyaux; preserve that human entry unchanged. The proposed reading sur \$S\$ et à supports quasi-finis at Corollaire 6.3.11.2, parenthèse sur les conoyaux changes a substantive mathematical assertion or bibliographic identification; no independent derivation of this exact replacement is established in this review.\par
\par\noindent\textbf{OM-SOL-1153 -- t3 XVII p.441 -- SOL}\par
Premier paragraphe, « où … est une telle indéterminée »\par
\emph{Reported reading}: \$j\$\par
\emph{Proposal}: \$T\$\par
This finding intersects Orgogozo's \textbackslash{}sisi material at Premier paragraphe, « où … est une telle indéterminée »; preserve that human entry unchanged. The proposed reading \$T\$ at Premier paragraphe, « où … est une telle indéterminée » changes a substantive mathematical assertion or bibliographic identification; no independent derivation of this exact replacement is established in this review.\par
\par\noindent\textbf{OM-SOL-1154 -- t3 XVII p.442 -- SOL}\par
Formule après « On sait qu'on a », première norme\par
\emph{Reported reading}: \textbackslash{}det u'=\allowbreak N\_\{X'/X\}\textbackslash{}det u\par
\emph{Proposal}: \textbackslash{}det u'=\allowbreak N\_\{X/X'\}\textbackslash{}det u\par
This finding intersects Orgogozo's \textbackslash{}sisi material at Formule après « On sait qu'on a », première norme; preserve that human entry unchanged. The replacement \textbackslash{}det u'=\allowbreak N\_\{X/X'\}\textbackslash{}det u at Formule après « On sait qu'on a », première norme identifies a particular indexed object or symbol; that identification is not independently established by the quoted local expression.\par
\par\noindent\textbf{OM-SOL-1155 -- t3 XVII p.442 -- SOL}\par
Même formule, norme dans l'expression de Q'\par
\emph{Reported reading}: Q'=\allowbreak \textbackslash{}Coker N\_\{X'/X\}(\textbackslash{}det u)\par
\emph{Proposal}: Q'=\allowbreak \textbackslash{}Coker N\_\{X/X'\}(\textbackslash{}det u)\par
This finding intersects Orgogozo's \textbackslash{}sisi material at Même formule, norme dans l'expression de Q'; preserve that human entry unchanged. The replacement Q'=\allowbreak \textbackslash{}Coker N\_\{X/X'\}(\textbackslash{}det u) at Même formule, norme dans l'expression de Q' identifies a particular indexed object or symbol; that identification is not independently established by the quoted local expression.\par
\par\noindent\textbf{OM-SOL-1156 -- t3 XVII p.442 -- SOL}\par
Égalité après « On a trivialement », second membre\par
\emph{Reported reading}: \{\textbackslash{}det\}\_\{S\}(\textbackslash{}varphi',P)\par
\emph{Proposal}: \{\textbackslash{}det\}\_\{S\}(\textbackslash{}varphi',P')\par
This finding intersects Orgogozo's \textbackslash{}sisi material at Égalité après « On a trivialement », second membre; preserve that human entry unchanged. The replacement \{\textbackslash{}det\}\_\{S\}(\textbackslash{}varphi',P') at Égalité après « On a trivialement », second membre identifies a particular indexed object or symbol; that identification is not independently established by the quoted local expression.\par
\par\noindent\textbf{OM-SOL-1157 -- t3 XVII p.442 -- SOL}\par
Dernier paragraphe, réduction au cas M =\allowbreak  O\_X\par
\emph{Reported reading}: et \$\textbackslash{}varphi\$ par \$\textbackslash{}det u\$\par
\emph{Proposal}: et \$u\$ par \$\textbackslash{}det u\$\par
The replacement et \$u\$ par \$\textbackslash{}det u\$ at Dernier paragraphe, réduction au cas M =\allowbreak  O\_X requires a complete check of the maps, their endpoints and the composition convention; this identification is not independently verified.\par
\par\noindent\textbf{OM-SOL-1158 -- t3 XVII p.444 -- SOL}\par
Définition de t, renvoi après « la section »\par
\emph{Reported reading}: (6.3.4.3)\par
\emph{Proposal}: (6.3.4.2)\par
This finding intersects Orgogozo's \textbackslash{}sisi material at Définition de t, renvoi après « la section »; preserve that human entry unchanged. The proposed reference (6.3.4.2) at Définition de t, renvoi après « la section » needs comparison with the statement actually invoked; that target has not been independently established here.\par
\par\noindent\textbf{OM-SOL-1160 -- t3 XVII p.444 -- SOL}\par
Formule définissant u\_n\par
\emph{Reported reading}: u\_\{n\}\textbackslash{}in \textbackslash{}Sym\textasciicircum{}\{n\}\_\{S\}(T)\par
\emph{Proposal}: u\_\{n\}\textbackslash{}in G(\textbackslash{}Sym\textasciicircum{}\{n\}\_\{S\}(T))\par
This finding intersects Orgogozo's \textbackslash{}sisi material at Formule définissant u\_n; preserve that human entry unchanged. The proposed reading u\_\{n\}\textbackslash{}in G(\textbackslash{}Sym\textasciicircum{}\{n\}\_\{S\}(T)) at Formule définissant u\_n changes a substantive mathematical assertion or bibliographic identification; no independent derivation of this exact replacement is established in this review.\par
\par\noindent\textbf{OM-SOL-1161 -- t3 XVII p.445 -- SOL}\par
Après f =\allowbreak  gh, noms des morphismes de rangs constants\par
\emph{Reported reading}: avec \$f\$ et \$g\$ finis localement libres de rangs constants \$n\$ et \$m\$\par
\emph{Proposal}: avec \$g\$ et \$h\$ finis localement libres de rangs constants \$n\$ et \$m\$\par
This finding intersects Orgogozo's \textbackslash{}sisi material at Après f =\allowbreak  gh, noms des morphismes de rangs constants; preserve that human entry unchanged. The replacement avec \$g\$ et \$h\$ finis localement libres de rangs constants \$n\$ et \$m\$ at Après f =\allowbreak  gh, noms des morphismes de rangs constants requires a complete check of the maps, their endpoints and the composition convention; this identification is not independently verified.\par
\par\noindent\textbf{OM-SOL-1162 -- t3 XVII p.446 -- SOL}\par
Diagramme (6.3.15.2), les trois étiquettes de trace\par
\emph{Reported reading}: \textbackslash{}Tr\_g(u)\_n ; \textbackslash{}Tr\_f\textbackslash{}Tr\_g(u) ; \textbackslash{}Tr\_\{fg\}(u)\par
\emph{Proposal}: \textbackslash{}Tr\_h(u)\_n ; \textbackslash{}Tr\_g\textbackslash{}Tr\_h(u) ; \textbackslash{}Tr\_\{gh\}(u)\par
This finding intersects Orgogozo's \textbackslash{}sisi material at Diagramme (6.3.15.2), les trois étiquettes de trace; preserve that human entry unchanged. The replacement \textbackslash{}Tr\_h(u)\_n ; \textbackslash{}Tr\_g\textbackslash{}Tr\_h(u) ; \textbackslash{}Tr\_\{gh\}(u) at Diagramme (6.3.15.2), les trois étiquettes de trace identifies a particular indexed object or symbol; that identification is not independently established by the quoted local expression.\par
\par\noindent\textbf{OM-SOL-1163 -- t3 XVII p.446 -- SOL}\par
Première phrase après le diagramme\par
\emph{Reported reading}: car S' est plat sur S\par
\emph{Proposal}: car \$S''\$ est plat sur \$S\$\par
This finding intersects Orgogozo's \textbackslash{}sisi material at Première phrase après le diagramme; preserve that human entry unchanged. The replacement car \$S''\$ est plat sur \$S\$ at Première phrase après le diagramme requires a complete check of the maps, their endpoints and the composition convention; this identification is not independently verified.\par
\par\noindent\textbf{OM-SOL-1165 -- t3 XVII p.446 -- SOL}\par
6.3.16, définition de l'image directe à support propre\par
\emph{Reported reading}: f\_\{!\}T\par
\emph{Proposal}: f\_\{!\}\ensuremath{\mathcal{F}}\par
This finding intersects Orgogozo's \textbackslash{}sisi material at 6.3.16, définition de l'image directe à support propre; preserve that human entry unchanged. The proposed reading f\_\{!\}\ensuremath{\mathcal{F}} at 6.3.16, définition de l'image directe à support propre changes a substantive mathematical assertion or bibliographic identification; no independent derivation of this exact replacement is established in this review.\par
\par\noindent\textbf{OM-SOL-1166 -- t3 XVII p.446 -- SOL}\par
6.3.16, référence après « limite inductive locale »\par
\emph{Reported reading}: (IV 14. )\par
\emph{Proposal}: \textbackslash{}hyperrefext[V][V.subsec8.2]\{8.2\}\par
This finding intersects Orgogozo's \textbackslash{}sisi material at 6.3.16, référence après « limite inductive locale »; preserve that human entry unchanged. The proposed reference \textbackslash{}hyperrefext[V][V.subsec8.2]\{8.2\} at 6.3.16, référence après « limite inductive locale » needs comparison with the statement actually invoked; that target has not been independently established here.\par
\par\noindent\textbf{OM-SOL-1169 -- t3 XVII p.449 -- SOL}\par
Après « l'élément canonique »\par
\emph{Reported reading}: (6.3.4.3)\par
\emph{Proposal}: \textbackslash{}eqhyperref[XVII.eq6.3.4.2]\{6.3.4.2\}\par
This finding intersects Orgogozo's \textbackslash{}sisi material at Après « l'élément canonique »; preserve that human entry unchanged. The proposed reference \textbackslash{}eqhyperref[XVII.eq6.3.4.2]\{6.3.4.2\} at Après « l'élément canonique » needs comparison with the statement actually invoked; that target has not been independently established here.\par
\par\noindent\textbf{OM-SOL-1171 -- t3 XVII p.450 -- SOL}\par
Proposition 6.3.19 (ii), second membre de la formule\par
\emph{Reported reading}: \textbackslash{}Tr\_\{(u\_\{\textbackslash{}ast\}E)/S\}(G)\par
\emph{Proposal}: \textbackslash{}Tr\_\{(u\textasciicircum{}\{\textbackslash{}ast\}E)/S\}(g)\par
This finding intersects Orgogozo's \textbackslash{}sisi material at Proposition 6.3.19 (ii), second membre de la formule; preserve that human entry unchanged. The proposed reading \textbackslash{}Tr\_\{(u\textasciicircum{}\{\textbackslash{}ast\}E)/S\}(g) at Proposition 6.3.19 (ii), second membre de la formule changes a substantive mathematical assertion or bibliographic identification; no independent derivation of this exact replacement is established in this review.\par
\par\noindent\textbf{OM-SOL-1172 -- t3 XVII p.450 -- SOL}\par
Preuve (i), terme \ensuremath{\sigma}(D)* (u* g\_n)\par
\emph{Reported reading}: \textbackslash{}sigma(D)\textasciicircum{}\{\textbackslash{}ast\}(u\textasciicircum{}\{\textbackslash{}ast\}g\_n)\par
\emph{Proposal}: \textbackslash{}sigma(D)\textasciicircum{}\{\textbackslash{}ast\}(\textbackslash{}Sym\textasciicircum{}\{n\}\_\{S\}(u)\textasciicircum{}\{\textbackslash{}ast\}g\_\{n\})\par
This finding intersects Orgogozo's \textbackslash{}sisi material at Preuve (i), terme \ensuremath{\sigma}(D)* (u* g\_n); preserve that human entry unchanged. The proposed reading \textbackslash{}sigma(D)\textasciicircum{}\{\textbackslash{}ast\}(\textbackslash{}Sym\textasciicircum{}\{n\}\_\{S\}(u)\textasciicircum{}\{\textbackslash{}ast\}g\_\{n\}) at Preuve (i), terme \ensuremath{\sigma}(D)* (u* g\_n) changes a substantive mathematical assertion or bibliographic identification; no independent derivation of this exact replacement is established in this review.\par
\par\noindent\textbf{OM-SOL-1173 -- t3 XVII p.450 -- SOL}\par
Preuve (i), terme suivant dans la chaîne d'égalités\par
\emph{Reported reading}: u\textbackslash{}circ\textbackslash{}sigma(D)\textasciicircum{}\{\textbackslash{}ast\}(g\_n)\par
\emph{Proposal}: (\textbackslash{}Sym\textasciicircum{}\{n\}\_\{S\}(u)\textbackslash{}circ\textbackslash{}sigma(D))\textasciicircum{}\{\textbackslash{}ast\}(g\_\{n\})\par
This finding intersects Orgogozo's \textbackslash{}sisi material at Preuve (i), terme suivant dans la chaîne d'égalités; preserve that human entry unchanged. The proposed reading (\textbackslash{}Sym\textasciicircum{}\{n\}\_\{S\}(u)\textbackslash{}circ\textbackslash{}sigma(D))\textasciicircum{}\{\textbackslash{}ast\}(g\_\{n\}) at Preuve (i), terme suivant dans la chaîne d'égalités changes a substantive mathematical assertion or bibliographic identification; no independent derivation of this exact replacement is established in this review.\par
\par\noindent\textbf{OM-SOL-1174 -- t3 XVII p.450 -- SOL}\par
Preuve (i), dernier terme de la chaîne d'égalités\par
\emph{Reported reading}: \textbackslash{}Tr\_\{u\_\{\textbackslash{}ast\}D\}(g\_n)\par
\emph{Proposal}: \textbackslash{}Tr\_\{u\_\{\textbackslash{}ast\}D\}(g)\par
The replacement \textbackslash{}Tr\_\{u\_\{\textbackslash{}ast\}D\}(g) at Preuve (i), dernier terme de la chaîne d'égalités identifies a particular indexed object or symbol; that identification is not independently established by the quoted local expression.\par
\par\noindent\textbf{OM-SOL-1175 -- t3 XVII p.451 -- SOL}\par
Proposition 6.3.21, clause specifying g\par
\emph{Reported reading}: « g section locale (au sens fppf) de G sur Y »\par
\emph{Proposal}: \$g\$ section locale (au sens fppf) de \$u\_\{\textbackslash{}ast\}G\_\{X\}\$ sur \$Y\$\par
The proposed reading \$g\$ section locale (au sens fppf) de \$u\_\{\textbackslash{}ast\}G\_\{X\}\$ sur \$Y\$ at Proposition 6.3.21, clause specifying g changes a substantive mathematical assertion or bibliographic identification; no independent derivation of this exact replacement is established in this review.\par
\par\noindent\textbf{OM-SOL-1176 -- t3 XVII p.451 -- SOL}\par
Paragraph beginning « Le foncteur trace »\par
\emph{Reported reading}: « Le foncteur trace 6.3.20 »\par
\emph{Proposal}: 6.3.21\par
This finding intersects Orgogozo's \textbackslash{}sisi material at Paragraph beginning « Le foncteur trace »; preserve that human entry unchanged. The proposed reference 6.3.21 at Paragraph beginning « Le foncteur trace » needs comparison with the statement actually invoked; that target has not been independently established here.\par
\par\noindent\textbf{OM-SOL-1177 -- t3 XVII p.452 -- SOL}\par
Lemme 6.3.22.1, clause specifying k\par
\emph{Reported reading}: \$u\_\{\textbackslash{}ast\}K\$\par
\emph{Proposal}: \$f\_\{\textbackslash{}ast\}K\$\par
This finding intersects Orgogozo's \textbackslash{}sisi material at Lemme 6.3.22.1, clause specifying k; preserve that human entry unchanged. The replacement \$f\_\{\textbackslash{}ast\}K\$ at Lemme 6.3.22.1, clause specifying k requires a complete check of the maps, their endpoints and the composition convention; this identification is not independently verified.\par
\par\noindent\textbf{OM-SOL-1178 -- t3 XVII p.453 -- SOL}\par
6.3.23.2, first image of the point z\par
\emph{Reported reading}: \$x(Z)\$\par
\emph{Proposal}: \$x(z)\$\par
This finding intersects Orgogozo's \textbackslash{}sisi material at 6.3.23.2, first image of the point z; preserve that human entry unchanged. The replacement \$x(z)\$ at 6.3.23.2, first image of the point z identifies a particular indexed object or symbol; that identification is not independently established by the quoted local expression.\par
\par\noindent\textbf{OM-SOL-1179 -- t3 XVII p.453 -- SOL}\par
6.3.23.2, second image of the point z\par
\emph{Reported reading}: \$y(Z)\$\par
\emph{Proposal}: \$y(z)\$\par
This finding intersects Orgogozo's \textbackslash{}sisi material at 6.3.23.2, second image of the point z; preserve that human entry unchanged. The replacement \$y(z)\$ at 6.3.23.2, second image of the point z identifies a particular indexed object or symbol; that identification is not independently established by the quoted local expression.\par
\par\noindent\textbf{OM-SOL-1180 -- t3 XVII p.453 -- SOL}\par
6.3.23.2, conclusion « est étale en »\par
\emph{Reported reading}: \$x(Z)\$\par
\emph{Proposal}: \$x(z)\$\par
This finding intersects Orgogozo's \textbackslash{}sisi material at 6.3.23.2, conclusion « est étale en »; preserve that human entry unchanged. The replacement \$x(z)\$ at 6.3.23.2, conclusion « est étale en » identifies a particular indexed object or symbol; that identification is not independently established by the quoted local expression.\par
\par\noindent\textbf{OM-SOL-1181 -- t3 XVII p.453 -- SOL}\par
Proof of 6.3.23.2, change of base\par
\emph{Reported reading}: \$S'\textbackslash{}longrightarrow S\$\par
\emph{Proposal}: \$T'\textbackslash{}rightarrow T\$\par
This finding intersects Orgogozo's \textbackslash{}sisi material at Proof of 6.3.23.2, change of base; preserve that human entry unchanged. The proposed reading \$T'\textbackslash{}rightarrow T\$ at Proof of 6.3.23.2, change of base changes a substantive mathematical assertion or bibliographic identification; no independent derivation of this exact replacement is established in this review.\par
\par\noindent\textbf{OM-SOL-1182 -- t3 XVII p.453 -- SOL}\par
Proof of 6.3.23.2, finite presentation\par
\emph{Reported reading}: « de présentation finie sur S »\par
\emph{Proposal}: finie sur \$T\$\par
This finding intersects Orgogozo's \textbackslash{}sisi material at Proof of 6.3.23.2, finite presentation; preserve that human entry unchanged. The proposed reading finie sur \$T\$ at Proof of 6.3.23.2, finite presentation changes a substantive mathematical assertion or bibliographic identification; no independent derivation of this exact replacement is established in this review.\par
\par\noindent\textbf{OM-SOL-1183 -- t3 XVII p.453 -- SOL}\par
Proof of 6.3.23.2, flatness\par
\emph{Reported reading}: « plat sur S »\par
\emph{Proposal}: plat sur \$T\$\par
This finding intersects Orgogozo's \textbackslash{}sisi material at Proof of 6.3.23.2, flatness; preserve that human entry unchanged. The proposed reading plat sur \$T\$ at Proof of 6.3.23.2, flatness changes a substantive mathematical assertion or bibliographic identification; no independent derivation of this exact replacement is established in this review.\par
\par\noindent\textbf{OM-SOL-1184 -- t3 XVII p.453 -- SOL}\par
Proof of 6.3.23.2, reduction to a field\par
\emph{Reported reading}: « au cas S spectre d'un corps »\par
\emph{Proposal}: au cas \$T\$ spectre d'un corps\par
This finding intersects Orgogozo's \textbackslash{}sisi material at Proof of 6.3.23.2, reduction to a field; preserve that human entry unchanged. The proposed reading au cas \$T\$ spectre d'un corps at Proof of 6.3.23.2, reduction to a field changes a substantive mathematical assertion or bibliographic identification; no independent derivation of this exact replacement is established in this review.\par
\par\noindent\textbf{OM-SOL-1185 -- t3 XVII p.453 -- SOL}\par
Proof of 6.3.23.2, first completion\par
\emph{Reported reading}: « Le complété de Sym\_T\textasciicircum{}n(X) en Z »\par
\emph{Proposal}: en \$z\$\par
This finding intersects Orgogozo's \textbackslash{}sisi material at Proof of 6.3.23.2, first completion; preserve that human entry unchanged. The replacement en \$z\$ at Proof of 6.3.23.2, first completion identifies a particular indexed object or symbol; that identification is not independently established by the quoted local expression.\par
\par\noindent\textbf{OM-SOL-1186 -- t3 XVII p.454 -- SOL}\par
Final reference in the proof of Proposition 6.3.23\par
\emph{Reported reading}: 6.3.2.3\par
\emph{Proposal}: 6.3.23\par
This finding intersects Orgogozo's \textbackslash{}sisi material at Final reference in the proof of Proposition 6.3.23; preserve that human entry unchanged. The proposed reading 6.3.23 at Final reference in the proof of Proposition 6.3.23 changes a substantive mathematical assertion or bibliographic identification; no independent derivation of this exact replacement is established in this review.\par
\par\noindent\textbf{OM-SOL-1187 -- t3 XVII p.454 -- SOL}\par
6.3.24.0, target of c\_v\par
\emph{Reported reading}: \$\textbackslash{}Sym\textasciicircum{}\{n\}\_\{T\}(v)\textasciicircum{}\{\textbackslash{}ast\}K\_\{n\}\$\par
\emph{Proposal}: \$\textbackslash{}Sym\textasciicircum{}\{n\}\_\{T\}(v)\textasciicircum{}\{\textbackslash{}ast\}(u\textasciicircum{}\{\textbackslash{}ast\}K)\_\{n\}\$\par
This finding intersects Orgogozo's \textbackslash{}sisi material at 6.3.24.0, target of c\_v; preserve that human entry unchanged. The proposed reading \$\textbackslash{}Sym\textasciicircum{}\{n\}\_\{T\}(v)\textasciicircum{}\{\textbackslash{}ast\}(u\textasciicircum{}\{\textbackslash{}ast\}K)\_\{n\}\$ at 6.3.24.0, target of c\_v changes a substantive mathematical assertion or bibliographic identification; no independent derivation of this exact replacement is established in this review.\par
\par\noindent\textbf{OM-SOL-1188 -- t3 XVII p.458 -- SOL}\par
Paragraph beginning « La commutativité », second reference before the composition formula\par
\emph{Reported reading}: (6.3.25.1) (6.3.25.4)\par
\emph{Proposal}: (\textbackslash{}hyperref[XVII.par6.3.25.1]\{6.3.25.1\}) (\textbackslash{}hyperref[XVII.par6.3.25.5]\{6.3.25.5\})\par
The proposed reference (\textbackslash{}hyperref[XVII.par6.3.25.1]\{6.3.25.1\}) (\textbackslash{}hyperref[XVII.par6.3.25.5]\{6.3.25.5\}) at Paragraph beginning « La commutativité », second reference before the composition formula needs comparison with the statement actually invoked; that target has not been independently established here.\par
\par\noindent\textbf{OM-SOL-1190 -- t3 XVII p.459 -- SOL}\par
Proposition 6.3.28, « Alors le complexe… »\par
\emph{Reported reading}: complexe semi-simplicial d'algèbres\par
\emph{Proposal}: complexe co-semi-simplicial d'algèbres\par
This finding intersects Orgogozo's \textbackslash{}sisi material at Proposition 6.3.28, « Alors le complexe… »; preserve that human entry unchanged. The proposed reading complexe co-semi-simplicial d'algèbres at Proposition 6.3.28, « Alors le complexe… » changes a substantive mathematical assertion or bibliographic identification; no independent derivation of this exact replacement is established in this review.\par
\par\noindent\textbf{OM-SOL-1191 -- t3 XVII p.459 -- SOL}\par
Proposition 6.3.28, « en tant que… »\par
\emph{Reported reading}: \$TS\textasciicircum{}\{n\}\_\{R\}(A)\$-module différentiel semi-simplicial\par
\emph{Proposal}: \$TS\textasciicircum{}\{n\}\_\{R\}(A)\$-module différentiel co-semi-simplicial\par
This finding intersects Orgogozo's \textbackslash{}sisi material at Proposition 6.3.28, « en tant que… »; preserve that human entry unchanged. The proposed reading \$TS\textasciicircum{}\{n\}\_\{R\}(A)\$-module différentiel co-semi-simplicial at Proposition 6.3.28, « en tant que… » changes a substantive mathematical assertion or bibliographic identification; no independent derivation of this exact replacement is established in this review.\par
\par\noindent\textbf{OM-SOL-1192 -- t3 XVII p.459 -- SOL}\par
Proof of Proposition 6.3.28, opening sentence\par
\emph{Reported reading}: associé au module semi-simplicial\par
\emph{Proposal}: associé au module co-semi-simplicial\par
This finding intersects Orgogozo's \textbackslash{}sisi material at Proof of Proposition 6.3.28, opening sentence; preserve that human entry unchanged. The proposed reading associé au module co-semi-simplicial at Proof of Proposition 6.3.28, opening sentence changes a substantive mathematical assertion or bibliographic identification; no independent derivation of this exact replacement is established in this review.\par
\par\noindent\textbf{OM-SOL-1193 -- t3 XVII p.461 -- SOL}\par
Fin de la preuve de 6.3.31\par
\emph{Reported reading}: \textbackslash{}Sym\textasciicircum{}\{n\}(X)\par
\emph{Proposal}: \textbackslash{}Sym\textasciicircum{}\{n\}\_\{T\}(X)\par
This finding intersects Orgogozo's \textbackslash{}sisi material at Fin de la preuve de 6.3.31; preserve that human entry unchanged. The proposed reading \textbackslash{}Sym\textasciicircum{}\{n\}\_\{T\}(X) at Fin de la preuve de 6.3.31 changes a substantive mathematical assertion or bibliographic identification; no independent derivation of this exact replacement is established in this review.\par
\par\noindent\textbf{OM-SOL-1194 -- t3 XVII p.462 -- SOL}\par
7.0.1, phrase « Ceci posé »\par
\emph{Reported reading}: \textbackslash{}sheaf\{\textbackslash{}Mod\}(\textbackslash{}Ee\textasciicircum{}\{\textbackslash{}circ\}\_\{R\},\textbackslash{}mathcal\{A\})\par
\emph{Proposal}: \textbackslash{}sheaf\{\textbackslash{}Mod\}(\textbackslash{}sheafGamma(\textbackslash{}Ee\textasciicircum{}\{\textbackslash{}circ\}\_\{R\}),\textbackslash{}mathcal\{A\})\par
This finding intersects Orgogozo's \textbackslash{}sisi material at 7.0.1, phrase « Ceci posé »; preserve that human entry unchanged. The proposed reading \textbackslash{}sheaf\{\textbackslash{}Mod\}(\textbackslash{}sheafGamma(\textbackslash{}Ee\textasciicircum{}\{\textbackslash{}circ\}\_\{R\}),\textbackslash{}mathcal\{A\}) at 7.0.1, phrase « Ceci posé » changes a substantive mathematical assertion or bibliographic identification; no independent derivation of this exact replacement is established in this review.\par
\par\noindent\textbf{OM-SOL-1195 -- t3 XVII p.463 -- SOL}\par
Lemme 7.0.3, flèche canonique sur les complexes\par
\emph{Reported reading}: j\_\{!\}\textbackslash{}circ R\textasciicircum{}\{+\}f'(K\textasciicircum{}\{\textbackslash{}bigdot\})\textbackslash{}longrightarrow R\textasciicircum{}\{+\}f\textbackslash{}circ j\_\{!\}(K\textasciicircum{}\{\textbackslash{}bigdot\})\par
\emph{Proposal}: j\_\{!\}\textbackslash{}circ R\textasciicircum{}\{+\}f'\_\{*\}(K\textasciicircum{}\{\textbackslash{}bigdot\})\textbackslash{}longrightarrow R\textasciicircum{}\{+\}f\_\{*\}\textbackslash{}circ j'\_\{!\}(K\textasciicircum{}\{\textbackslash{}bigdot\})\par
This finding intersects Orgogozo's \textbackslash{}sisi material at Lemme 7.0.3, flèche canonique sur les complexes; preserve that human entry unchanged. The replacement j\_\{!\}\textbackslash{}circ R\textasciicircum{}\{+\}f'\_\{*\}(K\textasciicircum{}\{\textbackslash{}bigdot\})\textbackslash{}longrightarrow R\textasciicircum{}\{+\}f\_\{*\}\textbackslash{}circ j'\_\{!\}(K\textasciicircum{}\{\textbackslash{}bigdot\}) at Lemme 7.0.3, flèche canonique sur les complexes identifies a particular indexed object or symbol; that identification is not independently established by the quoted local expression.\par
\par\noindent\textbf{OM-SOL-1196 -- t3 XVII p.463 -- SOL}\par
Formule (7.0.3.2)\par
\emph{Reported reading}: j\_\{!\}\textbackslash{}circ R\textasciicircum{}\{i\}f'(F)\textbackslash{}longrightarrow R\textasciicircum{}\{i\}f\textbackslash{}circ j\_\{!\}(F)\par
\emph{Proposal}: j\_\{!\}\textbackslash{}circ R\textasciicircum{}\{i\}f'\_\{*\}(F)\textbackslash{}longrightarrow R\textasciicircum{}\{i\}f\_\{*\}\textbackslash{}circ j'\_\{!\}(F)\par
This finding intersects Orgogozo's \textbackslash{}sisi material at Formule (7.0.3.2); preserve that human entry unchanged. The replacement j\_\{!\}\textbackslash{}circ R\textasciicircum{}\{i\}f'\_\{*\}(F)\textbackslash{}longrightarrow R\textasciicircum{}\{i\}f\_\{*\}\textbackslash{}circ j'\_\{!\}(F) at Formule (7.0.3.2) requires a complete check of the maps, their endpoints and the composition convention; this identification is not independently verified.\par
\par\noindent\textbf{OM-SOL-1197 -- t3 XVII p.463 -- SOL}\par
Fin de la preuve du lemme 7.0.3, flèche sur les fibres géométriques\par
\emph{Reported reading}: (j\_\{!\}\textbackslash{}circ R\textasciicircum{}\{i\}f'(F))\_\{\textbackslash{}xi\}\textbackslash{}rightarrow (R\textasciicircum{}\{i\}f\textbackslash{}circ j\_\{!\}(F))\_\{\textbackslash{}xi\}\par
\emph{Proposal}: (j\_\{!\}\textbackslash{}circ R\textasciicircum{}\{i\}f'\_\{*\}(F))\_\{\textbackslash{}xi\}\textbackslash{}rightarrow (R\textasciicircum{}\{i\}f\_\{*\}\textbackslash{}circ j'\_\{!\}(F))\_\{\textbackslash{}xi\}\par
This finding intersects Orgogozo's \textbackslash{}sisi material at Fin de la preuve du lemme 7.0.3, flèche sur les fibres géométriques; preserve that human entry unchanged. The replacement (j\_\{!\}\textbackslash{}circ R\textasciicircum{}\{i\}f'\_\{*\}(F))\_\{\textbackslash{}xi\}\textbackslash{}rightarrow (R\textasciicircum{}\{i\}f\_\{*\}\textbackslash{}circ j'\_\{!\}(F))\_\{\textbackslash{}xi\} at Fin de la preuve du lemme 7.0.3, flèche sur les fibres géométriques identifies a particular indexed object or symbol; that identification is not independently established by the quoted local expression.\par
\par\noindent\textbf{OM-SOL-1198 -- t3 XVII p.464 -- SOL}\par
Construction 7.0.4, sommet supérieur gauche du diagramme\par
\emph{Reported reading}: D(\textbackslash{}sheafGamma(\textbackslash{}bar\{X\}),\textbackslash{},)\par
\emph{Proposal}: D(\textbackslash{}sheafGamma(\textbackslash{}bar\{X\}),\textbackslash{}Aa)\par
This finding intersects Orgogozo's \textbackslash{}sisi material at Construction 7.0.4, sommet supérieur gauche du diagramme; preserve that human entry unchanged. The proposed reading D(\textbackslash{}sheafGamma(\textbackslash{}bar\{X\}),\textbackslash{}Aa) at Construction 7.0.4, sommet supérieur gauche du diagramme changes a substantive mathematical assertion or bibliographic identification; no independent derivation of this exact replacement is established in this review.\par
\par\noindent\textbf{OM-SOL-1199 -- t3 XVII p.464 -- SOL}\par
Construction 7.0.4, les deux flèches verticales du diagramme\par
\emph{Reported reading}: R e\textasciicircum{}\{+\}\_\{\textbackslash{}alpha\}\par
\emph{Proposal}: R e\textasciicircum{}\{*\}\_\{\textbackslash{}alpha\}\par
This finding intersects Orgogozo's \textbackslash{}sisi material at Construction 7.0.4, les deux flèches verticales du diagramme; preserve that human entry unchanged. The proposed reading R e\textasciicircum{}\{*\}\_\{\textbackslash{}alpha\} at Construction 7.0.4, les deux flèches verticales du diagramme changes a substantive mathematical assertion or bibliographic identification; no independent derivation of this exact replacement is established in this review.\par
\par\noindent\textbf{OM-SOL-1200 -- t3 XVII p.464 -- SOL}\par
Construction 7.0.5, isomorphisme canonique\par
\emph{Reported reading}: i\_\{!\}\textbackslash{}isomto R\textasciicircum{}\{+\}\textbackslash{}bar\{u\}\_\{f+\}\textbackslash{}circ j\_\{!\}\textbackslash{}circ L\textasciicircum{}\{+\}\textbackslash{}bar\{u\}\_\{f'\}\par
\emph{Proposal}: i\_\{!\}\textbackslash{}isomto R\textasciicircum{}\{+\}\textbackslash{}bar\{u\}\_\{f*\}\textbackslash{}circ j\_\{!\}\textbackslash{}circ L\textasciicircum{}\{+\}\textbackslash{}bar\{u\}\_\{f'\}\textasciicircum{}\{*\}\par
This finding intersects Orgogozo's \textbackslash{}sisi material at Construction 7.0.5, isomorphisme canonique; preserve that human entry unchanged. The replacement i\_\{!\}\textbackslash{}isomto R\textasciicircum{}\{+\}\textbackslash{}bar\{u\}\_\{f*\}\textbackslash{}circ j\_\{!\}\textbackslash{}circ L\textasciicircum{}\{+\}\textbackslash{}bar\{u\}\_\{f'\}\textasciicircum{}\{*\} at Construction 7.0.5, isomorphisme canonique identifies a particular indexed object or symbol; that identification is not independently established by the quoted local expression.\par
\par\noindent\textbf{OM-SOL-1201 -- t3 XVII p.465 -- SOL}\par
Phrase « Dès lors », premier facteur\par
\emph{Reported reading}: L\textasciicircum{}\{+\}\textbackslash{}bar\{u\}\_\{f\}\textasciicircum{}\{+\}\textbackslash{}circ i\_\{!\}\par
\emph{Proposal}: L\textasciicircum{}\{+\}\textbackslash{}bar\{u\}\_\{f\}\textasciicircum{}\{*\}\textbackslash{}circ i\_\{!\}\par
This finding intersects Orgogozo's \textbackslash{}sisi material at Phrase « Dès lors », premier facteur; preserve that human entry unchanged. The replacement L\textasciicircum{}\{+\}\textbackslash{}bar\{u\}\_\{f\}\textasciicircum{}\{*\}\textbackslash{}circ i\_\{!\} at Phrase « Dès lors », premier facteur identifies a particular indexed object or symbol; that identification is not independently established by the quoted local expression.\par
\par\noindent\textbf{OM-SOL-1202 -- t3 XVII p.466 -- SOL}\par
Diagramme (7.1.1.0), flèche horizontale\par
\emph{Reported reading}: La flèche X\_1 \ensuremath{\to} T\_1 n'est pas étiquetée.\par
\emph{Proposal}: X\_1\textbackslash{}xrightarrow\{j\_1\}T\_1\par
This finding intersects Orgogozo's \textbackslash{}sisi material at Diagramme (7.1.1.0), flèche horizontale; preserve that human entry unchanged. The proposed reading X\_1\textbackslash{}xrightarrow\{j\_1\}T\_1 at Diagramme (7.1.1.0), flèche horizontale changes a substantive mathematical assertion or bibliographic identification; no independent derivation of this exact replacement is established in this review.\par
\par\noindent\textbf{OM-SOL-1203 -- t3 XVII p.466 -- SOL}\par
7.1.2, définition de R\textasciicircum{}\{+(1)\}f\_!\par
\emph{Reported reading}: L\textasciicircum{}\{+\}\textbackslash{}bar\{u\}\_\{p\_\{1\}\}\textbackslash{}cdot j\_\{1!\}\textbackslash{}circ R\textasciicircum{}\{+\}\textbackslash{}bar\{u\}\_\{f\_\{1\}\}\par
\emph{Proposal}: R\textasciicircum{}\{+(1)\}f\_\{!\}=\allowbreak R\textasciicircum{}\{+\}\textbackslash{}bar\{u\}\_\{f\_\{1\}*\}\textbackslash{}circ j\_\{1!\}\textbackslash{}circ L\textasciicircum{}\{+\}\textbackslash{}bar\{u\}\_\{p\_\{1\}\}\textasciicircum{}\{*\}\par
The proposed reading R\textasciicircum{}\{+(1)\}f\_\{!\}=\allowbreak R\textasciicircum{}\{+\}\textbackslash{}bar\{u\}\_\{f\_\{1\}*\}\textbackslash{}circ j\_\{1!\}\textbackslash{}circ L\textasciicircum{}\{+\}\textbackslash{}bar\{u\}\_\{p\_\{1\}\}\textasciicircum{}\{*\} at 7.1.2, définition de R\textasciicircum{}\{+(1)\}f\_! changes a substantive mathematical assertion or bibliographic identification; no independent derivation of this exact replacement is established in this review.\par
\par\noindent\textbf{OM-SOL-1204 -- t3 XVII p.467 -- SOL}\par
Après le premier diagramme, référence au procédé de calcul\par
\emph{Reported reading}: VI b (2.3)\par
\emph{Proposal}: \textbackslash{}hyperrefext[Vbis][Vbis.subsec2.3]\{2.3\}\par
This finding intersects Orgogozo's \textbackslash{}sisi material at Après le premier diagramme, référence au procédé de calcul; preserve that human entry unchanged. The proposed reference \textbackslash{}hyperrefext[Vbis][Vbis.subsec2.3]\{2.3\} at Après le premier diagramme, référence au procédé de calcul needs comparison with the statement actually invoked; that target has not been independently established here.\par
\par\noindent\textbf{OM-SOL-1205 -- t3 XVII p.467 -- SOL}\par
Formule définissant \ensuremath{\alpha}\_1, dernier facteur\par
\emph{Reported reading}: L\textasciicircum{}\{+\}\textbackslash{}bar\{u\}\_\{p\_\{1\}p\_\{2\}\}\textasciicircum{}\{+\}\par
\emph{Proposal}: L\textasciicircum{}\{+\}\textbackslash{}bar\{u\_\{p\_\{1\}p\_\{2\}\}\}\textasciicircum{}\{\textbackslash{}ast\}\par
This finding intersects Orgogozo's \textbackslash{}sisi material at Formule définissant \ensuremath{\alpha}\_1, dernier facteur; preserve that human entry unchanged. The proposed reading L\textasciicircum{}\{+\}\textbackslash{}bar\{u\_\{p\_\{1\}p\_\{2\}\}\}\textasciicircum{}\{\textbackslash{}ast\} at Formule définissant \ensuremath{\alpha}\_1, dernier facteur changes a substantive mathematical assertion or bibliographic identification; no independent derivation of this exact replacement is established in this review.\par
\par\noindent\textbf{OM-SOL-1206 -- t3 XVII p.467 -- SOL}\par
7.1.3 c), formule définissant \ensuremath{\alpha}\_2\par
\emph{Reported reading}: R\textasciicircum{}\{+\}\textbackslash{}bar\{u\}\_\{f\_\{1\}f\_\{2\}+\}\par
\emph{Proposal}: R\textasciicircum{}\{+\}\textbackslash{}bar\{u\_\{f\_\{1\}f\_\{2\}\}\}\_\{*\}\par
This finding intersects Orgogozo's \textbackslash{}sisi material at 7.1.3 c), formule définissant \ensuremath{\alpha}\_2; preserve that human entry unchanged. The replacement R\textasciicircum{}\{+\}\textbackslash{}bar\{u\_\{f\_\{1\}f\_\{2\}\}\}\_\{*\} at 7.1.3 c), formule définissant \ensuremath{\alpha}\_2 identifies a particular indexed object or symbol; that identification is not independently established by the quoted local expression.\par
\par\noindent\textbf{OM-SOL-1207 -- t3 XVII p.468 -- SOL}\par
Diagramme initial, sommet inférieur gauche\par
\emph{Reported reading}: R\textasciicircum{}\{+(2)\}f\_!\par
\emph{Proposal}: R\textasciicircum{}\{+(3)\}f\_!\par
This finding intersects Orgogozo's \textbackslash{}sisi material at Diagramme initial, sommet inférieur gauche; preserve that human entry unchanged. The replacement R\textasciicircum{}\{+(3)\}f\_! at Diagramme initial, sommet inférieur gauche identifies a particular indexed object or symbol; that identification is not independently established by the quoted local expression.\par
\par\noindent\textbf{OM-SOL-1208 -- t3 XVII p.468 -- SOL}\par
Diagramme initial, deux flèches vers le sommet inférieur R\textasciicircum{}\{+(1,3)\}f\_!\par
\emph{Reported reading}: Les deux flèches obliques venant des sommets latéraux inférieurs sont absentes.\par
\emph{Proposal}: R\textasciicircum{}\{+(3)\}f\_\{!\}\textbackslash{}longrightarrow R\textasciicircum{}\{+(1,3)\}f\_\{!\}\textbackslash{}longleftarrow R\textasciicircum{}\{+(1)\}f\_\{!\}\par
This finding intersects Orgogozo's \textbackslash{}sisi material at Diagramme initial, deux flèches vers le sommet inférieur R\textasciicircum{}\{+(1,3)\}f\_!; preserve that human entry unchanged. The replacement R\textasciicircum{}\{+(3)\}f\_\{!\}\textbackslash{}longrightarrow R\textasciicircum{}\{+(1,3)\}f\_\{!\}\textbackslash{}longleftarrow R\textasciicircum{}\{+(1)\}f\_\{!\} at Diagramme initial, deux flèches vers le sommet inférieur R\textasciicircum{}\{+(1,3)\}f\_! requires a complete check of the maps, their endpoints and the composition convention; this identification is not independently verified.\par
\par\noindent\textbf{OM-SOL-1209 -- t3 XVII p.468 -- SOL}\par
Construction 7.1.4, isomorphisme annoncé\par
\emph{Reported reading}: R\textasciicircum{}\{+(3)\}fg\_\{!\}\textbackslash{}isomfrom R\textasciicircum{}\{+(2)\}g\_\{!\}\textbackslash{}circ R\textasciicircum{}\{+(1)\}f\_\{!\}\par
\emph{Proposal}: R\textasciicircum{}\{+(3)\}(fg)\_\{!\}\textbackslash{}isomfrom R\textasciicircum{}\{+(1)\}f\_\{!\}\textbackslash{}circ R\textasciicircum{}\{+(2)\}g\_\{!\}\par
The proposed reading R\textasciicircum{}\{+(3)\}(fg)\_\{!\}\textbackslash{}isomfrom R\textasciicircum{}\{+(1)\}f\_\{!\}\textbackslash{}circ R\textasciicircum{}\{+(2)\}g\_\{!\} at Construction 7.1.4, isomorphisme annoncé changes a substantive mathematical assertion or bibliographic identification; no independent derivation of this exact replacement is established in this review.\par
\par\noindent\textbf{OM-SOL-1210 -- t3 XVII p.469 -- SOL}\par
Après (7.1.4.1.3), « on notera »\par
\emph{Reported reading}: R\textasciicircum{}\{+(12)\}\textbackslash{}int g\_\{!\}\par
\emph{Proposal}: R\textasciicircum{}\{+(12)\}(fg)\_\{!\}\par
This finding intersects Orgogozo's \textbackslash{}sisi material at Après (7.1.4.1.3), « on notera »; preserve that human entry unchanged. The proposed reading R\textasciicircum{}\{+(12)\}(fg)\_\{!\} at Après (7.1.4.1.3), « on notera » changes a substantive mathematical assertion or bibliographic identification; no independent derivation of this exact replacement is established in this review.\par
\par\noindent\textbf{OM-SOL-1211 -- t3 XVII p.469 -- SOL}\par
7.1.4.2, membre gauche de l'isomorphisme\par
\emph{Reported reading}: R\textasciicircum{}\{+(12)\}\textbackslash{}int g\_\{!\}\par
\emph{Proposal}: R\textasciicircum{}\{+(12)\}(fg)\_\{!\}\par
This finding intersects Orgogozo's \textbackslash{}sisi material at 7.1.4.2, membre gauche de l'isomorphisme; preserve that human entry unchanged. The proposed reading R\textasciicircum{}\{+(12)\}(fg)\_\{!\} at 7.1.4.2, membre gauche de l'isomorphisme changes a substantive mathematical assertion or bibliographic identification; no independent derivation of this exact replacement is established in this review.\par
\par\noindent\textbf{OM-SOL-1212 -- t3 XVII p.470 -- SOL}\par
Lemme 7.1.4.2.2, hypothèse sur j et j\_1\par
\emph{Reported reading}: j, j\_1 sont les immersions\par
\emph{Proposal}: \$j\$, \$j\_\{1\}\$ sont des immersions.\par
The lemma calls j and j\_1 les immersions after displaying their diagram; changing the definite article to des is plausible, but the supplied wording does not independently force that editorial choice.\par
\par\noindent\textbf{OM-SOL-1213 -- t3 XVII p.470 -- SOL}\par
Lemme 7.1.4.2.2, cible de la flèche canonique\par
\emph{Reported reading}: R\textasciicircum{}\{+(1)\}(hh')\_\{!\}\par
\emph{Proposal}: R\textasciicircum{}\{+(1)\}f\_\{!\}\par
This finding intersects Orgogozo's \textbackslash{}sisi material at Lemme 7.1.4.2.2, cible de la flèche canonique; preserve that human entry unchanged. The replacement R\textasciicircum{}\{+(1)\}f\_\{!\} at Lemme 7.1.4.2.2, cible de la flèche canonique requires a complete check of the maps, their endpoints and the composition convention; this identification is not independently verified.\par
\par\noindent\textbf{OM-SOL-1214 -- t3 XVII p.471 -- SOL}\par
Chaîne d'isomorphismes, membre initial\par
\emph{Reported reading}: R\textasciicircum{}\{+(3)\}fg\_\{!\}\par
\emph{Proposal}: R\textasciicircum{}\{+(12)\}(fg)\_\{!\}\par
This finding intersects Orgogozo's \textbackslash{}sisi material at Chaîne d'isomorphismes, membre initial; preserve that human entry unchanged. The replacement R\textasciicircum{}\{+(12)\}(fg)\_\{!\} at Chaîne d'isomorphismes, membre initial identifies a particular indexed object or symbol; that identification is not independently established by the quoted local expression.\par
\par\noindent\textbf{OM-SOL-1215 -- t3 XVII p.471 -- SOL}\par
Chaîne d'isomorphismes, première ligne, après j'\_\{1!\}\par
\emph{Reported reading}: j'\_\{1!\}\textbackslash{}circ j'\_\{2!\}\par
\emph{Proposal}: j'\_\{1!\}\textbackslash{}circ j''\_\{2!\}\par
This finding intersects Orgogozo's \textbackslash{}sisi material at Chaîne d'isomorphismes, première ligne, après j'\_\{1!\}; preserve that human entry unchanged. The replacement j'\_\{1!\}\textbackslash{}circ j''\_\{2!\} at Chaîne d'isomorphismes, première ligne, après j'\_\{1!\} identifies a particular indexed object or symbol; that identification is not independently established by the quoted local expression.\par
\par\noindent\textbf{OM-SOL-1216 -- t3 XVII p.471 -- SOL}\par
Chaîne d'isomorphismes, deuxième et troisième lignes, termes en g'\_2\par
\emph{Reported reading}: R\textasciicircum{}\{+\}g'\_\{2\} ; R\textasciicircum{}\{+\}g'\_\{2+\}\par
\emph{Proposal}: R\textasciicircum{}\{+\}(g'\_2p)\_*\par
This finding intersects Orgogozo's \textbackslash{}sisi material at Chaîne d'isomorphismes, deuxième et troisième lignes, termes en g'\_2; preserve that human entry unchanged. The proposed reading R\textasciicircum{}\{+\}(g'\_2p)\_* at Chaîne d'isomorphismes, deuxième et troisième lignes, termes en g'\_2 changes a substantive mathematical assertion or bibliographic identification; no independent derivation of this exact replacement is established in this review.\par
\par\noindent\textbf{OM-SOL-1217 -- t3 XVII p.471 -- SOL}\par
Chaîne d'isomorphismes : première ligne pour u\_\{f\_1\}, troisième et quatrième lignes pour r\par
\emph{Reported reading}: R\textasciicircum{}\{+\}\textbackslash{}bar\{u\}\_\{f\_1\} ; L\textasciicircum{}\{+\}\textbackslash{}bar\{r\} ; R\textasciicircum{}\{+\}\textbackslash{}bar\{r\}\par
\emph{Proposal}: R\textasciicircum{}\{+\}\textbackslash{}bar\{u\}\_\{f\_\{1\}*\},\textbackslash{}qquad L\textasciicircum{}\{+\}\textbackslash{}bar\{r\}\textasciicircum{}\{*\},\textbackslash{}qquad R\textasciicircum{}\{+\}\textbackslash{}bar\{r\}\_*\par
This finding intersects Orgogozo's \textbackslash{}sisi material at Chaîne d'isomorphismes : première ligne pour u\_\{f\_1\}, troisième et quatrième lignes pour r; preserve that human entry unchanged. The proposed reading R\textasciicircum{}\{+\}\textbackslash{}bar\{u\}\_\{f\_\{1\}*\},\textbackslash{}qquad L\textasciicircum{}\{+\}\textbackslash{}bar\{r\}\textasciicircum{}\{*\},\textbackslash{}qquad R\textasciicircum{}\{+\}\textbackslash{}bar\{r\}\_* at Chaîne d'isomorphismes : première ligne pour u\_\{f\_1\}, troisième et quatrième lignes pour r changes a substantive mathematical assertion or bibliographic identification; no independent derivation of this exact replacement is established in this review.\par
\par\noindent\textbf{OM-SOL-1218 -- t3 XVII p.471 -- SOL}\par
Notation récurrente des images directes et inverses, également p. 474\par
\emph{Reported reading}: R\textasciicircum{}\{+\}\textbackslash{}bar\{u\}\_\{f\_\{1\}+\} ; L\textasciicircum{}\{+\}\textbackslash{}bar\{r\}\textasciicircum{}\{+\} ; \textbackslash{}pi\_\{1\}\textasciicircum{}\{p+\} (p. 474)\par
\emph{Proposal}: R\textasciicircum{}\{+\}\textbackslash{}bar\{u\}\_\{f\_\{1\}*\},\textbackslash{}qquad L\textasciicircum{}\{+\}\textbackslash{}bar\{r\}\textasciicircum{}\{*\},\textbackslash{}qquad \textbackslash{}pi\_\{1\}\textasciicircum{}\{p*\}\par
This finding intersects Orgogozo's \textbackslash{}sisi material at Notation récurrente des images directes et inverses, également p. 474; preserve that human entry unchanged. The proposed reading R\textasciicircum{}\{+\}\textbackslash{}bar\{u\}\_\{f\_\{1\}*\},\textbackslash{}qquad L\textasciicircum{}\{+\}\textbackslash{}bar\{r\}\textasciicircum{}\{*\},\textbackslash{}qquad \textbackslash{}pi\_\{1\}\textasciicircum{}\{p*\} at Notation récurrente des images directes et inverses, également p. 474 changes a substantive mathematical assertion or bibliographic identification; no independent derivation of this exact replacement is established in this review.\par
\par\noindent\textbf{OM-SOL-1219 -- t3 XVII p.472 -- SOL}\par
7.1.4.4, référence après «et un diagramme»\par
\emph{Reported reading}: (7.1.4.3)\par
\emph{Proposal}: \textbackslash{}eqhyperref[XVII.eq7.1.4.1.3]\{7.1.4.1.3\}\par
This finding intersects Orgogozo's \textbackslash{}sisi material at 7.1.4.4, référence après «et un diagramme»; preserve that human entry unchanged. The proposed reference \textbackslash{}eqhyperref[XVII.eq7.1.4.1.3]\{7.1.4.1.3\} at 7.1.4.4, référence après «et un diagramme» needs comparison with the statement actually invoked; that target has not been independently established here.\par
\par\noindent\textbf{OM-SOL-1220 -- t3 XVII p.472 -- SOL}\par
Après le premier carré, référence après «du type»\par
\emph{Reported reading}: (7.1.4.4)\par
\emph{Proposal}: \textbackslash{}eqhyperref[XVII.eq7.1.4.2.1]\{7.1.4.2.1\}\par
This finding intersects Orgogozo's \textbackslash{}sisi material at Après le premier carré, référence après «du type»; preserve that human entry unchanged. The proposed reference \textbackslash{}eqhyperref[XVII.eq7.1.4.2.1]\{7.1.4.2.1\} at Après le premier carré, référence après «du type» needs comparison with the statement actually invoked; that target has not been independently established here.\par
\par\noindent\textbf{OM-SOL-1222 -- t3 XVII p.473 -- SOL}\par
Démonstration de 7.1.6, identification dans la première phrase\par
\emph{Reported reading}: R\textasciicircum{}\{+\}fg\_\{!\}\textbackslash{}longrightarrow Rf\_\{!\}\textbackslash{}circ Rg\_\{!\}\par
\emph{Proposal}: R\textasciicircum{}\{+\}fg\_\{!\}\textbackslash{}xrightarrow\{c\_\{fg\}\}R\textasciicircum{}\{+\}f\_\{!\}\textbackslash{}circ R\textasciicircum{}\{+\}g\_\{!\}\par
The replacement R\textasciicircum{}\{+\}fg\_\{!\}\textbackslash{}xrightarrow\{c\_\{fg\}\}R\textasciicircum{}\{+\}f\_\{!\}\textbackslash{}circ R\textasciicircum{}\{+\}g\_\{!\} at Démonstration de 7.1.6, identification dans la première phrase identifies a particular indexed object or symbol; that identification is not independently established by the quoted local expression.\par
\par\noindent\textbf{OM-SOL-1223 -- t3 XVII p.473 -- SOL}\par
Diagramme (7.1.6.1), sommet gauche de la deuxième ligne\par
\emph{Reported reading}: T\_3, alors qu'un autre sommet T\_3 figure déjà sur la première ligne\par
\emph{Proposal}: T\_6\par
The proposed reading T\_6 at Diagramme (7.1.6.1), sommet gauche de la deuxième ligne changes a substantive mathematical assertion or bibliographic identification; no independent derivation of this exact replacement is established in this review.\par
\par\noindent\textbf{OM-SOL-1224 -- t3 XVII p.473 -- SOL}\par
Dernière phrase, référence après «du type»\par
\emph{Reported reading}: (7.1.4.4)\par
\emph{Proposal}: \textbackslash{}eqhyperref[XVII.eq7.1.4.2.1]\{7.1.4.2.1\}\par
The proposed reference \textbackslash{}eqhyperref[XVII.eq7.1.4.2.1]\{7.1.4.2.1\} at Dernière phrase, référence après «du type» needs comparison with the statement actually invoked; that target has not been independently established here.\par
\par\noindent\textbf{OM-SOL-1225 -- t3 XVII p.474 -- SOL}\par
Parenthèse après (7.1.9.2), indice de bar u\par
\emph{Reported reading}: L\textasciicircum{}\{+\}\textbackslash{}bar\{u\}\_\{p\_1\}\textasciicircum{}\{+\}(K\textasciicircum{}\{\textbackslash{}bigdot\})\par
\emph{Proposal}: L\textasciicircum{}\{+\}\textbackslash{}bar\{u\}\_\{\textbackslash{}pi\_\{1\}\}\textasciicircum{}\{*\}(K\textasciicircum{}\{\textbackslash{}bigdot\})\par
This finding intersects Orgogozo's \textbackslash{}sisi material at Parenthèse après (7.1.9.2), indice de bar u; preserve that human entry unchanged. The proposed reading L\textasciicircum{}\{+\}\textbackslash{}bar\{u\}\_\{\textbackslash{}pi\_\{1\}\}\textasciicircum{}\{*\}(K\textasciicircum{}\{\textbackslash{}bigdot\}) at Parenthèse après (7.1.9.2), indice de bar u changes a substantive mathematical assertion or bibliographic identification; no independent derivation of this exact replacement is established in this review.\par
\par\noindent\textbf{OM-SOL-1226 -- t3 XVII p.475 -- SOL}\par
Définition de K\textasciicircum{}n\_\{\textbackslash{}mathcal U\}, indice de la somme\par
\emph{Reported reading}: |P|=\allowbreak n+1\par
\emph{Proposal}: |P|=\allowbreak -n+1\par
This finding intersects Orgogozo's \textbackslash{}sisi material at Définition de K\textasciicircum{}n\_\{\textbackslash{}mathcal U\}, indice de la somme; preserve that human entry unchanged. The proposed reading |P|=\allowbreak -n+1 at Définition de K\textasciicircum{}n\_\{\textbackslash{}mathcal U\}, indice de la somme changes a substantive mathematical assertion or bibliographic identification; no independent derivation of this exact replacement is established in this review.\par
\par\noindent\textbf{OM-SOL-1227 -- t3 XVII p.475 -- SOL}\par
Suite spectrale de la démonstration de 7.2.1, cardinalité de P\par
\emph{Reported reading}: |P|=\allowbreak p+1\par
\emph{Proposal}: |P|=\allowbreak -p+1\par
This finding intersects Orgogozo's \textbackslash{}sisi material at Suite spectrale de la démonstration de 7.2.1, cardinalité de P; preserve that human entry unchanged. The proposed reading |P|=\allowbreak -p+1 at Suite spectrale de la démonstration de 7.2.1, cardinalité de P changes a substantive mathematical assertion or bibliographic identification; no independent derivation of this exact replacement is established in this review.\par
\par\noindent\textbf{OM-SOL-1228 -- t3 XVII p.475 -- SOL}\par
Suite spectrale, exposant du foncteur appliqué à j\_P\textasciicircum{}*(F)\par
\emph{Reported reading}: R\textasciicircum{}\{p\}(fj\_P)\_\{!\}j\_P\textasciicircum{}\{*\}(F)\par
\emph{Proposal}: R\textasciicircum{}\{q\}(fj\_\{P\})\_\{!\}j\textasciicircum{}\{\textbackslash{}ast\}\_\{P\}(F)\par
The replacement R\textasciicircum{}\{q\}(fj\_\{P\})\_\{!\}j\textasciicircum{}\{\textbackslash{}ast\}\_\{P\}(F) at Suite spectrale, exposant du foncteur appliqué à j\_P\textasciicircum{}*(F) identifies a particular indexed object or symbol; that identification is not independently established by the quoted local expression.\par
\par\noindent\textbf{OM-SOL-1229 -- t3 XVII p.475 -- SOL}\par
Référence aux conventions de complexes doubles\par
\emph{Reported reading}: VI b.(2.3.2)\par
\emph{Proposal}: \textbackslash{}hyperrefext[Vbis][Vbis.subsubsec2.3.2]\{2.3.2\}\par
This finding intersects Orgogozo's \textbackslash{}sisi material at Référence aux conventions de complexes doubles; preserve that human entry unchanged. The proposed reading \textbackslash{}hyperrefext[Vbis][Vbis.subsubsec2.3.2]\{2.3.2\} at Référence aux conventions de complexes doubles changes a substantive mathematical assertion or bibliographic identification; no independent derivation of this exact replacement is established in this review.\par
\par\noindent\textbf{OM-SOL-1230 -- t3 XVII p.476 -- SOL}\par
Triangle distingué, sens des trois flèches\par
\emph{Reported reading}: Le cycle imprimé va de K\textasciicircum{}\{n+1,\ensuremath{\cdot}\} à \ensuremath{\sigma}\ensuremath{\leq}n, puis à \ensuremath{\sigma}\ensuremath{\leq}n+1, puis à K\textasciicircum{}\{n+1,\ensuremath{\cdot}\}.\par
\emph{Proposal}: \textbackslash{}xymatrix@C=\allowbreak 1em\{ \& K\textasciicircum{}\{n+1,\textbackslash{}bigdot\}\_\{\textbackslash{}Uu\}\textbackslash{}ar[rd] \& \textbackslash{}\textbackslash{} (\textbackslash{}sigma\_\{\textbackslash{}leq n\}(K\textasciicircum{}\{\textbackslash{}bigdot\textbackslash{}bigdot\}\_\{\textbackslash{}Uu\}))\_s\textbackslash{}ar[ur] \&\& (\textbackslash{}sigma\_\{\textbackslash{}leq n+1\}(K\textasciicircum{}\{\textbackslash{}bigdot\textbackslash{}bigdot\}\_\{\textbackslash{}Uu\}))\_s\textbackslash{}ar[ll] \}\par
This finding intersects Orgogozo's \textbackslash{}sisi material at Triangle distingué, sens des trois flèches; preserve that human entry unchanged. The replacement \textbackslash{}xymatrix@C=\allowbreak 1em\{ \& K\textasciicircum{}\{n+1,\textbackslash{}bigdot\}\_\{\textbackslash{}Uu\}\textbackslash{}ar[rd] \& \textbackslash{}\textbackslash{} (\textbackslash{}sigma\_\{\textbackslash{}leq n\}(K\textasciicircum{}\{\textbackslash{}bigdot\textbackslash{}bigdot\}\_\{\textbackslash{}Uu\}))\_s\textbackslash{}ar[ur] \&\& (\textbackslash{}sigma\_\{\textbackslash{}leq n+1\}(K\textasciicircum{}\{\textbackslash{}bigdot\textbackslash{}bigdot\}\_\{\textbackslash{}Uu\}))\_s\textbackslash{}ar[ll] \} at Triangle distingué, sens des trois flèches requires a complete check of the maps, their endpoints and the composition convention; this identification is not independently verified.\par
\par\noindent\textbf{OM-SOL-1231 -- t3 XVII p.476 -- SOL}\par
Remarque 7.2.3, renvoi final\par
\emph{Reported reading}: (cf. (3.7))\par
\emph{Proposal}: (7.3.7)\par
This finding intersects Orgogozo's \textbackslash{}sisi material at Remarque 7.2.3, renvoi final; preserve that human entry unchanged. The proposed reference (7.3.7) at Remarque 7.2.3, renvoi final needs comparison with the statement actually invoked; that target has not been independently established here.\par
\par\noindent\textbf{OM-SOL-1232 -- t3 XVII p.477 -- SOL}\par
Carré de suites spectrales, second \ensuremath{\pi} du terme inférieur gauche\par
\emph{Reported reading}: \textbackslash{}pi'\textasciicircum{}\{p\}\_1(g'\textasciicircum{}\{\textbackslash{}ast\}(K\textasciicircum{}\{\textbackslash{}bigdot\}))\par
\emph{Proposal}: \textbackslash{}pi'\textasciicircum{}\{p\textbackslash{}ast\}\_1(g'\textasciicircum{}\{\textbackslash{}ast\}(K\textasciicircum{}\{\textbackslash{}bigdot\}))\par
The proposed reading \textbackslash{}pi'\textasciicircum{}\{p\textbackslash{}ast\}\_1(g'\textasciicircum{}\{\textbackslash{}ast\}(K\textasciicircum{}\{\textbackslash{}bigdot\})) at Carré de suites spectrales, second \ensuremath{\pi} du terme inférieur gauche changes a substantive mathematical assertion or bibliographic identification; no independent derivation of this exact replacement is established in this review.\par
\par\noindent\textbf{OM-SOL-1233 -- t3 XVII p.477 -- SOL}\par
Carré de suites spectrales, exposant de R dans le terme inférieur droit\par
\emph{Reported reading}: R\textasciicircum{}\{+\}f'\_!\par
\emph{Proposal}: R\textasciicircum{}\{p+q\}f'\_!\par
The proposed reading R\textasciicircum{}\{p+q\}f'\_! at Carré de suites spectrales, exposant de R dans le terme inférieur droit changes a substantive mathematical assertion or bibliographic identification; no independent derivation of this exact replacement is established in this review.\par
\par\noindent\textbf{OM-SOL-1234 -- t3 XVII p.478 -- SOL}\par
Formule (7.3.1.1), premier différentiel\par
\emph{Reported reading}: \textbackslash{}Xx\textasciicircum{}\{0\}(F)\textbackslash{}xrightarrow\{S\textasciicircum{}\{1\}(F)\}\textbackslash{}Xx\textasciicircum{}\{1\}(F)\par
\emph{Proposal}: \textbackslash{}Xx\textasciicircum{}\{0\}(F)\textbackslash{}xrightarrow\{S\textasciicircum{}\{0\}(F)\}\textbackslash{}Xx\textasciicircum{}\{1\}(F)\par
The replacement \textbackslash{}Xx\textasciicircum{}\{0\}(F)\textbackslash{}xrightarrow\{S\textasciicircum{}\{0\}(F)\}\textbackslash{}Xx\textasciicircum{}\{1\}(F) at Formule (7.3.1.1), premier différentiel identifies a particular indexed object or symbol; that identification is not independently established by the quoted local expression.\par
\par\noindent\textbf{OM-SOL-1235 -- t3 XVII p.478 -- SOL}\par
Démonstration du lemme 7.3.2, première colonne du diagramme\par
\emph{Reported reading}: \textbackslash{}Ker S\textasciicircum{}\{2d\}(F'), \textbackslash{}Ker S\textasciicircum{}\{2d\}(F), \textbackslash{}Ker S\textasciicircum{}\{2d\}(F'')\par
\emph{Proposal}: \textbackslash{}Ker S\textasciicircum{}\{2d+1\}(F'),\textbackslash{}quad \textbackslash{}Ker S\textasciicircum{}\{2d+1\}(F),\textbackslash{}quad \textbackslash{}Ker S\textasciicircum{}\{2d+1\}(F'')\par
The replacement \textbackslash{}Ker S\textasciicircum{}\{2d+1\}(F'),\textbackslash{}quad \textbackslash{}Ker S\textasciicircum{}\{2d+1\}(F),\textbackslash{}quad \textbackslash{}Ker S\textasciicircum{}\{2d+1\}(F'') at Démonstration du lemme 7.3.2, première colonne du diagramme identifies a particular indexed object or symbol; that identification is not independently established by the quoted local expression.\par
\par\noindent\textbf{OM-SOL-1236 -- t3 XVII p.478 -- SOL}\par
Diagramme de la démonstration du lemme 7.3.2, case inférieure droite\par
\emph{Reported reading}: \textbackslash{}Xx\textasciicircum{}\{2d+2\}(F''')\par
\emph{Proposal}: \textbackslash{}Xx\textasciicircum{}\{2d+2\}(F'')\par
The third prime requires identifying the row's exact-sequence object throughout the diagram.\par
\par\noindent\textbf{OM-SOL-1237 -- t3 XVII p.479 -- SOL}\par
Démonstration du lemme 7.3.3, renvoi à la Remarque 2.3.2.2\par
\emph{Reported reading}: VI b (2.3.2.2). Remarque\par
\emph{Proposal}: V bis (2.3.2.2), Remarque\par
This finding intersects Orgogozo's \textbackslash{}sisi material at Démonstration du lemme 7.3.3, renvoi à la Remarque 2.3.2.2; preserve that human entry unchanged. The proposed reference V bis (2.3.2.2), Remarque at Démonstration du lemme 7.3.3, renvoi à la Remarque 2.3.2.2 needs comparison with the statement actually invoked; that target has not been independently established here.\par
\par\noindent\textbf{OM-SOL-1238 -- t3 XVII p.479 -- SOL}\par
7.3.5, renvoi après « grâce à »\par
\emph{Reported reading}: 7.2.3\par
\emph{Proposal}: 7.3.3\par
This finding intersects Orgogozo's \textbackslash{}sisi material at 7.3.5, renvoi après « grâce à »; preserve that human entry unchanged. The proposed reference 7.3.3 at 7.3.5, renvoi après « grâce à » needs comparison with the statement actually invoked; that target has not been independently established here.\par
\par\noindent\textbf{OM-SOL-1239 -- t3 XVII p.479 -- SOL}\par
7.3.5, renvoi après « En vertu de »\par
\emph{Reported reading}: 7.2.4\par
\emph{Proposal}: 7.3.4\par
This finding intersects Orgogozo's \textbackslash{}sisi material at 7.3.5, renvoi après « En vertu de »; preserve that human entry unchanged. The proposed reference 7.3.4 at 7.3.5, renvoi après « En vertu de » needs comparison with the statement actually invoked; that target has not been independently established here.\par
\par\noindent\textbf{OM-SOL-1240 -- t3 XVII p.480 -- SOL}\par
Bibliographie, Gabriel–Zisman, titre\par
\emph{Reported reading}: Calculus of fraction and homotopy theory\par
\emph{Proposal}: Calculus of fractions and homotopy theory\par
The cited book title needs verification against the publication.\par
\par\noindent\textbf{OM-SOL-1241 -- t3 XVIII p.488 -- SOL}\par
Paragraph beginning « Soient les courbes complètes »\par
\emph{Reported reading}: courbes complètes non lisses\par
\emph{Proposal}: lisses\par
This finding intersects Orgogozo's \textbackslash{}sisi material at Paragraph beginning « Soient les courbes complètes »; preserve that human entry unchanged. The proposed reading lisses at Paragraph beginning « Soient les courbes complètes » changes a substantive mathematical assertion or bibliographic identification; no independent derivation of this exact replacement is established in this review.\par
\par\noindent\textbf{OM-SOL-1242 -- t3 XVIII p.489 -- SOL}\par
Opening Picard/cohomology diagram\par
\emph{Reported reading}: \textbackslash{}Pic(X)/n, H\textasciicircum{}2\_c(X,\textbackslash{}ZZ/n(1)); \textbackslash{}Pic(Y)/n, H\textasciicircum{}2\_c(Y,\textbackslash{}ZZ/n(1)); \textbackslash{}Tr\_u\par
\emph{Proposal}: \textbackslash{}Pic(\textbackslash{}bar\{X\})/n, H\textasciicircum{}2\_c(\textbackslash{}bar\{X\},\textbackslash{}ZZ/n(1)); \textbackslash{}Pic(\textbackslash{}bar\{Y\})/n, H\textasciicircum{}2\_c(\textbackslash{}bar\{Y\},\textbackslash{}ZZ/n(1)); \textbackslash{}Tr\_\{\textbackslash{}bar\{u\}\}\par
The replacement \textbackslash{}Pic(\textbackslash{}bar\{X\})/n, H\textasciicircum{}2\_c(\textbackslash{}bar\{X\},\textbackslash{}ZZ/n(1)); \textbackslash{}Pic(\textbackslash{}bar\{Y\})/n, H\textasciicircum{}2\_c(\textbackslash{}bar\{Y\},\textbackslash{}ZZ/n(1)); \textbackslash{}Tr\_\{\textbackslash{}bar\{u\}\} at Opening Picard/cohomology diagram requires a complete check of the maps, their endpoints and the composition convention; this identification is not independently verified.\par
\par\noindent\textbf{OM-SOL-1243 -- t3 XVIII p.489 -- SOL}\par
Reference concerning compatibility of the norm\par
\emph{Reported reading}: EGA IV 21.10.7\par
\emph{Proposal}: \textbackslash{}EGA IV 21.10.17\par
This finding intersects Orgogozo's \textbackslash{}sisi material at Reference concerning compatibility of the norm; preserve that human entry unchanged. The proposed reference \textbackslash{}EGA IV 21.10.17 at Reference concerning compatibility of the norm needs comparison with the statement actually invoked; that target has not been independently established here.\par
\par\noindent\textbf{OM-SOL-1244 -- t3 XVIII p.489 -- SOL}\par
Paragraph beginning « Soit p », morphism obtained from Kummer theory\par
\emph{Reported reading}: \textbackslash{}ZZ/n\textbackslash{}longrightarrow R\textasciicircum{}\{1\}p\_\{\textbackslash{}ast\}\textbackslash{}ZZ/n(1)\par
\emph{Proposal}: \textbackslash{}ZZ/n\textbackslash{}longrightarrow R\textasciicircum{}\{2\}p\_\{\textbackslash{}ast\}\textbackslash{}ZZ/n(1)\par
The replacement \textbackslash{}ZZ/n\textbackslash{}longrightarrow R\textasciicircum{}\{2\}p\_\{\textbackslash{}ast\}\textbackslash{}ZZ/n(1) at Paragraph beginning « Soit p », morphism obtained from Kummer theory requires a complete check of the maps, their endpoints and the composition convention; this identification is not independently verified.\par
\par\noindent\textbf{OM-SOL-1245 -- t3 XVIII p.489 -- SOL}\par
Last paragraph, target of the quasi-finite flat morphism u\par
\emph{Reported reading}: \textbackslash{}PP'\_S\par
\emph{Proposal}: \textbackslash{}PP\textasciicircum{}1\_S\par
The proposed reading \textbackslash{}PP\textasciicircum{}1\_S involves a printed glyph or a missing stroke at Last paragraph, target of the quasi-finite flat morphism u; the textual record does not independently establish that glyph.\par
\par\noindent\textbf{OM-SOL-1246 -- t3 XVIII p.490 -- SOL}\par
Paragraph beginning « Pour un X général », morphism induced by the dense open immersion\par
\emph{Reported reading}: R\textasciicircum{}\{2\}(fj\_\{!\})\textbackslash{}ZZ/n(1)\par
\emph{Proposal}: R\textasciicircum{}\{2\}(fj)\_\{!\}\textbackslash{}ZZ/n(1)\par
The proposed reading R\textasciicircum{}\{2\}(fj)\_\{!\}\textbackslash{}ZZ/n(1) at Paragraph beginning « Pour un X général », morphism induced by the dense open immersion changes a substantive mathematical assertion or bibliographic identification; no independent derivation of this exact replacement is established in this review.\par
\par\noindent\textbf{OM-SOL-1248 -- t3 XVIII p.496 -- SOL}\par
1.2.7, condition sur \ensuremath{\phi}\par
\emph{Reported reading}: \textbackslash{}varphi\textbackslash{}in\textbackslash{}Hom(U,G)\par
\emph{Proposal}: \textbackslash{}varphi\textbackslash{}in \textbackslash{}Hom(U\_\{m\},G)\par
The replacement \textbackslash{}varphi\textbackslash{}in \textbackslash{}Hom(U\_\{m\},G) at 1.2.7, condition sur \ensuremath{\phi} identifies a particular indexed object or symbol; that identification is not independently established by the quoted local expression.\par
\par\noindent\textbf{OM-SOL-1250 -- t3 XVIII p.500 -- SOL}\par
Variantes 1.2.15, première explication après l'énumération\par
\emph{Reported reading}: Dans le cas (i), les flèches […] sont des isomorphismes d'après le théorème de changement de base\par
\emph{Proposal}: Dans le cas (ii)\par
The proposed reference Dans le cas (ii) at Variantes 1.2.15, première explication après l'énumération needs comparison with the statement actually invoked; that target has not been independently established here.\par
\par\noindent\textbf{OM-SOL-1251 -- t3 XVIII p.500 -- SOL}\par
Variantes 1.2.15, seconde explication après l'énumération\par
\emph{Reported reading}: Dans le cas (ii), que les flèches […] soient des isomorphismes est standard […] pour la topologie de Zariski\par
\emph{Proposal}: Dans le cas (i)\par
The proposed reading Dans le cas (i) at Variantes 1.2.15, seconde explication après l'énumération changes a substantive mathematical assertion or bibliographic identification; no independent derivation of this exact replacement is established in this review.\par
\par\noindent\textbf{OM-SOL-1253 -- t3 XVIII p.501 -- SOL}\par
1.3.1.2 (c), référence finale\par
\emph{Reported reading}: XVII 6.3.3.1\par
\emph{Proposal}: XVII\textasciitilde{}6.3.23.1\par
The proposed reference XVII\textasciitilde{}6.3.23.1 at 1.3.1.2 (c), référence finale needs comparison with the statement actually invoked; that target has not been independently established here.\par
\par\noindent\textbf{OM-SOL-1254 -- t3 XVIII p.503 -- SOL}\par
Calcul après « se réécrit », début de la dernière ligne\par
\emph{Reported reading}: p\textasciicircum{}\{\textbackslash{}ast\}\_\{12\}\textbackslash{}langle\textbackslash{}Ll\_\{1\},K]+p\textasciicircum{}\{\textbackslash{}ast\}\_\{12\}\textbackslash{}langle\textbackslash{}Ll\_\{2\},K]+\textbackslash{}cdots\par
\emph{Proposal}: =\allowbreak p\textasciicircum{}\{\textbackslash{}ast\}\_\{12\}\textbackslash{}langle\textbackslash{}Ll\_\{1\},K]+p\textasciicircum{}\{\textbackslash{}ast\}\_\{12\}\textbackslash{}langle\textbackslash{}Ll\_\{2\},K]+\textbackslash{}cdots\par
The proposed reading =\allowbreak p\textasciicircum{}\{\textbackslash{}ast\}\_\{12\}\textbackslash{}langle\textbackslash{}Ll\_\{1\},K]+p\textasciicircum{}\{\textbackslash{}ast\}\_\{12\}\textbackslash{}langle\textbackslash{}Ll\_\{2\},K]+\textbackslash{}cdots at Calcul après « se réécrit », début de la dernière ligne changes a substantive mathematical assertion or bibliographic identification; no independent derivation of this exact replacement is established in this review.\par
\par\noindent\textbf{OM-SOL-1256 -- t3 XVIII p.508 -- SOL}\par
1.3.8.6, deuxième tiret\par
\emph{Reported reading}: Pour \$\textbackslash{}Ll\$ faisceau inversible sur \$X\$ et \$K\$ un \$G\_X\$-torseur sur \$X\$\par
\emph{Proposal}: Pour \$\textbackslash{}Ll\$ faisceau inversible sur \$Y\$ et \$K\$ un \$G\_\{X\}\$-torseur sur \$X\$, on a\par
The proposed reading Pour \$\textbackslash{}Ll\$ faisceau inversible sur \$Y\$ et \$K\$ un \$G\_\{X\}\$-torseur sur \$X\$, on a at 1.3.8.6, deuxième tiret changes a substantive mathematical assertion or bibliographic identification; no independent derivation of this exact replacement is established in this review.\par
\par\noindent\textbf{OM-SOL-1258 -- t3 XVIII p.512 -- SOL}\par
Dernier paragraphe, première référence\par
\emph{Reported reading}: L'isomorphisme (1.3.12.3) est compatible à la biadditivité\par
\emph{Proposal}: (1.3.12.4)\par
The proposed reference (1.3.12.4) at Dernier paragraphe, première référence needs comparison with the statement actually invoked; that target has not been independently established here.\par
\par\noindent\textbf{OM-SOL-1259 -- t3 XVIII p.516 -- SOL}\par
Premier isomorphisme en haut de page, membre de droite\par
\emph{Reported reading}: \textbackslash{}Ll\textasciicircum{}\{\textbackslash{}otimes\textbackslash{}deg(\textbackslash{}Mm)\}\par
\emph{Proposal}: \textbackslash{}langle f\textasciicircum{}\{\textbackslash{}ast\}\textbackslash{}Ll\_0,\textbackslash{}Mm\textbackslash{}rangle\textbackslash{}xrightarrow\{\textbackslash{}sim\}\textbackslash{}Ll\_0\textasciicircum{}\{\textbackslash{}otimes\textbackslash{}deg(\textbackslash{}Mm)\}\par
The replacement \textbackslash{}langle f\textasciicircum{}\{\textbackslash{}ast\}\textbackslash{}Ll\_0,\textbackslash{}Mm\textbackslash{}rangle\textbackslash{}xrightarrow\{\textbackslash{}sim\}\textbackslash{}Ll\_0\textasciicircum{}\{\textbackslash{}otimes\textbackslash{}deg(\textbackslash{}Mm)\} at Premier isomorphisme en haut de page, membre de droite identifies a particular indexed object or symbol; that identification is not independently established by the quoted local expression.\par
\par\noindent\textbf{OM-SOL-1260 -- t3 XVIII p.516 -- SOL}\par
Isomorphisme « resp. » en haut de page, membre de droite\par
\emph{Reported reading}: \textbackslash{}Mm\textasciicircum{}\{\textbackslash{}otimes\textbackslash{}deg(\textbackslash{}Ll)\}\par
\emph{Proposal}: \textbackslash{}langle\textbackslash{}Ll,f\textasciicircum{}\{\textbackslash{}ast\}\textbackslash{}Mm\_0\textbackslash{}rangle\textbackslash{}xrightarrow\{\textbackslash{}sim\}\textbackslash{}Mm\_0\textasciicircum{}\{\textbackslash{}otimes\textbackslash{}deg(\textbackslash{}Ll)\}\par
The replacement \textbackslash{}langle\textbackslash{}Ll,f\textasciicircum{}\{\textbackslash{}ast\}\textbackslash{}Mm\_0\textbackslash{}rangle\textbackslash{}xrightarrow\{\textbackslash{}sim\}\textbackslash{}Mm\_0\textasciicircum{}\{\textbackslash{}otimes\textbackslash{}deg(\textbackslash{}Ll)\} at Isomorphisme « resp. » en haut de page, membre de droite identifies a particular indexed object or symbol; that identification is not independently established by the quoted local expression.\par
\par\noindent\textbf{OM-SOL-1261 -- t3 XVIII p.519 -- SOL}\par
Formule en haut de page, avant la dernière égalité\par
\emph{Reported reading}: ((f\_1f\_2)\textasciicircum{}\{v\_x(g\_2)\}\textbackslash{}cdot(g\_1g\_2)\textasciicircum{}\{-v\_x(f\_1)\}=\allowbreak \textbackslash{}langle f\_1f\_2,g\_1g\_2\textbackslash{}rangle\textasciicircum{}\{\textbackslash{}ast\}\par
\emph{Proposal}: ((f\_\{1\}f\_\{2\})\textasciicircum{}\{v\_\{x\}(g\_\{2\})\}\textbackslash{}cdot(g\_\{1\}g\_\{2\})\textasciicircum{}\{-v\_\{x\}(f\_\{1\})\})(x)=\allowbreak \textbackslash{}langle f\_\{1\}f\_\{2\},g\_\{1\}g\_\{2\}\textbackslash{}rangle\textasciicircum{}\{\textbackslash{}ast\}\par
The proposed reading ((f\_\{1\}f\_\{2\})\textasciicircum{}\{v\_\{x\}(g\_\{2\})\}\textbackslash{}cdot(g\_\{1\}g\_\{2\})\textasciicircum{}\{-v\_\{x\}(f\_\{1\})\})(x)=\allowbreak \textbackslash{}langle f\_\{1\}f\_\{2\},g\_\{1\}g\_\{2\}\textbackslash{}rangle\textasciicircum{}\{\textbackslash{}ast\} at Formule en haut de page, avant la dernière égalité changes a substantive mathematical assertion or bibliographic identification; no independent derivation of this exact replacement is established in this review.\par
\par\noindent\textbf{OM-SOL-1263 -- t3 XVIII p.519 -- SOL}\par
Axiome (Ass), définition de a\_\{g,h\}, premier F après la flèche\par
\emph{Reported reading}: a\_\{g,h\}:\textbackslash{}underline\{F\}(gh)\textbackslash{}xrightarrow\{\textbackslash{}sim\}\textbackslash{}underline\{F\}(\textbackslash{}underline\{F\}(g),\textbackslash{}underline\{F\}(h))\par
\emph{Proposal}: a\_\{g,h\}:\ensuremath{\mathcal{F}}(gh)\textbackslash{}xrightarrow\{\textbackslash{}sim\}F(\ensuremath{\mathcal{F}}(g),\ensuremath{\mathcal{F}}(h))\par
The proposed reading a\_\{g,h\}:\ensuremath{\mathcal{F}}(gh)\textbackslash{}xrightarrow\{\textbackslash{}sim\}F(\ensuremath{\mathcal{F}}(g),\ensuremath{\mathcal{F}}(h)) at Axiome (Ass), définition de a\_\{g,h\}, premier F après la flèche changes a substantive mathematical assertion or bibliographic identification; no independent derivation of this exact replacement is established in this review.\par
\par\noindent\textbf{OM-SOL-1264 -- t3 XVIII p.521 -- SOL}\par
Lemme 1.4.3, target of the isomorphism a\_\{n,m\}\par
\emph{Reported reading}: \textbackslash{}Sigma(n)+\textbackslash{}Sigma(m)\par
\emph{Proposal}: \textbackslash{}Sigma(\textbackslash{}underline\{n\})+\textbackslash{}Sigma(\textbackslash{}underline\{m\})\par
The replacement \textbackslash{}Sigma(\textbackslash{}underline\{n\})+\textbackslash{}Sigma(\textbackslash{}underline\{m\}) at Lemme 1.4.3, target of the isomorphism a\_\{n,m\} requires a complete check of the maps, their endpoints and the composition convention; this identification is not independently verified.\par
\par\noindent\textbf{OM-SOL-1266 -- t3 XVIII p.523 -- SOL}\par
1.4.8(a), final reference\par
\emph{Reported reading}: les compatibilités de 1.4.3.\par
\emph{Proposal}: 1.4.6\par
The proposed reference 1.4.6 at 1.4.8(a), final reference needs comparison with the statement actually invoked; that target has not been independently established here.\par
\par\noindent\textbf{OM-SOL-1269 -- t3 XVIII p.528 -- SOL}\par
1.4.14, first occurrence of the category name\par
\emph{Reported reading}: \textbackslash{}ch\textasciicircum{}\{\textbackslash{}flat\}(\textbackslash{}Ss)\par
\emph{Proposal}: \textbackslash{}Ch\textasciicircum{}\{\textbackslash{}flat\}(\textbackslash{}Ss)\par
The replacement \textbackslash{}Ch\textasciicircum{}\{\textbackslash{}flat\}(\textbackslash{}Ss) at 1.4.14, first occurrence of the category name identifies a particular indexed object or symbol; that identification is not independently established by the quoted local expression.\par
\par\noindent\textbf{OM-SOL-1270 -- t3 XVIII p.530 -- SOL}\par
Construction 1.4.20 (b), source and target of k\_i\par
\emph{Reported reading}: k\_i:\textbackslash{}ch(K'\_i)\textbackslash{}longrightarrow\textbackslash{}ch(K\_i)\par
\emph{Proposal}: k\_\{i\}:K'\_\{i\}\textbackslash{}rightarrow K\_\{i\}\par
The replacement k\_\{i\}:K'\_\{i\}\textbackslash{}rightarrow K\_\{i\} at Construction 1.4.20 (b), source and target of k\_i requires a complete check of the maps, their endpoints and the composition convention; this identification is not independently verified.\par
\par\noindent\textbf{OM-SOL-1271 -- t3 XVIII p.535 -- SOL}\par
Suite exacte (1.5.4.2), terme central\par
\emph{Reported reading}: R'f\_\{\textbackslash{}ast\}G\par
\emph{Proposal}: R\textasciicircum{}\{1\}f\_\{\textbackslash{}ast\}G\par
The replacement R\textasciicircum{}\{1\}f\_\{\textbackslash{}ast\}G at Suite exacte (1.5.4.2), terme central identifies a particular indexed object or symbol; that identification is not independently established by the quoted local expression.\par
\par\noindent\textbf{OM-SOL-1273 -- t3 XVIII p.537 -- SOL}\par
1.5.8, définition de \textbackslash{}GG\_m(\textbackslash{}mathfrak m)\par
\emph{Reported reading}: \textbackslash{}GG\_\{mX\}\par
\emph{Proposal}: \textbackslash{}GG\_\{m\textbackslash{}bar\{X\}\}\par
The replacement \textbackslash{}GG\_\{m\textbackslash{}bar\{X\}\} at 1.5.8, définition de \textbackslash{}GG\_m(\textbackslash{}mathfrak m) identifies a particular indexed object or symbol; that identification is not independently established by the quoted local expression.\par
\par\noindent\textbf{OM-SOL-1274 -- t3 XVIII p.538 -- SOL}\par
Fin de la preuve du lemme 1.5.9, groupe d'obstruction\par
\emph{Reported reading}: H\textasciicircum{}\{1\}(X,\textbackslash{}Nn(n))\par
\emph{Proposal}: H\textasciicircum{}\{1\}(\textbackslash{}bar\{X\},\textbackslash{}Nn(n))\par
The replacement H\textasciicircum{}\{1\}(\textbackslash{}bar\{X\},\textbackslash{}Nn(n)) at Fin de la preuve du lemme 1.5.9, groupe d'obstruction identifies a particular indexed object or symbol; that identification is not independently established by the quoted local expression.\par
\par\noindent\textbf{OM-SOL-1275 -- t3 XVIII p.538 -- SOL}\par
Conditions (1.5.10.1) et (1.5.10.2), domaine des sections\par
\emph{Reported reading}: \textbackslash{}Ll\textbackslash{}otimes k(s) sur X\_s, dans les deux conditions\par
\emph{Proposal}: \textbackslash{}Ll\textbackslash{}otimes k(s)\textbackslash{}text\{ sur \}\textbackslash{}bar\{X\}\_\{s\}\par
The replacement \textbackslash{}Ll\textbackslash{}otimes k(s)\textbackslash{}text\{ sur \}\textbackslash{}bar\{X\}\_\{s\} at Conditions (1.5.10.1) et (1.5.10.2), domaine des sections identifies a particular indexed object or symbol; that identification is not independently established by the quoted local expression.\par
\par\noindent\textbf{OM-SOL-1276 -- t3 XVIII p.538 -- SOL}\par
Début du lemme 1.5.11\par
\emph{Reported reading}: un faisceau inversible sur X, trivialisé le long de Y\par
\emph{Proposal}: un faisceau inversible sur \$\textbackslash{}bar\{X\}\$, trivialisé le long de \$Y\$\par
The replacement un faisceau inversible sur \$\textbackslash{}bar\{X\}\$, trivialisé le long de \$Y\$ at Début du lemme 1.5.11 identifies a particular indexed object or symbol; that identification is not independently established by the quoted local expression.\par
\par\noindent\textbf{OM-SOL-1277 -- t3 XVIII p.541 -- SOL}\par
Proposition 1.5.14, identification of the inverse functors\par
\emph{Reported reading}: les foncteurs e et j\par
\emph{Proposal}: les foncteurs \$e\$ et \$j\textasciicircum{}\{\textbackslash{}ast\}\$\par
The proposed reading les foncteurs \$e\$ et \$j\textasciicircum{}\{\textbackslash{}ast\}\$ at Proposition 1.5.14, identification of the inverse functors changes a substantive mathematical assertion or bibliographic identification; no independent derivation of this exact replacement is established in this review.\par
\par\noindent\textbf{OM-SOL-1278 -- t3 XVIII p.541 -- SOL}\par
1.5.15, the line bundle preceding the final displayed isomorphism\par
\emph{Reported reading}: faisceau inversible sur X\_2\par
\emph{Proposal}: faisceau inversible sur \$\textbackslash{}bar\{X\}\_\{2\}\$\par
The replacement faisceau inversible sur \$\textbackslash{}bar\{X\}\_\{2\}\$ at 1.5.15, the line bundle preceding the final displayed isomorphism identifies a particular indexed object or symbol; that identification is not independently established by the quoted local expression.\par
\par\noindent\textbf{OM-SOL-1279 -- t3 XVIII p.542 -- SOL}\par
The line bundle in the hypothesis preceding the norm isomorphism\par
\emph{Reported reading}: faisceau inversible sur X\_1\par
\emph{Proposal}: faisceau inversible sur \$\textbackslash{}bar\{X\}\_\{1\}\$\par
The replacement faisceau inversible sur \$\textbackslash{}bar\{X\}\_\{1\}\$ at The line bundle in the hypothesis preceding the norm isomorphism identifies a particular indexed object or symbol; that identification is not independently established by the quoted local expression.\par
\par\noindent\textbf{OM-SOL-1280 -- t3 XVIII p.544 -- SOL}\par
Paragraph after Definition 1.6.4, argument of the presheaf\par
\emph{Reported reading}: à chaque section S sur k\par
\emph{Proposal}: à chaque schéma \$S\$ sur \$k\$\par
The proposed reading à chaque schéma \$S\$ sur \$k\$ at Paragraph after Definition 1.6.4, argument of the presheaf changes a substantive mathematical assertion or bibliographic identification; no independent derivation of this exact replacement is established in this review.\par
\par\noindent\textbf{OM-SOL-1281 -- t3 XVIII p.544 -- SOL}\par
Paragraph after Definition 1.6.4, the products supporting the modules and their trivializations\par
\emph{Reported reading}: S followed by a blank product-sign position with subscript k and then X; likewise before Y\par
\emph{Proposal}: \$\textbackslash{}fprod\{S\}\{X\}\{k\}\$ et \$\textbackslash{}fprod\{S\}\{Y\}\{k\}\$\par
The record documents a lacuna at Paragraph after Definition 1.6.4, the products supporting the modules and their trivializations; it does not uniquely establish the proposed completion \$\textbackslash{}fprod\{S\}\{X\}\{k\}\$ et \$\textbackslash{}fprod\{S\}\{Y\}\{k\}\$.\par
\par\noindent\textbf{OM-SOL-1282 -- t3 XVIII p.544 -- SOL}\par
Formula (1.6.4.1), subscript of Pic\par
\emph{Reported reading}: \{\}\_\{n\}\textbackslash{}sheafPic\textasciicircum{}\{\textbackslash{}tau\}\_\{\textbackslash{}mathfrak\{m\}\_\{X/k\}\}(k)\par
\emph{Proposal}: \{\}\_\{n\}\textbackslash{}sheafPic\textasciicircum{}\{\textbackslash{}tau\}\_\{\textbackslash{}mathfrak\{m\},X/k\}(k)\par
The proposed reading \{\}\_\{n\}\textbackslash{}sheafPic\textasciicircum{}\{\textbackslash{}tau\}\_\{\textbackslash{}mathfrak\{m\},X/k\}(k) at Formula (1.6.4.1), subscript of Pic changes a substantive mathematical assertion or bibliographic identification; no independent derivation of this exact replacement is established in this review.\par
\par\noindent\textbf{OM-SOL-1283 -- t3 XVIII p.545 -- SOL}\par
1.6.5, beginning of the paragraph « Le groupe… »\par
\emph{Reported reading}: \textbackslash{}sheafPic\_\{\textbackslash{}mathfrak\{m\},X/k\}\par
\emph{Proposal}: \textbackslash{}sheafPic\_\{\textbackslash{}mathfrak\{m\},\textbackslash{}bar\{X\}/k\}\par
The replacement \textbackslash{}sheafPic\_\{\textbackslash{}mathfrak\{m\},\textbackslash{}bar\{X\}/k\} at 1.6.5, beginning of the paragraph « Le groupe… » requires a complete check of the maps, their endpoints and the composition convention; this identification is not independently verified.\par
\par\noindent\textbf{OM-SOL-1285 -- t3 XVIII p.546 -- SOL}\par
Construction of the lifted sections\par
\emph{Reported reading}: section \$s'\_\{i\}\$ de \$\textbackslash{}bar\{X\}\$ sur \$W(k)\$\par
\emph{Proposal}: section \$s'\_\{i\}\$ de \$\textbackslash{}bar\{X\}\_\{1\}\$ sur \$W(k)\$\par
The proposed reading section \$s'\_\{i\}\$ de \$\textbackslash{}bar\{X\}\_\{1\}\$ sur \$W(k)\$ at Construction of the lifted sections changes a substantive mathematical assertion or bibliographic identification; no independent derivation of this exact replacement is established in this review.\par
\par\noindent\textbf{OM-SOL-1286 -- t3 XVIII p.546 -- SOL}\par
Last paragraph, first higher direct image\par
\emph{Reported reading}: \$R\textasciicircum{}\{1\}f\_\{1\}(\textbackslash{}ZZ/n)\$\par
\emph{Proposal}: \$R\textasciicircum{}\{1\}f\_\{1\textbackslash{}ast\}(\textbackslash{}ZZ/n)\$\par
The proposed reading \$R\textasciicircum{}\{1\}f\_\{1\textbackslash{}ast\}(\textbackslash{}ZZ/n)\$ at Last paragraph, first higher direct image changes a substantive mathematical assertion or bibliographic identification; no independent derivation of this exact replacement is established in this review.\par
\par\noindent\textbf{OM-SOL-1287 -- t3 XVIII p.546 -- SOL}\par
Last paragraph, second higher direct image\par
\emph{Reported reading}: \$R\textasciicircum{}\{1\}f'\_\{1\}(\textbackslash{}ZZ/n)\$\par
\emph{Proposal}: \$R\textasciicircum{}\{1\}f'\_\{1\textbackslash{}ast\}(\textbackslash{}ZZ/n)\$\par
The proposed reading \$R\textasciicircum{}\{1\}f'\_\{1\textbackslash{}ast\}(\textbackslash{}ZZ/n)\$ at Last paragraph, second higher direct image changes a substantive mathematical assertion or bibliographic identification; no independent derivation of this exact replacement is established in this review.\par
\par\noindent\textbf{OM-SOL-1289 -- t3 XVIII p.551 -- SOL}\par
Lemme 2.5, première énumération des systèmes de paramètres\par
\emph{Reported reading}: \textbackslash{}underline\{w\}\_\{0\},\textbackslash{}ldots,w\_\{n\}\par
\emph{Proposal}: \textbackslash{}underline\{w\}\_\{0\},\textbackslash{}ldots,\textbackslash{}underline\{w\}\_\{n\}\par
The proposed reading \textbackslash{}underline\{w\}\_\{0\},\textbackslash{}ldots,\textbackslash{}underline\{w\}\_\{n\} at Lemme 2.5, première énumération des systèmes de paramètres changes a substantive mathematical assertion or bibliographic identification; no independent derivation of this exact replacement is established in this review.\par
\par\noindent\textbf{OM-SOL-1290 -- t3 XVIII p.561 -- SOL}\par
Opening reference, before the first displayed morphism\par
\emph{Reported reading}: (XVII 6.2.7.2)\par
\emph{Proposal}: \textbackslash{}eqhyperrefext[XVII][XVII.eq6.2.7.4]\{6.2.7.4\}\par
This finding intersects Orgogozo's \textbackslash{}sisi material at Opening reference, before the first displayed morphism; preserve that human entry unchanged. The proposed reference \textbackslash{}eqhyperrefext[XVII][XVII.eq6.2.7.4]\{6.2.7.4\} at Opening reference, before the first displayed morphism needs comparison with the statement actually invoked; that target has not been independently established here.\par
\par\noindent\textbf{OM-SOL-1291 -- t3 XVIII p.562 -- SOL}\par
Definition of t, term immediately after the first arrow\par
\emph{Reported reading}: \textbackslash{}sigma\_\{\textbackslash{}geq 2d\}Rf'\_V\textbackslash{}ZZ/n(d)\par
\emph{Proposal}: \textbackslash{}sigma\_\{\textbackslash{}geq 2d\}Rf'\_\{V!\}\textbackslash{}ZZ/n(d)\par
The proposed reading \textbackslash{}sigma\_\{\textbackslash{}geq 2d\}Rf'\_\{V!\}\textbackslash{}ZZ/n(d) at Definition of t, term immediately after the first arrow changes a substantive mathematical assertion or bibliographic identification; no independent derivation of this exact replacement is established in this review.\par
\par\noindent\textbf{OM-SOL-1292 -- t3 XVIII p.562 -- SOL}\par
Definition of t, argument of H\textasciicircum{}\{2d\}\par
\emph{Reported reading}: H\textasciicircum{}\{2d\}(Rf'\_V\textbackslash{}ZZ/n(d))[-2d]\par
\emph{Proposal}: H\textasciicircum{}\{2d\}(Rf'\_\{V!\}\textbackslash{}ZZ/n(d))[-2d]\par
The proposed reading H\textasciicircum{}\{2d\}(Rf'\_\{V!\}\textbackslash{}ZZ/n(d))[-2d] at Definition of t, argument of H\textasciicircum{}\{2d\} changes a substantive mathematical assertion or bibliographic identification; no independent derivation of this exact replacement is established in this review.\par
\par\noindent\textbf{OM-SOL-1294 -- t3 XVIII p.564 -- SOL}\par
Sentence immediately below the two factorisations, second trace\par
\emph{Reported reading}: \textbackslash{}Tr\_\{g'\}\par
\emph{Proposal}: \textbackslash{}Tr\_\{g'\_0\}\par
The replacement \textbackslash{}Tr\_\{g'\_0\} at Sentence immediately below the two factorisations, second trace identifies a particular indexed object or symbol; that identification is not independently established by the quoted local expression.\par
\par\noindent\textbf{OM-SOL-1295 -- t3 XVIII p.564 -- SOL}\par
Definition of f\_W at the end of the base-change sentence\par
\emph{Reported reading}: f\_w=\allowbreak g\_0h\_0\par
\emph{Proposal}: f\_\{W\}=\allowbreak g\_\{0\}h\_\{0\}\par
The replacement f\_\{W\}=\allowbreak g\_\{0\}h\_\{0\} at Definition of f\_W at the end of the base-change sentence identifies a particular indexed object or symbol; that identification is not independently established by the quoted local expression.\par
\par\noindent\textbf{OM-SOL-1296 -- t3 XVIII p.564 -- SOL}\par
First commutative diagram, upper-left object\par
\emph{Reported reading}: Rf'\textbackslash{}ZZ/n(d)\par
\emph{Proposal}: Rf'\_\{!\}\textbackslash{}ZZ/n(d)\par
The proposed reading Rf'\_\{!\}\textbackslash{}ZZ/n(d) at First commutative diagram, upper-left object changes a substantive mathematical assertion or bibliographic identification; no independent derivation of this exact replacement is established in this review.\par
\par\noindent\textbf{OM-SOL-1297 -- t3 XVIII p.564 -- SOL}\par
First commutative diagram, lower-left object\par
\emph{Reported reading}: Rg'\_\{0!\}Rh'\_1\textbackslash{}ZZ/n(d)\par
\emph{Proposal}: Rg'\_\{0!\} Rh'\_\{1!\}\textbackslash{}ZZ/n(d)\par
The proposed reading Rg'\_\{0!\} Rh'\_\{1!\}\textbackslash{}ZZ/n(d) at First commutative diagram, lower-left object changes a substantive mathematical assertion or bibliographic identification; no independent derivation of this exact replacement is established in this review.\par
\par\noindent\textbf{OM-SOL-1298 -- t3 XVIII p.565 -- SOL}\par
First sentence after the diagram, target of the canonical morphism\par
\emph{Reported reading}: Rf\_W\textbackslash{}ZZ/n(d)\par
\emph{Proposal}: Rf\_\{W!\}\textbackslash{}ZZ/n(d)\par
The proposed reading Rf\_\{W!\}\textbackslash{}ZZ/n(d) at First sentence after the diagram, target of the canonical morphism changes a substantive mathematical assertion or bibliographic identification; no independent derivation of this exact replacement is established in this review.\par
\par\noindent\textbf{OM-SOL-1299 -- t3 XVIII p.565 -- SOL}\par
Vanishing assertion, target of the morphisms on cohomology\par
\emph{Reported reading}: R\textasciicircum{}i f\_W\textbackslash{}ZZ/n(d)\par
\emph{Proposal}: R\textasciicircum{}\{i\}f\_\{W!\}\textbackslash{}ZZ/n(d)\par
The proposed reading R\textasciicircum{}\{i\}f\_\{W!\}\textbackslash{}ZZ/n(d) at Vanishing assertion, target of the morphisms on cohomology changes a substantive mathematical assertion or bibliographic identification; no independent derivation of this exact replacement is established in this review.\par
\par\noindent\textbf{OM-SOL-1301 -- t3 XVIII p.570 -- SOL}\par
Remarque 3.1.5, formule (3.1.5.1)\par
\emph{Reported reading}: Rf\textasciicircum{}\{!\}:D(S,\textbackslash{}mathcal\{A\})\textbackslash{}longrightarrow D(X,\textbackslash{}mathcal\{A\})\par
\emph{Proposal}: Rf\textasciicircum{}\{!\}:D(S,\textbackslash{}mathcal\{A\})\textbackslash{}longrightarrow D(X,f\textasciicircum{}\{\textbackslash{}ast\}\textbackslash{}mathcal\{A\})\par
The proposed reading Rf\textasciicircum{}\{!\}:D(S,\textbackslash{}mathcal\{A\})\textbackslash{}longrightarrow D(X,f\textasciicircum{}\{\textbackslash{}ast\}\textbackslash{}mathcal\{A\}) at Remarque 3.1.5, formule (3.1.5.1) changes a substantive mathematical assertion or bibliographic identification; no independent derivation of this exact replacement is established in this review.\par
\par\noindent\textbf{OM-SOL-1302 -- t3 XVIII p.572 -- SOL}\par
3.1.9, sentence immediately before (3.1.9.1)\par
\emph{Reported reading}: Le complexe \$\textbackslash{}Hom(K,L)\$ est à composantes flasques\par
\emph{Proposal}: \textbackslash{}Hhom(K,L)\par
The replacement \textbackslash{}Hhom(K,L) at 3.1.9, sentence immediately before (3.1.9.1) identifies a particular indexed object or symbol; that identification is not independently established by the quoted local expression.\par
\par\noindent\textbf{OM-SOL-1303 -- t3 XVIII p.572 -- SOL}\par
3.1.9, display following « de sorte que », right-hand side\par
\emph{Reported reading}: \textbackslash{}Hhom\textasciicircum{}\{\textbackslash{}bigdot\}(f\textasciicircum{}\{\textbackslash{}bigdot\}\_\{V!\}(k\textasciicircum{}\{\textbackslash{}ast\}\_\{X\}K),f\textasciicircum{}\{\textbackslash{}bigdot\}\_\{V!\}(k\textasciicircum{}\{\textbackslash{}ast\}\_\{X\}L))\par
\emph{Proposal}: \textbackslash{}Hom\textasciicircum{}\{\textbackslash{}bigdot\}(f\textasciicircum{}\{\textbackslash{}bigdot\}\_\{V!\}(k\textasciicircum{}\{\textbackslash{}ast\}\_\{X\}K),f\textasciicircum{}\{\textbackslash{}bigdot\}\_\{V!\}(k\textasciicircum{}\{\textbackslash{}ast\}\_\{X\}L))\par
The replacement \textbackslash{}Hom\textasciicircum{}\{\textbackslash{}bigdot\}(f\textasciicircum{}\{\textbackslash{}bigdot\}\_\{V!\}(k\textasciicircum{}\{\textbackslash{}ast\}\_\{X\}K),f\textasciicircum{}\{\textbackslash{}bigdot\}\_\{V!\}(k\textasciicircum{}\{\textbackslash{}ast\}\_\{X\}L)) at 3.1.9, display following « de sorte que », right-hand side identifies a particular indexed object or symbol; that identification is not independently established by the quoted local expression.\par
\par\noindent\textbf{OM-SOL-1304 -- t3 XVIII p.573 -- SOL}\par
(3.1.10.2), final term\par
\emph{Reported reading}: \textbackslash{}Hhom(k\textasciicircum{}\{\textbackslash{}ast\}Rf\_\{!\}K,k\textasciicircum{}\{\textbackslash{}ast\}L)\par
\emph{Proposal}: \textbackslash{}Hom(k\textasciicircum{}\{\textbackslash{}ast\}Rf\_\{!\}K,k\textasciicircum{}\{\textbackslash{}ast\}L)\par
The replacement \textbackslash{}Hom(k\textasciicircum{}\{\textbackslash{}ast\}Rf\_\{!\}K,k\textasciicircum{}\{\textbackslash{}ast\}L) at (3.1.10.2), final term identifies a particular indexed object or symbol; that identification is not independently established by the quoted local expression.\par
\par\noindent\textbf{OM-SOL-1305 -- t3 XVIII p.574 -- SOL}\par
Lemma 3.1.10.3, lower-left node, pullback of L\par
\emph{Reported reading}: Rf\textasciicircum{}\{!\}\_Vk\textasciicircum{}\{\textbackslash{}ast\}\_XL\par
\emph{Proposal}: Rf\textasciicircum{}\{!\}\_Vk\textasciicircum{}\{\textbackslash{}ast\}L\par
The proposed reading Rf\textasciicircum{}\{!\}\_Vk\textasciicircum{}\{\textbackslash{}ast\}L at Lemma 3.1.10.3, lower-left node, pullback of L changes a substantive mathematical assertion or bibliographic identification; no independent derivation of this exact replacement is established in this review.\par
\par\noindent\textbf{OM-SOL-1306 -- t3 XVIII p.574 -- SOL}\par
Lemma 3.1.10.3, lower-right node\par
\emph{Reported reading}: \textbackslash{}Hom(Rf\_!k\textasciicircum{}\{\textbackslash{}ast\}K,k\textasciicircum{}\{\textbackslash{}ast\}L)\par
\emph{Proposal}: \textbackslash{}Hom(Rf\_\{V!\}k\textasciicircum{}\{\textbackslash{}ast\}\_XK,k\textasciicircum{}\{\textbackslash{}ast\}L)\par
The replacement \textbackslash{}Hom(Rf\_\{V!\}k\textasciicircum{}\{\textbackslash{}ast\}\_XK,k\textasciicircum{}\{\textbackslash{}ast\}L) at Lemma 3.1.10.3, lower-right node requires a complete check of the maps, their endpoints and the composition convention; this identification is not independently verified.\par
\par\noindent\textbf{OM-SOL-1307 -- t3 XVIII p.574 -- SOL}\par
3.1.11, definition of the bimodule complex M\par
\emph{Reported reading}: \$(f\textasciicircum{}\{\textbackslash{}ast\}\textbackslash{}mathcal\{A\},\textbackslash{}mathcal\{B\})\$-bimodules\par
\emph{Proposal}: (g\textasciicircum{}\{\textbackslash{}ast\}\textbackslash{}mathcal\{A\},\textbackslash{}mathcal\{B\})\par
The proposed reading (g\textasciicircum{}\{\textbackslash{}ast\}\textbackslash{}mathcal\{A\},\textbackslash{}mathcal\{B\}) at 3.1.11, definition of the bimodule complex M changes a substantive mathematical assertion or bibliographic identification; no independent derivation of this exact replacement is established in this review.\par
\par\noindent\textbf{OM-SOL-1308 -- t3 XVIII p.574 -- SOL}\par
(3.1.11.4), right-hand side\par
\emph{Reported reading}: Rf\textasciicircum{}\{!\}Rg\_\{\textbackslash{}ast\}R\textbackslash{}Hom\_\{\textbackslash{}mathcal\{B\}\}(M,L)\par
\emph{Proposal}: Rf\textasciicircum{}\{!\}Rg\_\{\textbackslash{}ast\}R\textbackslash{}Hhom\_\{\textbackslash{}mathcal\{B\}\}(M,L)\par
The replacement Rf\textasciicircum{}\{!\}Rg\_\{\textbackslash{}ast\}R\textbackslash{}Hhom\_\{\textbackslash{}mathcal\{B\}\}(M,L) at (3.1.11.4), right-hand side identifies a particular indexed object or symbol; that identification is not independently established by the quoted local expression.\par
\par\noindent\textbf{OM-SOL-1311 -- t3 XVIII p.574 -- SOL}\par
Large final diagram, lower-left node, tensor product\par
\emph{Reported reading}: \textbackslash{}Hom(f'\textasciicircum{}\{\textbackslash{}ast\}M\textbackslash{}otimes\_\{\textbackslash{}Aa\}g'\textasciicircum{}\{\textbackslash{}ast\}K,Rf'\textasciicircum{}\{!\}L)\par
\emph{Proposal}: \textbackslash{}Hom(f'\textasciicircum{}\{\textbackslash{}ast\}M\textbackslash{}otimes\textasciicircum{}\{L\}\_\{\textbackslash{}Aa\}g'\textasciicircum{}\{\textbackslash{}ast\}K,Rf'\textasciicircum{}\{!\}L)\par
The replacement \textbackslash{}Hom(f'\textasciicircum{}\{\textbackslash{}ast\}M\textbackslash{}otimes\textasciicircum{}\{L\}\_\{\textbackslash{}Aa\}g'\textasciicircum{}\{\textbackslash{}ast\}K,Rf'\textasciicircum{}\{!\}L) at Large final diagram, lower-left node, tensor product requires a complete check of the maps, their endpoints and the composition convention; this identification is not independently verified.\par
\par\noindent\textbf{OM-SOL-1312 -- t3 XVIII p.575 -- SOL}\par
Corollary 3.1.12.3, parenthetical specialization\par
\emph{Reported reading}: \textbackslash{}mathcal\{B\}=\allowbreak f\textasciicircum{}\{\textbackslash{}ast\}\textbackslash{}mathcal\{A\}=\allowbreak M\par
\emph{Proposal}: \textbackslash{}mathcal\{B\}=\allowbreak g\textasciicircum{}\{\textbackslash{}ast\}\textbackslash{}mathcal\{A\}=\allowbreak M\par
The proposed reading \textbackslash{}mathcal\{B\}=\allowbreak g\textasciicircum{}\{\textbackslash{}ast\}\textbackslash{}mathcal\{A\}=\allowbreak M at Corollary 3.1.12.3, parenthetical specialization changes a substantive mathematical assertion or bibliographic identification; no independent derivation of this exact replacement is established in this review.\par
\par\noindent\textbf{OM-SOL-1313 -- t3 XVIII p.577 -- SOL}\par
3.1.14 a), isomorphism for the inverse-image functors\par
\emph{Reported reading}: Rg\_\{\textbackslash{}ast\}Rf\textasciicircum{}\{!\}\textbackslash{}xrightarrow\{\textbackslash{}sim\}Rf\textasciicircum{}\{!\}Rg\_\{\textbackslash{}ast\}\par
\emph{Proposal}: Rg'\_\{\textbackslash{}ast\}R\{f'\}\textasciicircum{}\{!\}\textbackslash{}xrightarrow\{\textbackslash{}sim\}Rf\textasciicircum{}\{!\}Rg\_\{\textbackslash{}ast\}\par
The replacement Rg'\_\{\textbackslash{}ast\}R\{f'\}\textasciicircum{}\{!\}\textbackslash{}xrightarrow\{\textbackslash{}sim\}Rf\textasciicircum{}\{!\}Rg\_\{\textbackslash{}ast\} at 3.1.14 a), isomorphism for the inverse-image functors requires a complete check of the maps, their endpoints and the composition convention; this identification is not independently verified.\par
\par\noindent\textbf{OM-SOL-1314 -- t3 XVIII p.579 -- SOL}\par
Last sentence, reference after « morphisme d'objets »\par
\emph{Reported reading}: (3.1.11.1)\par
\emph{Proposal}: \textbackslash{}eqhyperref[XVIII.par3.1.16.1]\{3.1.16.1\}\par
The proposed reference \textbackslash{}eqhyperref[XVIII.par3.1.16.1]\{3.1.16.1\} at Last sentence, reference after « morphisme d'objets » needs comparison with the statement actually invoked; that target has not been independently established here.\par
\par\noindent\textbf{OM-SOL-1315 -- t3 XVIII p.580 -- SOL}\par
(3.1.17.2), abutment\par
\emph{Reported reading}: \textbackslash{}underline\{H\}\textasciicircum{}\{n\}(K\textasciicircum{}\{!\}L)\_\{x\}\par
\emph{Proposal}: \ensuremath{\mathcal{H}}\textasciicircum{}\{n\}(RK\textasciicircum{}\{!\}L)\_\{x\}\par
The replacement \ensuremath{\mathcal{H}}\textasciicircum{}\{n\}(RK\textasciicircum{}\{!\}L)\_\{x\} at (3.1.17.2), abutment identifies a particular indexed object or symbol; that identification is not independently established by the quoted local expression.\par
\par\noindent\textbf{OM-SOL-1316 -- t3 XVIII p.580 -- SOL}\par
3.1.17(iii), reference after « morphisme d'objets »\par
\emph{Reported reading}: (3.1.11.1)\par
\emph{Proposal}: \textbackslash{}eqhyperref[XVIII.par3.1.16.1]\{3.1.16.1\}\par
The proposed reference \textbackslash{}eqhyperref[XVIII.par3.1.16.1]\{3.1.16.1\} at 3.1.17(iii), reference after « morphisme d'objets » needs comparison with the statement actually invoked; that target has not been independently established here.\par
\par\noindent\textbf{OM-SOL-1320 -- t3 XVIII p.587 -- SOL}\par
Bibliographie, second item by J.L. Verdier, « A duality theorem… »\par
\emph{Reported reading}: J.L. VERDIER [1] A duality theorem in the etale cohomology of schemes.\par
\emph{Proposal}: J.L. \textbackslash{}textsc\{Verdier\} [2]\par
The proposed reading J.L. \textbackslash{}textsc\{Verdier\} [2] at Bibliographie, second item by J.L. Verdier, « A duality theorem… » changes a substantive mathematical assertion or bibliographic identification; no independent derivation of this exact replacement is established in this review.\par
\par\noindent\textbf{OM-SOL-1321 -- t3 XIX p.589 -- SOL}\par
Resolution hypothesis, definition of Y\par
\emph{Reported reading}: \$Y=\allowbreak X-U\$\par
\emph{Proposal}: \$Y=\allowbreak X'-U\$\par
The replacement \$Y=\allowbreak X'-U\$ at Resolution hypothesis, definition of Y identifies a particular indexed object or symbol; that identification is not independently established by the quoted local expression.\par
\par\noindent\textbf{OM-SOL-1324 -- t3 XIX p.596 -- SOL}\par
First paragraph, definition of B\textasciicircum{}0\_j\par
\emph{Reported reading}: le normalisé ... de \$A[Y]/(f\textasciicircum{}\{0\}\_\{j\})\$\par
\emph{Proposal}: \$A\textasciicircum{}\{0\}[Y]/(f\textasciicircum{}\{0\}\_\{j\})\$\par
The replacement \$A\textasciicircum{}\{0\}[Y]/(f\textasciicircum{}\{0\}\_\{j\})\$ at First paragraph, definition of B\textasciicircum{}0\_j identifies a particular indexed object or symbol; that identification is not independently established by the quoted local expression.\par
\par\noindent\textbf{OM-SOL-1325 -- t3 XIX p.600 -- SOL}\par
First paragraph after the opening diagram\par
\emph{Reported reading}: Pour démontrer que \$c'\$ est surjectif\par
\emph{Proposal}: \$c\$\par
The replacement \$c\$ at First paragraph after the opening diagram requires a complete check of the maps, their endpoints and the composition convention; this identification is not independently verified.\par
\par\noindent\textbf{OM-SOL-1327 -- t3 XIX p.605 -- SOL}\par
Leray spectral sequence following (2.6), abutment\par
\emph{Reported reading}: H\textasciicircum{}\{p+q\}(U',\textbackslash{}mu\_\{n\})\par
\emph{Proposal}: H\textasciicircum{}\{p+q\}(\textbackslash{}hat\{U\},\ensuremath{𝛍}\_\{n\})\par
The replacement H\textasciicircum{}\{p+q\}(\textbackslash{}hat\{U\},\ensuremath{𝛍}\_\{n\}) at Leray spectral sequence following (2.6), abutment identifies a particular indexed object or symbol; that identification is not independently established by the quoted local expression.\par
\par\noindent\textbf{OM-SOL-1328 -- t3 XIX p.605 -- SOL}\par
Display after « donnera des isomorphismes », target\par
\emph{Reported reading}: H\textasciicircum{}\{q\}(U',\textbackslash{}mu\_\{n\})\par
\emph{Proposal}: H\textasciicircum{}\{q\}(\textbackslash{}hat\{U\},\ensuremath{𝛍}\_\{n\})\par
The replacement H\textasciicircum{}\{q\}(\textbackslash{}hat\{U\},\ensuremath{𝛍}\_\{n\}) at Display after « donnera des isomorphismes », target identifies a particular indexed object or symbol; that identification is not independently established by the quoted local expression.\par
\par\noindent\textbf{OM-SOL-1329 -- t3 XIX p.605 -- SOL}\par
Exact sequence after « et une suite exacte », middle cohomology group\par
\emph{Reported reading}: H\textasciicircum{}\{N+1\}(U',\textbackslash{}mu\_\{n\})\par
\emph{Proposal}: H\textasciicircum{}\{N+1\}(\textbackslash{}hat\{U\},\ensuremath{𝛍}\_\{n\})\par
The replacement H\textasciicircum{}\{N+1\}(\textbackslash{}hat\{U\},\ensuremath{𝛍}\_\{n\}) at Exact sequence after « et une suite exacte », middle cohomology group identifies a particular indexed object or symbol; that identification is not independently established by the quoted local expression.\par
\par\noindent\textbf{OM-SOL-1331 -- t3 XIX p.609 -- SOL}\par
3.3, ordre des groupes cycliques\par
\emph{Reported reading}: groupes cycliques d'ordres n\par
\emph{Proposal}: groupes cycliques d'ordre \$n\$\par
The proposed reading groupes cycliques d'ordre \$n\$ at 3.3, ordre des groupes cycliques changes a substantive mathematical assertion or bibliographic identification; no independent derivation of this exact replacement is established in this review.\par
\par\noindent\textbf{OM-SOL-1334 -- t3 XIX p.619 -- SOL}\par
Corollaire 6.3, définition de R(A)\par
\emph{Reported reading}: Soit \$R(A)\$ le corps des fonctions rationnelles sur \$\textbackslash{}Spec A\$.\par
\emph{Proposal}: Soit \$R(A)\$ l'anneau des fonctions rationnelles sur \$\textbackslash{}Spec A\$.\par
The proposed reading Soit \$R(A)\$ l'anneau des fonctions rationnelles sur \$\textbackslash{}Spec A\$. at Corollaire 6.3, définition de R(A) changes a substantive mathematical assertion or bibliographic identification; no independent derivation of this exact replacement is established in this review.\par
\par\noindent\textbf{OM-SOL-1335 -- t3 XIX p.623 -- SOL}\par
Suite spectrale en tête de page, terme d'aboutissement\par
\emph{Reported reading}: \$P\textasciicircum{}\{p+q\}(\textbackslash{}EE\textasciicircum{}\{1\},F)\$\par
\emph{Proposal}: H\textasciicircum{}\{p+q\}(\textbackslash{}EE\textasciicircum{}\{1\},F)\par
The proposed reading H\textasciicircum{}\{p+q\}(\textbackslash{}EE\textasciicircum{}\{1\},F) at Suite spectrale en tête de page, terme d'aboutissement changes a substantive mathematical assertion or bibliographic identification; no independent derivation of this exact replacement is established in this review.\par
\par\noindent\textbf{OM-SOL-1336 -- t3 XIX p.623 -- SOL}\par
Même phrase, schéma sur lequel la section x est prise\par
\emph{Reported reading}: \$x\textbackslash{}in\textbackslash{}Gamma(X,O\_\{X\})\$\par
\emph{Proposal}: x\textbackslash{}in\textbackslash{}Gamma(\textbackslash{}widetilde\{X\},O\_\{\textbackslash{}widetilde\{X\}\})\par
The replacement x\textbackslash{}in\textbackslash{}Gamma(\textbackslash{}widetilde\{X\},O\_\{\textbackslash{}widetilde\{X\}\}) at Même phrase, schéma sur lequel la section x est prise identifies a particular indexed object or symbol; that identification is not independently established by the quoted local expression.\par
\par\noindent\textbf{OM-SOL-1338 -- t3 XIX p.625 -- SOL}\par
Dernier paragraphe, extension de corps engendrée par b\par
\emph{Reported reading}: \$R(A)\$ sur \$k((x))\$\par
\emph{Proposal}: R(B)\par
The proposed reading R(B) at Dernier paragraphe, extension de corps engendrée par b changes a substantive mathematical assertion or bibliographic identification; no independent derivation of this exact replacement is established in this review.\par
\par\noindent\textbf{OM-SOL-1339 -- t3 XIX p.627 -- SOL}\par
“applicable au morphisme”, before the vanishing formula\par
\emph{Reported reading}: au morphisme f\textasciicircum{}\{0\}\par
\emph{Proposal}: au morphisme \$f\textasciicircum{}\{0\}\_\{\textbackslash{}alpha\}\$\par
The proposed reading au morphisme \$f\textasciicircum{}\{0\}\_\{\textbackslash{}alpha\}\$ at “applicable au morphisme”, before the vanishing formula changes a substantive mathematical assertion or bibliographic identification; no independent derivation of this exact replacement is established in this review.\par
\par\noindent\textbf{OM-SOL-1340 -- t3 XIX p.632 -- SOL}\par
Dernière inégalité affichée, avant « Par conséquent »\par
\emph{Reported reading}: (R\textasciicircum{}\{q\}i\_\{\textbackslash{}ast\}\textbackslash{}ZZ/\textbackslash{}ell)\textbackslash{}leq d-q\par
\emph{Proposal}: \textbackslash{}delta(R\textasciicircum{}\{q\}i\_\{\textbackslash{}ast\}\textbackslash{}ZZ/\textbackslash{}ell)\textbackslash{}leq d-q\par
The proposed reading \textbackslash{}delta(R\textasciicircum{}\{q\}i\_\{\textbackslash{}ast\}\textbackslash{}ZZ/\textbackslash{}ell)\textbackslash{}leq d-q at Dernière inégalité affichée, avant « Par conséquent » changes a substantive mathematical assertion or bibliographic identification; no independent derivation of this exact replacement is established in this review.\par
\par\noindent\textbf{OM-SOL-1341 -- t3 XIX p.633 -- SOL}\par
Après (7.7), description du fermé C à l'infini\par
\emph{Reported reading}: et s'identifie dans l'autre ouvert Spec B[x/y] de Z à l'ouvert C=\allowbreak V(x/y)\par
\emph{Proposal}: qui s'identifie dans l'autre ouvert \$\textbackslash{}Spec B[x/y]\$ de \$Z\$ au fermé \$C=\allowbreak V(x/y)\$\par
The proposed reading qui s'identifie dans l'autre ouvert \$\textbackslash{}Spec B[x/y]\$ de \$Z\$ au fermé \$C=\allowbreak V(x/y)\$ at Après (7.7), description du fermé C à l'infini changes a substantive mathematical assertion or bibliographic identification; no independent derivation of this exact replacement is established in this review.\par
\par\noindent\textbf{OM-SOL-1342 -- t3 XIX p.637 -- SOL}\par
Seconde occurrence de k((y,z,x)), après « extension génériquement étale de »\par
\emph{Reported reading}: extension génériquement étale de k((y,z,x))\par
\emph{Proposal}: extension génériquement étale de \$k[[y,z,x]]\$\par
The proposed reading extension génériquement étale de \$k[[y,z,x]]\$ at Seconde occurrence de k((y,z,x)), après « extension génériquement étale de » changes a substantive mathematical assertion or bibliographic identification; no independent derivation of this exact replacement is established in this review.\par
\par\noindent\textbf{OM-SOL-1343 -- t3 XIX p.638 -- SOL}\par
Bibliographie [2], nom de l'auteur\par
\emph{Reported reading}: ABYANKHAR, S.\par
\emph{Proposal}: \textbackslash{}textsc\{Abhyankar\}, S.\par
The proposed reading \textbackslash{}textsc\{Abhyankar\}, S. at Bibliographie [2], nom de l'auteur changes a substantive mathematical assertion or bibliographic identification; no independent derivation of this exact replacement is established in this review.\par
\endgroup
